%%%%%%%%%%%%%%%%%%%%%%%%%%%%%%%%%%%%%%%%%
% The Legrand Orange Book
% LaTeX Template
% Version 2.3 (8/8/17)
%
% This template has been downloaded from:
% http://www.LaTeXTemplates.com
%
% Original author:
% Mathias Legrand (legrand.mathias@gmail.com) with modifications by:
% Vel (vel@latextemplates.com)
%
% License:
% CC BY-NC-SA 3.0 (http://creativecommons.org/licenses/by-nc-sa/3.0/)
%
% Compiling this template:
% This template uses biber for its bibliography and makeindex for its index.
% When you first open the template, compile it from the command line with the 
% commands below to make sure your LaTeX distribution is configured correctly:
%
% 1) pdflatex main
% 2) makeindex main.idx -s StyleInd.ist
% 3) biber main
% 4) pdflatex main x 2
%
% After this, when you wish to update the bibliography/index use the appropriate
% command above and make sure to compile with pdflatex several times 
% afterwards to propagate your changes to the document.
%
% This template also uses a number of packages which may need to be
% updated to the newest versions for the template to compile. It is strongly
% recommended you update your LaTeX distribution if you have any
% compilation errors.
%
% Important note:
% Chapter heading images should have a 2:1 width:height ratio,
% e.g. 920px width and 460px height.
%
%%%%%%%%%%%%%%%%%%%%%%%%%%%%%%%%%%%%%%%%%

%----------------------------------------------------------------------------------------
%	PACKAGES AND OTHER DOCUMENT CONFIGURATIONS
%----------------------------------------------------------------------------------------

\documentclass[11pt]{book} % Default font size and left-justified equations

%----------------------------------------------------------------------------------------

%%%%%%%%%%%%%%%%%%%%%%%%%%%%%%%%%%%%%%%%%
% The Legrand Orange Book
% Structural Definitions File
% Version 2.0 (9/2/15)
%
% Original author:
% Mathias Legrand (legrand.mathias@gmail.com) with modifications by:
% Vel (vel@latextemplates.com)
% 
% This file has been downloaded from:
% http://www.LaTeXTemplates.com
%
% License:
% CC BY-NC-SA 3.0 (http://creativecommons.org/licenses/by-nc-sa/3.0/)
%
%%%%%%%%%%%%%%%%%%%%%%%%%%%%%%%%%%%%%%%%%

%----------------------------------------------------------------------------------------
%	VARIOUS REQUIRED PACKAGES AND CONFIGURATIONS
%----------------------------------------------------------------------------------------

\usepackage[top=3cm,bottom=3cm,left=3cm,right=3cm,headsep=10pt,a4paper]{geometry} % Page margins

\usepackage{graphicx} % Required for including pictures
\graphicspath{{Pictures/}} % Specifies the directory where pictures are stored

\usepackage{lipsum} % Inserts dummy text
\usepackage{yhmath}
\usepackage{tikz} % Required for drawing custom shapes

\usepackage{tikz-cd}
\tikzset{
  subseteq/.style={
    draw=none,
    edge node={node [sloped, allow upside down, auto=false]{$\subseteq$}}},
  Subseteq/.style={
    draw=none,
    every to/.append style={
      edge node={node [sloped, allow upside down, auto=false]{$\subseteq$}}}
    },
    Subsetneq/.style={
    draw=none,
    every to/.append style={
      edge node={node [sloped, allow upside down, auto=false]{$\subsetneq$}}}
    },
  Supseteq/.style={
    draw=none,
    every to/.append style={
      edge node={node [sloped, allow upside down, auto=false]{$\supseteq$}}}
  }
}

\usepackage[english]{babel} % English language/hyphenation

\usepackage{enumitem} % Customize lists
\setlist{nolistsep} % Reduce spacing between bullet points and numbered lists

\usepackage{booktabs} % Required for nicer horizontal rules in tables

\usepackage{xcolor} % Required for specifying colors by name
\definecolor{ocre}{RGB}{51, 153, 255} % Define the orange color used for highlighting throughout the book

%----------------------------------------------------------------------------------------
%	FONTS
%----------------------------------------------------------------------------------------


\usepackage{avant} % Use the Avantgarde font for headings
%\usepackage{times} % Use the Times font for headings
\DeclareMathAlphabet{\mathpzc}{OT1}{pzc}{m}{it}
\usepackage{mathptmx} % Use the Adobe Times Roman as the default text font together with math symbols from the Sym­bol, Chancery and Com­puter Modern fonts

\usepackage{microtype} % Slightly tweak font spacing for aesthetics
\usepackage[utf8]{inputenc} % Required for including letters with accents
\usepackage[T1]{fontenc} % Use 8-bit encoding that has 256 glyphs

%----------------------------------------------------------------------------------------
%	BIBLIOGRAPHY AND INDEX
%----------------------------------------------------------------------------------------

\usepackage[style=numeric,citestyle=numeric,sorting=nyt,sortcites=true,autopunct=true,babel=hyphen,hyperref=true,abbreviate=false,backref=true,backend=biber]{biblatex}
\addbibresource{bibliography.bib} % BibTeX bibliography file
\defbibheading{bibempty}{}

\usepackage{calc} % For simpler calculation - used for spacing the index letter headings correctly
\usepackage{makeidx} % Required to make an index
\makeindex % Tells LaTeX to create the files required for indexing

%----------------------------------------------------------------------------------------
%	MAIN TABLE OF CONTENTS
%----------------------------------------------------------------------------------------

\usepackage{titletoc} % Required for manipulating the table of contents

\contentsmargin{0cm} % Removes the default margin

% Part text styling
\titlecontents{part}[0cm]
{\addvspace{20pt}\centering\large\bfseries}
{}
{}
{}

% Chapter text styling
\titlecontents{chapter}[1.25cm] % Indentation
{\addvspace{12pt}\large\sffamily\bfseries} % Spacing and font options for chapters
{\color{ocre!60}\contentslabel[\Large\thecontentslabel]{1.25cm}\color{ocre}} % Chapter number
{\color{ocre}}  
{\color{ocre!60}\normalsize\;\titlerule*[.5pc]{.}\;\thecontentspage} % Page number

% Section text styling
\titlecontents{section}[1.25cm] % Indentation
{\addvspace{3pt}\sffamily\bfseries} % Spacing and font options for sections
{\contentslabel[\thecontentslabel]{1.25cm}} % Section number
{}
{\hfill\color{black}\thecontentspage} % Page number
[]

% Subsection text styling
\titlecontents{subsection}[1.25cm] % Indentation
{\addvspace{1pt}\sffamily\small} % Spacing and font options for subsections
{\contentslabel[\thecontentslabel]{1.25cm}} % Subsection number
{}
{\ \titlerule*[.5pc]{.}\;\thecontentspage} % Page number
[]

% List of figures
\titlecontents{figure}[0em]
{\addvspace{-5pt}\sffamily}
{\thecontentslabel\hspace*{1em}}
{}
{\ \titlerule*[.5pc]{.}\;\thecontentspage}
[]

% List of tables
\titlecontents{table}[0em]
{\addvspace{-5pt}\sffamily}
{\thecontentslabel\hspace*{1em}}
{}
{\ \titlerule*[.5pc]{.}\;\thecontentspage}
[]

%----------------------------------------------------------------------------------------
%	MINI TABLE OF CONTENTS IN PART HEADS
%----------------------------------------------------------------------------------------

% Chapter text styling
\titlecontents{lchapter}[0em] % Indenting
{\addvspace{15pt}\large\sffamily\bfseries} % Spacing and font options for chapters
{\color{ocre}\contentslabel[\Large\thecontentslabel]{1.25cm}\color{ocre}} % Chapter number
{}  
{\color{ocre}\normalsize\sffamily\bfseries\;\titlerule*[.5pc]{.}\;\thecontentspage} % Page number

% Section text styling
\titlecontents{lsection}[0em] % Indenting
{\sffamily\small} % Spacing and font options for sections
{\contentslabel[\thecontentslabel]{1.25cm}} % Section number
{}
{}

% Subsection text styling
\titlecontents{lsubsection}[.5em] % Indentation
{\normalfont\footnotesize\sffamily} % Font settings
{}
{}
{}

%----------------------------------------------------------------------------------------
%	PAGE HEADERS
%----------------------------------------------------------------------------------------

\usepackage{fancyhdr} % Required for header and footer configuration

\pagestyle{fancy}
\renewcommand{\chaptermark}[1]{\markboth{\sffamily\normalsize\bfseries\chaptername\ \thechapter.\ #1}{}} % Chapter text font settings
\renewcommand{\sectionmark}[1]{\markright{\sffamily\normalsize\thesection\hspace{5pt}#1}{}} % Section text font settings
\fancyhf{} \fancyhead[LE,RO]{\sffamily\normalsize\thepage} % Font setting for the page number in the header
\fancyhead[LO]{\rightmark} % Print the nearest section name on the left side of odd pages
\fancyhead[RE]{\leftmark} % Print the current chapter name on the right side of even pages
\renewcommand{\headrulewidth}{0.5pt} % Width of the rule under the header
\addtolength{\headheight}{2.5pt} % Increase the spacing around the header slightly
\renewcommand{\footrulewidth}{0pt} % Removes the rule in the footer
\fancypagestyle{plain}{\fancyhead{}\renewcommand{\headrulewidth}{0pt}} % Style for when a plain pagestyle is specified

% Removes the header from odd empty pages at the end of chapters
\makeatletter
\renewcommand{\cleardoublepage}{
\clearpage\ifodd\c@page\else
\hbox{}
\vspace*{\fill}
\thispagestyle{empty}
\newpage
\fi}

%----------------------------------------------------------------------------------------
%	THEOREM STYLES
%----------------------------------------------------------------------------------------
\usepackage{stmaryrd}
\usepackage{amsmath,amsfonts,amssymb,amsthm} % For math equations, theorems, symbols, etc


\newcommand{\intoo}[2]{\mathopen{]}#1\,;#2\mathclose{[}}
\newcommand{\ud}{\mathop{\mathrm{{}d}}\mathopen{}}
\newcommand{\intff}[2]{\mathopen{[}#1\,;#2\mathclose{]}}
\newtheorem{notation}{Notation}[chapter]

% Boxed/framed environments
\newtheoremstyle{ocrenumbox}% % Theorem style name
{0pt}% Space above
{0pt}% Space below
{\normalfont}% % Body font
{}% Indent amount
{\small\bf\sffamily\color{ocre}}% % Theorem head font
{\;}% Punctuation after theorem head
{0.25em}% Space after theorem head
{\small\sffamily\color{ocre}\thmname{#1}\nobreakspace\thmnumber{\@ifnotempty{#1}{}\@upn{#2}}% Theorem text (e.g. Theorem 2.1)
\thmnote{\nobreakspace\the\thm@notefont\sffamily\bfseries\color{black}---\nobreakspace#3.}} % Optional theorem note

\DeclareRobustCommand{\coprod}{\mathop{\text{\fakecoprod}}}
\newcommand{\fakecoprod}{%
  \sbox0{$\prod$}%
  \smash{\raisebox{\dimexpr.9625\depth-\dp0}{\scalebox{1}[-1]{$\prod$}}}%
  \vphantom{$\prod$}%
}

%\makeatletter
\newcommand{\colim@}[2]{%
  \vtop{\m@th\ialign{##\cr
    \hfil$#1\operator@font colim$\hfil\cr
    \noalign{\nointerlineskip\kern1.5\ex@}#2\cr
    \noalign{\nointerlineskip\kern-\ex@}\cr}}%
}
\newcommand{\colim}{%
  \mathop{\mathpalette\colim@{\rightarrowfill@\scriptscriptstyle}}\nmlimits@
}
\renewcommand{\varprojlim}{%
  \mathop{\mathpalette\varlim@{\leftarrowfill@\scriptscriptstyle}}\nmlimits@
}

\renewcommand{\varinjlim}{%
  \mathop{\mathpalette\varlim@{\rightarrowfill@\scriptscriptstyle}}\nmlimits@
}
\newcommand{\dirlim}{\varinjlim}
\newcommand{\invlim}{\varprojlim}

\renewcommand{\qedsymbol}{$\blacksquare$}% Optional qed square
\renewcommand{\baselinestretch}{1.1}
\newcommand{\affn}{\mathbb A}
\newcommand{\reals}{\mathbb R}
\newcommand{\cplx}{\mathbb C}
\newcommand{\proj}{\mathbb P}
\newcommand{\intg}{\mathbb Z}
\newcommand{\bbf}{\mathbb F}
\newcommand{\ratl}{\mathbb Q}
\newcommand{\torus}{\mathbb T}
\newcommand{\sca}{{\mathfrak a}}
\newcommand{\scb}{{\mathfrak b}}
\newcommand{\scc}{{\mathfrak c}}
\newcommand{\scm}{{\mathfrak m}}
\newcommand{\scn}{{\mathfrak n}}
\newcommand{\scp}{{\mathfrak p}}
\newcommand{\scq}{\mathfrak q}
\newcommand{\frakg}{{\mathfrak g}}
\newcommand{\frakd}{{\mathfrak d}}
\newcommand{\cala}{{\cal A}}
\newcommand{\calb}{{\cal B}}
\newcommand{\calc}{{\cal C}}
\newcommand{\cale}{{\cal E}}
\newcommand{\calf}{{\cal F}}
\newcommand{\calg}{{\cal G}}
\newcommand{\calh}{{\cal H}}
\newcommand{\cali}{{\cal I}}
\newcommand{\call}{{\cal L}}
\newcommand{\caln}{{\cal N}}
\newcommand{\calo}{{\cal O}}
\newcommand{\pd}{{\partial}}
\newcommand{\calr}{{\cal R}}
\newcommand{\mathbold}{\bf}
\newcommand{\Proj}{\text{ Proj\,}}
\newcommand{\res}{\text{res}}
\newcommand{\sym}{\text{ Sym\,}}
\newcommand{\spec}{\text{ Spec\,}}
\newcommand{\tor}{\text{ Tor\,}}
\newcommand{\ext}{\text{ Ext\,}}
\renewcommand{\char}{\text{char\,}}
\newcommand{\mor}{\text{ Mor}}
\renewcommand{\hom}{\text{ Hom}}
\newcommand{\phom}{\mathpzc{Hom\,}}
\newcommand{\pzf}{{\mathpzc F}}
\newcommand{\pzg}{{\mathpzc G}}
\newcommand{\pzm}{{\mathpzc M}}
\newcommand{\dims}{\text{dim\,}}
\newcommand{\coker}{\text{coker\,}}
\newcommand{\im}{{\text im\thinspace}}
\newcommand{\supp}{{\text Supp\thinspace}}
\newcommand{\cinf}{C^{\infty}}
\newcommand{\row}[2]{#1_1,\dots ,#1_{#2}}
\newcommand{\dbyd}[2]{{\partial #1\over\partial #2}}
\newcommand{\Space}{{\bf Space}}
\newcommand{\alg}{{\mathbold Alg}}
\newcommand{\notsubset}{\not \subset}
\newcommand{\notsupset}{\not \supset}
\newcommand{\pois}{{\mathbold Pois}}
\newcommand{\pitilde}{\tilde{\pi}}
\newcommand{\rta}{\rightarrow}
\newcommand{\Lrta}{\Longrightarrow}
\newcommand{\lrta}{\longrightarrow}
\newcommand{\llta}{\longleftarrow}
\newcommand{\llrta}{\longleftrightarrow}
\newcommand{\Llta}{\Longleftarrow}
\newcommand{\Llrta}{\Longleftrightarrow}
\newcommand{\lgl}{\langle}
\newcommand{\rgl}{\rangle}
\newcommand{\inj}{\hookrightarrow}
\newcommand{\Spec}{{\text Spec}}
\newcommand{\rad}{{\text rad}}
\newcommand{\Jac}{{\text Jac}}
\newcommand{\surj}{\twoheadrightarrow}
\newcommand{\cmark}{\ding{51}}%
\newcommand{\xmark}{\ding{55}}%
\newcommand{\downmapsto}{\rotatebox[origin=c]{-90}{$\scriptstyle\mapsto$}\mkern2mu}

\newtheoremstyle{blacknumex}% Theorem style name
{5pt}% Space above
{5pt}% Space below
{\normalfont}% Body font
{} % Indent amount
{\small\bf\sffamily}% Theorem head font
{\;}% Punctuation after theorem head
{0.25em}% Space after theorem head
{\small\sffamily{\tiny\ensuremath{\blacksquare}}\nobreakspace\thmname{#1}\nobreakspace\thmnumber{\@ifnotempty{#1}{}\@upn{#2}}% Theorem text (e.g. Theorem 2.1)
\thmnote{\nobreakspace\the\thm@notefont\sffamily\bfseries---\nobreakspace#3.}}% Optional theorem note

\newtheoremstyle{blacknumbox} % Theorem style name
{0pt}% Space above
{0pt}% Space below
{\normalfont}% Body font
{}% Indent amount
{\small\bf\sffamily}% Theorem head font
{\;}% Punctuation after theorem head
{0.25em}% Space after theorem head
{\small\sffamily\thmname{#1}\nobreakspace\thmnumber{\@ifnotempty{#1}{}\@upn{#2}}% Theorem text (e.g. Theorem 2.1)
\thmnote{\nobreakspace\the\thm@notefont\sffamily\bfseries---\nobreakspace#3.}}% Optional theorem note

% Non-boxed/non-framed environments
\newtheoremstyle{ocrenum}% % Theorem style name
{5pt}% Space above
{5pt}% Space below
{\normalfont}% % Body font
{}% Indent amount
{\small\bf\sffamily\color{ocre}}% % Theorem head font
{\;}% Punctuation after theorem head
{0.25em}% Space after theorem head
{\small\sffamily\color{ocre}\thmname{#1}\nobreakspace\thmnumber{\@ifnotempty{#1}{}\@upn{#2}}% Theorem text (e.g. Theorem 2.1)
\thmnote{\nobreakspace\the\thm@notefont\sffamily\bfseries\color{black}---\nobreakspace#3.}} % Optional theorem note
\renewcommand{\qedsymbol}{$\blacksquare$}% Optional qed square
\makeatother

% Defines the theorem text style for each type of theorem to one of the three styles above
\newcounter{dummy} 
\numberwithin{dummy}{section}
\theoremstyle{ocrenumbox}
\newtheorem{theoremeT}[dummy]{Theorem}
\newtheorem{problem}{Problem}[chapter]
\newtheorem{exerciseT}{Exercise??}[section]
\theoremstyle{blacknumex}
\newtheorem{exampleT}{Example}[section]
\theoremstyle{blacknumbox}
\newtheorem{vocabulary}{Vocabulary}[chapter]
\newtheorem{definitionT}{Definition}[section]
\newtheorem{corollaryT}[dummy]{Corollary}
\theoremstyle{ocrenum}
\newtheorem{proposition}[dummy]{Proposition}
\newtheorem{lemma}[dummy]{Lemma}

\renewcommand*{\theexerciseT}{\thesection.\Alph{exerciseT}}
%----------------------------------------------------------------------------------------
%	DEFINITION OF COLORED BOXES
%----------------------------------------------------------------------------------------

\RequirePackage[framemethod=default]{mdframed} % Required for creating the theorem, definition, exercise and corollary boxes

% Theorem box
\newmdenv[skipabove=7pt,
skipbelow=7pt,
backgroundcolor=black!5,
linecolor=ocre,
innerleftmargin=5pt,
innerrightmargin=5pt,
innertopmargin=5pt,
leftmargin=0cm,
rightmargin=0cm,
innerbottommargin=5pt]{tBox}

% Exercise box	  
\newmdenv[skipabove=7pt,
skipbelow=7pt,
rightline=false,
leftline=true,
topline=false,
bottomline=false,
backgroundcolor=ocre!10,
linecolor=ocre,
innerleftmargin=5pt,
innerrightmargin=5pt,
innertopmargin=5pt,
innerbottommargin=5pt,
leftmargin=0cm,
rightmargin=0cm,
linewidth=4pt]{eBox}	

% Definition box
\newmdenv[skipabove=7pt,
skipbelow=7pt,
rightline=false,
leftline=true,
topline=false,
bottomline=false,
linecolor=ocre,
innerleftmargin=5pt,
innerrightmargin=5pt,
innertopmargin=0pt,
leftmargin=0cm,
rightmargin=0cm,
linewidth=4pt,
innerbottommargin=0pt]{dBox}	

% Corollary box
\newmdenv[skipabove=7pt,
skipbelow=7pt,
rightline=false,
leftline=true,
topline=false,
bottomline=false,
linecolor=gray,
backgroundcolor=black!5,
innerleftmargin=5pt,
innerrightmargin=5pt,
innertopmargin=5pt,
leftmargin=0cm,
rightmargin=0cm,
linewidth=4pt,
innerbottommargin=5pt]{cBox}

% Creates an environment for each type of theorem and assigns it a theorem text style from the "Theorem Styles" section above and a colored box from above
\newenvironment{theorem}{\begin{tBox}\begin{theoremeT}}{\end{theoremeT}\end{tBox}}
\newenvironment{exr}{\begin{eBox}\begin{exerciseT}}{\hfill{\color{ocre}\tiny\ensuremath{\blacksquare}}\end{exerciseT}\end{eBox}}				  
\newenvironment{definition}{\begin{dBox}\begin{definitionT}}{\end{definitionT}\end{dBox}}	
\newenvironment{example}{\begin{exampleT}}{\hfill{\tiny\ensuremath{\blacksquare}}\end{exampleT}}		
\newenvironment{corollary}{\begin{cBox}\begin{corollaryT}}{\end{corollaryT}\end{cBox}}	

%----------------------------------------------------------------------------------------
%	REMARK ENVIRONMENT
%----------------------------------------------------------------------------------------

\newenvironment{remark}{\par\vspace{10pt}\small % Vertical white space above the remark and smaller font size
\begin{list}{}{
\leftmargin=35pt % Indentation on the left
\rightmargin=25pt}\item\ignorespaces % Indentation on the right
\makebox[-2.5pt]{\begin{tikzpicture}[overlay]
\node[draw=ocre!60,line width=1pt,circle,fill=ocre!25,font=\sffamily\bfseries,inner sep=2pt,outer sep=0pt] at (-15pt,0pt){\textcolor{ocre}{R}};\end{tikzpicture}} % Orange R in a circle
\advance\baselineskip -1pt}{\end{list}\vskip5pt} % Tighter line spacing and white space after remark

%----------------------------------------------------------------------------------------
%	SECTION NUMBERING IN THE MARGIN
%----------------------------------------------------------------------------------------

\makeatletter
\renewcommand{\@seccntformat}[1]{\llap{\textcolor{ocre}{\csname the#1\endcsname}\hspace{1em}}}                    
\renewcommand{\section}{\@startsection{section}{1}{\z@}
{-4ex \@plus -1ex \@minus -.4ex}
{1ex \@plus.2ex }
{\normalfont\large\sffamily\bfseries}}
\renewcommand{\subsection}{\@startsection {subsection}{2}{\z@}
{-3ex \@plus -0.1ex \@minus -.4ex}
{0.5ex \@plus.2ex }
{\normalfont\sffamily\bfseries}}
\renewcommand{\subsubsection}{\@startsection {subsubsection}{3}{\z@}
{-2ex \@plus -0.1ex \@minus -.2ex}
{.2ex \@plus.2ex }
{\normalfont\small\sffamily\bfseries}}                        
\renewcommand\paragraph{\@startsection{paragraph}{4}{\z@}
{-2ex \@plus-.2ex \@minus .2ex}
{.1ex}
{\normalfont\small\sffamily\bfseries}}

%----------------------------------------------------------------------------------------
%	PART HEADINGS
%----------------------------------------------------------------------------------------

% numbered part in the table of contents
\newcommand{\@mypartnumtocformat}[2]{%
\setlength\fboxsep{0pt}%
\noindent\colorbox{ocre!20}{\strut\parbox[c][.7cm]{\ecart}{\color{ocre!70}\Large\sffamily\bfseries\centering#1}}\hskip\esp\colorbox{ocre!40}{\strut\parbox[c][.7cm]{\linewidth-\ecart-\esp}{\Large\sffamily\centering#2}}}%
%%%%%%%%%%%%%%%%%%%%%%%%%%%%%%%%%%
% unnumbered part in the table of contents
\newcommand{\@myparttocformat}[1]{%
\setlength\fboxsep{0pt}%
\noindent\colorbox{ocre!40}{\strut\parbox[c][.7cm]{\linewidth}{\Large\sffamily\centering#1}}}%
%%%%%%%%%%%%%%%%%%%%%%%%%%%%%%%%%%
\newlength\esp
\setlength\esp{4pt}
\newlength\ecart
\setlength\ecart{1.2cm-\esp}
\newcommand{\thepartimage}{}%
\newcommand{\partimage}[1]{\renewcommand{\thepartimage}{#1}}%
\def\@part[#1]#2{%
\ifnum \c@secnumdepth >-2\relax%
\refstepcounter{part}%
\addcontentsline{toc}{part}{\texorpdfstring{\protect\@mypartnumtocformat{\thepart}{#1}}{\partname~\thepart\ ---\ #1}}
\else%
\addcontentsline{toc}{part}{\texorpdfstring{\protect\@myparttocformat{#1}}{#1}}%
\fi%
\startcontents%
\markboth{}{}%
{\thispagestyle{empty}%
\begin{tikzpicture}[remember picture,overlay]%
\node at (current page.north west){\begin{tikzpicture}[remember picture,overlay]%	
\fill[ocre!20](0cm,0cm) rectangle (\paperwidth,-\paperheight);
\node[anchor=north] at (4cm,-3.25cm){\color{ocre!40}\fontsize{220}{100}\sffamily\bfseries\thepart}; 
\node[anchor=south east] at (\paperwidth-1cm,-\paperheight+1cm){\parbox[t][][t]{8.5cm}{
\printcontents{l}{0}{\setcounter{tocdepth}{1}}%
}};
\node[anchor=north east] at (\paperwidth-1.5cm,-3.25cm){\parbox[t][][t]{15cm}{\strut\raggedleft\color{white}\fontsize{30}{30}\sffamily\bfseries#2}};
\end{tikzpicture}};
\end{tikzpicture}}%
\@endpart}
\def\@spart#1{%
\startcontents%
\phantomsection
{\thispagestyle{empty}%
\begin{tikzpicture}[remember picture,overlay]%
\node at (current page.north west){\begin{tikzpicture}[remember picture,overlay]%	
\fill[ocre!20](0cm,0cm) rectangle (\paperwidth,-\paperheight);
\node[anchor=north east] at (\paperwidth-1.5cm,-3.25cm){\parbox[t][][t]{15cm}{\strut\raggedleft\color{white}\fontsize{30}{30}\sffamily\bfseries#1}};
\end{tikzpicture}};
\end{tikzpicture}}
\addcontentsline{toc}{part}{\texorpdfstring{%
\setlength\fboxsep{0pt}%
\noindent\protect\colorbox{ocre!40}{\strut\protect\parbox[c][.7cm]{\linewidth}{\Large\sffamily\protect\centering #1\quad\mbox{}}}}{#1}}%
\@endpart}
\def\@endpart{\vfil\newpage
\if@twoside
\if@openright
\null
\thispagestyle{empty}%
\newpage
\fi
\fi
\if@tempswa
\twocolumn
\fi}

%----------------------------------------------------------------------------------------
%	CHAPTER HEADINGS
%----------------------------------------------------------------------------------------

% A switch to conditionally include a picture, implemented by  Christian Hupfer
\newif\ifusechapterimage
\usechapterimagetrue
\newcommand{\thechapterimage}{}%
\newcommand{\chapterimage}[1]{\ifusechapterimage\renewcommand{\thechapterimage}{#1}\fi}%
\newcommand{\autodot}{.}
\def\@makechapterhead#1{%
{\parindent \z@ \raggedright \normalfont
\ifnum \c@secnumdepth >\m@ne
\if@mainmatter
\begin{tikzpicture}[remember picture,overlay]
\node at (current page.north west)
{\begin{tikzpicture}[remember picture,overlay]
\node[anchor=north west,inner sep=0pt] at (0,0) {\ifusechapterimage\includegraphics[width=\paperwidth]{\thechapterimage}\fi};
\draw[anchor=west] (\Gm@lmargin,-9cm) node [line width=2pt,rounded corners=15pt,draw=ocre,fill=white,fill opacity=0.5,inner sep=15pt]{\strut\makebox[22cm]{}};
\draw[anchor=west] (\Gm@lmargin+.3cm,-9cm) node {\huge\sffamily\bfseries\color{black}\thechapter\autodot~#1\strut};
\end{tikzpicture}};
\end{tikzpicture}
\else
\begin{tikzpicture}[remember picture,overlay]
\node at (current page.north west)
{\begin{tikzpicture}[remember picture,overlay]
\node[anchor=north west,inner sep=0pt] at (0,0) {\ifusechapterimage\includegraphics[width=\paperwidth]{\thechapterimage}\fi};
\draw[anchor=west] (\Gm@lmargin,-9cm) node [line width=2pt,rounded corners=15pt,draw=ocre,fill=white,fill opacity=0.5,inner sep=15pt]{\strut\makebox[22cm]{}};
\draw[anchor=west] (\Gm@lmargin+.3cm,-9cm) node {\huge\sffamily\bfseries\color{black}#1\strut};
\end{tikzpicture}};
\end{tikzpicture}
\fi\fi\par\vspace*{270\p@}}}

%-------------------------------------------

\def\@makeschapterhead#1{%
\begin{tikzpicture}[remember picture,overlay]
\node at (current page.north west)
{\begin{tikzpicture}[remember picture,overlay]
\node[anchor=north west,inner sep=0pt] at (0,0) {\ifusechapterimage\includegraphics[width=\paperwidth]{\thechapterimage}\fi};
\draw[anchor=west] (\Gm@lmargin,-9cm) node [line width=2pt,rounded corners=15pt,draw=ocre,fill=white,fill opacity=0.5,inner sep=15pt]{\strut\makebox[22cm]{}};
\draw[anchor=west] (\Gm@lmargin+.3cm,-9cm) node {\huge\sffamily\bfseries\color{black}#1\strut};
\end{tikzpicture}};
\end{tikzpicture}
\par\vspace*{270\p@}}
\makeatother

%----------------------------------------------------------------------------------------
%	HYPERLINKS IN THE DOCUMENTS
%----------------------------------------------------------------------------------------

\usepackage{hyperref}
\hypersetup{hidelinks,backref=true,pagebackref=true,hyperindex=true,colorlinks,linkcolor = ocre,breaklinks=true,urlcolor= ocre,bookmarks=true,bookmarksopen=false,pdftitle={Title},pdfauthor={Author}}
\usepackage{bookmark}
\bookmarksetup{
open,
numbered,
addtohook={%
\ifnum\bookmarkget{level}=0 % chapter
\bookmarksetup{bold}%
\fi
\ifnum\bookmarkget{level}=-1 % part
\bookmarksetup{color=ocre,bold}%
\fi
}
}
 % Insert the commands.tex file which contains the majority of the structure behind the template

\begin{document}

%----------------------------------------------------------------------------------------
%	TITLE PAGE
%----------------------------------------------------------------------------------------

\begingroup
\thispagestyle{empty}
\begin{tikzpicture}[remember picture,overlay]
\node[inner sep=0pt] (background) at (current page.center) {\includegraphics[width=\paperwidth]{background}};
\draw (current page.center) node [fill=ocre!30!white,fill opacity=0.6,text opacity=1,inner sep=1cm]{\Huge\centering\bfseries\sffamily\parbox[c][][t]{\paperwidth}{\centering Evanescent Tide\\[15pt] % Book title
{\Large  A solution manual by and for stupid student}\\[20pt] % Subtitle
{\huge Vector\_Cat }}}; % Author name
\end{tikzpicture}
\vfill
\endgroup

%----------------------------------------------------------------------------------------
%	COPYRIGHT PAGE
%----------------------------------------------------------------------------------------

\newpage
~\vfill
\thispagestyle{empty}

\noindent Copyright \copyright\ 2013 John Smith\\ % Copyright notice

\noindent \textsc{Published by Publisher}\\ % Publisher

\noindent \textsc{book-website.com}\\ % URL

\noindent Licensed under the Creative Commons Attribution-NonCommercial 3.0 Unported License (the ``License''). You may not use this file except in compliance with the License. You may obtain a copy of the License at \url{http://creativecommons.org/licenses/by-nc/3.0}. Unless required by applicable law or agreed to in writing, software distributed under the License is distributed on an \textsc{``as is'' basis, without warranties or conditions of any kind}, either express or implied. See the License for the specific language governing permissions and limitations under the License.\\ % License information

\noindent \textit{First printing, March 2013} % Printing/edition date

%----------------------------------------------------------------------------------------
%	TABLE OF CONTENTS
%----------------------------------------------------------------------------------------

%\usechapterimagefalse % If you don't want to include a chapter image, use this to toggle images off - it can be enabled later with \usechapterimagetrue

\chapterimage{chapter_head_1.pdf} % Table of contents heading image

\pagestyle{empty} % No headers

\tableofcontents % Print the table of contents itself

\cleardoublepage % Forces the first chapter to start on an odd page so it's on the right

\pagestyle{fancy} % Print headers again

%----------------------------------------------------------------------------------------
%	PART
%----------------------------------------------------------------------------------------

\part{Preliminaries}

%----------------------------------------------------------------------------------------
%	CHAPTER 1
%----------------------------------------------------------------------------------------
\chapterimage{orange_cat.jpg}
\chapter{Some category theory}

\section{Motivation}
\begin{exr}
A category in which each morphism is an isomorphism is called a groupoid. (This notion is not important in what we will discuss. The point of this exercise is to give you some practice with categories, by relating them to an object you know well.)
\begin{enumerate}[label=(\alph*)]
\item A perverse definition of a \textit{group} is: a groupoid with one object. Make sense of this.
\item Describe a groupoid that is not a group.
\end{enumerate}
\end{exr}
\begin{proof}
\begin{enumerate}[label=(\alph*)]
\item 
\end{enumerate}
\end{proof}
\section{Categories and functors}

\section{Universal properties determine an object up to unique isomorphism}

\section{Limis and colimits}

\section{Adjoints}

\section{An introduction to Abelian categroies}

\section{Spectral sequences}
%----------------------------------------------------------------------------------------
%	CHAPTER 2
%----------------------------------------------------------------------------------------

\chapterimage{sheaves.jpg}
\chapter{Sheaves}
\section{Motivating example}
\begin{exr}
Prove that $\scm_p:=$\{germs that vanishing at $p$\} is the only maximal ideal in $\calo_p$(Hint: prove that $\calo_p-\scm_p\subseteq \calo_p^{\times}$)
\end{exr}
\begin{proof}In general setting, consider a ring $A$.
The set of non-unit can be written as $\cup\scm_j$, where each $\scm_j$ is a maximal ideal in $A$. $A-\scm\subseteq A^\times\Lrta A-\scm\subseteq A-\cup_j\scm_j\Lrta \scm\supseteq\cup_j\scm_j$, then there is only one maximal ideal. It suffices to prove each element in $\calo_p-\scm_p$ is invertible. 

Use $(f,U)_p$ to denote the equivalent class of $(f,U)$.
$(f,U)_p\in \calo_p-\scm_p\Lrta f(p)\neq 0$, because $f$ is continuous on $U$, there exists an open set $V\subseteq U$, s.t. $f(q)\neq0,\forall q\in V\Lrta (1/f,V)$ is an inverse.
\end{proof}

\begin{exr}
Prove that $\frac{\scm_p}{\scm_p^2}\cong T^*_pX$, $T^*_p X$ is the cotangent bundle of differentiable manifold $X$ at $p$. 
\end{exr}
\begin{proof}Construct explicitly the isomorphism:
$$
(f,V)_p+\scm_p^2\longmapsto df_p=\sum_i\left.\frac{\partial f\circ \varphi^{-1}}{\partial x^i}\right|_p (dx^i)_p,
$$ 
where for a given tangent vector $\delta_p\in T_p X\cong Der_p(C^{1}(X))$ at $p$, $df_p(\delta_p)=\delta_p f\in \reals$ is an $\reals$-linear map. It is well-defined because $h\in \scm_p^2\Lrta d h_p=0$ and $(g,U)\sim (f,V)$, then there exists an open $p\in W\subseteq U\cap V$, s.t. $f|_W-g|_W=0$, $(df_p-dg_p)(\delta)=\delta_p(f)-\delta_p(g)$. Consider a bump function $\phi(x)=1,\forall x\in W'\ni p$ and $\phi(x)=0,\forall x\in U\cap V-W$, then $0=\delta_p(f|_{W'}-g|_{W'})=\delta_p((f-g)\phi|_{W'})=(f(p)-g(p))\delta_p(\phi)+\delta_p(f-g)\phi(p)\Lrta \delta_p(f-g)=0$.

It is surjective because the basis $(dx^j)_p$ can be generated by this map.

Construct the inverse map. Consider a chart of $p$, $(V,\varphi)$ s.t. $\varphi(p)=0\in \reals^n$, and define $x_i$ to be the $i$-th entry of $\phi$
$$
\sum_i a_i (dx^i)_p\longmapsto \left(\sum_i a_i x_i, V \right)+\scm^2_p.
$$
To check this is indeed a inverse, we only need to check (The other direction is tricial)
$$
(f,V)_p+\scm_p^2=\left(\sum_i\left.\frac{\partial f\circ \varphi^{-1}}{\partial x^i}\right|_0 x^i, V'\right)+\scm_p^2.
$$
But, assuming $\varphi(V)$ to be a star-shaped domain in $\reals^n$: 
$$
f\circ\varphi^{-1}(x)= \left(\int_0^1d t\frac{\partial f\circ\varphi^{-1}(t x)}{\partial t}\right)=\sum_ix^i\underbrace{\int_0^1d t\frac{\partial f\circ\varphi^{-1}}{\partial x^i}(t x)}_{:=h_i(x)}.
$$
Though $h_i(x)$ is at most continuous, $\sum_i x^i h_i(x)$ is differentiable and $$
h_i(0)=\left.\frac{\partial f\circ \varphi^{-1}}{\partial x^i}\right|_0
\Lrta h_i(x)-\left.\frac{\partial f\circ \varphi^{-1}}{\partial x^i}\right|_0\in \scm_p
$$
then
$$
\sum_i x^i \left(h_i(x)-\left.\frac{\partial f\circ \varphi^{-1}}{\partial x^i}\right|_0\right)\in \scm_p^2
$$
\end{proof}
\section{Definition of presheaf and sheaf}
\begin{exr}
Verify that the data of a presheaf is precisely the data of a contravariant functor from the category of open sets of $X$ to the category of sets.
\end{exr}
\begin{proof}We use a diagram to replace all the words:
\begin{center}
\begin{tikzcd}
U\ar[loop left] \arrow[d, "\calf"] \arrow[r, hook] & V \arrow[d, "\calf"] \arrow[r, hook] & W \arrow[d, "\calf"] \\
\calf(U)\ar[loop left, "res_{U,U}"] & \calf(V) \arrow[l, "{res_{V,U}}"'] & \calf(W) \arrow[l, "{res_{W,V}}"'] \arrow[ll, "{res_{W,V}}", bend left]
\end{tikzcd}
\end{center}
\end{proof}

\begin{exr}
Show that the following are presheaves on $\cplx$ (with the classical topology), but not sheaves: (a) bounded functions, (b) holomorphic functions admitting a holomorphic square root.
\end{exr}
\begin{proof}
\begin{enumerate}[label=(\alph*)]
\item Bounded functions $B$ is a presheaf, with the restriction map just the ordinary restriction of functions. But $B$ is not a sheaf because there could be an open cover $\{U_i\}_{i\in I}$(necessarily infinite), for example $U_n:=\{z\in\cplx: n-1< \text{Re}z<n+1\}$, and we can choose $f_n(z):=z$. each is bounded function on the open set, but there is no global bounded function that restricts to them.
\item Holomorphic functions admitting a square roots, in other words, they are of the form $h^2(z)$ where $h(z)$ is a holomorphic function. It is a presheaf with ordinary restrictions and it is not a sheaf because it violate the gluability axiom. Consider $U$ to  be a strip $\{2<|z|<3 \}$ and $U$ is covered by open disks $D_n:=\{|z-5/2e^{i n2\pi/10}|<1\}$, on each open disk, there is a holomorphic square root corresponding to $\sqrt z$. But there is no $\sqrt z$ defined on the strip $U$.
\end{enumerate}
\end{proof}
\begin{exr}
The identity and gluability axioms may be interpreted as saying that $\calf(\cup_{i\in I}U_i)$ is a certain limit. What is that limit?
\end{exr}
\begin{proof}
The index category $\cali$ is the ``subsets index by $I$'' with inclusion being the morphism it is a partially ordered set. Or we can simply denote it by $I$ with $i\leq j$ iff $U_i\inj U_j$. Then $\calf(\cup_{i\in I}U_i)$ can be interpreted as a limit with morphism $\calf(\cup_{i\in I}U_i)\lrta \calf(U_k)$ being the restriction $\res_{\cup_i U_i, U_k}$
$$
\calf(\cup_{i\in I}U_i)=\invlim_{i\in I}\calf(U_i)
$$
\end{proof}
\begin{exr}\ 
\begin{enumerate}[label=(\alph*)]
\item
 Verify that the the following examples are indeed sheaves ( differentiable functions, or continuous functions, or smooth functions, or functions on a manifold or $\reals^n$).
 \item  Show that real-valued continuous functions on (open sets of) a topological space X form a sheaf.
\end{enumerate}
\end{exr}
\begin{proof}
\begin{enumerate}[label=(\alph*)]
\item All of them are easily seen to be presheaves, we only need to verify they satisfies the identity axiom and gluability axiom of sheaf.

Identity axiom holds for all of them, $\res_{U,U_i}f_1=\res_{U,U_i}f_2,\forall U_i$ means $f_1(p)=f_2(p)\forall p\in U_i,\forall U_i$. Hence $f_1$ and $f_2$ coincide on each point in $U$, they are equal. (This is not true for functions on schemes)

Manifolds are Hausdorff and paracompact, hence endowed with partition of unity. For each over cover $\{U_i\}_{i\in I}$, there is a partition of unity $\{\rho_i\}_{i\in I}$ subordinate to it (each $\rho_i$ is non-negative and continuous, and we can choose it to be smooth if the manifold is smooth). For $\{f_i\in C^{k}(M)\}$ compatible on the intersections, we can define $f=\sum_i \rho_i(x)f_i(x)\in C^{k}(M)$. Such that it restrict to each $f_i$. $C^k(M)$ can be continuous, smooth, or even just plain functions.
(In particular the partition of unity property is included in the notion of fine sheaves. In fact, we don't need partition of unity to prove $(a)$ here.)

\item Continuous functions on a topological space $C(X)$. It is a presheaf. We don't have to worry about the fact that function may be not determined by it's value at points, because here the continuous function means continuous map from $X$ to $\reals$ and is indeed determined valuewise. 

For identity axiom, $f_1(p)=f_2(p)$ at every point in $U$. Then they are equal. 

For gluability axiom, we define $f(p)=f_i(p)$, where $f_i\in C(U_i)$ and $p\in U_i$. Then it remains to prove $f$ is a continuous function on $U$. Recall the definition of continuous function in terms of neighborhoods. Consider $f(p)\in \reals$, for any neighborhood $V$ of $f(p)=f_i(p)$. Because $f_i$ is continuous, there is a neighborhood $p\in W\subset U_i$ such that $f_i(W)\subset V$. $W$ is also an open subset in $U$. $\res_{U,W}f=\res_{U_s,W}\circ \res_{U,U_i}f=\res_{U_i,W}f_i$, hence $f(W)\subset V$. This means $f$ is continuous at every point $p$.
\end{enumerate}
\end{proof}

\begin{exr}
Now let $\calf(U)$ be the maps to $S$ that are locally constant, i.e., for any point $p$ in U, there is an open neighborhood of $p$ where the function is constant. Show that this is a sheaf. (A better description is this: endow $S$ with the discrete topology, and let $\calf (U)$ be the continuous maps $U \lrta S$.) This is called the constant sheaf (associated to $S$); do not confuse it with the constant presheaf. We denote this sheaf $\underline{S}$.
\end{exr}
\begin{proof}
The restriction maps are just the ordinary restrictions. The constant sheaf is indeed a sheaf because it is continuous functions from $U$ to $S$ where $S$ is with discrete topology. The proof is identical. Ae will be explained in the next exercise, in fact by \href{https://en.wikipedia.org/wiki/Pasting_lemma}{pasting lemma}, every presheaf of continuous maps $C(X,Y)$ is in  fact a sheaf. 
\end{proof}

\begin{exr}\label{chap2exr:sheaf_of_continuous_maps}
Suppose $Y$ is a topological space. Show that ``continuous maps to $Y$'' form a sheaf of sets on $X$. More precisely, to each open set $U$ of $X$, we associate the set of continuous maps of $U$ to $Y$. Show that this forms a sheaf.
\end{exr}
\begin{proof}
It is obviously a presheaf and the identity holds true.

As for the gluability. Choose $V$ an open set in $Y$, and the preimage of $f\in C(X,Y)$ defined point wise by $f_i\in C(U_i,Y)$, where $U_i$ is in the subset topology.
$$
f^{-1}(V)=\cup_i f_i^{-1}(V).
$$
Each $f_i^{-1}(V)$ is open set in $U_i$ hence is an open set in $X$.
\end{proof}

\begin{exr}\ 
\begin{enumerate}[label=(\alph*)]
\item 
(sheaf of sections of a map) Suppose we are given a continuous map $\mu : Y \lrta X$. Show that sections of $\mu$ form a sheaf. More precisely, to each open set $U$ of $X$, associate the set of continuous maps $s: U \lrta  Y$ such that $\mu\circ s   = id|_U$. Show that this forms a sheaf. 
\item Suppose that $Y$ is a topological group. Show that continuous maps to $Y$ form a sheaf of groups.
\end{enumerate}
\end{exr}
\begin{proof}\ 
\begin{enumerate}[label=(\alph*)]
\item It is a presheaf, and with obvious restriction maps. $s_i:=\res_{U,U_i}s$ and $\mu\circ s_i=id_{U_i}$. 

Given an open cover $\{U_i\}_{i\in I}$ of $U$.

Identity Axiom:  if $s_1,s_2$ agree when restricted to any $U_i$, then they agree point wisely.

Gluability Axiom: if $\{s_i\}_{i\in I}$  are compatible on each intersection. We define $s$ point wisely, we have to check $\mu\circ s=id|_U$. And this can be checked also value-wisely.
\item A topological group is a topological space together with  compatible group structures.

It is presheaf of groups, $s_1,s_2:U\lrta Y$ the multiplication is defined as $s_1 \cdot s_2(p)=s_1(p)\times s_2(p)$. The identity and inverse is also defined point-wisely. 

Given an open cover $\{U_i\}_{i\in I}$.

Identity axiom: if $s_1,s_2$ agree upon restricted to each $U_i$, then they agrees pint-wisely.

Gluability axiom: if $\{s_i\}_{i\in I}$  are compatible on each intersection. We define $s$ point wisely, we have to check $s$ is indeed continuous, and this has been checked in Exercise~\ref{chap2exr:sheaf_of_continuous_maps}.

 In fact, the group structure has nothing to do with whether it is a sheaf of sets.
 What we have to do is just to verify that the sheaf of sets possesses the addition structure of sheaf of groups.
\end{enumerate}
\end{proof}

\begin{exr}
 THE PUSHFORWARD SHEAF OR DIRECT IMAGE SHEAF. Suppose $\pi: X \lrta Y$ is a continuous map, and $\calf$ is a presheaf on $X$. Then define $\pi_*\calf$ by $\pi_*\calf(V)=\calf(\pi^{-1}(V))$, where $V $ is an open subset of $Y$. Show that $\pi_*\calf$ is a presheaf on $Y$, and is a sheaf if $\calf{F}$ is. This is called the \textbf{pushforward} or \textbf{direct image} of $\calf$. More precisely, $\pi_*\calf$ is called the pushforward of $\calf$ by $\pi$.
\end{exr}
\begin{proof}
It is a presheaf:
The new restriction map is induced as 
$$
\res'_{V,W}=\res_{\pi^{-1}V,\pi^{-1}W}.
$$
The functorial property holds
$$
\res'_{V,V}=\res_{\pi^{-1}V,\pi^{-1}V}=id_{\calf(\pi^{-1}(V))}:=id_{\pi_*\calf(V)}
$$
and
$$
\res'_{U,V}\circ \res'_{V,W}=\res_{\pi^{-1}U,\pi^{-1}V}\circ \res_{\pi^{-1}V,\pi^{-1}W}=\res_{\pi^{-1}U,\pi^{-1}W}=\res'_{V,W}.
$$

Now assume $\calf$ is a sheaf.
The open cover $\{V_i\}_{i\in I}$ of $V$ pullback to an open cover $\pi^{-1}(V_i)$ of $\pi^{-1}(V)$.

Identity axiom: $f_1,f_2$ in $\pi_*\calf(V)$ agrees when restricted to any $V_i$, then $f_1,f_2$ as element in $\calf(\pi^{-1}{V})$ agrees when restricted to each $\pi^{-1}(V_i)$. Because $\calf$ is a sheaf, $f_1,f_2$ has to agree in $\calf(\pi^{-1}(V))$, which is equivalent to $f_1,f_2$ agrees in $\pi_*\calf(V)$.

Gluability axiom: if $\{f_i\}_{i\in I}$  are compatible on each intersection. They are also compatible as sections in $\calf(\pi^{-1}V_i)$, because $\calf$ is a sheaf, there is an element $\tilde{f}\in \calf(\pi^{-1}V)$ such that it restricts to each $f_i$. $\tilde{f}$ can be interpreted as an element $f\in \pi_*\calf(V)$.
\end{proof}

\begin{exr}\label{chap2exr:Pushforward_induces_morphism_stalks}
Suppose $\pi: X \lrta Y $ is a continuous map, and $\calf$ is a sheaf of sets (or rings or $A$-modules) on $X$. If $ \pi(p) = q$, describe the natural morphism of stalks $(\pi_*\calf)_q \lrta \calf_p$.
\end{exr}
\begin{proof}
First we define the morphism of stalk using the universal property of colimits.
$$
\calf_p=\dirlim\calf(U),
$$
 where the index category is the ``open sets that containing $p$'' and each $\pi_*\calf(V)\cong \calf(\pi^{-1}(V))$
\begin{center}
\begin{tikzcd}
\pi_*\calf_q \arrow[rr, "\exists!", dashed] &  & \calf_p \\
 & \calf(\pi^{-1}V_i) \arrow[ru, "g_i"] \arrow[r] & \calf(\pi^{-1}V_j) \arrow[u, "g_j"] \\
\pi_*\calf(V_i) \arrow[r] \arrow[uu, "f_i"] \arrow[ru, "="] & \pi_*\calf(V_j) \arrow[luu, "f_j"'] \arrow[ru, "="] &
\end{tikzcd}
\end{center}
Because the index category of either colimits is filtered set. We can equivalently define the at the level of disjoint union modulo equivalence. 
$$
\calf_p:=\left\{(f_i; U_i)\in\coprod_{i\in\cali} \calf(U_i):\right\}\left/
\left(
\begin{matrix}
& (f_i;U_i)\sim (f_j;U_j)\text{iff there are $\res_{U_i\rta U_k}:\calf(U_i)\rta \calf(U_k)$}\\
&\text{$\res_{U_j,U_k}:\calf(U_j)\rta \calf(U_k)$ s.t., $f_i|_{U_k}=f_j|_{U_k}\in \calf(U_k)$}
\end{matrix}
\right)
\right.
$$
$$
[(f_i;V_i)]\mapsto [(f_i;\pi^{-1}(V_i))]
$$
It is well defined map because
$$
\res'_{V,W}=\res_{\pi^{-1}V,\pi^{-1}W}.
$$
\end{proof}

\begin{definition}
Ringed spaces, and $\calo_X$-modules. Suppose $\calo_X$ is a sheaf of rings on a topological space $X$. Then $(X,\calo_X)$ is called a \textbf{ringed space}. The sheaf of rings is often denoted by $\calo_X$. This sheaf is called the structure sheaf of the ringed space. Sections of the structure sheaf $\calo_X$ over an open subset $U$ are called \textbf{functions} on $U$. Functions in this sense is not defined point wise, and is not determined by their value at points.
\end{definition}

\begin{exr}
If $(X,\calo_X)$ is a ringed space, and $\calf$ is an $\calo_X$-module, describe how for each $p\in X$, $\calf_p$ is an $\calo_{X,p}$-module.
\end{exr}
\begin{proof} The crucial print is that colimit of product is product of colimit.
\begin{center}
\begin{tikzcd}
 & \calo_X(V)\times \calf(V) \arrow[rr, "\text{action}"] \arrow[ld, "{\res_{V,U}\times\res_{V,U}}"'] \arrow[rdd] &  & \calf(V) \arrow[ld, "{\res_{V,U}}"] \arrow[rdd] &  \\
\calo_X(U)\times \calf(U) \arrow[rr, "\text{action}"] \arrow[rrd] &  & \calf(U) \arrow[rrd] &  &  \\
 &  & \calo_{X,p}\times\calf_p \arrow[rr, "\text{action}"] &  & \calf_p
\end{tikzcd}
\end{center}
At the explicit level,
$$
\calo_{X,p}\times \calf_p\ni([(a_V;V)],[(m_U;U)])=[((a_{W}, m_W);W)]\in (\calo_X\times \calf)_p,
$$
where $W\subset V\cap U$.
and 
$$
\calo_{X,p}\times \calf_p\ni([(a_V;V)],[(m_U;U)])\overset{\text{action}}{\longmapsto}[(a_w\cdot m_W;W)]\in\calf_p.
$$
And it is routine to verify that map defined above indeed gives $\calf_p$ the structure of a $\calo_p$ module.


\end{proof}
\section{Morphism of presheaves and sheaves}
\begin{exr}\label{chap2exr:morphism_of_sheaves_induce_on_stalks}
MORPHISMS OF (PRE)SHEAVES INDUCE MORPHISMS OF STALKS. If $\phi: \calf\lrta \calg $ is a morphism of presheaves on $X$, and $p \in X$, describe an induced morphism of stalks $\phi_p : \calf_p \lrta \calg_p$. Translation: taking the stalk at $p$ induces a functor $Sets_X \lrta  Sets$.
\end{exr}
\begin{proof}
Using the language of colimits,
\begin{center}
\begin{tikzcd}
\calf_q \arrow[rr, "\exists!\phi_p", dashed] &  & \calg_p \\
 & \calg(V_i) \arrow[ru, "g_i"] \arrow[r,"\res_{V_i,V_j}"] & \calg(V_j) \arrow[u, "g_j"] \\
\calf(V_i) \arrow[r,"\res_{V_i,V_j}"] \arrow[uu, "f_i"] \arrow[ru, "\phi(V_i)"] & \calf(V_j) \arrow[luu, "f_j"'] \arrow[ru, "\phi(V_j)"] &
\end{tikzcd}
\end{center}

Using the explicit construction, we can check the map 
$$
\phi_p:[(s;V)]\mapsto [(\phi(V)s;V)]
$$
is well-defined. Assume $(s;V)\sim (s';V')$, there exists $(t;W)$ so that $W\subset V\cap V'\ni p$ and $t=s|_W=s'|_W$.
 We can check $(\res_{V,W}\phi(V)s;W)=(\res_{V'W}\phi(V')s';W)=(\phi(W)t;W)$. Hence $[(\phi(V)s;V)]=[\phi(V')s';V']$.
\end{proof}

\begin{exr}
Suppose $\pi:X\lrta Y$ is a continuous map of topological spaces (i.e., a morphism in the category of topological spaces). Show that pushforward gives a functor $\pi_* : Sets_X \lrta Sets_Y$ . Here $Sets $ can be replaced by other categories. 
\end{exr}
\begin{proof}
We want ot verufy the diagram 
\begin{center}
\begin{tikzcd}
\calf \arrow[r, "\phi",blue] \arrow[d, "\pi_*"',red] & \calg \arrow[d, "\pi_*",red] \\
\pi_*\calf \arrow[r, "\pi_*(\phi)",blue] & \pi_*\calg,
\end{tikzcd}
\end{center}
where $\phi:\calf\lrta \calg$ is a morphism of sheaves and $\pi_*\phi(V)=\phi(\pi^{-1}V)$. 
In detail, it means each square is commutative in the following diagram
\begin{center}
\begin{tikzcd}
 &  & \calf(\pi^{-1}U) \arrow[lld,"\pi_*",red] \arrow[dd, bend left,"\phi", blue] &  \\
\pi_*\calf (U) \arrow[dd, bend right,"\pi_*(\phi)"', blue] &  & \pi^{-1}U \arrow[u] \arrow[rd, hook] \arrow[d] & \calf(\pi^{-1}V) \arrow[lu] \arrow[lld,"\pi_*"', red] \arrow[dd, bend left,"\phi", blue] \\
U \arrow[u] \arrow[rd, hook] \arrow[d] \arrow[rru] & \pi_*\calf(V) \arrow[lu] \arrow[dd, bend right,"\pi_*(\phi)"', blue] & \calg(\pi^{-1}U) \arrow[lld,"\pi_*", red] & \pi^{-1}V \arrow[u] \arrow[d] \\
\pi_*\calg(U) & V \arrow[u] \arrow[d] \arrow[rru] &  & \calg(\pi^{-1}V) \arrow[lu] \arrow[lld,"\pi_*", red] \\
 & \pi_*\calg(V) \arrow[lu] &  & 
\end{tikzcd}
\end{center}
\end{proof}

\begin{exr}
Suppose $\calf$ and $\calg$ are two sheaves of sets on $X$. (In fact, it will suffice that $\calf$ is a presheaf.) Let $\mathpzc{Hom} (\calf , \calf)$ be the collection of data
$$
\mathpzc{Hom}(\calf,\calg)(U):=\mor(\calf|_U,\calg|_U).
$$
Show that this is a sheaf of sets on $X$. (To avoid a common confusion: the right side does not say $\mor(\calf(U),\calg(U))$.) This sheaf is called ``sheaf Hom''. (Strictly speaking, we should reserve Hom for when we are in an additive category, so this should possibly be called ``sheaf Mor''. But the terminology ``sheaf Hom'' is too established to uproot.) It will be clear from your construction that, like Hom, $\mathpzc{Hom}$ is a contravariant functor in its first argument and a covariant functor in its second argument.
\end{exr}
\begin{proof}
Given $\calf$ a presheaf and $\calg$ a sheaf. We ahve to check that $\phom(\calf,\calg)$ is a sheaf of sets on $X$.

First, we check that it is a presheaf:
\begin{center}
\begin{tikzcd}
U \arrow[d] & V \arrow[d] \arrow[l, hook] \\
\phom(\calf,\calg)(U) \arrow[r,"\res_{U,V}"] & \phom(\calf,\calg)(V)
\end{tikzcd}
\end{center}
The $\res_{U,V}$ is induced by 
$$
\res_{U,V}:\phom(\calf,\calg)(U)=\mor(\calf|_U,\calg|_U)\lrta\mor(\calf|_V,\calg|_V)= \phom(\calf,\calg)
$$
$$
\mor(\calf|_U,\calg|_U)\ni\phi\longmapsto \phi|_V,
$$
where $\phi|_V:\calf|_V\lrta \calg|_V$, $\phi|_V(V_i)=\phi(V_i)$. In the later term $V_i$ is regarded as an open subset in $U$. It is easy to see that $\res_{U,U}$ is identity. And it suffices to check that $\res_{V,W}\circ \res_{U,V}=\res_{U,W}$. For an element $\phi$ in in $\phom(\calf,\calg)$, and for $W_i$ open subset in $W$:
$$
\left(\phi|_V\right)|_W(W_i)=\phi|_V(W_i)=\phi(W_i)=\phi|_W(W_i).
$$
This equality is true for arbitrary open sets in $W$ and notice that morphism $\varphi$ of (pre)sheaves is defined by the behavior of $\varphi(W_i)$. We can conclude that it satisfies the functorial property of a presheaf. 

Then given an open cover $\{U_i\}_{i\in I}$ of $U$, we check that sheaf Hom satisfies the two axioms of sheaf.

Identity axiom: Let $\phi_1,\phi_2$ be two morphisms in $\phom(\calf,\calg)(U)$. They agree when restricted to each $U_i$.

\underline{Know}:
$$
\phi_1|_{U_i}=\phi_2|_{U_i},\forall i\in I
$$
which means $\phi_1|_{U_i}(U_{i,j})=\phi_2|_{U_i}(U_{i,j})$ for each open $U_{i,j}\subset U_i$

\underline{Want}:$\phi_1(V)=\phi_2(V)$ for all $V$ open sets in $U$.

In particular, each open set $V$ is covered by $\{U_i\cap V\}_{i\in I}$, then we know 
$$
\phi_1|_{U_i}(U_i\cap V)=\phi_1(U_i\cap V)=\phi_2(U_i\cap V)=\phi_2|_{U_i}(U_i\cap V).
$$
Then we have reduced the problem to verifying that, for $\{V_i\}_{i\in I}$ an open cover of $V$, $\phi_1(V)=\phi_2(V)$ if $\phi_1(V_i)$ agrees with $\phi_2(V_i)$ as morphism of sets for all $V_i$.

Given $f\in\calf(V)$.

\underline{Want}: $g_1:=\phi_1(V)(f)=g_2:=\phi_2(V)(f)$. 

 They restrict to $f_i\in \calf(V_i)$ and $g_{1,i},g_{2,i}\in\calg(V_i)$. Set $g_{1,i},g_{2,i}\in \calg(V_i)$ and
$\phi_1(V_i):f_i\mapsto g_{1,i}$ and $\phi_2(V_i):f_i\mapsto g_{2,i}$ by the definition of morphism of (pre)sheaves. $g_{1,i}=g_{2,i}$ for all $i$ because $\phi_1(V_i)$ and $\phi_2(V_i)$ agree as morphism of sets.

Now, \textbf{because $\calg$ is a sheaf}, $g_1=g_2$. \textbf{It suffices that $\calf$ being only a presheaf rather than sheaf}.

Gluability axiom: Let $\{\phi_i\}_{i\in I}, \phi_i\in \phom(\calf,\calg)(U_i)$ be compatible when they are restricted to intersections. 

\underline{Want}: There exists a $\phi\in \phom(\calf,\calg)(U)$ such that $\phi|_{U_i}=\phi_i$, which means there exists a morphism of sets $\phi(V)$ for each open set $V\subset U$ such that $\phi(U_{i,j})=\phi|_{U_i}(U_{i,j})=\phi_i(U_{i,j})$ for all $U_{i,j}$ open subsets in $U_i$ and the defined morphism of sets $\phi(V)$ indeed determines a morphism of (pre)sheaves.

We construct each $\phi(V)$ explicitly. $V$ is covered by $\{V_i:=U_i\cap V\}_{i\in I}$ and $V_i\cap V_j=U_i\cap U_j\cap V$.

\underline{Know}: $\phi_i|_{U_i\cap U_j}=\phi_j|_{U_i\cap U_j}$. In particular 
$$
\phi_i|_{U_i\cap U_j}(V_i\cap V_j)=\phi_i(V_i\cap V_j)=\phi_j(V_i\cap V_j)=\phi_j|_{U_i\cap U_j}(V_i\cap V_j).
$$
Given $f\in\calf(V)$. $f$ restricts to $f_i\in \calf(V_i)$ and set
$$
g_i:=\phi_i(V_i)(f_i)\in\calg(V_i)
$$
$$
g_i|_{V_i\cap V_j}=\phi_i(V_i)(f_i)|_{V_i\cap V_j}=\phi_i(V_i\cap V_j)(f_i|_{V_i\cap V_j}),
$$
where the second equality holds because each $\phi_i$ is a morphism of (pre)sheaves. Then because $\phi_i|_{U_i\cap U_j}(V_i\cap V_j)=\phi_j|_{U_i\cap U_j}(V_i\cap V_j)
$ and $f_i|_{V_i\cap V_j}=f_j|_{V_i\cap V_j}=f|_{V_i\cap V_j}$ (because $\calf$ is a presheaf)
$$
g_i|_{V_i\cap V_j}=g_j|_{V_i\cap V_j}
$$
$g_i$ are compatible when restricted to intersections. Then, again,  \textbf{because $\calg$ is a sheaf, by gluability}, there is a unique element $g$ to glue to in $\calg(V)$. 

We define $\phi(V):f\mapsto g$, $g$ is constructed as above. $\phi(V)$ is a well-defined morphism in category $\mathsf{Set}$. \footnote{This value-wisely defined map is well defined morphism is any category of structured sets. For general category we should define $\phi(V)$ as the limit of $\phi(U_i)$s in the arrow category.}
Then there two things left: 1. verify that $\phi$ is indeed a morphism of (pre)sheaves. 2. verify that $\phi$ to restricts to $\phi_i$.
\begin{center}
\begin{tikzcd}
\calf(V) \arrow[r, "\phi(V)"] \arrow[d, "{\res_{V,W}}"'] & \calg(V) \arrow[d, "{\res_{V,W}}"] \\
\calf(W) \arrow[r, "\phi(W)"] & \calg(W)
\end{tikzcd}
\end{center}
Given an element $f\in\calf(V)$, $g=\phi(V)(f)\in \calg(V)$ is the element to which $\phi_i(V_i)(f|_{V_i})$ glue up. $f$ restricts to $f|_W$, and $g'=\phi(W)(f|_W)\in\calg(W)$ is defined to be the element $g\in \calg(W)$ to which $\phi(W_i)(f|_{W_i})$ glue up, where $W_i=U_i\cap W$. $W_i$ is an open subset of $V_i$ and $\{W_i\}_{i\in I}$ is an open cover of $W$. \underline{Want}: $g'=g|_W$, but we have
$$
g'|_{W_i}=\phi(W_i)(f|_{W_i})=\phi(V_i)(f|_{V_i})|_{W_i}=(g|_{V_i})|_{W_i}
=(g|_W)|_{W_i}.
$$
Then \textbf{because $\calg$ is a sheaf, by identity axiom}, $g'=g|_W$. The above diagram commutes and $\phi$ is a morphism of (pre)sheaves.

On the other hand, by definition of $\phi(V)$,
$$
\phi|_{U_i}(V_i)(f_i)=\phi_i(V_i)(f_i),
$$
but our $V$ is chosen arbitrarily, hence $V_i=V\cap U$ is arbitrary open subsets in $U_i$, which means $\phi|_{U_i}=\phi_i$.

Now $\phom(\_,\calg)$ is a contravariant functor from the the category of sheaves of something to the category of category of sheaves of sets and $\phom(\calf,\_)$ is a covariant functor...
\end{proof}
\begin{remark}
Warning: $\mathpzc{Hom}$ does not commute with taking stalks. More precisely: it is not true that $\mathpzc{Hom} (\calf , \calg)_p$ is isomorphic to $\mathpzc{Hom}(\calf_p , \calg_p )$. 
\begin{center}
\begin{tikzcd}
 & \mor(\calf_p,\calg_p) &  \\
\phom(\calf,\calg)(U) \arrow[d] \arrow[ru] \arrow[rrd] &  &  \\
\phom(\calf,\calg)(V) \arrow[ruu] \arrow[rr] &  & \phom(\calf,\calg)_p \arrow[luu, "\exists!"', dashed]
\end{tikzcd}
\end{center}
There is always a natural morphism from $\mathpzc{Hom} (\calf , \calg)_p$ to $\mathpzc{Hom}(\calf_p , \calg_p )$, because by~\ref{chap2exr:morphism_of_sheaves_induce_on_stalks}, each morphism of sheaves induces a morphism of stalks. The map $\phi\mapsto \phi_p$ is an arrow from $\mathpzc{Hom} (\calf , \calg)(U)$ to $\mathpzc{Hom}(\calf_p , \calg_p )$ and compatible with restriction of $\mathpzc{Hom} (\calf , \calg)$. Then by the definition of colimit, there exists a unique morphism from $\mathpzc{Hom} (\calf , \calg)_p$ to $\mathpzc{Hom}(\calf_p , \calg_p )$. But in general, the morphism is neither injective nor surjective.

Non-examples: Let $H$ to be any nontrivial Abelian group, $X$ be any topological space. Take $\calg$
 to be the constant sheaf $\underline{H}$ and $\calf$ to be the skyscraper sheaf at a point a point $p$ with stalk $H$ at $p$.

\textbf{Then $\phom(\calf,\calg)_p\lrta \hom(\calf_p,\calg_p)$, is not surjective.}
$\phom(\calf,\calg)_p=\{0\}$ and we have $\hom(\calf_p,\calg_p)$ $=\hom(H,H)\neq \{0\}$.

Because $\phom(\calf,\calg)(V)=\{0\}, p\notin V$. Given $f\in\calf(V)$, and $\phi\in\phom(\calf,\calg)$. $\phi(V)(f)\in\calg(V)$ is the section that $\phi(V_i)(f|_{V_i})$ glues to. $\phi(V_i)(f|_{V_i})$ is trivial when $V\not\ni p$. $\calg$ is locally constant, and hence they must glue up to trivial element in $\calg(V)$. 

On the other hand, let $Y=X\backslash p$, and let $Ab_X\ni\calf=\calg$ be the extension by zero of the constant sheaf $\underline{H}$ on $Y$ (i.e. it's the constant sheaf $\underline{H}$ on $V$ and $\calf(U)=0$ if $U\notsubset V$). Then $\calf_p=\calg_p=0$, but $\phom(\calf,\calg)$ contains a natural copy of $\hom(\underline{H},\underline{H})\neq0$ for any $U$, so $\hom(\calf,\calg)_p$ contains a natural copy of $\hom(\underline{H},\underline{H})$, so it's not zero. Hence \textbf{$\phom(\calf,\calg)_p\lrta \hom(\calf_p,\calg_p)$, is not surjective.}
 \end{remark}

 \begin{exr}\label{chap2exr:OX-moduel-hom_from_special_Ox-modules}\ 
 \begin{enumerate}[label=(\alph*)]
 	\item
If $\calf$ is a sheaf of sets on $X$, then show that $\phom(\underline{\{p\}},\calf) \cong \calf$, where \underline{\{p\}} is the constant sheaf associated to the one element set $\{p\}$.
\item 
 If $\calf$ is a sheaf of abelian groups on $X$, then show that $\phom_{Ab_X}(\underline{\intg},\calf)\cong \calf$ (an isomorphism of sheaves of abelian groups).
\item  If $\calf$ is an $\calo_X$-module, then show that $\phom_{Mod_{\calo_X}}(\calo_X,\calf) \cong\calf$ (an isomorphism of $\calo_X$-modules).
\end{enumerate}
 \end{exr}
\begin{proof}
\begin{enumerate}[label=(\alph*)]
\item The one element constant sheaf is just the corresponding constant presheaf, $\underline{\{p\}}(V)=\{p\}$ for all open sets. For each morphism of sheaves $\phi$ in $\phom(\underline{\{p\}},\calf)$, the corresponding morphism of sets $\phi(V):\{p\}\lrta\calf(V)$ defines a unique element $\phi(V)(p)\in\calf(V)$. We can define the morphism of sheaves $ev$
$$
ev(V):\phom(\underline{\{p\}},\calf)(V)\lrta\calf(V) 
$$
$$
\phi(V)\longmapsto \phi(V)(p)
$$
and the following diagram commutes
\begin{center}
\begin{tikzcd}
\phom(\underline{\{p\}},\calf)(U) \arrow[d] \arrow[r, "ev(U)"] & \calf(U) \arrow[d] \\
\phom(\underline{\{p\}},\calf)(V) \arrow[r, "ev(V)"] & \calf(V).
\end{tikzcd}
\end{center}
It has an inverse
$$
ve(V):\calf(V) \lrta \phom(\underline{\{p\}},\calf)(V)
$$
$$
\calf(V)\ni f|_{V}\longmapsto \phi_f(V):[p\mapsto f|_V].
$$
$\phi_f$ determines a morphism of sheaves because $\res_{U,V}(\phi_f)$
$ve$ is well-defined morphism of sheaves because
$p\mapsto p\mapsto f|_U=p\mapsto f|_V\mapsto f|_U$
\begin{center}
\begin{tikzcd}
\phom(\underline{\{p\}},\calf)(U) \arrow[ddd] &  &  & \calf(U) \arrow[ddd] \arrow[lll, "ve(U)"'] \\
 & \phi_f(U)=[f|_U\mapsfrom p] \arrow[d, maps to] & f|_U \arrow[l, maps to] \arrow[d, maps to] &  \\
 & \phi_f|_V(V)=\phi_f(V)=[f|_V\mapsfrom p] & f|_V \arrow[l, maps to] &  \\
\phom(\underline{\{p\}},\calf)(V) &  &  & \calf(V) \arrow[lll, "ve(V)"']
\end{tikzcd}
\end{center}
They are indeed mutually inverse because
$$
ev(V)\circ ve(V):f|_V\mapsto [p\mapsto f|_V]\mapsto [p\mapsto f|_V](p)= f|_V
$$
and
$$
ve(V)\circ ev(V):\phi(V)\mapsto \phi(V)(p)\mapsto [p\mapsto \phi(V)(p)]=\phi(V)
$$

\item  $\underline{\intg}(V) $ is the set of locally constant maps from $V$ to $\intg$. (continuous maps from $V$ to $\intg$ with discrete topology. For $V$ connected $\underline{\intg}(V)= \intg$ and for $V$ with $n$ connected components, $\underline{\intg}(V)=\oplus_i^n\intg$ and the restriction map is just ordinary restriction of functions.) Similarly, we can define a morphism of sheaves of Abelian groups: $ev:\phom_{Ab_X}(\underline{\intg},\calf)\lrta\calf$
$$
ev(V):\phom(\underline{\intg},\calf)(V)\lrta\calf(V) 
$$
$$
\phi(V)\mapsto \phi(V)(1|_V),
$$
where $1|_V$ is the constant map that maps each point in $V$ to $1$. The explicit form of $1|_V$ can be $(1,1...,1)$. Notice that $\phi(V)$ is only morphism of Abelian groups and it is quite legal to map $1$ to $0$. There is a inverse of morphism of sheaves $ve:\calf\lrta \phom(\underline{\intg},\calf)$
$$
ve(V):\calf(V) \lrta \phom(\underline{\intg},\calf)(V)
$$
$$
\calf(V)\ni f|_V\longmapsto \phi_f(V):[\underline{\intg}(V)\ni 1|_V\mapsto f|_V].
$$
We have to mention that $1|_V$ does not generate $\underline{\intg}(V)$ in general but the morphism $\phi_f(V)$ is totally determined by its image of $1|_V$. For example $a\in\underline{\intg}(V)$, we can find an open cover $\{V_i\}_{i\in I}$ such that $a|_{V_i}$ is constant on each $V_i$. $1|_{V_i}$ is still the constant map to $1$. Then $a|_{V_i}=n_i\times 1|_{V_i}$ for some integer $n_i$.
$$
\phi_f(V_i):a|_{V_i}\longmapsto n_i\times \phi_f(V)(1|_{V})|_{V_i}=n_i\times f|_{V_i}\in \calf(V_i)
$$
If $V_i\cap V_j\neq \emptyset$, $a|_{V_i}$ and $a|_{V_j}$ has to agree on $V_i\cap V_j$, and hence $n_i=n_j$, hence
$$
\phi_f(V_i)(a|_{V_i})|_{V_i\cap V_j}=\phi_f(V_j)(a|_{V_j})|_{V_i\cap V_j}.
$$
Then because $\calf$ is  a sheaf, there is a unique section in $\calf(V)$ to glue to, we define it as $\phi_f(V)(a|_V)$. Then it is routine to check 
\begin{center}
\begin{tikzcd}
\phom(\underline{\intg},\calf)(U) \arrow[d] \arrow[r, "ev(U)"] & \calf(U) \arrow[d] \\
\phom(\underline{\intg},\calf)(V) \arrow[r, "ev(V)"] & \calf(V).
\end{tikzcd}
\end{center}
then diagram commutes and $ev$, $ve$ is mutually inverse.
\item $(b)$ is a special case of $(c)$.  Interpret $\phom_{Mod_{\calo_X}}(\calo_X,\calf)(V)$ as a $\calo_X(V)$-module by considering the $\calo_X(V)$-module structure of its image. $\calo_X(V)$ is  ring with unit element $1|_V$.

We can define a morphism of $\calo_X$-modules, $ev$
$$
ev(V):\phom_{Mod_{\calo_X}}(\calo_X,\calf)(V)\lrta\calf(V)
$$
$$
\phi(V)\mapsto \phi(V)(1|_V),
$$
each $ev(V)$ is indeed a morphism of $\calo_X(V)$-modules and it is routine to check $ev$ forms a morphism of $\calo_X$-modules. 
\begin{center}
\tiny
\begin{tikzcd}
\calo_X(U)\times \phom(\calo_X,\calf)(U) \arrow[rd, "id\times ev(U)"] \arrow[dd] \arrow[rrr, "action'"] &  &  & \phom(\calo_X,\calf)(U) \arrow[rd, "ev(U)"] \arrow[dd] &  \\
 & \calo_X(U)\times \calf(U) \arrow[dd] \arrow[rrr, "action"] &  &  & \calf(U) \arrow[dd] \\
\calo_X(V)\times\phom(\calo_X,\calf)(V) \arrow[rd, "id\times ev(V)"] \arrow[rrr, "action'"] &  &  & \phom(\calo_X,\calf)(V) \arrow[rd, "ev(V)"] &  \\
 & \calo_X(V)\times \calf(V) \arrow[rrr, "action"] &  &  & \calf(V),
\end{tikzcd}
\end{center}
where the vertical arrows a restrictions. The $action'$
 is  defined as 
 $$
 a|_V\times \phi(V)\mapsto [1|_V\longmapsto a|_V\phi(V)(1|_V)]
 $$
 and above diagram indeed commutes.

 Then, we  construct inverse of $ev$ by 
$$
ve(V):\calf(V) \lrta \phom_{Mod_{\calo_X}}(\calo_X,\calf)(V)
$$
$$
\calf(V)\ni f|_V\longmapsto \phi_f(V):[\calo_X(V)\ni 1|_V\mapsto f|_V].
$$
$\phi_f(V)$ is uniquely determined by it's image of $1|_V$.
Then it is routine to check 
\begin{center}
\begin{tikzcd}
\phom(\calo_X,\calf)(U) \arrow[d] \arrow[r, "ev(U)"] & \calf(U) \arrow[d] \\
\phom(\calo_X,\calf)(V) \arrow[r, "ev(V)"] & \calf(V).
\end{tikzcd}
\end{center}
then diagram commutes and $ev$, $ve$ is mutually inverse.
\end{enumerate}
\end{proof}

\begin{exr}
Show that $\ker_{pre}\phi$ is a presheaf. (Hint: if $U  \inj V$, define the restriction map by chasing the following diagram:
\begin{center}
\begin{tikzcd}
0 \arrow[r] & \ker_{pre}\phi(V) \arrow[r, hook] \arrow[d, "\exists !"', dotted] & \calf(V) \arrow[r] \arrow[d, "{\res_{V,U}}"] & \calg(V) \arrow[d, "{\res_{V,U}}"] \\
0 \arrow[r] & \ker_{pre}\phi(U) \arrow[r, hook] & \calf(U) \arrow[r] & \calg(U)
\end{tikzcd}
\end{center}
\end{exr}
\begin{proof}
$f|_V\in\ker_{pre}\phi(V)\subset \calf(V)$, $\res_{V,U}\phi(V)(f|_V)=\phi(U)(f|_U)=0$. Hence $f|_U=\res_{V,U}f|_U$ is in the kernel. It must be a \textbf{restriction of restriction map}, because $\ker_{pre}\phi(V)\inj\calf(V)$. Then it suffices to check only the dashed triangle commutes in the following diagram
\begin{center}
\begin{tikzcd}
0 \arrow[r] & \ker_{pre}(\phi)(V) \arrow[r, hook] \arrow[d, "\exists !"', dotted] \arrow[rdd, dotted] & \calf(V) \arrow[r] \arrow[d, "{\res_{V,U}}"] \arrow[rdd, "{\res_{V,W}}"] & \calg(V) \arrow[d, "{\res_{V,U}}"] \arrow[rdd, "{\res_{V,W}}"] &  \\
0 \arrow[r] & \ker_{pre}\phi(U) \arrow[r, hook] \arrow[rd, dotted] & \calf(U) \arrow[r] \arrow[rd, "{\res_{U,W}}"] & \calg(U) \arrow[rd, "{\res_{U,W}}"] &  \\
 & 0 \arrow[r] & \ker_{pre}\phi(W) \arrow[r, hook] & \calf(W) \arrow[r] & \calg(W)
\end{tikzcd}
\end{center}
It commutes because it can be regraded as part of the solid triangle of $\calf(V),\calf(W),\calf(U)$.
\end{proof}

\begin{exr}
Show that the presheaf cokernel satisfies the universal property of cokernels
\end{exr}
\begin{proof}, Given a morphism $\phi:\calf\lrta \calg$ and a morphism of presheaves $\sigma$ from $\calg$ to $\calc$ such that
We need to check that, the composition $\phi\circ \sigma=0$ there is a unique morphism of presheaves from $\coker \phi$ to $\calc$
\begin{center}
\begin{tikzcd}
\calc &  &  \\
\coker \phi  \arrow[u, "\exists!", dashed] & \calg \arrow[l, "\pi"] \arrow[lu] & \calf \arrow[l, "\phi"'] \arrow[llu, "0"'] \arrow[ll, "0", bend left]
\end{tikzcd}
\end{center}
By definition of $\coker \phi$, $(\coker \phi)(V):=\coker (\phi(V))$. Because each $\phi(V)$ is morphism of Abelian groups, each $\coker \phi(V)$ has the universal property and there exists a unique morphism $\psi(V):\coker\phi(V)\lrta \calc(V)$ such that the triangles in the below diagram commute.
\begin{center}
\begin{tikzcd}
\calc(V) &  & \calc(U) \arrow[ll,red] &  &  \\
\coker \phi(V) \arrow[u, "\exists!\psi(V)", dashed,red] &  & \coker\phi(U) \arrow[u, "\exists!\psi(U)", dashed,red] \arrow[ll,red] &  &  \\
 & \calg(V) \arrow[lu, "\pi(V)"] \arrow[luu, "\sigma(V)"'] &  & \calg(U) \arrow[lu, "\pi(U)"] \arrow[ll] \arrow[luu, "\sigma(U)"'] &  \\
 &  & \calf(V) \arrow[lu, "\phi(V)"] \arrow[lluuu, "0"'] \arrow[lluu, "0", bend left] &  & \calf(U) \arrow[lu, "\phi(U)"] \arrow[ll] \arrow[lluuu, "0"'] \arrow[lluu, "0", bend left]
\end{tikzcd}
\end{center}
It remains to check the red square commutes such that $\psi$ indeed determines a morphism of presheaves. Because $\pi(U)$ is an epimorphism. Surjective morphism of when regarded as morphism of sets are always epimorphism in small category. Then the red square has to be commutative.
\end{proof}

\begin{exr}\label{chap2:exr_presheave_section_funcotr_exact}
Show (orobserve) that for a topological space $X$with open set $U$, $\calf  \lrta \calf(U)$ gives a functor from presheaves of abelian groups on $X$, $Ab^{pre}_X$, to abelian groups, \textit{Ab}. Then show that this functor is exact.
\end{exr}
\begin{proof}
The diagram commutes by the definition of morphism of presheaves
\begin{center}
\begin{tikzcd}
\calf \arrow[r, "\phi"] \arrow[d, "\Gamma"'] & \calg \arrow[d, "\Gamma"] \arrow[r, "\psi"] & \calh \arrow[d, "\Gamma"] \\
\calf(U) \arrow[r, "\phi(U)"] & \calg(U) \arrow[r, "\psi(U)"] & \calh(U)
\end{tikzcd}
\end{center}
Taking the section of $U$ gives a functor because

Given a short exact sequence of presheaves
$$
0\lrta \calf\overset{\phi}{\lrta} \calg\overset{\psi}{\lrta} \calh\lrta 0
$$

The section functor $\Gamma(U,\_)$ maps $\calf\mapsto\calf(U)$. $0$ the the $0$-object in the $Ab_X^{pre}$
By definition of $\ker \phi$, $\ker\phi=0$ iff $\ker \phi(V)=0$ for each open set $V$. Similarly $\im \phi=\ker\psi$ iff $\im \phi(V)=\ker\psi(V)$ for all open sets $V$. $\coker \psi=0$ iff $\coker \psi(V)=0$.
This means in particular
$$
0\lrta \calf(U)\overset{\phi(U)}{\lrta} \calg(U)\overset{\psi(U)}{\lrta} \calh(U)\lrta 0
$$
is exact, which means $\Gamma(U,\_)$ is exact functor.
\end{proof}

\begin{exr}\label{chap2:exr2.3.H}
Show that a sequence of presheaves $$
0 \lrta \calf_1 \lrta \calf_2 \lrta ... \lrta \calf_n \lrta0$$
 is exact if and only if
$$0 \lrta \calf_1(U) \lrta \calf_2(U) \lrta ... \lrta \calf_n(U) \lrta0$$ is exact for all $U$.
\end{exr}
\begin{proof}
In~\ref{chap2:exr_presheave_section_funcotr_exact}, we have proved that the short sequence of presheaves
$$
0\lrta \calf\overset{\phi}{\lrta} \calg\overset{\psi}{\lrta} \calh\lrta 0
$$
is exact iff only
$$
0\lrta \calf(U)\overset{\phi(U)}{\lrta} \calg(U)\overset{\psi(U)}{\lrta} \calh(U)\lrta 0
$$ 
is exact. Because each long exact sequence splits into short exact sequences, we are done.

\end{proof}


\begin{exr}\label{chap2exr:kernel_sheaf}
Suppose $\phi: \calf\lrta\calg $ is a morphism of sheaves. Show that the presheaf kernel $\ker_{pre} \phi$ is in fact a sheaf. Show that it satisfies the universal property of kernels.
\end{exr}
\begin{proof}
The $0$-object in sheaves of abelian groups is just the $0$-presheaf (It is a constant presheaf of sets with one element set.)

Given an open set $U$, and an open cover $\{U_i\}_{i\in I}$.

Identity axiom: $k_1,k_2\in\ker_{pre}\phi(U)$ agree when restricted to any  $U_i$. They are regarded as section in $\calf(U)$, hence $k_1=k_2$, because $\calf$ is a sheaf.

Gluability axiom: $\{k_i\in \ker_{pre}(U_i)\}_{i\in I}$ compatible when restricted to intersections $U_i\cap U_j$. Again there is one element $k\in\calf(U)$ to glue to because $\calf$ is a sheaf. It remains to check that $k\in\ker_{pre}(U)$. $\phi(U)(k)|_{U_i}=\phi(U_i)(k_i)=0$, then by $\calg$ being a sheaf, there is a unique section to glue to in $\calg(U)$, which is $0$, therefore $k\in\ker_{pre}(U)$.

As for the universal property, it suffices to check that $\ker_{pre}\phi$ has the universal property in the category of presheaves, because category of sheaf is a full subcategory of presheaves and we have verified that $\ker_{pre} \phi$ is a sheaf.
\begin{center}
\tiny
\begin{tikzcd}
\calc(V) \arrow[d, "\exists!", dashed,red] \arrow[rdd] \arrow[rrddd, "0"] &  & \calc(U) \arrow[ll,red] \arrow[d, "\exists!", dashed,red] \arrow[rdd] \arrow[rrddd, "0"] &  &  \\
\ker\phi(V) \arrow[rd, "\iota(V)", hook] \arrow[rrdd, "0"', bend right] &  & \ker\phi(U) \arrow[ll,red] \arrow[rd, "\iota(U)", hook] \arrow[rrdd, "0"', bend right] &  &  \\
 & \calf(V) \arrow[rd, "\phi(V)"] &  & \calf(U) \arrow[ll] \arrow[rd, "\phi(U)"] &  \\
 &  & \calg(V) &  & \calg(U) \arrow[ll]
\end{tikzcd}
\end{center}
This time the red square commutes because $\iota(V)$ is injective (injective morphism is always monomorphism in small category)
\end{proof}

\begin{exr}
Let $X$ be $\cplx$ with the classical topology, let $\underline{\intg}$ be the constant sheaf on $X$ associated to $\intg$, $\calo_X$ the sheaf of holomorphic functions, and $\calf$ the presheaf of functions admitting a holomorphic logarithm. Describe an exact sequence of presheaves on $X$:
$$0 \lrta \underline{\intg}     \lrta \calo_X   \lrta \calf  \lrta 0$$
where $\underline{\intg} \lrta \calo_X$ is the natural inclusion and $\calo_X \lrta \calf$ is given by $f  \mapsto exp(2\pi if)$. (Be sure to verify exactness.) Show that $\calf$ is not a sheaf.
\end{exr}
\begin{proof}
By~\ref{chap2:exr2.3.H}, we can check 
$$
0 \lrta \underline{\intg}(U)     \lrta \calo_X (U)  \lrta \calf (U) \lrta 0
$$
is exact for any open subset of $\cplx$.

$\lrta \underline{\intg}(U)\lrta \calo_X(U)$ is exact because it is just inclusion.

$\calo_X (U)  \lrta \calf (U) \lrta$ is exact: Choose $g\in \calf(U)$, then $\frac{1}{2\pi i}\log g=h\in\calo_X(U)$ , then $\exp(2\pi h)=g$. This map is thus surjective.

$\lrta \calo_X (U)  \lrta $ is exact. $\ker(f\mapsto \exp(2\pi if))=\im (\underline{\intg}(U)\lrta\calo_X(U))$. Notice here $\calf(U)$ is an Abelian group under multiplication while the other two is under addition. $p\in\ker(f\mapsto \exp(2\pi if))$ means in fact $\exp(2\pi i p)=1$.

$\calf$ is not a sheaf, because it violates the gluability axiom. 

Consider the annulus $U$ around the $0$, covered by open disks. In each disks, we consider the constant functions, we can choose branch cuts according to each disks such that each constant function can have holomorphic logarithms. But they can not glue to an element in $\calf(U)$


\end{proof}
\section{Properties determined at the level of stalks, and sheafification}
\begin{exr}\label{chap2exr2.4.A}
Prove that a section of a sheaf of sets is determined by its germs, i.e., the natural map

$$
\calf (U) \lrta  \prod_{p\in U}\calf_p$$
is injective.
\end{exr}
\begin{proof}
$f_1,f_2\in \calf(U)$, assume they agree when restricted to each points. In particular, for each point $p$, there are open neighborhood $U_p,V_p$ such that $(f_1|_{U_P};U_p)\sim (f_2|_{V_p};V_p)$, which means there is an open neighborhood $W_p$  such that $f_1|_{W_p}=f_2|_{W_p}$. We choose such $W_p$ for each point. $\{W_p\}_{p\in U}$ form an open cover of $U$ and $f_1, f_2$ agree when restricted to each $W_p$. By identity axiom, $f_1=f_2$. We don't need the gluability axiom, thus this claim is also true for separated sheaves.
\end{proof}

\begin{exr}
Show that $\supp(s)$ is a closed subset of X.
\end{exr}
\begin{proof}
For $p\notin \supp (s)$, $s_p=0\in\calf_p$. This means $[(s|_V;V)]=[(0;U)]$, which means there is an open neighborhood $W$ of $p$ such that $s|_W=0$. Then $W\in X-\supp(s)$. This has effectively proved the complement of $\supp (s)$ is  open.
\end{proof}

\begin{exr}\label{chap2:exr2_compatible germs_glue_to_section}
Prove that any choice of compatible germs for a sheaf of sets $\calf$ over $U$ is the image of a section of $\calf$ over $U$.
\end{exr}
\begin{proof}
Suppose $\prod_{p\in U}s_p$ consists of compatible germs. There is an open cover $\{U_i\}_{i\in I}$ and choice of sections $f_i\in\calf(U_i)$ such that $(f_i)_p=s_p$ for all $p\in U_i$, In particular $f_i|_{U_i\cap U_j}=f_j|_{U_i\cap U_j}$ because they agree when restricted to stalks in the intersection. Then by gluability, there is a section $f\in \calf(U)$ such that $f|_p=s_p$ 
\end{proof}

\begin{exr}\label{chap2:exr2.4.D}
If $\phi_1$ and $\phi_2$ are morphisms from a \textbf{presheaf} of sets $\calf$ to a \textbf{sheaf} of sets $\calg$ that induce the same maps on each stalk, show that $\phi_1=\phi_2$.
\end{exr}
\begin{proof}
We temporarily regard $\calg$ as a presheaf. Consider the diagram
\begin{center}
\begin{tikzcd}
\calf(U) \arrow[d, hook] \arrow[r, "\phi(U)"] & \calg(U) \arrow[d, hook] \\
\prod_{p\in U}\calf_p \arrow[r, "\prod \phi_p"'] & \prod_{p\in U}\calg_p
\end{tikzcd}
\end{center}
The diagram by two points: 

1. by the universal property of $\calf_p$ (colimits), each $\phi_p$ makes $\phi_p\circ [\calf(U)\lrta \calf_p]=[\calg(U)\lrta \calg_p]\circ \phi(U)$.

2. By universal property of $\prod_p\calg_p$(product), $\prod_p\phi_p$ makes the above diagram commutes.

We only need to check $\phi(U)$ is uniquely determined by $\phi_p$. Consider an element $f\in \calf(U)$

\underline{Claim}: the image $\prod_p\phi_p(f_p)$ of $\prod\phi_p$ consists of compatible germs. This is true because $\calg$ is a (pre)sheaf and the diagram commutes.

Then there is a unique section $g$ to glue to in $\calg(U)$, and by identity axiom $\phi(U)f=g$. Each $\phi(U)$ is uniquely determined by $\phi_p$ and hence the morphism of presheaves(from a presheaf to a sheaf) $\phi$ is uniquely determined by $\phi_p$ for all points. 

We used only identity axiom here, and this statement is also true when we loosen $\calg$ to be separated presheaves
\end{proof}
\begin{exr}\label{chap2exr2.4.E}
Show that a morphism of sheaves of sets is an isomorphism if and only if it induces an isomorphism of all stalks.
\end{exr}
\begin{proof}
The ``only if'' part is obvious because of the universal property of colimits.

Then we focus on the ``if'' direction.

Previous exercise~\ref{chap2:exr2.4.D} has told us that morphism from a presheaf to a sheaf is uniquely determined by the induced morphism on stalks. We still consider the following diagram but then the problem reduces to the case when $\phi_p$ are isomorphisms and $\calf$ is also a sheaf.
\begin{center}
\begin{tikzcd}
\calf(U) \arrow[d,hook] \arrow[r, "\phi(U)"] & \calg(U) \arrow[d, hook] \\
\prod_{p\in U}\calf_p \arrow[r, "\prod \phi_p"'] & \prod_{p\in U}\calg_p
\end{tikzcd}
\end{center}
Assume each $\phi_p$ is isomorphism, want to prove each $\phi(U)$ is isomorphism.

\underline{Want}: $\phi(U)$ is injective: $f_1,f_2\in\calf(U)$, $\phi(U)f_1=\phi(U)f_2$.

 $f_1,f_2$ restrict to $\prod_p f_{1,p}$, $\prod_p f_{2,p}$ and then map to $\prod_p\phi_p(f_{1,p})$, $\prod_p\phi_p(f_{2,p})$ respectively.

 $\phi(U)f_1=\phi(U)f_2$ also restricts to $\prod_p\phi_p(f_{1,p})$, $\prod_p\phi_p(f_{2,p})$ respectively. $\prod_p\phi_p(f_{1,p})=\prod_p\phi_p(f_{2,p})$, because $\calg(U)\lrta \prod_p\calg_p$ is injective.

Then, because each $\phi_p$ is isomorphism, we have $\prod_p f_{1,p}=\prod_p f_{2,p}$. Then $f_1=f_2$ because $\calf(U)\lrta \prod_p\calf_p$ is injective.

\underline{Want}: $\phi(U)$ is surjective: $g\in\calg(U)$, want to find a preimage $f\in\calf(U)$.

$g\mapsto\prod_pg_p\in\prod_p\calg_p$ consists of compatible germs. And because $\phi_p$ is isomorphism, there is $\prod_pf_p\in \prod_{p\in U}\calf_p$ that maps to $\prod_pg_p$.

\underline{Claim}: $\prod_pf_p$ consists of compatible germs.
$$
\prod_p f_p\mapsto \prod_p g_p
$$
$$
[(f'_{(p)}\in\calf(U_p);U_p)]\mapsto [(\phi(U_p)f'_{(p)};U_p)]
$$
We choose $\{U_p\}_{p\in U}$ a cover fine enough such that $(\phi(U_p)f'_{(p)};U_p)$ form a representative of compatible germs $g_p$. By definition of compatible germs, for every $q\in U_p$, 
$$
\phi(U_p)f'_{(p)}|_q=[(\phi(U_q)f'_{(q)};U_q)]=g_q
$$
which means there exists open $W_{p,q}\subset U_p\cap U_q$ such that
$$
\phi(W_{p,q})f'_{(p)}|_{W_{p,q}}=\phi(W_{p,q})f'_{(q)}|_{W_{p,q}}.
$$
But We already proved that each $\phi(U)$ is injective, we have
$$
f'_{(p)}|_{W_{p,q}}=f'_{(q)}|_{W_{p,q}}.
$$
In particular,
$$
f'_{(p)}|_q=f_q,
$$
which means our choice of representative $[(f'_{(p)}\in\calf(U_p);U_p)]$ satisfies the requirement of compatible germs.

Then by exercise~\ref{chap2:exr2_compatible germs_glue_to_section}, we know there is a section $f\in\calf(U)$ such that $f$ to restricts to $\prod_p f_p$ and such $\phi(U)f=g$ because they agree on stalks.


This exercise are restricted to the case we know the morphism between sheaves.  It does not say that if two sheaves have isomorphic stalks, then they are isomorphic.
\end{proof}

\begin{exr}\ 
\begin{enumerate}[label=(\alph*)]
\item
Show that Exercise~\ref{chap2exr2.4.A} is false for general presheaves.
\item Show that Exercise~\ref{chap2:exr2.4.D} is false for general presheaves.
\item Show that Exercise~\ref{chap2exr2.4.E} is false for general presheaves.
\end{enumerate}
\end{exr}
\begin{proof}
\begin{enumerate}[label=(\alph*)]
\item Consider a two point set with discrete topology $X=\{p,q\}$. Assume $\calf$ is the presheaf such that  $\calf(X)=\{a,b\}$ and $\calf_p=\{a\},\calf_q=\{b\}$, $\calf_p\times\calf_q=\{(a,b)\}$. There is no injection from $\{a,b\}$ to $\{(a,b)\}$.
\item Still consider a same topological space  $X=\{p,q\}$. $\calf(X)=\{a,b,c\}$ and $\calf_p=\{a\},\calf_q=\{a\}$ and $\calg(X)=\{a,b,c\}$. Consider $\phi_1(X)=id$ and $\phi_2(X)$ switch $b,c$ but fixes $a$. They would  agree on stalks. 
\item Still consider the topological space. $X=\{p,q\}$. $\calf(X)=\{a,b,c\}$ and $\calf_p=\{a\},\calf_q=\{a\}$ and $\calg(X)=\{a,b,c\}$. This time consider the map $\phi(X):a\mapsto a,b\mapsto b,c\mapsto b$, it induces isomorphisms on the stalks but itself is not an isomorphism.
\end{enumerate}
\end{proof}

\begin{exr}
Show that sheafification is unique up to unique isomorphism, assuming it exists. Show that if $\calf$ is a sheaf, then the sheafification is $id:\calf\lrta\calf$.
\end{exr}
\begin{proof}
By universal property, nothing special here.
\end{proof}

\begin{exr}
Assume for now that sheafification exists. Use the universal property to show that for any morphism of presheaves $\phi\calf\lrta\calg$ , we get a natural induced morphism of sheaves $\phi^{sh}:\calf^{sh}\lrta\calg^{sh}$. Show that sheafification is a functor from presheaves on $X$ to sheaves on $X$.
\end{exr}
\begin{proof}
There is a morphism of sheaves form $\calf$ to $\calg^{sh}$,
\begin{center}
\begin{tikzcd}
\calf \arrow[d, "\phi"] \arrow[r] & \calf^{sh} \arrow[d, dashed,"\exists!\phi^{sh}"'] \arrow[dd, dashed, bend left,"\exists!(\psi\circ \phi)^{sh}"] \\
\calg \arrow[r] \arrow[d, "\psi"] & \calg^{sh} \arrow[d, dashed,"\exists!\psi^{sh}"'] \\
\calh \arrow[r] & \calh^{sh}
\end{tikzcd}
\end{center} 
then there exists a unique morphism from $\calf^{sh}$ to $\calg^{sh}$ such that the diagram commutes. Sheafification forms a functor because the morphism from $\calf^{sh}$ to $\calh^{sh}$ is unique, hence $\psi^{sh}\circ \phi^{sh}=(\psi\circ \phi)^{sh}$. 
\end{proof}

\begin{exr}\label{chap2exr.sheafication_is_a_sheaf}
Show that $\calf^{sh}$ (using the tautological restriction maps) forms a sheaf.
\end{exr}
\begin{proof}
$$
\calf^{sh}(U)\subset\prod_{p\in U}\calf_p
$$
as the set of compatible germs. It is a presheaf and the restriction map is induced by
$$
\res_{U,V}:\prod_{p\in U}\calf_p\lrta\prod_{q\in V}\calf_q.
$$
Given an open covering $\{U_i\}_{i\in I}$ of $U$.

Identity axiom: $f_1=\prod_{p\in U}f_{1,p},f_2=\prod_{q\in U}f_{2,q}\in \calf^{sh}(U)$. They agree when restricted to each open. Tautologically, this means $f_{1,q}=f_{2,q}\forall q\in U_i$. It means $f_{1,p}=f_{2,p}\forall p\in U$. And by definition this means $f_1=f_2$.

Gluability axiom: $\{f_i\in \calf^{sh}(U_i)\}$ has the property $f_i|_{U_i\cap U_j}=f_j|_{U_j\cap U_i}$. We can choose $f_p\in\calf_p$ such that $f_p=f_{i,p}\forall p\in U_i\subset U$. These are compatible germs, because $f_{i,p}$ are compatible germs. We can find $s_i\in\calf(U_i)$ such that $s_{i,p}=f_{i,p}=f_p$. Then by definition $\prod_{p\in U}f_p\in \calf^{sh}$. We define the bunch of $\prod_{p\in U}f_p$ to be $f$ and one can check that $f|_{U_i}=f_i$.

In this proof $f_p$ means the $p$-entry of $f$ but not the germ of $f$ at $p$. (We will see they are equal afterwards)
\end{proof}

\begin{exr}
Describe a natural map of presheaves $sh:\calf\lrta \calf^{sh}$.
\end{exr}
\begin{proof}
Consider the natural map, (we extend to presheaves, where it is no longer injective but it is still natural.)
$$
\calf(U)\lrta \prod_{p\in U}\calf_p
$$
$$
\calf(U)\ni s\lrta \prod_{p\in U}[(s;U)]_p,
$$
where $[(s;U)]_p$ means taking the equivalent class of $(s;U)$ at point $p$.
The image of $\calf(U)$ consists of compatible germs, hence lies in $\calf^{sh}(U)$. 
$$
\calf(U)\overset{sh(U)}{\lrta}\calf^{sh}(U)\inj \prod_{p\in U}\calf_p,
$$
we define the map to be $sh(U)$. $sh(U)$ is a well-defined morphism of sets and it preserves the extra structures on the sets.
We just need to prove that $sh(U)$ indeed determines a morphism of presheaf, which reduce to checking the commutativity of the left square in the diagram
\begin{center}
\begin{tikzcd}
\calf(U) \arrow[d] \arrow[r, "sh(U)"] & \calf^{sh}(U) \arrow[d] \arrow[r, hook] & \prod_{p\in U}\calf_p \arrow[d] \\
\calf(V) \arrow[r, "sh(V)"] & \calf^{sh}(V) \arrow[r, hook] & \prod_{q\in V}\calf_q.
\end{tikzcd}
\end{center}
The restriction map $\calf^{sh}(U)\lrta\calf^{sh}(V)$ is just induced by the restriction map $\prod_{p\in U}\calf_p\lrta \prod_{q\in V}\calf_q$. Hence the left square commutes since the outer larger square commutes.
\end{proof}
\begin{exr}
Show that the map $sh$ satisfies the universal property of sheafification.
\end{exr}
\begin{proof}
The unique morphism $\psi$ from $\calf^{sh}$ to $\calg$ is induced at the level of stalks.
\begin{center}
\begin{tikzcd}
\calf \arrow[r, "sh"] \arrow[d, "\phi"'] & \calf^{sh} \arrow[r] \arrow[d, "sh(\phi)"] & \prod_p\calf_p \arrow[d, "\prod_p\phi_p"] \\
\calg \arrow[r, "sh"] & \calg^{sh} \arrow[r] & \prod_p\calg_p
\end{tikzcd}
\end{center}
$sh(\phi)$ is induced by $\prod_{p}\phi_p$. Uniquely determined by the morphism on stalks. We need to prove compatible germs are sent to compatible germs.
$$
\prod_p\phi_p:\prod_pf_p\mapsto\prod_pg_p
$$
Know:$f_p=[(\tilde{s}_p;V_p)]$ such that $\forall q\in V_p$ $f_q=\tilde{s}_p|_q$.

Then, $g_p=\phi_p f_p=[(\phi(V_p)\tilde{s}_p;V_p)]$ any $q\in V_p$ $g_q=\phi_qf_q=[(\phi(V_q)\tilde{s}_q;V_q)]_p=[(\phi(V_q\cap V_p)\tilde{s}_p;V_q\cap V_p)]_q=(\phi(V_p)\tilde{s}_p)_q$.
As a result, we know the above diagram commutes and it remains to check $\calg^{sh}=\calg$ if $\calg$ is a sheaf (We have to check it without the universal property).

Notice that $\calg_p=\calg^{sh}_p$ (by~\ref{chap2:exr_shfification_iso_stalks}), the $sh$ induce isomorphism on stalks, and by Exercise~\ref{chap2exr2.4.E}. $\calg\cong\calg^{sh}$ as sheaves.
\end{proof}

\begin{exr}
Show that the sheafification functor is left-adjoint to the forgetful functor from sheaves on $X$ to presheaves on $X$.
\end{exr}
\begin{proof}$\calc$ is a small category and we denote the category of presheaves of objects in $\calc$ on a topological space $X$ by $\calc^{pre}_X$ and the corresponding category of sheaves by $\calc_X$.
Denote the forgetful  functor $i:\calc_X\lrta \calc_X^{pre}$. Still, we look at the diagram.
\begin{center}
\begin{tikzcd}
\calf \arrow[r, "sh"] \arrow[d, "\phi"'] & \calf^{sh} \arrow[r] \arrow[d, "sh(\phi)"] & \prod_p\calf_p \arrow[d, "\prod_p\phi_p"] \\
\calg \arrow[r, "sh"] & \calg^{sh} \arrow[r] & \prod_p\calg_p
\end{tikzcd}
\end{center}
It says for each morphism of presheaves $\phi:\calf\lrta i(\calg)$, there is a unique morphism of sheaves $sh(\phi):\calf^{sh}\lrta\calg^{sh}$. And in return, for each morphism of sheaves $\psi:\calf^{sh}\lrta\calg$, there is a morphism $i(\psi)$ of presheaves. Hence they establish a bijection
$$
\hom_{\calc^{pr}_X}(\calf,i(\calg))\cong\hom_{\calc_X}(\calf^{sh},\calg)
$$
{\color{red} It is also true that we have the following isomorphism of presheaves of sets, where we have extended sheaf hom to presheaf hom.
$$
\phom(\calf,i(\calg))\cong\phom(\calf^{sh},\calg)
$$
}

\end{proof}

\begin{exr}\label{chap2:exr_shfification_iso_stalks}
Show $\calf\lrta \calf^{sh}$ induces an isomorphism of stalks.
\end{exr}
\begin{proof}
By definition,
$$
\calf^{sh}(U)=\left(\prod_{p\in U}f_p: \text{ for each $p\in U$, $\exists(\tilde{s}_p;V_p)$ representing $f_p$ s.t. $f_q=\tilde{s}_p|_q\forall q\in V_q$}\right)
$$
Consider $\bar{f}=\prod_{p\in U}f_p\in \calf^{sh}(U)$, for each $p\in U$, $\exists(\tilde{s}_p;V_p)$ representing $f_p$ s.t. $f_q=\tilde{s}_p|_q\forall q\in V_q$. And the $\res_{U,V_p}\bar{f}=\prod_{q\in V_p}f_q$. We can define the map
$$
\calf^{sh}_p\ni[(\bar{f}|_{V_p};V_p)]\longmapsto f_p\in \calf_p.
$$
Then map is well-defined. (It is just the projection of the product to $p$-entry.) It has a well-defined inverse $sh_p$:
$$
sh_p:f_p=[(\tilde{s}_p;V_p)]\longmapsto [(\prod_{q\in V_p}\tilde{s}_p|_q;V_p)]
$$
\end{proof}

\begin{exr}
Suppose $\phi: \calf\lrta\calg$ is a morphism of sheaves of sets on a topological space $X$. Show that the following are equivalent.
\begin{enumerate}[label=(\alph*)]
\item
 $\phi$ is a monomorphism in the category of sheaves.
\item $\phi$ is injective on the level of stalks: $\phi_p:\calf_p \lrta \calg_p$ is injective for all $p\in X$.
\item $\phi$ is injective on the level of open sets:$\phi(U):\calf(U)\lrta\calg(U)$  is injective for all open $U\subset X$.
\end{enumerate}
\end{exr}
\begin{proof}
$(a)\Lrta (c)$: Consider the indicator sheaf $\cali_U$. ($\cali_U(V)=\{*\},\forall V\subset U$ and $\cali_U(V)=\emptyset,\forall V\notsubset U$). Assume there are two morphisms of sheafs from $\cali_U$ to $\calf$.
\begin{center}
\begin{tikzcd}
\cali_U \arrow[r, "\psi_2"', bend right] \arrow[r, "\psi_1", bend left] & \calf \arrow[r, "\phi"] & \calg
\end{tikzcd}
\end{center}
$\phi\circ \psi_1=\phi\circ \psi_2$ $\Llrta$ $\phi(V)\circ\psi_1(V)=\phi(V)\circ \psi_2(V),\forall V\subset X$.

$\phi:\calf\lrta \calg$ being a monomorphism implies $\cali_U$ $\Llrta$ $\phi\circ \psi_1=\phi\circ \psi_2\Lrta \psi_1=\psi_2$.

$\Llrta$ $\phi(V)\circ\psi_1(V)=\phi(V)\circ \psi_2(V),\forall V\subset X\Lrta \psi_1(V)=\psi_2(V)\forall V\subset X$.

Notice that in the special case of indicator sheaf. $\psi_i(V)$ is empty function for $V\notsubset U$, then the condition $\phi(V)\circ\psi_1(V)=\phi(V)\circ \psi_2(V)$ is vacuous. The only efficient conditions are $\phi(V)\circ\psi_1(V)=\phi(V)\circ \psi_2(V),\forall V\subset U$, which is equivalent to $\phi(U)\circ\psi_1(U)=\phi(U)\circ \psi_2(U)$ alone, (by definition of morphism of presheaves).

Then the origin statement means
$\phi(U)\circ\psi_1(U)=\phi(U)\circ \psi_2(U)\Lrta \psi_1(U)=\psi_2(U)$. $\psi_i(U)$ are just an injection of $\{*\}$ into $\calf(U)$, which are uniquely determined by their images $f_i:=\psi_i(U)(*)$. This means in particular, $\phi(U)(f_1)=\phi(U)(f_2)\Lrta f_1=f_2$. $f_i$ are initially arbitrary, we know $\phi(U)$ is injective. And $U$ is arbitrarily chosen, we know $\phi(U)$ is injective for all open $U\subset X$.

(This part of proof is also true for presheaves)

$(c)\Lrta (b)$: True because \textbf{filtered colimits preserve finite limits}. Or we can prove it by hand:
$\phi_p(f_{p})=\phi_p(f'_p)$ 
$$
\Llrta \phi(V)f|_V=\phi(V)f'|_V,\exists V\ni p
$$
$$
\Llrta f|_V=f'|_V \exists V\ni p\text{ because $\phi(V)$ is injective}
$$
$$
\Lrta f_p=f_p'
$$

$(b)\Lrta (a)$. Assume $\phi\circ \psi_1=\phi\circ \psi_2$, and this implies $\phi_p\circ \psi_{1,p}=\phi_p\circ \psi_{2,p}$ and we know in the category of sets injection is monomorphism, $\Lrta \psi_{1,p}=\psi_{2,p}$ and then $\psi_{1}=\psi_{2}$ because morphism from a presheaves to a sheaf are determined at the level of stalks~\ref{chap2:exr2.4.D}. 
\end{proof}

\begin{remark}
In the category of presheaves, $(a)\Llrta (c)$ and $(c)\Lrta (b)$.

If furthermore, we restrict to separated presheaves $(a)\Llrta (c)\Llrta (b)$
\end{remark}

\begin{exr}\label{chap2:exr.surj_stalk_epimorphism_sheaf}
Continuing the notation of the previous exercise, show that the following are equivalent.
\begin{enumerate}[label=(\alph*)]
\item
 $\phi$ is an epimorphism in the category of sheaves.
\item $\phi$ is surjective at the level of stalks: $\phi_p$ is surjective for each $p\in X$.
\end{enumerate}
\end{exr}
\begin{proof}
$(a)\Lrta (b)$: Consider the skyscraper sheaf $i_{p,*}S$. 
\begin{center}
\begin{tikzcd}
\calf \arrow[r, "\phi"] & \calg \arrow[r, "\psi_1", bend left] \arrow[r, "\psi_2"', bend right] & i_{*,p}S
\end{tikzcd}
\end{center}
$\phi$ is epimorphism: $\psi_1\circ \phi=\psi_2\circ \phi\Lrta \psi_1=\psi_2$

Suppose $\phi_p$ is not surjective.  There is an element $k_p\in\calg_p$ such that $k_p\notin im(\phi_p)$. Then we can construct $\psi_{1,p}$ and $\psi_{2,p}$ such that they agrees on the image of $\phi_p$ but differs on $k_p$. Then we have
$$
\psi_{1,p}\phi_p=\psi_{2,p}\phi_p,\text{ but } \psi_{1,p}\neq \psi_{2,p}.
$$
Notice that for skyscraper sheaf $q\neq p$, $\psi_{i,q}$ is the unique morphism to final object $\{*\}$ and  hence $\psi_{1,q}\phi_q=\psi_{2,q}\phi_q\forall q\in X$.
Remember that morphism of sheaves are determined at the level of stalks. This means $\phi$ can not be a epimorphism in the category of sheaves.

$(b)\Lrta (a)$: Assume $ \psi_1\circ \phi= \psi_2\circ \phi$, and this implies $\psi_{1,p}\circ \phi_p= \psi_{2,p}\circ \phi_p,\forall p$ and we know in the category of sets surjection is monomorphism, $\Lrta \psi_{1,p}=\psi_{2,p}$ and then $\psi_{1}=\psi_{2}$ because morphism from a presheaves to a sheaf are determined at the level of stalks~\ref{chap2:exr2.4.D}. 
\end{proof}
\begin{remark}
In the category of presheaves of small category, being a epimorphism is equivalent to being surjective at open set level. See \href{https://stacks.math.columbia.edu/tag/00V5}{Stack Project Tag 00V5}.
\end{remark}
\begin{exr}
Show that $\exp: \calo_X \lrta \calo_X^*$ describes $\calo_X^*$ as a quotient sheaf of $\calo_X$. Find an open set on which this map is not surjective.
\end{exr}
\begin{proof}
Exercise~\ref{chap2:exr.surj_stalk_epimorphism_sheaf} says $\calo_X^*$ is a quotient sheaf of $\calo_X$ iff $\exp_p$ is surjective. Locally logarithm is well defined for nonzero holomorphic functions. 
$$
\log f(z)=\int_\gamma\frac{f'(z)}{f(z)} dz,
$$
where $\gamma$ is path from $0$ to $z$ and the definition is not dependent on the choice of paths. And we have 
$$
\exp(\log f)=f
$$
Hence $\exp_p$ is surjective. 

Again, we can not define logarithm on annulus $U$. There is no holomorphic function $f\in\calo_X(U)$ such that $\exp(f)=z\in\calo_X^{*}(U)$.
\end{proof}


\section{Recovering sheaves form a ``sheaf on a base''}
\begin{exr}
How can you recover a sheaf $\calf$ from this partial information on a base of topology?
$$
(\{\calf(B_i)\},\{\res_{B_i,B_j}:B_i\rta B_j\})
$$
\end{exr}
\begin{proof}
As defined in the below text, a sheaf on a base $\{B_j\}$ of $X$, is denoted by $\pzf$. And we can recover a sheaf on $X$ by 
$$
\calf(U):=\left\{\prod_{p\in U}(f_p\in\pzf_p):\forall p\in U,\exists B, p\in B\subset U,\text{ and } s\in \pzf(B), \text{ s.t.} s_q=f_q\forall q\in B\right\}
$$
We will prove $\calf$ is a sheaf and the proof is identical to~\ref{chap2exr.sheafication_is_a_sheaf}
$$
\calf(U)\subset\prod_{p\in U}\pzf_p
$$
as the set of compatible germs. It is a presheaf and the restriction map is induced by
$$
\res_{U,V}:\prod_{p\in U}\pzf_p\lrta\prod_{q\in V}\pzf_q.
$$
Given an open covering $\{U_i\}_{i\in I}$ of $U$.

Identity axiom: $f_1=\prod_{p\in U}f_{1,p},f_2=\prod_{q\in U}f_{2,q}\in \calf(U)$. They agree when restricted to each open. Tautologically, this means $f_{1,q}=f_{2,q}\forall q\in U_i$. It means $f_{1,p}=f_{2,p}\forall p\in U$. And by definition this means $f_1=f_2$.

Gluability axiom: $\{f_i\in \calf(U_i)\}$ has the property $f_i|_{U_i\cap U_j}=f_j|_{U_j\cap U_i}$. We can choose $f_p\in\calf_p$ such that $f_p=f_{i,p}\forall p\in U_i\subset U$. We can find, $p\in B_i\subset U_i$ and  $s_i\in\pzf(B_i)$ such that $s_{i,q}=f_{i,q}=f_q$. Then by definition $\prod_{p\in U}f_p\in \calf$. We define the bunch of $\prod_{p\in U}f_p$ to be $f$ and one can check that $f|_{U_i}=f_i$.

In this proof $f_p$ means the $p$-entry of $f$ but not the germ of $f$ at $p$. (We will see they are equal afterwards)
\end{proof}
\begin{exr}
Verify that $\calf(B) \lrta \pzf(B)$ is an isomorphism, likely by showing that it is injective and surjective (or else by describing the inverse map and verifying that it is indeed inverse).
\end{exr}
\begin{proof}
We can not use exercise~\ref{chap2exr2.4.E} because $\pzf$ is not a sheaf. But we can construct the inverse at the level of stalks.
$$
\calf(B)\lrta \prod_{p\in B}\calf_p
$$
is injective by~\ref{chap2exr2.4.A}. Similarly, 
$$
\pzf(B)\lrta \prod_{p\in B}\pzf_p
$$
is injective, because it has base identity axiom. Elements of $\calf(B)$ are determined by stalks, as are elements of $\pzf(B)$. 

We define the morphism directly
$$
ex:\pzf(B)\lrta \calf(B)
$$
$$
s\mapsto \prod_{p\in B}[(s;B)]_p,
$$
where $[(s;U)]_p$ means taking the equivalence class at point $p$, by definition $\prod_{p\in B}[(s;B)]_p\in \calf(B)$. 

It is injective because this map is identical to $\pzf(B)\lrta \prod_{p\in B}\pzf(B)$. 

It is surjective. Consider $f:=\prod_{p\in B}f_p\in \calf(B)$. By definition, for a point $p$, there exists $B_i$, $s_i\in \pzf(B_i)$ such that $s_i|_q=f_q\forall q\in B_i$. Such $\{B_i\}$ covers $B$ and by base gluability, $\exists s\in\pzf(B)$ such that $s|_{B_i}=s_i$. Then we have $\prod_{p\in B}f_p=\prod_{p\in U}[(s;U)]_p$. 
\end{proof}
\begin{remark}
Similar to sheafification, the sheaf $\calf$ corresponding to a sheaf on a base $\pzf$ is unique up to unique isomorphism in the sense $\calf$ is the initial object in category of sheaves that extend a sheaf on base $B_i$, $\pzf$.
\begin{center}
\begin{tikzcd}
\pzf \arrow[rd, "\gamma"'] \arrow[r, "ex"] & \calf \arrow[d, "\exists!", dashed] \\
 & \calg
\end{tikzcd}
\end{center}
where $\phi,\gamma$ are morphism of ``sheaf on a base $B_i$''(there is a forgetful functor from category of sheaves to category of sheaves on a base $B_i$)

Again, $ex$ is a functor and it is left-adjoint to the forget functor from category of sheaves to category of sheaves on a base $B_i$.

We also know $\calf_p\cong \pzf_p\forall p\in X$. The proof is identical to~\ref{chap2:exr_shfification_iso_stalks}.
\end{remark}

\begin{exr}\label{chap2exr:equivalence_of_sheaves_on_base_and_sheaves}
MORPHISMS OF SHEAVES CORRESPOND TO MORPHISMS OF SHEAVES ON A BASE. Suppose $\{B_i\}$ is a base for the topology of $X$. A morphism $\pzf\lrta\pzg$ of sheaves on the base is a collection of maps $\pzf(B_k) \lrta \pzg(B_k)$ such that the diagram commutes:
\begin{center}
\begin{tikzcd}
\pzf(B_i) \arrow[r] \arrow[d] & \pzg(B_j) \arrow[d] \\
\pzf(B_j) \arrow[r] & \pzg(B_j)
\end{tikzcd}
\end{center}
\begin{enumerate}[label=(\alph*)]
\item Verify that a morphism of sheaves is determined by the induced morphism of sheaves on the base
\item Show that a morphism of sheaves on the base gives a morphism of the induced sheaves
\end{enumerate}
\end{exr}
\begin{proof}
\begin{enumerate}[label=(\alph*)]
\item  There is a forgetful functor $\iota$ form  $\calc_{X}$ (category of $\calc$-sheaves on $X$) to $\calc_{X}^{B}$ (category of sheaves on a base $B_i$ of the topology of $X$).
Each morphism of sheaves $\psi:\calf\lrta\calg$ induces a morphism $\iota(\psi):\pzf\lrta \pzg$. Recall that $\pzf_p\cong\calf_p$ $\iota(\psi)$ induces same morphism on stalks as $\psi_p$. And a morphism of sheaves is determined by the morphism it induces on stalks~\ref{chap2:exr2.4.D}.
\item
The unique morphism $ex(\phi)$ from $\calf$ to $\calg$ is induced at the level of stalks.
\begin{center}
\begin{tikzcd}
\pzf \arrow[r, "ex"] \arrow[d, "\phi"'] & \calf \arrow[r,hook] \arrow[d, "ex(\phi)"] & \prod_p\pzf_p \arrow[d, "\prod_p\phi_p"] \\
\pzg \arrow[r, "ex"] & \calg \arrow[r,hook] & \prod_p\pzg_p
\end{tikzcd}
\end{center}
$ex(\phi)$ is induced by $\prod_{p}\phi_p$. Uniquely determined by the morphism on stalks. We need to prove compatible germs are sent to compatible germs.
$$
\prod_p\phi_p:\prod_pf_p\mapsto\prod_pg_p
$$
Know:$f_p=[(\tilde{s}_p;B_p)], \tilde{s}_p\in \pzf(B_p)$ such that $\forall q\in B_p$ $f_q=\tilde{s}_p|_q$.

Then, $g_p=\phi_p f_p=[(\phi(V_p)\tilde{s}_p;B_p)]$ any $q\in B_p$ $g_q=\phi_qf_q=[(\phi(B_q)\tilde{s}_q;B_q)]_p=[(\phi(B_q\cap B_p)\tilde{s}_p;B_q\cap B_p)]_q=(\phi(B_p)\tilde{s}_p)_q$.
As a result, we know the above diagram commutes.
\end{enumerate}
\end{proof}
\begin{remark}
The above exercise basically means $ex$ is left-adjoint to $\iota$
$$
\mor_{\calc_X^{B}}(\pzf,\iota(\calg))\cong\mor_{\calc_X}(ex(\pzf),\calg)
$$
\end{remark}

\begin{exr}\label{chap3exr:glue_sheaf}
Suppose $X = \cup_{i\in I}U_i$ is an open cover of $X$, and we have sheaves $\calf_i$ on $U_i$ along with isomorphisms $\phi_{ij}:\calf_i|_{U_i\cap U_j} \lrta \calf_j|_{U_i\cap U_j}$ (with $\phi_{ii}$ the identity) that agree on triple overlaps, i.e.,$ \phi_{jk} \circ \phi_{ij} = \phi_{ik}$ on $U_i\cap U_j \cap U_k$ (this is called the \textbf{cocycle condition}, for reasons we ignore). Show that these sheaves can be glued together into a sheaf $\calf$ on $X$ (unique up to unique isomorphism), such that $\calf_i \cong\calf|_{U_i}$ , and the isomorphisms over $U_i \cap U_j$ are the obvious ones. (Thus we can ``glue sheaves together'', using limited patching information.)
\end{exr}
\begin{proof}
Following the hint, we first construct a base of topology associated to the open cover $\{U_i\}_{i\in I}$.
$$
\omega:=\left\{(B_{i,\alpha})_{\alpha\in A_i,i\in I}\text{: open set in $U_i$}\right\}/\sim,
$$
where $B_{i,\alpha}\sim B_{j,\beta}$ iff $B_{i,\alpha}=B_{j,\beta}$ as open sets in $X$.

It is routine to verify it is really a base of topology. In addition each $\{B_{i,\alpha}\}_{\alpha\in A_i}$ forms a base of $U_i$.

Then we apply the forgetful functor $\iota:\calc_X\lrta\calc_{X}^B$ such that we get $\pzf_i$ as a sheaf on base $\{B_{i,\alpha}\}_{\alpha\in A_i}$ of $U_i$, with the induced restriction map from $\calf_i$.
 $$
\res^B_{B_{i,\alpha},B_{i,\beta}}:=\res^{(i)}_{B_{i,\alpha},B_{i,\beta}},
$$

\underline{Claim}: The ``Atlas'' of the data can form a sheaf on the base $\omega$, we denote it by $\pzf$. To be more concrete, we need to check the compatibility of these objects and morphisms.


\begin{center}
\tiny
\begin{tikzpicture}
\draw (2,2) circle (2cm);
\draw (3,2) circle (2cm);
\draw (2.5,1.43) circle (2cm);
\draw (2.5,1.7) circle (1.2cm);
\draw (2.5,1.7) circle (1cm);
\draw (2.5,1.7) circle (0.8cm);
\end{tikzpicture}
\end{center}

Assume $B_\gamma\subset B_\beta\subset B_\alpha\subset U_i\cap U_j\cap U_k$, we need to check the following diagram commutes (the horizontal arrows are restriction maps in $\pzf_i,\pzf_j,\pzf_k$).
\begin{center}
\tiny
\begin{tikzcd}
\pzf_i(B_\alpha) \arrow[dd, "\phi_{ik}(B_\alpha)"'] \arrow[rd, "\phi_{ij}(B_\alpha)"] \arrow[rr] \arrow[rrrr, bend left] &  & \pzf_i(B_\beta) \arrow[dd, "\phi_{ik}(B_\beta)"'] \arrow[rd, "\phi_{ij}(B_\beta)"] \arrow[rr] &  & \pzf_i(B_\gamma) \arrow[dd, "\phi_{ik}(B_\gamma)"'] \arrow[rd, "\phi_{ij}(B_\gamma)"] &  \\
 & \pzf_j(B_\alpha) \arrow[ld, "\phi_{jk}(B_\alpha)"] \arrow[rr] \arrow[rrrr, bend left] &  & \pzf_j(B_\beta) \arrow[ld, "\phi_{jk}(B_\beta)"] \arrow[rr] &  & \pzf_j(B_\gamma) \arrow[ld, "\phi_{jk}(B_\gamma)"] \\
\pzf_k(B_\alpha) \arrow[rr] \arrow[rrrr, bend right] &  & \pzf_k(B_\beta) \arrow[rr] &  & \pzf_k(B_\gamma) & 
\end{tikzcd}
\end{center}
It commutes because of the cocycle condition.

Then for each equivalence class $[B_{i,\alpha}]\in \omega$, we can define $\pzf([B_{i,\alpha}])=\pzf_i(B_{i,\alpha})/\sim$, where the equivalence relation is just the isomorphism induced by $\phi_{ij}$. The restriction maps are also induced. 



It obviously forms a presheaf on the base $\omega$.

Base identity axiom:
For open sets $B$ in the base, there is $U_i$ containing it. $B=\cup_{i,\alpha\in A'_i}[B_{i,\alpha}]$. The base identity axiom holds because $\pzf_i$ is a sheaf on base of $U_i$. 

Base gluability holds similarly.

By previous theorem and exercises, there is a unique sheaf $\calf$ extending the sheaf on base $\pzf$.

Then it remains to check
$\calf|_{U_i}\cong\calf_i$. Notice that $\calf|_{U_i}([B_{i,\alpha}])=\calf([B_{i,\alpha}])=\pzf_i(B_{i,\alpha})=\calf_i(B_{i,\alpha})$. Recall the forgetful functor $\iota$ from category of sheaves to category of sheaves on a base, $\iota(\calf|_{U_i})\cong \pzf_i=\iota(\calf_i)$. Then $\calf|_{U_i}\cong \calf_i$ due to  exercise~\ref{chap2exr:equivalence_of_sheaves_on_base_and_sheaves}. (category of sheafs on $X$ is equivalent to the category of sheaves on a base of $X$)
\end{proof}


\begin{exr}\label{chap2exr:surjectivity_morphism_of_sheaves_checked_on_base}
Suppose a morphism of sheaves $\pzf\lrta\pzg$ on a base $\{B_i\}$ is surjective for all $B_i$ (i.e., $\pzf(B_i) \lrta \pzg(B_i)$ is surjective for all $i$). Show that the corresponding morphism of sheaves (not on the base) is surjective (or more precisely: an epimorphism).
\end{exr}
\begin{proof}
Recall~\ref{chap2:exr.surj_stalk_epimorphism_sheaf}, we can prove that surjective at the level of a base implies surjective at the level of stalks, which is true.
\end{proof}




\section{Sheaves of Abelian group, and $\calo_X$-modules, form Abelian categories}
\begin{exr}
Show that the stalk of the kernel is the kernel of the stalks: for all$ p \in  X$, there is a natural isomorphism
$$
(\ker(\calf\lrta \calg))_p\cong \ker (\calf_p\lrta \calg_p)
$$
\end{exr}
\begin{proof}
In exercise~\ref{chap2exr:kernel_sheaf}, we showed that the kernel of morphism of sheaves of abelian group is indeed a sheaf. The punchline is \textbf{filtered colimits preserve finite limits} (see \cite[p.~216]{mac1998categories}). There is no simple derivation using only the universal property. We have to prove it by hand. 

Consider the sheaf morphism $\phi:\calf\lrta \calg$.
Assume $[(f,U)]_p\in \ker(\calf\lrta\calg)_p$, where $(f,U)$ is a representative so that $f\in \ker (\phi(U))\in \calf(U)$.  Also $(\phi(U)f)_p=\phi_pf_p=0$, which means $f_p\in \ker(\phi_p)$. This gives a morphism from $\ker(\phi)_p$ to $\ker(\phi_p)$.

Then we construct the inverse. Given $f_p\in \ker(\phi_p)$, $\phi_pf_p=0$. Notice this means $\phi_p[(f';V)]=[(\phi(V)f';V)]=[(0;V)]$. Then there exists an open $W\ni p$ such that $\phi(W)f'|_W=0$. There is a representative of $f_p$ such that $f'|_W\in\ker \phi(W)$. This would mean $f'_p\in\ker(\phi)_p$, and we get a morphism from $\ker(\phi_p)$ to $\ker(\phi_p)$. Then it is routine to check that the two morphisms are mutually inverse.
\end{proof}

\begin{exr}
Show that the stalk of the cokernel is naturally isomorphic to the cokernel of the stalk.
\end{exr}
\begin{proof}
A cokernel in category of sheaves is the sheafification of cokernel in the category of presheaves. By Exercise~\ref{chap2:exr_shfification_iso_stalks}, sheafification induces isomorphism on stalks. then we only need to verify that the stalk of cokernel presheaves is isomorphic to cokernel of the stalks. But notice that 
$$
(\coker_{pre}(\phi))(U)=\coker(\phi(U)).
$$
Then because colimit commute with colimit we have 
$$
\coker_{pre}(\phi)_p\cong \coker(\phi_p).
$$
\end{proof}

\begin{exr}
Suppose $\phi:\calf\lrta \calg$ is a morphism of sheaves of Abelian groups. Show that the image sheaf $\im \phi$ is the sheafification of the image presheaf. Show that the stalk of the image is the image of the stalk.
\end{exr}
\begin{proof}
Recall the definition of images in Abelian category.

The image of a morphism $\phi : \calf \lrta \calg$ is defined as $\im(\phi) = \ker(\coker \phi)$ whenever it exists.

In the category of presheaves of Abelian groups, the kernel presheaves and cokernel presheaves do exists. And we already know $\coker(\phi)=\coker_{pre}(\phi)^{sh}$. We want to prove
$$
\ker_{pre}(\coker_{pre}(\phi))^{sh}=\ker(\coker_{pre}(\phi)^{sh})
$$
The punchline is \textbf{Sheafification commutes with finite limits} and ultimately, it relies on \textbf{filtered colimit commutes with finite limit}

\end{proof}
\textbf{As a consequnce, the exactness of a sequence of sheaves may be checked art the level of stalks.}
\begin{remark}
There is no direct simple minded proof based on abstract nonsense because it is not tautology. In the above exercise, sheafification commutes with finite limit in the category we care but its not tautology. If we want to have a try: We have two options, one is to check that $\ker_{pre}(\coker_{pre}(\phi))^{sh}$ satisfies the universal property of $\ker(\coker(\phi))$, another is to check that $\ker(\coker_{pre}(\phi)^{sh})$ satisfies the universal property of sheafification. For example, we try to prove that the sheafification has the universal property of being a kernel.
\begin{center}
\begin{tikzcd}
 & \calh \arrow[rrr, dashed] \arrow[ddd, dashed] &  &  & \calg \arrow[ddd, "\pi"'] \arrow[ldd, "\pi_{pre}"] \\
Z \arrow[rrrru] \arrow[rdd] \arrow[rr, dashed, "does\ not\ exists"'] &  & \calh_{pre} \arrow[ldd] \arrow[rru] \arrow[lu, "sh"] &  &  \\
 &  &  & \coker_{pre}(\phi) \arrow[rd, "sh"] &  \\
 & 0 \arrow[rrr] \arrow[rru] &  &  & \coker(\phi)
\end{tikzcd}
\end{center}
We denote $\calh_{pre}$ as $\ker(\coker_{pre}(\phi))$ and $\calh$ as its sheafification.
$Z\lrta \calg$ compose with $\calg\lrta \coker(\phi)$ to get $0$.
What we need is an arrow from $Z$ to $\calh$, but the universal property of sheafification at best can give us an arrow from $\calh$ to $Z$ and nothing else. We can neither construct the morphism from $Z$ to $\calh_{pre}$ because it would imply $\calh_{pre}=\calh$, which 
\end{remark}
\begin{exr}
Show that taking the stalk of a sheaf of abelian groups is an exact functor. More precsely, if $X$ is a toplogical space and $p\in X$ is a point, show that taking the stalk at $p$ defines an exact functor $Ab_X\lrta Ab$.
\end{exr}

\begin{exr}
Check that the exponential exact sequence is exact.
\end{exr}

\begin{exr}\label{chap2exr:left_exact_global_section}
Suppose $U\subset X$ is an open set, and $0\lrta \calf\lrta \calg\lrta \calh$ is an exact sequence of sheaves of abelian groups. Show that 
$$
0\lrta \calf(U)\lrta\calg(U)\lrta \calh(U)
$$
is exact. Show that the section functor need not be exact: show that if $0\lrta \calf\lrta \calg\lrta \calh\lrta 0$ is an exact sequence of sheaves of abelian groups, then
$$
0\lrta\calf(U)\lrta\calg(U)\lrta\calh(U)\lrta 0
$$
need not be exact.
\end{exr}

\begin{exr}
Suppose $\calf$ is a sheaf of abelian groups on a topological space $X$. Show that $\phom(\calf,\cot) $ is a left-exct covariant functor $Ab_X\lrta Ab_X$. Show that $\phom(\cdot,\calf)$ is a left excat contravaiant functor $Ab_X\lrta Ab_X$
\end{exr}

\begin{exr}
Show that given $(X,\calo_X)$ is a ringed space, then $\calo_X$-modules form an abelian category. 
\end{exr}
\begin{exr}\label{chap2exr:sheaf_section_functor_left_excact}
\ \begin{enumerate}[label=(\alph*)]
\item Suppose $\calo_X$ is a sheaf of rings on $X$. Define (categorically) what we should mean by \textbf{ tensor product of two $\calo_X$-modules.} Given an explicit construction, and show that it satisfies your categorical definition.
\item Show that the tensor product of stlks is the stalk of the tensor product. (colimits commute with tensor products)
\end{enumerate}
\end{exr}
\section{The inverse image sheaf}
\begin{exr}
Show that the map
$$
\pi^{-1}_{pre}\calg(U)=\colim_{}V\supset \pi(U)\calg(V).
$$
 defines a presheaf on $X$. Show that it needn’t form a sheaf. 
\end{exr}

\begin{exr}\label{chap2exr:adjoint_inverse_image_direct_image}
If $\pi:X\lrta Y$ is a continuous map, and $\calf$ is a shef on $X$ and $\calg$ is a sheaf on $Y$. describe a bijection 
$$
\mor(\pi^{-1}\calg,\calf)\llrta\mor(\calg,\pi_*\calf).
$$
\end{exr}

\begin{exr}
Show that the stalks of $\pi^{-1}\calg$ are the same as the stalks of $\calg$. More precisely, if $\pi(p)=q$, describe a natural isomorphism $\calg_q\lrta (\pi^{-1}\calg)_p$.
\end{exr}

\begin{exr}
If $U$ is an open subset of $Y$, $i:U\lrta Y$ is the inclusion, and $\calg$ is a sheaf on $Y$, show that $i^{-1}Y$ is naturally isomorphic to $\calg|_U$
\end{exr}




%------------------------------------------------





%------------------------------------------------


%------------------------------------------------


%------------------------------------------------



%----------------------------------------------------------------------------------------
%	PART
%----------------------------------------------------------------------------------------

\part{Schemes}

%----------------------------------------------------------------------------------------
%	CHAPTER 3-5
%----------------------------------------------------------------------------------------
\chapterimage{scheme.jpg} % Chapter heading image
\chapter{Toward affine schemes: the underlying set, and topological space}
%\documentclass[11pt,fleqn]{book}
%%%%%%%%%%%%%%%%%%%%%%%%%%%%%%%%%%%%%%%%%%
% The Legrand Orange Book
% Structural Definitions File
% Version 2.0 (9/2/15)
%
% Original author:
% Mathias Legrand (legrand.mathias@gmail.com) with modifications by:
% Vel (vel@latextemplates.com)
% 
% This file has been downloaded from:
% http://www.LaTeXTemplates.com
%
% License:
% CC BY-NC-SA 3.0 (http://creativecommons.org/licenses/by-nc-sa/3.0/)
%
%%%%%%%%%%%%%%%%%%%%%%%%%%%%%%%%%%%%%%%%%

%----------------------------------------------------------------------------------------
%	VARIOUS REQUIRED PACKAGES AND CONFIGURATIONS
%----------------------------------------------------------------------------------------

\usepackage[top=3cm,bottom=3cm,left=3cm,right=3cm,headsep=10pt,a4paper]{geometry} % Page margins

\usepackage{graphicx} % Required for including pictures
\graphicspath{{Pictures/}} % Specifies the directory where pictures are stored

\usepackage{lipsum} % Inserts dummy text
\usepackage{yhmath}
\usepackage{tikz} % Required for drawing custom shapes

\usepackage{tikz-cd}
\tikzset{
  subseteq/.style={
    draw=none,
    edge node={node [sloped, allow upside down, auto=false]{$\subseteq$}}},
  Subseteq/.style={
    draw=none,
    every to/.append style={
      edge node={node [sloped, allow upside down, auto=false]{$\subseteq$}}}
    },
    Subsetneq/.style={
    draw=none,
    every to/.append style={
      edge node={node [sloped, allow upside down, auto=false]{$\subsetneq$}}}
    },
  Supseteq/.style={
    draw=none,
    every to/.append style={
      edge node={node [sloped, allow upside down, auto=false]{$\supseteq$}}}
  }
}

\usepackage[english]{babel} % English language/hyphenation

\usepackage{enumitem} % Customize lists
\setlist{nolistsep} % Reduce spacing between bullet points and numbered lists

\usepackage{booktabs} % Required for nicer horizontal rules in tables

\usepackage{xcolor} % Required for specifying colors by name
\definecolor{ocre}{RGB}{51, 153, 255} % Define the orange color used for highlighting throughout the book

%----------------------------------------------------------------------------------------
%	FONTS
%----------------------------------------------------------------------------------------


\usepackage{avant} % Use the Avantgarde font for headings
%\usepackage{times} % Use the Times font for headings
\DeclareMathAlphabet{\mathpzc}{OT1}{pzc}{m}{it}
\usepackage{mathptmx} % Use the Adobe Times Roman as the default text font together with math symbols from the Sym­bol, Chancery and Com­puter Modern fonts

\usepackage{microtype} % Slightly tweak font spacing for aesthetics
\usepackage[utf8]{inputenc} % Required for including letters with accents
\usepackage[T1]{fontenc} % Use 8-bit encoding that has 256 glyphs

%----------------------------------------------------------------------------------------
%	BIBLIOGRAPHY AND INDEX
%----------------------------------------------------------------------------------------

\usepackage[style=numeric,citestyle=numeric,sorting=nyt,sortcites=true,autopunct=true,babel=hyphen,hyperref=true,abbreviate=false,backref=true,backend=biber]{biblatex}
\addbibresource{bibliography.bib} % BibTeX bibliography file
\defbibheading{bibempty}{}

\usepackage{calc} % For simpler calculation - used for spacing the index letter headings correctly
\usepackage{makeidx} % Required to make an index
\makeindex % Tells LaTeX to create the files required for indexing

%----------------------------------------------------------------------------------------
%	MAIN TABLE OF CONTENTS
%----------------------------------------------------------------------------------------

\usepackage{titletoc} % Required for manipulating the table of contents

\contentsmargin{0cm} % Removes the default margin

% Part text styling
\titlecontents{part}[0cm]
{\addvspace{20pt}\centering\large\bfseries}
{}
{}
{}

% Chapter text styling
\titlecontents{chapter}[1.25cm] % Indentation
{\addvspace{12pt}\large\sffamily\bfseries} % Spacing and font options for chapters
{\color{ocre!60}\contentslabel[\Large\thecontentslabel]{1.25cm}\color{ocre}} % Chapter number
{\color{ocre}}  
{\color{ocre!60}\normalsize\;\titlerule*[.5pc]{.}\;\thecontentspage} % Page number

% Section text styling
\titlecontents{section}[1.25cm] % Indentation
{\addvspace{3pt}\sffamily\bfseries} % Spacing and font options for sections
{\contentslabel[\thecontentslabel]{1.25cm}} % Section number
{}
{\hfill\color{black}\thecontentspage} % Page number
[]

% Subsection text styling
\titlecontents{subsection}[1.25cm] % Indentation
{\addvspace{1pt}\sffamily\small} % Spacing and font options for subsections
{\contentslabel[\thecontentslabel]{1.25cm}} % Subsection number
{}
{\ \titlerule*[.5pc]{.}\;\thecontentspage} % Page number
[]

% List of figures
\titlecontents{figure}[0em]
{\addvspace{-5pt}\sffamily}
{\thecontentslabel\hspace*{1em}}
{}
{\ \titlerule*[.5pc]{.}\;\thecontentspage}
[]

% List of tables
\titlecontents{table}[0em]
{\addvspace{-5pt}\sffamily}
{\thecontentslabel\hspace*{1em}}
{}
{\ \titlerule*[.5pc]{.}\;\thecontentspage}
[]

%----------------------------------------------------------------------------------------
%	MINI TABLE OF CONTENTS IN PART HEADS
%----------------------------------------------------------------------------------------

% Chapter text styling
\titlecontents{lchapter}[0em] % Indenting
{\addvspace{15pt}\large\sffamily\bfseries} % Spacing and font options for chapters
{\color{ocre}\contentslabel[\Large\thecontentslabel]{1.25cm}\color{ocre}} % Chapter number
{}  
{\color{ocre}\normalsize\sffamily\bfseries\;\titlerule*[.5pc]{.}\;\thecontentspage} % Page number

% Section text styling
\titlecontents{lsection}[0em] % Indenting
{\sffamily\small} % Spacing and font options for sections
{\contentslabel[\thecontentslabel]{1.25cm}} % Section number
{}
{}

% Subsection text styling
\titlecontents{lsubsection}[.5em] % Indentation
{\normalfont\footnotesize\sffamily} % Font settings
{}
{}
{}

%----------------------------------------------------------------------------------------
%	PAGE HEADERS
%----------------------------------------------------------------------------------------

\usepackage{fancyhdr} % Required for header and footer configuration

\pagestyle{fancy}
\renewcommand{\chaptermark}[1]{\markboth{\sffamily\normalsize\bfseries\chaptername\ \thechapter.\ #1}{}} % Chapter text font settings
\renewcommand{\sectionmark}[1]{\markright{\sffamily\normalsize\thesection\hspace{5pt}#1}{}} % Section text font settings
\fancyhf{} \fancyhead[LE,RO]{\sffamily\normalsize\thepage} % Font setting for the page number in the header
\fancyhead[LO]{\rightmark} % Print the nearest section name on the left side of odd pages
\fancyhead[RE]{\leftmark} % Print the current chapter name on the right side of even pages
\renewcommand{\headrulewidth}{0.5pt} % Width of the rule under the header
\addtolength{\headheight}{2.5pt} % Increase the spacing around the header slightly
\renewcommand{\footrulewidth}{0pt} % Removes the rule in the footer
\fancypagestyle{plain}{\fancyhead{}\renewcommand{\headrulewidth}{0pt}} % Style for when a plain pagestyle is specified

% Removes the header from odd empty pages at the end of chapters
\makeatletter
\renewcommand{\cleardoublepage}{
\clearpage\ifodd\c@page\else
\hbox{}
\vspace*{\fill}
\thispagestyle{empty}
\newpage
\fi}

%----------------------------------------------------------------------------------------
%	THEOREM STYLES
%----------------------------------------------------------------------------------------
\usepackage{stmaryrd}
\usepackage{amsmath,amsfonts,amssymb,amsthm} % For math equations, theorems, symbols, etc


\newcommand{\intoo}[2]{\mathopen{]}#1\,;#2\mathclose{[}}
\newcommand{\ud}{\mathop{\mathrm{{}d}}\mathopen{}}
\newcommand{\intff}[2]{\mathopen{[}#1\,;#2\mathclose{]}}
\newtheorem{notation}{Notation}[chapter]

% Boxed/framed environments
\newtheoremstyle{ocrenumbox}% % Theorem style name
{0pt}% Space above
{0pt}% Space below
{\normalfont}% % Body font
{}% Indent amount
{\small\bf\sffamily\color{ocre}}% % Theorem head font
{\;}% Punctuation after theorem head
{0.25em}% Space after theorem head
{\small\sffamily\color{ocre}\thmname{#1}\nobreakspace\thmnumber{\@ifnotempty{#1}{}\@upn{#2}}% Theorem text (e.g. Theorem 2.1)
\thmnote{\nobreakspace\the\thm@notefont\sffamily\bfseries\color{black}---\nobreakspace#3.}} % Optional theorem note

\DeclareRobustCommand{\coprod}{\mathop{\text{\fakecoprod}}}
\newcommand{\fakecoprod}{%
  \sbox0{$\prod$}%
  \smash{\raisebox{\dimexpr.9625\depth-\dp0}{\scalebox{1}[-1]{$\prod$}}}%
  \vphantom{$\prod$}%
}

%\makeatletter
\newcommand{\colim@}[2]{%
  \vtop{\m@th\ialign{##\cr
    \hfil$#1\operator@font colim$\hfil\cr
    \noalign{\nointerlineskip\kern1.5\ex@}#2\cr
    \noalign{\nointerlineskip\kern-\ex@}\cr}}%
}
\newcommand{\colim}{%
  \mathop{\mathpalette\colim@{\rightarrowfill@\scriptscriptstyle}}\nmlimits@
}
\renewcommand{\varprojlim}{%
  \mathop{\mathpalette\varlim@{\leftarrowfill@\scriptscriptstyle}}\nmlimits@
}

\renewcommand{\varinjlim}{%
  \mathop{\mathpalette\varlim@{\rightarrowfill@\scriptscriptstyle}}\nmlimits@
}
\newcommand{\dirlim}{\varinjlim}
\newcommand{\invlim}{\varprojlim}

\renewcommand{\qedsymbol}{$\blacksquare$}% Optional qed square
\renewcommand{\baselinestretch}{1.1}
\newcommand{\affn}{\mathbb A}
\newcommand{\reals}{\mathbb R}
\newcommand{\cplx}{\mathbb C}
\newcommand{\proj}{\mathbb P}
\newcommand{\intg}{\mathbb Z}
\newcommand{\bbf}{\mathbb F}
\newcommand{\ratl}{\mathbb Q}
\newcommand{\torus}{\mathbb T}
\newcommand{\sca}{{\mathfrak a}}
\newcommand{\scb}{{\mathfrak b}}
\newcommand{\scc}{{\mathfrak c}}
\newcommand{\scm}{{\mathfrak m}}
\newcommand{\scn}{{\mathfrak n}}
\newcommand{\scp}{{\mathfrak p}}
\newcommand{\scq}{\mathfrak q}
\newcommand{\frakg}{{\mathfrak g}}
\newcommand{\frakd}{{\mathfrak d}}
\newcommand{\cala}{{\cal A}}
\newcommand{\calb}{{\cal B}}
\newcommand{\calc}{{\cal C}}
\newcommand{\cale}{{\cal E}}
\newcommand{\calf}{{\cal F}}
\newcommand{\calg}{{\cal G}}
\newcommand{\calh}{{\cal H}}
\newcommand{\cali}{{\cal I}}
\newcommand{\call}{{\cal L}}
\newcommand{\caln}{{\cal N}}
\newcommand{\calo}{{\cal O}}
\newcommand{\pd}{{\partial}}
\newcommand{\calr}{{\cal R}}
\newcommand{\mathbold}{\bf}
\newcommand{\Proj}{\text{ Proj\,}}
\newcommand{\res}{\text{res}}
\newcommand{\sym}{\text{ Sym\,}}
\newcommand{\spec}{\text{ Spec\,}}
\newcommand{\tor}{\text{ Tor\,}}
\newcommand{\ext}{\text{ Ext\,}}
\renewcommand{\char}{\text{char\,}}
\newcommand{\mor}{\text{ Mor}}
\renewcommand{\hom}{\text{ Hom}}
\newcommand{\phom}{\mathpzc{Hom\,}}
\newcommand{\pzf}{{\mathpzc F}}
\newcommand{\pzg}{{\mathpzc G}}
\newcommand{\pzm}{{\mathpzc M}}
\newcommand{\dims}{\text{dim\,}}
\newcommand{\coker}{\text{coker\,}}
\newcommand{\im}{{\text im\thinspace}}
\newcommand{\supp}{{\text Supp\thinspace}}
\newcommand{\cinf}{C^{\infty}}
\newcommand{\row}[2]{#1_1,\dots ,#1_{#2}}
\newcommand{\dbyd}[2]{{\partial #1\over\partial #2}}
\newcommand{\Space}{{\bf Space}}
\newcommand{\alg}{{\mathbold Alg}}
\newcommand{\notsubset}{\not \subset}
\newcommand{\notsupset}{\not \supset}
\newcommand{\pois}{{\mathbold Pois}}
\newcommand{\pitilde}{\tilde{\pi}}
\newcommand{\rta}{\rightarrow}
\newcommand{\Lrta}{\Longrightarrow}
\newcommand{\lrta}{\longrightarrow}
\newcommand{\llta}{\longleftarrow}
\newcommand{\llrta}{\longleftrightarrow}
\newcommand{\Llta}{\Longleftarrow}
\newcommand{\Llrta}{\Longleftrightarrow}
\newcommand{\lgl}{\langle}
\newcommand{\rgl}{\rangle}
\newcommand{\inj}{\hookrightarrow}
\newcommand{\Spec}{{\text Spec}}
\newcommand{\rad}{{\text rad}}
\newcommand{\Jac}{{\text Jac}}
\newcommand{\surj}{\twoheadrightarrow}
\newcommand{\cmark}{\ding{51}}%
\newcommand{\xmark}{\ding{55}}%
\newcommand{\downmapsto}{\rotatebox[origin=c]{-90}{$\scriptstyle\mapsto$}\mkern2mu}

\newtheoremstyle{blacknumex}% Theorem style name
{5pt}% Space above
{5pt}% Space below
{\normalfont}% Body font
{} % Indent amount
{\small\bf\sffamily}% Theorem head font
{\;}% Punctuation after theorem head
{0.25em}% Space after theorem head
{\small\sffamily{\tiny\ensuremath{\blacksquare}}\nobreakspace\thmname{#1}\nobreakspace\thmnumber{\@ifnotempty{#1}{}\@upn{#2}}% Theorem text (e.g. Theorem 2.1)
\thmnote{\nobreakspace\the\thm@notefont\sffamily\bfseries---\nobreakspace#3.}}% Optional theorem note

\newtheoremstyle{blacknumbox} % Theorem style name
{0pt}% Space above
{0pt}% Space below
{\normalfont}% Body font
{}% Indent amount
{\small\bf\sffamily}% Theorem head font
{\;}% Punctuation after theorem head
{0.25em}% Space after theorem head
{\small\sffamily\thmname{#1}\nobreakspace\thmnumber{\@ifnotempty{#1}{}\@upn{#2}}% Theorem text (e.g. Theorem 2.1)
\thmnote{\nobreakspace\the\thm@notefont\sffamily\bfseries---\nobreakspace#3.}}% Optional theorem note

% Non-boxed/non-framed environments
\newtheoremstyle{ocrenum}% % Theorem style name
{5pt}% Space above
{5pt}% Space below
{\normalfont}% % Body font
{}% Indent amount
{\small\bf\sffamily\color{ocre}}% % Theorem head font
{\;}% Punctuation after theorem head
{0.25em}% Space after theorem head
{\small\sffamily\color{ocre}\thmname{#1}\nobreakspace\thmnumber{\@ifnotempty{#1}{}\@upn{#2}}% Theorem text (e.g. Theorem 2.1)
\thmnote{\nobreakspace\the\thm@notefont\sffamily\bfseries\color{black}---\nobreakspace#3.}} % Optional theorem note
\renewcommand{\qedsymbol}{$\blacksquare$}% Optional qed square
\makeatother

% Defines the theorem text style for each type of theorem to one of the three styles above
\newcounter{dummy} 
\numberwithin{dummy}{section}
\theoremstyle{ocrenumbox}
\newtheorem{theoremeT}[dummy]{Theorem}
\newtheorem{problem}{Problem}[chapter]
\newtheorem{exerciseT}{Exercise??}[section]
\theoremstyle{blacknumex}
\newtheorem{exampleT}{Example}[section]
\theoremstyle{blacknumbox}
\newtheorem{vocabulary}{Vocabulary}[chapter]
\newtheorem{definitionT}{Definition}[section]
\newtheorem{corollaryT}[dummy]{Corollary}
\theoremstyle{ocrenum}
\newtheorem{proposition}[dummy]{Proposition}
\newtheorem{lemma}[dummy]{Lemma}

\renewcommand*{\theexerciseT}{\thesection.\Alph{exerciseT}}
%----------------------------------------------------------------------------------------
%	DEFINITION OF COLORED BOXES
%----------------------------------------------------------------------------------------

\RequirePackage[framemethod=default]{mdframed} % Required for creating the theorem, definition, exercise and corollary boxes

% Theorem box
\newmdenv[skipabove=7pt,
skipbelow=7pt,
backgroundcolor=black!5,
linecolor=ocre,
innerleftmargin=5pt,
innerrightmargin=5pt,
innertopmargin=5pt,
leftmargin=0cm,
rightmargin=0cm,
innerbottommargin=5pt]{tBox}

% Exercise box	  
\newmdenv[skipabove=7pt,
skipbelow=7pt,
rightline=false,
leftline=true,
topline=false,
bottomline=false,
backgroundcolor=ocre!10,
linecolor=ocre,
innerleftmargin=5pt,
innerrightmargin=5pt,
innertopmargin=5pt,
innerbottommargin=5pt,
leftmargin=0cm,
rightmargin=0cm,
linewidth=4pt]{eBox}	

% Definition box
\newmdenv[skipabove=7pt,
skipbelow=7pt,
rightline=false,
leftline=true,
topline=false,
bottomline=false,
linecolor=ocre,
innerleftmargin=5pt,
innerrightmargin=5pt,
innertopmargin=0pt,
leftmargin=0cm,
rightmargin=0cm,
linewidth=4pt,
innerbottommargin=0pt]{dBox}	

% Corollary box
\newmdenv[skipabove=7pt,
skipbelow=7pt,
rightline=false,
leftline=true,
topline=false,
bottomline=false,
linecolor=gray,
backgroundcolor=black!5,
innerleftmargin=5pt,
innerrightmargin=5pt,
innertopmargin=5pt,
leftmargin=0cm,
rightmargin=0cm,
linewidth=4pt,
innerbottommargin=5pt]{cBox}

% Creates an environment for each type of theorem and assigns it a theorem text style from the "Theorem Styles" section above and a colored box from above
\newenvironment{theorem}{\begin{tBox}\begin{theoremeT}}{\end{theoremeT}\end{tBox}}
\newenvironment{exr}{\begin{eBox}\begin{exerciseT}}{\hfill{\color{ocre}\tiny\ensuremath{\blacksquare}}\end{exerciseT}\end{eBox}}				  
\newenvironment{definition}{\begin{dBox}\begin{definitionT}}{\end{definitionT}\end{dBox}}	
\newenvironment{example}{\begin{exampleT}}{\hfill{\tiny\ensuremath{\blacksquare}}\end{exampleT}}		
\newenvironment{corollary}{\begin{cBox}\begin{corollaryT}}{\end{corollaryT}\end{cBox}}	

%----------------------------------------------------------------------------------------
%	REMARK ENVIRONMENT
%----------------------------------------------------------------------------------------

\newenvironment{remark}{\par\vspace{10pt}\small % Vertical white space above the remark and smaller font size
\begin{list}{}{
\leftmargin=35pt % Indentation on the left
\rightmargin=25pt}\item\ignorespaces % Indentation on the right
\makebox[-2.5pt]{\begin{tikzpicture}[overlay]
\node[draw=ocre!60,line width=1pt,circle,fill=ocre!25,font=\sffamily\bfseries,inner sep=2pt,outer sep=0pt] at (-15pt,0pt){\textcolor{ocre}{R}};\end{tikzpicture}} % Orange R in a circle
\advance\baselineskip -1pt}{\end{list}\vskip5pt} % Tighter line spacing and white space after remark

%----------------------------------------------------------------------------------------
%	SECTION NUMBERING IN THE MARGIN
%----------------------------------------------------------------------------------------

\makeatletter
\renewcommand{\@seccntformat}[1]{\llap{\textcolor{ocre}{\csname the#1\endcsname}\hspace{1em}}}                    
\renewcommand{\section}{\@startsection{section}{1}{\z@}
{-4ex \@plus -1ex \@minus -.4ex}
{1ex \@plus.2ex }
{\normalfont\large\sffamily\bfseries}}
\renewcommand{\subsection}{\@startsection {subsection}{2}{\z@}
{-3ex \@plus -0.1ex \@minus -.4ex}
{0.5ex \@plus.2ex }
{\normalfont\sffamily\bfseries}}
\renewcommand{\subsubsection}{\@startsection {subsubsection}{3}{\z@}
{-2ex \@plus -0.1ex \@minus -.2ex}
{.2ex \@plus.2ex }
{\normalfont\small\sffamily\bfseries}}                        
\renewcommand\paragraph{\@startsection{paragraph}{4}{\z@}
{-2ex \@plus-.2ex \@minus .2ex}
{.1ex}
{\normalfont\small\sffamily\bfseries}}

%----------------------------------------------------------------------------------------
%	PART HEADINGS
%----------------------------------------------------------------------------------------

% numbered part in the table of contents
\newcommand{\@mypartnumtocformat}[2]{%
\setlength\fboxsep{0pt}%
\noindent\colorbox{ocre!20}{\strut\parbox[c][.7cm]{\ecart}{\color{ocre!70}\Large\sffamily\bfseries\centering#1}}\hskip\esp\colorbox{ocre!40}{\strut\parbox[c][.7cm]{\linewidth-\ecart-\esp}{\Large\sffamily\centering#2}}}%
%%%%%%%%%%%%%%%%%%%%%%%%%%%%%%%%%%
% unnumbered part in the table of contents
\newcommand{\@myparttocformat}[1]{%
\setlength\fboxsep{0pt}%
\noindent\colorbox{ocre!40}{\strut\parbox[c][.7cm]{\linewidth}{\Large\sffamily\centering#1}}}%
%%%%%%%%%%%%%%%%%%%%%%%%%%%%%%%%%%
\newlength\esp
\setlength\esp{4pt}
\newlength\ecart
\setlength\ecart{1.2cm-\esp}
\newcommand{\thepartimage}{}%
\newcommand{\partimage}[1]{\renewcommand{\thepartimage}{#1}}%
\def\@part[#1]#2{%
\ifnum \c@secnumdepth >-2\relax%
\refstepcounter{part}%
\addcontentsline{toc}{part}{\texorpdfstring{\protect\@mypartnumtocformat{\thepart}{#1}}{\partname~\thepart\ ---\ #1}}
\else%
\addcontentsline{toc}{part}{\texorpdfstring{\protect\@myparttocformat{#1}}{#1}}%
\fi%
\startcontents%
\markboth{}{}%
{\thispagestyle{empty}%
\begin{tikzpicture}[remember picture,overlay]%
\node at (current page.north west){\begin{tikzpicture}[remember picture,overlay]%	
\fill[ocre!20](0cm,0cm) rectangle (\paperwidth,-\paperheight);
\node[anchor=north] at (4cm,-3.25cm){\color{ocre!40}\fontsize{220}{100}\sffamily\bfseries\thepart}; 
\node[anchor=south east] at (\paperwidth-1cm,-\paperheight+1cm){\parbox[t][][t]{8.5cm}{
\printcontents{l}{0}{\setcounter{tocdepth}{1}}%
}};
\node[anchor=north east] at (\paperwidth-1.5cm,-3.25cm){\parbox[t][][t]{15cm}{\strut\raggedleft\color{white}\fontsize{30}{30}\sffamily\bfseries#2}};
\end{tikzpicture}};
\end{tikzpicture}}%
\@endpart}
\def\@spart#1{%
\startcontents%
\phantomsection
{\thispagestyle{empty}%
\begin{tikzpicture}[remember picture,overlay]%
\node at (current page.north west){\begin{tikzpicture}[remember picture,overlay]%	
\fill[ocre!20](0cm,0cm) rectangle (\paperwidth,-\paperheight);
\node[anchor=north east] at (\paperwidth-1.5cm,-3.25cm){\parbox[t][][t]{15cm}{\strut\raggedleft\color{white}\fontsize{30}{30}\sffamily\bfseries#1}};
\end{tikzpicture}};
\end{tikzpicture}}
\addcontentsline{toc}{part}{\texorpdfstring{%
\setlength\fboxsep{0pt}%
\noindent\protect\colorbox{ocre!40}{\strut\protect\parbox[c][.7cm]{\linewidth}{\Large\sffamily\protect\centering #1\quad\mbox{}}}}{#1}}%
\@endpart}
\def\@endpart{\vfil\newpage
\if@twoside
\if@openright
\null
\thispagestyle{empty}%
\newpage
\fi
\fi
\if@tempswa
\twocolumn
\fi}

%----------------------------------------------------------------------------------------
%	CHAPTER HEADINGS
%----------------------------------------------------------------------------------------

% A switch to conditionally include a picture, implemented by  Christian Hupfer
\newif\ifusechapterimage
\usechapterimagetrue
\newcommand{\thechapterimage}{}%
\newcommand{\chapterimage}[1]{\ifusechapterimage\renewcommand{\thechapterimage}{#1}\fi}%
\newcommand{\autodot}{.}
\def\@makechapterhead#1{%
{\parindent \z@ \raggedright \normalfont
\ifnum \c@secnumdepth >\m@ne
\if@mainmatter
\begin{tikzpicture}[remember picture,overlay]
\node at (current page.north west)
{\begin{tikzpicture}[remember picture,overlay]
\node[anchor=north west,inner sep=0pt] at (0,0) {\ifusechapterimage\includegraphics[width=\paperwidth]{\thechapterimage}\fi};
\draw[anchor=west] (\Gm@lmargin,-9cm) node [line width=2pt,rounded corners=15pt,draw=ocre,fill=white,fill opacity=0.5,inner sep=15pt]{\strut\makebox[22cm]{}};
\draw[anchor=west] (\Gm@lmargin+.3cm,-9cm) node {\huge\sffamily\bfseries\color{black}\thechapter\autodot~#1\strut};
\end{tikzpicture}};
\end{tikzpicture}
\else
\begin{tikzpicture}[remember picture,overlay]
\node at (current page.north west)
{\begin{tikzpicture}[remember picture,overlay]
\node[anchor=north west,inner sep=0pt] at (0,0) {\ifusechapterimage\includegraphics[width=\paperwidth]{\thechapterimage}\fi};
\draw[anchor=west] (\Gm@lmargin,-9cm) node [line width=2pt,rounded corners=15pt,draw=ocre,fill=white,fill opacity=0.5,inner sep=15pt]{\strut\makebox[22cm]{}};
\draw[anchor=west] (\Gm@lmargin+.3cm,-9cm) node {\huge\sffamily\bfseries\color{black}#1\strut};
\end{tikzpicture}};
\end{tikzpicture}
\fi\fi\par\vspace*{270\p@}}}

%-------------------------------------------

\def\@makeschapterhead#1{%
\begin{tikzpicture}[remember picture,overlay]
\node at (current page.north west)
{\begin{tikzpicture}[remember picture,overlay]
\node[anchor=north west,inner sep=0pt] at (0,0) {\ifusechapterimage\includegraphics[width=\paperwidth]{\thechapterimage}\fi};
\draw[anchor=west] (\Gm@lmargin,-9cm) node [line width=2pt,rounded corners=15pt,draw=ocre,fill=white,fill opacity=0.5,inner sep=15pt]{\strut\makebox[22cm]{}};
\draw[anchor=west] (\Gm@lmargin+.3cm,-9cm) node {\huge\sffamily\bfseries\color{black}#1\strut};
\end{tikzpicture}};
\end{tikzpicture}
\par\vspace*{270\p@}}
\makeatother

%----------------------------------------------------------------------------------------
%	HYPERLINKS IN THE DOCUMENTS
%----------------------------------------------------------------------------------------

\usepackage{hyperref}
\hypersetup{hidelinks,backref=true,pagebackref=true,hyperindex=true,colorlinks,linkcolor = ocre,breaklinks=true,urlcolor= ocre,bookmarks=true,bookmarksopen=false,pdftitle={Title},pdfauthor={Author}}
\usepackage{bookmark}
\bookmarksetup{
open,
numbered,
addtohook={%
\ifnum\bookmarkget{level}=0 % chapter
\bookmarksetup{bold}%
\fi
\ifnum\bookmarkget{level}=-1 % part
\bookmarksetup{color=ocre,bold}%
\fi
}
}
 
%\begin{document}
\section{Towards schemes}
\begin{exr}
Suppose that $\pi: X\longrightarrow Y$ is a continuous map of differentiable manifolds (as topological spaces — not a priori differentiable). Show that $\pi$ is differentiable if differentiable functions pull back to differentiable functions, i.e., if pullback by $\pi$ gives a map $\calo_Y \longrightarrow \pi_*\calo_X$. (Hint: check this on small patches. Once you figure out what you are trying to show, you will realize that the result is immediate.)
\end{exr}
\begin{proof}
Recall the definition of differentiable maps: At every point $p\in X$, local charts $(U,\varphi)\ni p$ and $(V,\phi)\ni \pi(p)$, where $\varphi: U\lrta \reals^n$ and $\phi:V\lrta \reals^m$ are homeomorphisms. We say $\pi$ is differentiable iff $\phi\circ\pi\circ\varphi^{-1}:\reals^n\lrta \reals^m$ is differentiable. And check that if $\forall g$, $g\circ \pi$ is differentiable at $p$. $g\circ \phi^{-1}\circ(\phi\circ \pi\circ \varphi^{-1})$ is differentiable for all $g$. We can choose $g=\phi_i$ such that $\phi_i\circ \phi^{-1}=pr_i:\reals^m\lrta\reals$ is the projection map to the $i$-th coordinate. Then the map $\phi\circ\pi\varphi^{-1}=(pr_1\circ\phi\circ\pi\varphi^{-1},...,pr_m\circ\phi\circ\pi\varphi^{-1})$ is differentiable.
\end{proof}

\begin{exr}
Show that a morphism of differentiable manifolds $\pi: X\lrta Y$ with $\pi(p)=q$ induces a morphism of stalks $\pi^\#:\calo_{Y,q}\lrta \calo_{X,p}$. Show that $\pi^\# (\scm_{Y,q})\subset \scm_{X,\scp}$ . In other words, if you pull back a function that vanishes at $q$, you get a function that vanishes at $p$. 
\end{exr}
\begin{proof}
$\calo_{X,p}$ is the germ of smooth functions at $p$. $\pi^{\#}:\calo_{Y,q}\lrta \calo_{X,p}; [(f,U)]\longmapsto [(f\circ \pi, \pi^{-1} U)]$. If $[(f,U)]\in \scm_{Y,q}$, $(f,U)\sim (0, V) $. In particulr, there is an open set $W\subseteq U\cap V$, so that $f|_W=0$. $f\circ\pi|_{\pi^{-1}W}=0\Lrta [f\circ\pi, \pi^{-1}W]\in \scm_{X,p}$.
\end{proof}
\section{The underlying sets of affine schemes}
\begin{exr}
A SMALL EXERCISE ABOUT SMALL SCHEMES
\begin{enumerate}[label=(\alph*)]
\item Describe the set $\Spec(k[\epsilon]/(\epsilon^2)) $. The ring $k[\epsilon]/(\epsilon^2)$ is called the ring of \textbf{dual numbers}, and will turn out to be quite useful. You should think of ε as a very small number, so small that its square is $0$ (although it itself is not $0$). It is a nonzero function whose value at all points is zero, thus giving our first example of functions not being determined by their values at points.
\item Describe the set $k[x]_{(x)}$
\end{enumerate}
\end{exr}
\begin{proof}\ 
\begin{enumerate}[label=(\alph*)]
\item Recall the isomorphism theorem, there is a one to one correspondence between \{$\scp\in \Spec(A),\scp\supseteq \sca$\} and \{$\scq\in \Spec(A/\sca)$\}. The only prime in $k[\epsilon]$ that contains $(\epsilon^2)$ is $(\epsilon)$. Hence, $\Spec(k[\epsilon]/(\epsilon^2))$ $=\{[(\epsilon)]\}$
\item  There is a one to one correspondence,
$$
\{\scp\in \Spec(A):\scp\cap S=\emptyset \}\llrta\{\scq\in Spec(S^{-1}A)\}.
$$
In this specific case, the multiplicative set is $S:=k[x]-(x)$. Then the primes in $k[x]_{(x)}$ corresponds to those primes ideals contained in $(x)$, which is just $(x)\cdot k[x]_{(x)},(0)$  

\end{enumerate}
\end{proof}

\begin{exr}
Show that for $\reals[X]/(p(X))\cong\cplx$, where $p(x)$ is an irreducible quadratic polynomial.
\end{exr}
\begin{proof}
For example, consider $p(X)=X^2+aX+b$, where $a^2-4b<0$. Then 
$$
A=\reals[X]/(X^2+aX+b)\cong \{f\in\reals[x]: x^2+ax+b=0\},
$$
hence every element in $A$ can be written as  $cx+d$, with $c,d\in\reals$. However we know $\cplx\cong\reals\oplus\reals i$ as $\reals$-algebra. Consider an $\reals$-linear morphism of vector space 
$$
\varphi:x\longmapsto -\frac{a}{2}+\frac{\sqrt{4b-a^2}}{2} i, 1\longmapsto 1
$$ 
is surjective and injective. Check that $\varphi$ is in fact a well-defined ring morphism. It suffices to check it on basis, we only need to check 
$$
-a\varphi(x)+b=\varphi(x^2)=\varphi(x)\varphi(x).
$$
\end{proof}

\begin{exr}
Describe the set $\affn^1_\ratl=\spec(\ratl[X])$.
\end{exr}
\begin{proof}
Each irreducible polynomial $p(X)\in\ratl[X]$ corresponds to a splitting field. And $\ratl[X]/(p(X))\cong \ratl[\alpha]$ where $\alpha\in\overline{Q}$ s.t. $p(\alpha)=0$. $\spec(\ratl[X])$ contains a generic point $(0)$, closed point corresponding to $(x-q)$, where $q\in\ratl$, and $p(X)$ irreducible, where $(p(X))$ corresponds an equivalent class of algebraic number and their Galois conjugates. 
\end{proof}

\begin{exr}
If $k$ is a field, show that $\spec k[X]$ has infinitely many points. (Hint: Euclid’s proof of the infinitude of primes of $\intg$.)
\end{exr}
\begin{proof}
Polynomial ring over a field is Euclidean domain, which means every polynomial has a unique factorization as product of irreducible polynomials. Assume there are only finitely many prime ideals in $k[X]$, which means there are only finitely many irreducible polynomials $p_1,...,p_n\in k[X]$. We can set $p=p_1\cdot\cdot\cdot p_n+1$. Assume $p$ is reducible, then it must contain some $p_i$, but then we have a contradiction that $p_i|1$. Then $p$ is irreducible and we have infinite prime ideals (points).
\end{proof}

\begin{exr}
Show that we have identified all the prime ideals of $\cplx[x,y]$. (Show that all prime ideals in $\cplx[x,y]$ are either of the form $(0)$, $(x-a,y-b)$,$(f(x,y))$ where $f(x,y)$ is irreducible)
\end{exr}
\begin{proof} $\cplx[x,y]$ is a UFD. Then, a nonzero principal ideal is prime if and only if it is generated by a irreducible element.

Then it suffices to check the  non-principal case. Suppose a prime $\scp\subseteq \cplx[x,y]$ is not principal ideal. We can find at least two elements $f(x,y), g(x,y)$ in $\scp$ with no common factor. We can regard $f(x,y)$ and $g(x,y)$ as polynomials of $y$ with coefficients in $c(x)$.
$$
f(x,y)=f_n(x)y^n+f_{n-1}(x)y^{n-1}+...+f_0
$$
$$
g(x,y)=g_m(x)y^m+g_{m-1}(x)y^{m-1}+...+g_0
$$
Consider the Euclidean algorithm in $\cplx(x)[y]$, we can calculate the greatest common divisor of $f,g$ when regard then as elements in $\cplx(x)[y]$. \\
\underline{Claim}: Let $R$ be a Euclidean domain, and $f,g\in\calr[y]$ such that $gcd(f,g)=1$ (in $R[y]$). They $gcd(f,g)=1$ in $K[y]$, where $K$ is the field of fractions of $R$.
 c
Let $h\in K[y]$ such that $h|f$ and $h|g$ in $K[y]$. We wan to show that $deg (h)=0$.

Let $d$ be the product of all denominators of the coefficients of $h$, and $k=dh\in R[y]$. Then $k|df$ and $k|dg$ in $K[y]$, so there are $a,b\in R\backslash\{0\}$ such that $k|(ad)f$ and $k|bdg$ in $R[y]$. Write $adf=kp$ and $bdg=kq$ with $p,q\in R[y]$.

In the following, one denotes by $c(r)$ the greatest common divisor of the coefficients of $r\in R[y]$, and write $r=c(r)r_1$ where $r_1$ is primitive, that is , the greatest common divisor of its coefficients is $1$.

From $adf=kp$ and $bdg=kq$ we get $ad c(f)=c(k)c(p)$ and $bd c(g)=c(k)c(q)$. But $(ad)c(f)f_1=c(k)c(p)k_1p_1$ and $bd c(g)g_1=c(k)c(q)k_1q_1$, so $f_1=k_1p_1$ and $g_1=k_1q_1$. Thus we get $k_1 |f_1|f$ and $k_1|g_1|g$, so $k_1=1$, and we are done.

$\cplx(x)$ is such a Euclidean domain. $f,g$ has no common factor in $\cplx[x][y]\Lrta gcd(f,g)=1\in \cplx(x)[y]$.
$\exists u,v\in \cplx(x)[y]$ s.t. 
$$
u f+vg=1.
$$
Multiplying the product of denominators of coefficients in $u,v$, we can get an equality in $\cplx[x,y]$
$$
u'f+v'g=h(x).
$$
$h(x)\in (f(x,y),g(x,y))\subseteq\scp$. $h(x)$ can split into product of linear polynomials in $\cplx[x]$, because $\scp$ is prime, at least one the linear factor $x-a$ is in $\scp$. Similarly we can prove that at least one linear factor $y-b\in \scp$. Thus, $(x-a,y-b)\in\scp$, but because $(x-a,y-b)$ is already maximal, we have every non-principal prime ideal is of the form $(x-a,y-b)$.
\end{proof}

\begin{exr}
Show that the Nullstellensatz(Zariski's Lemma) implies the Weak Nullstellensatz.
\end{exr}
\begin{proof} Set $A:=k[X_1,...,X_d]$, with $k$ being a algebraically closed field.
Let $\scm\subseteq {A}$ be any maximal ideal, then $ L={A}/\scm$ is a field.
\begin{center}
\begin{tikzcd}
k\ar[r]\ar[rr,bend right,"j"]&k[X_1,...,X_d]={A}\ar[r,two heads,"q"]& L={A}/\scm  
\end{tikzcd}
\end{center}
Note: $ L$ is a finitely generated $k$-algebra, generated by $q(X_1),...,q(X_d)$
$$
\begin{aligned}
\text{Zariski's Lemma}&\Lrta  L/j(k)\text{ is finite field extension}\\
 & \Lrta
 L\cong k \text{($k$ algebraically closed)}
\end{aligned}
$$
Set $x:=(j^{-1}(q(X_1)),...,j^{-1}(q(X_d)))\in \cplx^d$. Check $\scm=\scm_x:=(X_1-x_1,...,X_d-x_d)$.
We know $j$ is surjective because $q$ is, and $j$ is always injective because $k$ is a field (ring morpihsm of fields are injective). Suppose $P\in \scm\Lrta q(P)=0\Lrta j^{-1}(P(q(X)))=0\Lrta P(j^{-1}(q(X)))=P(x)=0$, hence $\scm_x\in \scm$, but we already know $\scm_x$ is maximal, then $\scm=\scm_x$.
\end{proof}

\begin{exr}\label{exr:chap3_fin-dim_k_domain_is_field}
Any integral domain $A$ which is a finite $k$-algebra (i.e., a $k$-algebra that is a finite-dimensional vector space over $k$) must be a field. 
\end{exr}
\begin{proof}
For a nonzero element $x\in A$, we $xA$ is a finite dimensional subspace of $A$. If $\dim xA\subsetneq \dim A$, then $x^2A\subsetneq xA$ otherwise $\forall z\in A, \exists y \in A xz=x^2 y$, $\Lrta x(z-xy)=0\Lrta z=xy$ because $A$ is an integral domain contradicting to $xA\neq A$. Then we have a descending chain of vector spaces $\cdots x^3A\subsetneq x^2 A\subsetneq xA\subsetneq A$. Because $A$ is finite dimensional, there is an $n\in \intg$ so that $x^nA=0\Lrta x^n=0$ contradiction. Therefore, multiplying $x$ must be an isomorphism.

In particular, the above argument is not true for general finitely generated $k$-algebras. $k[X]$ is not Artinian, there exists an infinitely descending chain of ideals $(X)\supsetneq (X^2)\supsetneq (X^3)\supsetneq ...$ Multiplying $X$ is not an isomorphism on $k[X]$.

Now, assume $A=k[X_1,...,X_d]/\scp$  for some prime ideal $\scp$ so that the residue ring $A$ is finite $k$-algebra. Such $\scp$ exists by Nullstellensatz. We already proved that $0\neq x\in A$ multiplies as an isomorphism, hence $x$ is not contained in any proper ideal in $A$. $\scp$ must be maximal ideal. Hence $A$ is a field.
\end{proof}

\begin{exr}
Describe the maximal ideal of  $\ratl[x,y]$ corresponding to points $(\sqrt 2,\sqrt 2)$
and $(-\sqrt2, -\sqrt 2)$. Describe the maximal ideal of $\ratl[x, y]$ corresponding to points $( \sqrt 2, -\sqrt 2)$ and $(-\sqrt2, \sqrt 2)$. What are the residue fields in each case?
\end{exr}
\begin{proof}
 $(x^2-2,y^2-2)$ is not a maximal ideal, geometrically it consists of two points.

$(\sqrt 2,\sqrt 2)\sim(-\sqrt 2,-\sqrt 2)$ correspond to $(x-y,x^2-2)=(x-y,y^2-2)$.

$(-\sqrt 2,\sqrt 2)\sim(\sqrt 2,-\sqrt 2)$ correspond to $(x+y,x^2-2)=(x+y,y^2-2)$.

The residue field in both case is $\ratl[x,y]/(x-y,y^2-2)=\ratl[y]/[y^2-2]\cong \ratl[\sqrt 2]\cong \ratl[x,y]/(x+y,y^2-2)$. And in particular, the quotient ring of $(x^2-2,y^2-2)$ is
$$
\ratl[x,y]/(x^2-2,y^2-2)\cong \ratl[\sqrt2]\times \ratl[\sqrt 2]
$$

\end{proof}

\begin{exr}
Consider the map of sets $\phi: \cplx^2\lrta \mathbb{A}_\ratl^{2}$ defined as follows. $(z_1,z_2)$ is sent to the prime ideal of $\ratl[x,y]$ consisting of polynomials vanishing at $(z_1,z_2)$.
\begin{enumerate}[label=(\alph*)]
\item What is the image of $(\pi,\pi^2)$
\item Show that $\phi$ is surjective.
\end{enumerate}
\end{exr}
\begin{proof}
\begin{enumerate}[label=(\alph*)]
\item 
There is an injective morphism from $\varphi:\ratl[x,y]\inj \cplx[x,y]$, it induces a morphism on $\spec \cplx[x,y]\lrta \spec \ratl[x,y]$. In particular, the point $(\pi,\pi^2)$ maps as 
$$
(x-\pi,y-\pi^2)\lrta (x-\pi,y-\pi^2)\cap \ratl[x,y]
$$
The problem reduces to find two polynomials $p(x,y),q(x,y)\in\cplx[x,y]$ what is the polynomial
$$
p(x,y)(x-\pi)+q(x,y)(y-\pi^2)\in \ratl[x,y].
$$
We can find 
$$
(x+\pi)(x-\pi)-(y-\pi^2)=x^2-y\in \ratl[x,y]
$$
$(x^2-y)\subseteq (x-\pi,y-\pi^2)\cap\ratl[x,y]$. The tricky part is the reverse inclusion. 

Suppose there is a polynomial 
$$
P(x,y)=\sum_{m,n}a_{m,n}x^{m}y^{n}, a_{m,n}\in \ratl
$$
so that $P(\pi,\pi^2)=0$. $P(\pi,\pi^2)$ is equal to rational polynomial $F$ of $\pi$ alone. 
$$
F(\pi)=\sum_k(\sum_{m+2n=k} a_{m,n})\pi^{k} 
$$
Because $\pi$ is transcendental, each coefficients of $F$ should be zero.
$$
\sum_n a_{k-2n,n}=0
$$
 Then we reorder the summation of $P$
$$
\begin{aligned}
P(x,y)&=\sum_{k}\left(\sum_{n}a_{k-2n,n}x^{k-2n}y^{n}\right)\\
&=\sum_{k}\left(\sum_{n}a_{k-2n,n}x^{k-2n}(y-x^2+x^2)^{n}\right)\\
&=\sum_{k}\left(\sum_{n}a_{k-2n,n}x^{k}\right)+\sum_{k}\left(\sum_{n}a_{k-2n,n}\sum_{j=1}\binom{n}{j} x^{k-2j}(y-x^2)^{j}\right)\\
&=0+\sum_{k}\left(\sum_{n}a_{k-2n,n}\sum_{j=1}\binom{n}{j} x^{k-2j}(y-x^2)^{j}\right)
\end{aligned}
$$

All such $P(x,y)\in (x^2-y)$.

\item $\phi$ is basically $\varphi^*$ for $\varphi:\ratl[x,y]\inj \cplx[x,y]$. 

\underline{Want}: $\varphi^*$ is surjective.

Consider the special case where $\scp$ is maximal. Then $\ratl[x,y]/\scp$ is a finitely generated field extension, hence is finite algebraic field extension by Nullstellensatz. It embeds into $\cplx$, with $a,b$ being the images of $x,y$.
$$
(x-a,y-b)\cap \ratl[x,y]=\scp
$$
$\scp\subseteq (x-a,y-b)\cap \ratl[x,y]$ is trivially true since $\scp\subseteq \varphi_*\varphi^*(\scp)$. And the reverse inclusion holds because $\scp$ is maximal.

For the prime ideal $(0)\subseteq \ratl[x,y]$, we can find two algebraically independent transcendental number $\xi,\delta$ s.t., $p(\xi,\delta)=0\lrta p=0$ given that $p\in \ratl[x,y]$.

For non-zero non-maximal ideal $\scp$, we consider the residue field
$$
K:=Frac\left(\ratl[x,y]/\scp\right)\cong \ratl[x,y]_\scp/\scp\ratl[x,y]_\scp,
$$
where $\scp\ratl[x,y]_\scp$ is the maximal ideal in the local ring $\ratl[x,y]_\scp$ and we have
$$
\scp=\scp\ratl[x,y]_\scp\cap \ratl[x,y]
$$
We have a chain of injections
$$
\ratl[x,y]\inj\ratl[x,y]_\scp\inj K\inj \cplx\inj\cplx[x,y].
$$

By the similar argument, let $\xi,\delta\in \cplx$ be the images of $\frac{x}{1},\frac{y}{1}$, we have 
$$
(x-\xi,y-\delta)\cap \scp\ratl[x,y]_\scp=\scp\ratl[x,y]_\scp\cap \ratl[x,y]
$$ 
because $\scp\ratl[x,y]_\scp\cap \ratl[x,y]$ is maximal. Hence,
$$
(x-\xi,y-\delta)\cap \ratl[x,y]=\scp\ratl[x,y]_\scp\cap \ratl[x,y]\cap \ratl[x,y]=\scp.
$$



\end{enumerate}
\end{proof}

\begin{exr}\label{exr:spec_quotient}
Suppose $A$ is a ring, and $I$ an ideal of $A$. Let $\varphi:A\lrta A/I$. Show that $\varphi^{-1}$ gives an inclusion-preserving bijection between prime ideals of $A/I$ and prime ideals of $A$ containing $I$. Thus we can picture $\spec A/I$ as a subset of $\spec A$.
\end{exr}
\begin{proof}
\underline{Claim1}: $\scp\subseteq A/I$ is prime iff 
$\varphi^{-1}(\scp)$ is prime.

$\varphi(x)\varphi(y)=\varphi(xy)\in \scp$, $xy\in\varphi^{-1}\scp$. If $\varphi^{-1}\scp$ is prime, at least one of $x,y$ is contained in $\varphi^{-1}\scp\Lrta $ at least one of $\varphi(x),\varphi(y)$ is contained in $\scp$.

$xy\in \varphi^{-1}\scp$ $\Lrta \varphi(xy)\in \scp$. If $\scp$ is prime  at least one of $\varphi(x)$ and $\varphi(y)$ is contained in $\scp$ $\Lrta x$ or $y$ is  contained in $\scp$.

\underline{Claim2}: $\varphi$ induces a inclusion-preserving bijection between ideals of $A/I$ and ideals of $A$ containing $I$.

Consider an ideal $\sca\in A/I$. $\varphi^{-1}\sca$ is  an ideal in $A$. In particular, because $\varphi(I)=[0]\in\sca$ $\Lrta I\subseteq \sca$.

$$
\sca\longmapsto \varphi^{-1}\sca
$$
is an bijection because we have the inverse, for $I\subseteq\scq\subseteq A$ $\varphi:\scq\longmapsto \scq/I$, where $\scq/I$ is a well-defined ideal in $A/I$.

And $\varphi^{-1}$ preserves the inclusion of ideals (it preserves the proper inclusion in fact). 

Consider $\sca\subsetneq \scb\subseteq A/I\Lrta \exists [x]\in \scb$ s.t. $[x]\notin \sca\Lrta$ the representative element $x\in \varphi^{-1}\scb, x\notin\varphi^{-1}\sca$.
\end{proof}

\begin{exr}\label{exr:spec_localization}
Suppose $S$ is a multiplicative subset of $A$. Describe an order-preserving bijection of the prime ideals of $S^{-1}A$ with the prime ideals of $A$ that don’t meet the multiplicative set $S$.

The claim should contain the following points
\begin{enumerate}[label=(\alph*)]
\item $\scp$ prime $\Llrta \iota_*\scp$ prime,
\item $\iota^*\iota_*\scp=\scp$, (True for only prime ideal $\scp$ which satisfies $\scp\cap S=\emptyset$. If $\scp\cap S\neq \emptyset$, $\iota^*\iota_*\scp={A}$)
\item $\iota_*(\sca)= S^{-1}{A}\Llrta\sca\cap S\neq \emptyset$, (True for any ideals)
\item  $\iota_*\iota^*\scq=\scq$ (True for any ideal of $S^{-1}{A}$, not necessarily prime)
\end{enumerate}
\end{exr}

\begin{proof}
\begin{enumerate}[label=(\alph*)]
   % \item Choose an ideal $\scq\subseteq S^{-1}{A}$, $\iota^*(\scq)=\iota^{-1}(\scq)$, $\iota(S)\subseteq (S^{-1}{A})^\times,\Lrta \iota(S)\cap\scq=\emptyset$ otherwise $1\in \scq$ (In fact this part of proof also works for other ideals.) and the remaining part for proof of (a) is postponed until the next lecture.
\item 
 As for point $(a)$, $\scp$ \text{ prime} $\overset{?}{\Lrta}\iota_*\scp$ prime. Consider $\frac{a}{s}\cdot \frac{b}{t}\in \iota_*\scp$, then $\frac{ab}{st}=\frac{c}{u},c\in\scp, u\in S$, then $\exists v\in S:abuv=cstv$, where $uv\in S$ $cstv\in\scp$, $uv\not\in \scp\Lrta ab\in\scp\Lrta$ at least one of $a, b\in\scp\Lrta$ at least one of $\frac{a}{s},\frac{b}{t}\in \iota_*\scp$.
\item 
 $\iota^*\iota_*\scp\supseteq\scp$ is a general fact.
 For the converse inclusion, $\iota^*\iota_*\scp=\iota^{-1}(\iota_*\scp)\overset{?}{\subseteq}\scp$, choose an $a\in \iota^{-1}(\iota_*\scp)$. $\iota(a)=\frac{a}{1}\in\iota_*\scp\Lrta\exists b\in \scp, s\in S$ s.t. $\frac{a}{1}=\frac{b}{s}\Lrta ast=bt\in \scp$ and $s,t \in S\subseteq {A}-\scp\Lrta a\in \scp$ because $\scp$ is a prime ideal.
\item
 $\iota_*(\sca)=S^{-1}{A}\Llrta \exists a\in \sca, s\in S$ s.t. $a/s=1/1$
$\Llrta\exists t\in S$ s.t. $\sca\ni ta=ts\in S$, then $\sca\cap S\neq\emptyset$. Conversely, $\sca\cap S\neq \emptyset$, exists an $a\in \sca, a=s\in S$, then $a/s=1/1$. 
\item 
$\iota_*(\iota^*(\scb))\subset \scb$ in general. For the converse inclusion, if $a/s\in \scb$, then $a/s\cdot s/1=a/1\in\scb$, which means $a\in \iota^{*}(\scb)\Lrta a/s\in \iota_{*}(\iota^*(\scb))$. This claims means every ideal in $S^{-1}{A}$ is extension of an ideal in ${A}$. 
\end{enumerate}
The fact that $\iota^*\iota_*\scp=\scp$ and $\iota_*\iota^*\scq$ means that the correspondence preserves the inclusion and proper inclusion.
\end{proof}

\begin{exr}\label{chap3exr:3.2.L}
Show that these two rings are isomorphic. 
$$
\left(\cplx[x,y]/(xy)\right)_{[x]}\cong \cplx[x]_x
$$
\end{exr}
\begin{proof}
$A:=\cplx[x,y]/(xy)$, $S:=\{[x]^n|n\geq 0,n\in\intg\}$. For the morphism $\iota:A\lrta S^{-1}A$,
$$
a\longmapsto \frac{a}{1}
$$
$$
\frac{a}{1}\sim \frac{0}{1}\Llrta \exists s\in S \text{ s.t. } sa=0
$$
 the kernel is $\{a\in A|\exists n\geq 0 \text{ s.t. } [x]^na=0\}$. $[y]\in \ker \iota$. A general element $\frac{a}{s}\in S^{-1}A$ should be 
 $$
\frac{a}{s}=\frac{p([x],[y])}{[x]^n}\sim \frac{1}{[x]^n}p\left(\frac{[x]}{1},\frac{[y]}{1}\right)=\frac{1}{[x]^n}p\left(\frac{[x]}{1},0\right)=\frac{q([x])}{[x]^n},
 $$
 where polynomial $q(X):=p(X,0)$.
And then we can define the isomorphism 
$$
\begin{aligned}
\phi: (\cplx[x,y]/(xy))_{[x]}&\lrta \cplx[x]_x\\
 \frac{q([x])}{[x]^n}&\longmapsto \frac{q(x)}{x^n}
\end{aligned}
$$
\end{proof}

\begin{exr}
If $\phi: B\lrta A$ is a map of rings, and $\scp$ is a prime ideal of $A$, show that $\phi^{-1}\scp$ is a prime ideal of $B$.
\end{exr}
\begin{proof}
Consider two elements $x,y\in B$ s.t. $xy\in \phi^{-1}\scp$. Then $\phi(xy)\in B\Lrta \phi(x)\phi(y)\in\scp$. Because $\scp$ is prime ideal, at least one of $\phi(x),\phi(y)\in\scp\Lrta $ at least one of $x,y$ is contained in $\phi^{-1}\scp$.
\end{proof}

\begin{exr}
Let $B$ be a ring.
\begin{enumerate}[label=(\alph*)]
	\item  Suppose $I\subseteq B$ is an ideal. Show that the map $\spec B/I\lrta \spec B$ is the inclusion of  prime ideals that containing $I$.
    \item  Suppose $S\subseteq B$ is a multiplicative set. Show that the map $\spec S^{-1}B\lrta \spec B$ is the inclusion of prime ideals that does not intersect with $S$.
\end{enumerate}
\end{exr}
\begin{proof}
The detailed proof of the correspondences are already contained in~\ref{exr:spec_quotient} and~\ref{exr:spec_localization}.
\end{proof}

\begin{exr}
Consider the map of complex manifolds sending $\cplx\lrta \cplx$ via $x\mapsto y=x^2$. We interpret the “source” $\cplx$ as the “$x$-line”, and the “target” $\cplx$ the “$y$-line”. You can picture it as the projection of the parabola $y=x^2$ in the $xy$-plane to the $y$-axis. Interpret the corresponding map of rings as given by ${\cplx[y] \lrta \cplx[x]}$ by  $y\mapsto x^2$. Verify that the preimage (the fiber) above the point $a\in\cplx$ is the point(s) $\pm\sqrt a\in\cplx$, using the definition given above.
\end{exr}
\begin{proof}
We have defined the morphism of rings
$$
\begin{aligned}
\phi:\cplx[y]&\lrta\cplx[x]\\
 p(y)&\longmapsto p(x^2)
\end{aligned}.
$$
The only prime ideals in $\cplx[x]$ are $(0)$ and $(x-b)$. 
$$
\phi^*(0)=\{p\in \cplx[y]: p(x^2)=0\}=(0)
$$
$$
\phi^*(x-b)=\{p\in\cplx[y]: p(x^2)\in (x-b)\}=\{p\in\cplx[y]: p(b^2)=0\}=(y-b^2).
$$
Then $\phi^*: \spec\cplx[x]\lrta\spec \cplx[y]$  is totally described. On each traditional point $y-a$, there are two points $(x-\pm \sqrt a)$ in the preimage.
\end{proof}

\begin{exr}
Suppose $k$ is a field,and $f_1,...,f_n \in k[x_1,...,x_m]$are given. Let $\phi: k[y_1,...,y_n] \lrta k[x_1,...,x_m]$ be the ring morphism defined by $y_i  \mapsto f_i$.
\begin{enumerate}[label=(\alph*)]
\item  Show that $\phi$ induces a map of sets $\spec k[x_1,...,x_m]/I \lrta \spec k[y_1,...,y_n]/J$ for any ideals $I \subseteq k[x_1,...,x_m]$ and $J \subseteq k[y_1,...,y_n]$ such that $\phi(J) \subseteq I$. 
\item Show that the map of part (a) sends the point $(a_1,...,a_m) \in k^m$ (or more
precisely, $[(x_1 -a_1,...,x_m - a_m)] \in \spec k[x_1,...,x_m])$ to
$$
(f_1(a_1,...,a_m),...,f_n(a_1,...,a_m)) \in k^n
$$
\end{enumerate}
\end{exr}
\begin{proof}
\begin{enumerate}[label=(\alph*)]
\item In fact, part (a) has nothing to do with $k$-algebras. We can choose any two rings $\phi:B\lrta A$ and two ideals $J\subseteq B, I\subseteq A$ and $\phi(J)\subseteq I$ and prove that $\phi$ induces a morphism $\spec (A/I)\lrta \spec (B/J)$. 

We direct define a ring morphism induced by $\phi$
$$
\begin{aligned}
\tilde{\phi}:B/J&\lrta A/I\\
b+J&\longmapsto \phi(b)+I.
\end{aligned}
$$
$\tilde{\phi}$ is well-defined because if we choose another representative $b'$ s.t. $b-b'\in J$, we have $\phi(b)-\phi(b')=\phi(b-b')\in I\Lrta \phi(b)+I=\phi(b')+I$. Then problem reduces to the basic case, where $\tilde{\phi}$ induces morphism
$$
\tilde{\phi}^*:\spec(A/I)\lrta \spec (B/J).
$$
\item The point (maximal ideal) $(x_1-a_1,...,x_m-a_m)$  is mapped to $\phi^{-1}(x_1-a_1,...,x_m-a_m)$
$$
\begin{aligned}
\phi^{-1}(x_1-a_1,...,x_m-a_m)&=\{p\in k[y_1,...,y_n]:\phi(p)=p[f_1,...,f_n]\in (x_1-a_1,...,x_m-a_m)\}\\
&=\{p\in k[y_1,...,y_n]: p\left(f_1(a_1,...,a_m),...,f_n(a_1,...a_m)\right)=0\}\\
&=(y_1-f_1(a_1,..,a_m),...,y_n-f_n(a_1,...,a_m)).
\end{aligned}
$$ 
\end{enumerate}
\end{proof}

\begin{exr}
Consider the map of sets $\pi: \mathbb{A}^n_\intg \lrta  \spec \intg$,
given by the ring map $\intg \lrta \intg[x_1,..., x_n]$. If $p$ is prime, describe a bijection between the fiber $\pi^{-1}([(p)])$ and $\mathbb{A}^n_{\bbf_p}$.
(Can you interpret the fiber over $[(0)]$ as $\mathbb{A}^n_{k}$ for some field $k$?)
\end{exr}
\begin{proof}
$\pi$ is induced by the canonical inclusion $\iota:\intg\inj \intg[x_1,..,x_n]$. In particular, for each prime ideal $\scq\in \mathbb{A}^n_\intg$
$$
\pi(\scq)=\scq\cap\intg.
$$
Hence $\mathbb{A}^{n}_\intg$ can be partitioned as 
$$\mathbb{A}^n_\intg
=\pi^{-1}([(0)])\cup \left(\bigcup_{p\text{ prime }}\pi^{-1}([(p)])\right).
$$
For each $\scp\in \mathbb{A}^n_\intg$, $\scp\in \pi^{-1}([(p)])\Llrta \scp\cap\intg=(p)$.\\
$\overset{(*)}{\Llrta}\scp\supseteq (p)\intg[x_1,...,x_n]=:\iota_*((p))$.
$$
\intg[x_1,...,x_n]/(p)\intg[x_1...,x_n]\cong \bbf_p[x_1,...,x_n].
$$
There is a one to one correspondence of prime ideals in $\bbf_p[x_1,...,x_n]$ and prime ideals in $\intg[x_1,...,x_n]$ that containing $(p)\intg[x_1,...,x_n]$. Hence $\pi^{-1}([(p)])=\spec \bbf_p[x_1,...,x_n]=\affn_{\bbf_p}^n$.

The $\Llta$ part of $(*)$ is questionable. We have to prove each prime ideal $\scq$ that contains $(p)\intg[x_1,..,x_n]$ must have the property that $\scq\cap\intg=(p)$. $\scq$ is such a prime ideal, $\iota_*\scq\supseteq\iota_*((p)\intg[x_1,..,x_n])$ but $(p)$ is already maximal in $\intg$, hence $\iota_*\scq=\scq\cap\intg=(p)$.

As for the fiber over $[(0)]$, $\scp\in\pi^{-1}([(0)])\Llrta \scp\cap \intg=0$.
 Considering the multiplicative set $S:=\intg^\times$, it is equivalent to $S\cap\scp=\emptyset$. There is a one to one correspondence between the primes $S^{-1}\intg[x_1,..,x_n]$ and the primes in $\intg[x_1,..,x_n]$ which does not intersect $S$ and 
 $
S^{-1}\intg[x_1,...,x_n]\cong \ratl[x_1,...,x_n],
 $   
 thus 
 $$
\pi^{-1}([(0)])\cong \affn^n_\ratl.
 $$

\end{proof}

\begin{exr}\label{exr:nilpotents_bijection}
Ring elements that have a power that is $0$ are called \textbf{nilpotents}. 
\begin{enumerate}[label=(\alph*)]
\item 
 Show that if $I$ is an ideal of nilpotents, then the inclusion $\spec B/I \lrta \spec B$ of Exercise~\ref{exr:spec_quotient} is a bijection. Thus nilpotents don’t affect the underlying set. 
\item Show that the nilpotents of a ring $B$ form an ideal. This ideal is called the \textbf{nilradical}, and is denoted $\mathfrak{N} = \mathfrak{N}(B)$.
\end{enumerate}
\end{exr}
\begin{proof}
(a) Each element of $I$ is nilpotent $x\in I$, $x^n=0\in\scp$. $\Lrta x\in\scp$ because $\scp$ is prime. $I\subseteq\scp$. Each prime ideals contain $I$, hence $\spec B/I\lrta \spec B$ is a bijection.

(b) Two nilpotents $x,y\in B$.
$x^n=0$ and $y^m=0$. $\forall a\in A, (ax)^n=0$ and $(ax+by)^{n+m}=0$. Hence, all nilpotents form an ideal.
\end{proof}

\begin{exr}
Prove : 
$$
\mathfrak{N}(A)=\bigcap_{\text{primes in }A}\scp
$$
\end{exr}
\begin{proof}
Denote by $\mathfrak{N}'$ the intersection of all prime ideals of ${A}$.
For any nilpotent element $f\in {A}$ with $n>0$ s.t. $f^n=0$,
We have $f^n\in \scp$ for every prime ideal $\scp$. Hence $f\in \scp\Lrta$ conclude $f\in \mathfrak{N}'$.

Conversely, suppose $f\in {A}$ is not nilpotent 
Define $\Sigma:=\{ \sca \subset A\text{ ideals} | \forall n>0: f^n\notin \sca\}$.
We will apply Zorn's lemma. 
We have 

    1.  $(0)\in \Sigma$, so $\Sigma$ is nonempty,

    2.  $\Sigma$ is partially ordered by inclusion.

    3.  For any chain $(\sca_i)_{i\in I}\subset \Sigma$, the set $\sca :=\cup_{i\in I}\sca_i$ is an ideal and 

    for all $n>0$, 
    we have $f^n \notin \sca$, hence $\sca \in \Sigma$.
By Zorn's lemma we conclude that there is a maximal element $\scp\in \Sigma$.

We show that $\scp$ is a prime ideal.
For any $x,y \notin \scp$, consider the ideals
$\scp+(x), \scp+(y)$. They strictly contain $\scp$ and are thus not in $\Sigma$.
Let $n,m>0$ s.t. $f^n\in \scp+(x), f^m\in \scp+(y)$.
We conclude that $f^{n+m}\in \scp+(xy)$, so $\scp+(xy) \notin \Sigma$.
Hence $xy \notin \scp$, which means, $\scp $ is a prime ideal so $f \notin \mathfrak{N}'$.
\end{proof}

\begin{exr}
Suppose we have a polynomial $ f(x) \in k[x]$. Instead, we work in $k[x, \epsilon]/(\epsilon^2)$. What then is $f(x + \epsilon)$?
\end{exr}
\begin{proof}
First we check what happens to simple examples like $x^2$ and $x^3$
$$
(x+\epsilon)^2=x^2+2x\epsilon +\epsilon^2=x^2+2x \epsilon
$$
and
$$
(x+\epsilon)^3=x^3+3x^2\epsilon.
$$
For general term $x^n$, we have $(x+\epsilon)^n=x^n+nx^{n-1}\epsilon$, which can be linearly extended to general polynomial $f$.
$$
f(x+\epsilon)=f(x)+f'(x)\epsilon
$$

\end{proof}

\section{Visualizing schemes I: generic points}
\section{The underlying topological space of an affine scheme}
\begin{exr}
Check that the $x$-axis is contained in $V(xy, yz)\subseteq \spec\cplx[x,y,z]$.
\end{exr}
\begin{proof}
$x$-axis is defined to be $(y-0,z-0)$. $\{xy,yz\}\subset (xy,yz)\subseteq (y,z)$ by definition $(y,z)\in V(xy,yz)$.
\end{proof}

\begin{exr}
Show that if $(S)$ is the ideal generated by $S$, then $V(S) =
V ((S))$.
\end{exr}
\begin{proof}
$V(S):=\{[\scp]\in \spec A: S\subset\scp\}$ and $V((S):=\{[\scp]\in \spec A: (S)\subset\scp\}$.

 If $\scq\in V(S)\Lrta S\subset\scq\Lrta (S)\subset\scq\Lrta \scq\in V((S))$, we have $V(S)\subseteq V((S))$.

 If $\scq\in V((S))\Lrta (S)\subset\scq\Lrta S\subseteq \scq\Lrta \scq\in V(S)$, we have $V((S))\subseteq V(S)$.
\end{proof}

\begin{exr}\ 
\begin{enumerate}[label=(\alph*)]
\item Show that $\emptyset$ and $\spec A$ are both open subsets of $\spec A$.
\item If $I_i$ is a collection of ideals (as $i$ runs over some index set), show that $\cap_i V(I_i) =
V(\sum_i I_i)$. Hence the union of any collection of open sets is open.
\item  Show that $V(I_1) \cup V(I_2) = V(I_1\cdot I_2)$. Hence the intersection of any finite number of open sets is open.
\end{enumerate}
\end{exr}
\begin{proof}
\begin{enumerate}[label=(\alph*)]
\item $\emptyset=\spec A-V((0))$ and $\spec A=\spec A-V((1))$, hence they are both open sets.
\item $\scp\in\cap_iV(I_i)\Lrta \scp\supseteq I_i\forall i$. Recall that $\sum_i I_i$ is defined to be $\sum_i a_i$, where $a_i\in I_i$ and only finitely many of they are non-zero. (this is basically the smallest ideal that contains $\cup_i I_i$). Hence, $\scp\supseteq \sum_i I_i$. Then, we have $\cap_i V(I_i)\subseteq V(\sum_i I_i)$.

The reverse inclusion is easier. $\scp\supseteq \sum_i I_i\supseteq \cup_iI_i$. We can conclude that $\cap_i V(I_i)=V(\sum_i I_i)$ and consider the de morgen law, union of arbitrary collection of open set is still open.
\item $\scp\in V(I_1)\cup V(I_2)\Lrta \scp\in V(I_1)$ or $\scp\in V(I_2)$. $\Lrta \scp\supseteq I_1$ or $\scp\supseteq I_2$. In either case $\scp\supseteq I_1I_2$ because $I_1I_2$ is a subset in $I_1\cap I_2$. $V(I_1)\cup V(I_2)\subseteq V(I_1\cdot I_2)$.

For the reverse inclusion, $\scp\in V(I_1\cdot I_2)$, we have $\scp\supseteq I_1\cdot I_2$.\\
\underline{Want}: $\scp\supseteq I_1$ or $\scp\supseteq I_2$.

Suppose it is not the case, $\scp\notsupset I_1$ and $\scp\notsupset I_2$. Then there are $x_1\in I_1$ and $x_2\in I_2$ such that both are not in $\scp$. $x_1x_2\in I_1\cdot I_2 \subseteq \scp$ but $\scp$ is prime ideal, contradiction.
\end{enumerate}
\end{proof}

\begin{exr}
Show that $\sqrt I$ is an ideal and
$$
\sqrt {\sqrt I}=\sqrt I.
$$
Show that prime ideals are radical
$$
\sqrt \scp=\scp
$$
\end{exr}
\begin{proof}
Two element $x,y\in \sqrt I$.
$x^n\in I$ and $y^m\in I$. $\forall a\in A, (ax)^n\in I$ and $(ax+by)^{n+m}\in I$. Hence, radical is an ideal.

$\sqrt{\sca}\supseteq \sca$ is trivially true. For the reverse inclusion, consider an element $x\in \sqrt{\sqrt I}$. $x^n\in \sqrt I$, then $(x^n)^m\in I\Lrta x\in \sqrt I$.

$x^n\in\scp\Lrta x\in \scp$ because $\scp$ is prime ideal. 
In fact we have $\sqrt{\scq^n}=\scq$ because 
$$
\scp\subseteq \sqrt{\scq^n}\subseteq\sqrt{\scp}=\scp.$$
\end{proof}

\begin{exr}
Prove that taking radical commutes with finite intersection of ideals.
\end{exr}

\begin{exr}
Prove that 
$$
\sqrt{\sca}=\bigcap_{\text{primes}\supseteq \sca }\scp
$$
\end{exr}

\begin{exr}
Describe the topological space $\affn^1_k$
\end{exr}
\begin{proof}
The points of the topological space is$[(0)]$ and $[(x-a)],\forall a\in k$. It is almost the cofinite topology on $k^1$ but there is a new point $[(0)]$. 

\underline{Claim}: Every point $[\scp]\in \affn^1_n$, every open neighborhood of $[\scp]$ contains $[(0)]$.

The claim holds for the case $\scp=(0)$. For $\scp=(x-a)$. An open neighborhood of $[\scp]=[(x-a)]$ is of the form $\affn^1_k-V(S)$. which means $[\scp]\notin V(S)\Lrta S\notsubset\scp$.

\underline{Want}: $\affn^1_k-V(S)\ni [(0)]\Llrta S\notsubset (0)$, which is a direct fact from $S\notsubset \scp$ and $(0)\subset \scp$

Then we know every open neighborhood of every point in $\spec k[x]$ contains $(0)$. In other words, the closure of $(0)$ is the whole affine line.
\end{proof}

\begin{exr}\label{exr:spec_functor} A ring morphism
$\phi: B\lrta A$ induces a map $\pi$ on the spectrum. By showing that closed sets pull back to closed sets, show that $\pi=\phi^*:\spec A\lrta \spec B$ is a continuous map. Interpret $\spec$ as a contravariant functor \textit{Rings}$\lrta$\textit{Top}.
\end{exr}
\begin{proof}
Assume $S$ is a subset in $B$. Consider a closed set $V(S)$ in $\spec B$, we will verify that $\pi^{-1}V(S)$ is also closed.
$$
\pi^{-1}V(S)=\{[\scp]\in \spec A: \pi(\scp)\supset S\}=\{[\scp]\in \spec A: \phi^*(\scp)\supset S\}
$$
\underline{Claim}: $\pi^{-1}V(S)=V(\phi(S))$

In fact $\scp\supset \phi(S)\Llrta \phi^{-1}(\scp)\supset S$, therefore,
$$
\pi^{-1}V(S)=\{[\scp]\in \spec A: \phi^*(\scp)\supset S\}=\{[\scp]\in \spec A: \scp\supset \phi(S)\}=V(\phi(S)).
$$
The preimage of a closed set is always closed. $\pi$ is continuous.

Now, we can interpret $\spec$ as a contravariant functor form $Rings$ to $Top$

\begin{center}
\begin{tikzcd}
B\ar[r, "\phi"]\ar[d, "\spec"] &A  \arrow[d,"\spec"] \\
\spec B & \spec A\ar[l,"\spec(\phi)=\pi"]       
\end{tikzcd}
\end{center}
\end{proof}


\begin{exr}\label{chap3exr:specA_f_and_specA/I}
Suppose that $I, S \subset B$ are an ideal and multiplicative subset respectively.
\begin{enumerate}[label=(\alph*)]
\item Show that $\spec B/I$ is naturally a closed subset of $\spec B$. If $S = {1,f,f^2,...} (f \in B)$, show that $\spec S^{-1}B$ is naturally an open subset of $\spec B$. Show that for arbitrary $S$, $\spec S^{-1}B$ need not be open or closed. (Hint:$ \spec \ratl \subset \spec \intg$, or possibly Figure 3.5.)
\item  Show that the Zariski topology on $\spec B/I$ (resp. $\spec S^{-1}B$) is the subspace topology induced by inclusion in $\spec B$. (Hint: compare closed subsets.)
\end{enumerate}
\end{exr}
\begin{proof}
\begin{enumerate}[label=(\alph*)]
\item Prime ideals in $B/I$ are the prime ideals in $B$ that  contains $I$ by~\ref{exr:spec_quotient}.
$$
\spec B/I=\{[\scp]\in B:\scp\supset I\}=V(I)
$$
which a closed set in $\spec B$.
For multiplicative set generated by one element $S:=\{1,f,f^2,...\}$. Recall~\ref{exr:spec_localization}. The prime ideals in $S^{-1}B$ are just the prime ideals in $B$ that do not intersect $S$. 
Because $\scp$ is prime, and $1\notin \scp$, $\scp\cap S\neq \emptyset\Lrta f^n\in \scp\Lrta f\in \scp$, thus $\scp\supset S-\{1\}$. Then we know
$$
\spec S^{-1}B=\spec B-\{[\scp]\in \spec B: \scp\ni f\}=D(f),
$$
which is an open set.

In general, $\spec S^{-1}B$ is neither open nor closed in $\spec B$. For example consider $\spec\ratl \subset\spec \intg$.
$\spec\ratl$ corresponds to the generic point $[(0)]$ in $\scp$, which is neither open nor closed. Or we can consider the morphism $\spec \cplx[x,y]_{(x,y)}$, which corresponds to prime ideals in $\cplx[x,y]$ that are contained in $(x,y)$, which is neither open nor closed.
\item Denote the projection $\phi:B\lrta B/I$. Consider a closed set in $\spec B/I$, call it $V_{B/I}(D)$, where $D$ is a subset in $B$.
$$
\begin{aligned}
V_{B/I}(D)&=\{[\scp]\in B/I:\scp\supset D\}\\
&\cong \{[\scq]\in B:\scq\supset I\text{ and }\scq\supset \phi^{-1}D\}\\
&=\{[\scq]\in B:\scq\supset I+(\phi^{-1}D)\}\\
&=\spec(B/I)\cap V_{B}(\phi^{-1}D).
\end{aligned}
$$
For the localization map $\iota:B\lrta S^{-1}B$. Consider a closed set in $\spec B$, call it $V_{S^{-1}B}(P)$, where $P$ is a subset in $S^{-1}B$.
$$
\begin{aligned}
V_{S^{-1}B}(P)&=\{[\scp]\in \spec (S^{-1}B):\scp\supset P\}\\
&\cong\{[\scq]\in B: \scq\cap S=\emptyset,\text{ and } \scq\supset \iota^{-1} P\}\\
&=\{[\scq]\in B: \scq\cap S=\emptyset\}\cap\{[\scq]\in B: \scq\supset \iota^{-1} P\}\\
&=\spec S^{-1}B\cap V_B(\iota^{-1}P).
\end{aligned}
$$
We have verified that Zariski topologies on $\spec B/I$ and $\spec S^{-1}B$ are the induced topology on by Zariski topology on $\spec B$.
\end{enumerate}
\end{proof}
\begin{remark}
In particular, if $I$ is an ideal in the nilradical, the $\spec B$ and $\spec B/I$ are homeomorphic. $\iota:\spec B/I\lrta \spec B$. It is continuous~\ref{exr:spec_functor} and bijective~\ref{exr:nilpotents_bijection}.
A continuous bijection is not necessarily homeomorphism. But we have a theorem that \textbf{bijective continuous map is homeomorphism iff it is closed or open}. In our case, $\iota$ is closed. 
\end{remark}
\begin{exr}\label{exr:vanishing_set_radical}
Suppose $ I \subset B$ is an ideal. Show that $f$ vanishes on $V(I)$ if and only if $f \in \sqrt I$
\end{exr}
\begin{proof}
Assume $f$ vanishes on $V(I)$, then
$$
f\mod \scp=0,\ \forall  \scp\supset I,
$$
which is equivalent to 
$$
f\in \bigcap_{\scp\supset I}\scp=\sqrt{I}
$$

\end{proof}

\begin{exr}
Describe the topological space $\spec k[x]_{(x)}$.
\end{exr}
\begin{proof}
The only two prime ideals in $k[x]_{(x)}$ are $(0)$ and $\scm:=(x)k[x]_{(x)}$, of which $[(0)]$ is generic point and $[\scm]$ is closed point.
\end{proof}


\section{A base of the Zariski topology on schemes: Distinguished open sets}
\begin{exr}
Show that the distinguished open sets form a base for the (Zariski) topology. (Hint: Given a subset $ S \subset A$, show that the complement of $V(S)$ is $\cup_{f\in S}D(f).$)
\end{exr}
\begin{proof}
Recall one of the equivalent definitions of base of topology:

Base of a topology is a subset $B$ of topology $\tau$ such that every open set in $\tau$ can be obtained by union of some elements of $B$.

Each open set in the Zariski topology is the complement of some closed set $V(S)$.

\underline{Claim}: $\spec A-V(S)=\cup_{f\in S} D(f)$\\
$[\scp]\in \cup_{f\in S} D(f)\Llrta [\scp]\in D(f)$ for some $f\in S$.\\
$ \Llrta\exists f\in S, f\notin \scp\Lrta \scp\notsupset S$ $\Llrta [\scp]\in\spec A-V(S)$.
\end{proof}

\begin{exr}\label{exr:open_cover_distinguished}
Suppose $f_i \in A$ as $i$ runs over some index set $J$. Show that $\cup_{i\in J}D(f_i) = \spec A$ if and only if $(\{f_i\}_{i\in J}) = A$, or equivalently and very usefully,
if there are $ \{a_i\}_{i\in J} $, all but finitely many $0$, such that 
$\sum_{i\in J}a_if_i = 1$.
\end{exr}
\begin{proof}
``$\Lrta$'': We know by definition $D(f_i)=\spec A-V((f_i))$.

$\cup_{i\in J}D(f_i)=A\Lrta \cap V((f_i))=\emptyset$. $\Lrta V(\sum_i(f_i))=\emptyset$.

In particular, for principal ideals 
$\sum_i(f_i)=(f_1,f_2,...f_i,..)$.

Any proper ideal is contained in some maximal ideal $\scm$, hence $\sum_i(f_i)$ can be a proper ideal, other wise there is at least one element $[\scm]\in V(\sum_i (f_i))$. We know $\sum_i(f_i)=(f_1,f_2,...f_i,..)=A=(1)$, which means there exists a summation such that $\sum_{i\in J}a_if_i = 1$ with finitely many $a_i$ nonzero.

``$\Llta$'': We know $\sum_{i\in J}a_if_i = 1$ and hence $\sum_{i\in J} (f_i)=A$. Then we know
$$
\spec A-\cup_i D(f_i)=\cap_i (\spec A-D(f_i))=\cap_i V((f_i))=V(\sum_{i\in J} (f_i))=\emptyset.
$$
\end{proof}


\begin{exr}\label{exr:distinguished_open_cover}
Show that if $\spec A $ is an infinite union of distinguished open sets $\cup_{i\in J}D(f_j)$, then in fact it is a union of a finite number of these, i.e., there is a finite subset $J'$ so that $\spec A = \cup_{j\in J'}  D(f_j)$. 
\end{exr}
\begin{proof}
Recall~\ref{exr:open_cover_distinguished},
$$
\cup_{i\in J}D(f_i) = \spec A\Llrta \sum_{j\in J}a_j f_i=1.
$$
There are only finitely many of $a_j$ nonzero. We can choose a finite subset $J'\subset J$ such that $\sum_{j\in J'}a_j f_i=1$. Then we have $\cup_{i\in J'}D(f_i) = \spec A$.
\end{proof}

\begin{exr}
Show that $D(f) \cap D(g) = D(fg)$.
\end{exr}
\begin{proof}
Pick $[\scp]\in D(f) \cap D(g)$, then $[\scp]\in D(f)$ and $D(g)$ $f\notin \scp$ and $g\notin \scp$, then by $\scp$ being prime, we have $fg\notin\scp\Lrta [\scp]\in D(fg)$. $D(f) \cap D(g)\subset D(fg)$.

For the reverse inclusion, pick $\scq\in D(fg)$, then $fg\notin \scp$, again by primality, we have $f\notin\scp$ and $g\notin\scp$, hence $[\scp]\in D(f)\cap D(g)$.
\end{proof}


\begin{exr}\label{chap3exr:inclusion_distinguished_open}
Show that $D(f) \subset D(g)$ if and
only if $f^n \in(g)$ for some $n\geq 1$, or equivalently, if and  only if the image of $g$ is an invertible element of $A_f$.
\end{exr}
\begin{proof}
$D(f)\subset D(g)\Llrta V((f))\supset V((g))$. $V((g))$ is the closed set where $g$ vanishes. The derived inclusion means is equivalent to saying that $f$ vanishes on $V((g))$. By Exercise~\ref{exr:vanishing_set_radical}, we know it is equivalent to 
$$
f\in \sqrt{(g)}
$$
$$
\Llrta (g)\cap S_f\neq \emptyset,
$$
where $S_f:=\{1,f,f^2,...,\}$.
$$
\Llrta S_f^{-1}(g)=A_f
$$
which is equivalent to ``the image of $g$ is invertible in $A_f$''.
\end{proof}

\begin{exr}
Show that $D(f) = \emptyset$ if and only if $f\in \mathfrak{N}(A)$
\end{exr}
\begin{proof}
$D(f) = \emptyset\Lrta \forall \scp\in \spec A$, $f\mod\scp=0\Lrta f\in \cap_{\text{prime}}\scp=\mathfrak{N}(A)$.

For $\Llta$ direction, consider an element $g\in \mathfrak{N}(A)$,  $(0)\subset (g)\subset \mathfrak{N}(A)$. 
$$
\spec A=V((0))\supset V((g))\supset V\left(\sqrt{(0)}\right)=\spec A
$$
 $D(g)=\spec A- V((g))=\emptyset$.
\end{proof}

\section{Topological (and Noetherian) properties}
\begin{exr}\label{exr:nonconnected_scheme}
$A = A_1 \times A_2 \times\cdots \times A_n$, describe $
\spec A_1\coprod \spec A_2\coprod \cdots \coprod\spec A_n\lrta \spec A$ as a homeomorphism for which each $\spec A_i $ is mapped onto a distinguished open subset $D(f_i)$ of $\spec A$. Thus $\spec \prod^n_{i=1}A_i=\coprod^n_{i=1}\spec A_i$ as topological spaces.
\end{exr}
\begin{proof}
We can induct on $n$ and reduce to the spacial case $n=2$. Consider $A=A_1\times A_2$. 

\underline{Claim}: $\sca\in A$ is an ideal iff $\sca=\sca_1\times \sca_2$, for some ideals $\sca_1\in A_1$ and $\sca_2\in A_2$.

``$\Llta$'' direction is trivial. Let's consider the ``$\Lrta$'' direction. $\sca$ is an ideal in $A$, therefore $\sca=S\times T$, where $S,T$ are general subset in $A_1,A_2$. $\sca$ should be closed under multiplication and linear summation of $A_1\times 0$, which implies that $S$ is an ideal in $A_1$, similarly we can prove that $T$ is an ideal in $A_2$.

\underline{Claim}: $\scp\in A$ is a prime ideal iff $\scp=\scp_1\times A_2$ or $\scp=A_1\times \scp_2$.

We already know that the ideal $\scp=\scq_1\times \scq_2$ and the quotient ring $A/\scp=A_1/\scq_1\times A_2/\scq_2$ is integral domain. Recall that $(1,0)\times (0,1)=(0,0)$, We must have $\scq_1=A_2$ or $\scq_2=A_2$. The reverse direction is trivial.

Now, set $f_1=(1,0)$ and $f_2=(0,1)$ and denote the two canonical projections $\phi_i:A_1\times A_2\lrta A_i$. We have
$$
\begin{aligned}
D(f_1)&=\{[\scp]\in \spec A: f_1=(1,0)\notin\scp\}\\
&=\{\scp_1\times A_2:\scp_1\in \spec A_1\}
\end{aligned}
$$ 
and 
$$
\begin{aligned}
\pi_1:=\spec(\phi_1):\spec A_i&\lrta \spec A\\
[\scp_1]&\longmapsto \phi_1^{-1}(\scp_1)=\scp_1\times A_2.\\
V(S)&\lrta V(S)\times A_2
\end{aligned}
$$
Hence,
$$
D(f_i)=\pi_i(\spec A_i).
$$
Each $\pi_i$ is continuous bijective. Furthermore, $\pi_i$ is closed map as shown above. Then $\pi_i$ is homeomorphism. 
Recall the universal property of coproduct (disjoint union in \textit{Top}), there is a unique morphism from $\spec A_1\coprod\spec A_2$ to $\spec A$, which is denoted by $\pi_1\coprod \pi_2$. This morphism is also continuous bijective and closed.

These proof can be generalized to finite product and finite coproduct without much difficulty.
\end{proof}
\begin{theorem}
$\spec A$ is not connected if and only if $A$ is isomorphic to the product of nonzero rings $A_1$ and $A_2$.
\end{theorem}
\begin{proof}
One direction has been proved in~\ref{exr:nonconnected_scheme}. Now we just focus on the other direction. Follow the hint, we call a pair of idempotents $a_1,a_2$ such that $a_1^2=a_1$, $a_2^2=a_2$, $a_1+a_2=1$ and $a_1a_2=0$ \textbf{complimentary idempotents}. 

We can see there are a pair of complimentary idempotents in $A_1\times A_2$, $(1,0)$ and $(0,1)$.

\underline{Claim}: If $\spec A$ not connected, there is a pair of complimentary idempotents in $A$.

$A$ is not connected, it can decompose into disjoint union of two open sets or equivalently two closed sets. Suppose they are $V(I_1)$ and $V(I_2)$, we have
$$
V(I_1)\cup V(I_2)=V(I_1I_2)=\spec A
$$
and 
$$
V(I_1)\cap V(I_2)=V(I_1+I_2)=\emptyset.
$$
At most can find a pair of element $f_1\in I_1$ and $f_2\in I_2$, such that 
$f_1+f_2=1$ and $f_1f_2$ is nilpotent. 

\underline{Lemma}: Every nontrivial idempotent in $A/\mathfrak{N}$ lifts to a unique nontrivial idempotents in $A$.

With this lemma, assume $e_1,e_2$ are the nontrivial idempotents in $A$. We have $(e_1)\cap (e_2)=(e_1)\cdot (e_2)$ and $(e_1)+(e_2)=(1)$, we can use the Chinese remainder theorem
$$
A=A/((e_1)\cdot(e_2))\cong\frac{A}{(e_1)}\times \frac{A}{(e_2)}\cong (e_2)\times (e_1)
$$

Now, we prove the lemma. $f\neq 0,1$ $f(1-f)$ is nilpotent. $\exists n\geq 0, f^n(1-f)^n=0$. Because $f^n$ and $(1-f)^n$ are coprime, we can still use the Chinese remainder theorem
$$
A\cong \frac{A}{(f^n)}\times \frac{A}{(1-f)^n}.
$$
The preimage of $(1,0)$ is unique and idempotent.
(This step in fact gives us the desired proof that $A$ should be a product of rings but we want to reduce it to the simplest case)


Or simpler, start with $x+y=1$ and $xy$ nilpotent. $(xy)^m=0$.
$$
1 = (x + y)^{2m} = \underbrace{x^{2m} + \cdots + \binom{2m}{m+1} x^{m+1}y^{m-1}}_{e_1} + \underbrace{\binom{2m}{m}x^my^m + \cdots + y^{2m}}_{e_2}
$$
and these $e_1,e_2$ satisfies
$$
e_1e_2 = 0, \quad e_1^2 = e_1, \quad e_2^2 = e_2.
$$
\end{proof}

\begin{exr}\label{chap3exr:irreducible_subset_dense_open_irreducible_closure}\ 
\begin{enumerate}[label=(\alph*)]
\item Show that in an irreducible topological space, any nonempty open set is dense. (For this reason, you will see that unlike in the classical topology, in the Zariski topology, nonempty open sets are all “huge”.)
\item If $X$ is a topological space, and $Z$ (with the subspace topology) is an irreducible subset, then the closure $\overline{Z}$ in $X$ is irreducible as well.
\end{enumerate}
\end{exr}
\begin{proof}
\begin{enumerate}[label=(\alph*)]
\item Recall that in topology $\overline{U\cup V}=\overline{U}\cup \overline{V}$. (The analogous statement for intersection of closure is not true.) In an irreducible topological space $X$, consider an open set $U$, if $\overline{U}$ is not the whole space $X$, we have $V:=X-\overline{U}\neq \emptyset$, then we have $\overline{U}\cup \overline{V}=X$, contradiction. 

Also an open subset $U$ in a irreducible space $X$ is irreducible with the subset topology. If $U=A\cup B$ with $A,B$ closed in $U$, then taking the closure in $X$ yields $X=\overline{U}=\overline{A}\cup \overline{B}$ because $U$ is dense in $X$. This forces $\overline{A}=X$. By assumption $A$ is closed in $U$ and this means $A=\overline{A}\cap U=X\cap U=U$. It follows that $U$ is irreducible in subspace topology on it.


In the special case $X=\spec A$. We only need to verify that every nonempty distinguished open set is dense in Zariski's topology.

Consider a point $[\scp]\in V((f))=\spec A-D(f)$, then without loss of generality, there is  an open neighborhood of the form $D(g)\ni[\scp]$. 

\underline{Claim}: $D(f)\cap D(g)\neq \emptyset$.

Assume $D(f)\cap D(g)=\emptyset$, then $\spec A-D(f)\cap D(g)=(\spec A-D(f))\cup(\spec A-D(g))=V((f))\cup V((g))=\spec A$ contradicting to our assumption that $\spec A$ is irreducible.

A by-product of this fact is that every product of non-nilpotent is non-nilpotent.

 \item If the closure of $Z$ in $X$ is reducible, $\overline{Z}=U\cup V$, where $U,V$ are closed set with induced topology in $\overline{Z}$ (They are not necessarily closed in $X$. Denote the closure of $U$, $V$ in $X$ by $\overline{U}$, $\overline{V}$. We can write $Z=(\overline{U}\cap Z)\cup (\overline{V}\cap Z)$, where $\overline{U}\cap Z$ and $\overline{V}\cap Z$ are closed set because $Z$ is endowed with subset topology.
\end{enumerate}
\end{proof}

\begin{exr}\label{chap3exr:Integral_domain_irreducible}
If $A$ is an integral domain, show that $\spec A$ is irreducible. (Hint: pay attention to the generic point $[(0)]$.)
\end{exr}
\begin{proof}
Assume $\spec A$ is reducible and can be written as $\spec A=V(I_1)\cup V(I_2)$.
$\spec A=V(I_1)\cup V(I_2)=V(I_1 I_2)$, which means all $I_1I_2$ vanishes on each point $[\scp]$, i.e. $I_1I_2\subseteq\cap \scp=\mathfrak{N}$. $V(I_1)$ and $V(I_2)$ are proper closed subsets, none of $I_1$, $I_2$ is contained in $\mathfrak{N}$. In particular they are non-zero ideals. There exist non-nilpotents $x,y\in I_1,I_2$ such that $xy$ is nilpotent. In particular $x^ny^n=0$, where $x^n, y^n$ are non-zero zero-divisors. Which contradicts the hypothesis that $A$ is integral domain.
\end{proof}

\begin{exr}\label{chap3exr:irreducible_topological_space_connected}
Show that an irreducible topological space is connected.
\end{exr}
\begin{proof}
Assume a topological space $X$ is not connected. Then $X=U\coprod V$ where $U, V$ are both closed and open. Hence we also have $X=U\cup V$, which means $X$ is reducible. 
\end{proof}

\begin{exr}
Give (with proof!) an example of a ring $A$ where $\spec A$ is connected but reducible.
\end{exr}
\begin{proof}
Follow the hint, consider the ring 
$$
A:=\cplx[x,y]/(xy)
$$
$\spec A$ is connected because $A$ is not of the form $A_1\times A_2$. The only idempotents in $A$ are $0$ and $1$.

It is reducible. Indeed,
$$
\spec A=V((x))\cup V((y)),
$$
each is proper closed subset.
\end{proof}

\begin{exr}\ 
\begin{enumerate}[label=(\alph*)]
	\item 
Suppose $I = (wz-xy, wy-x^2 , xz-y^2 ) \subset k[w, x, y, z]$. Show that $\spec k[w, x, y, z]/I$ is irreducible, by showing that $k[w, x, y, z]/I$ is an integral domain.
\item Note that the generators of the ideal of part (a) may be rewritten as the equations ensuring that
$$
\text{rank}
\begin{pmatrix}
w & x & y\\
x & y & z
\end{pmatrix}
\leq 1
$$
i.e., as the determinants of submatrices. Generalize this to $2\times n$ variables.
\end{enumerate}
\end{exr}
\begin{proof}
\begin{enumerate}[label=(\alph*)]
	\item 
Follow the hint, consider an morphism
$$
\begin{aligned}
&k[w,x,y,z]\lrta k[a,b] \\
& w\longmapsto a^3, y\longmapsto ab^2,\\
& x\longmapsto a^2b, z\longmapsto b^3.\\
\end{aligned}
$$
The kernel is just $I$, and the image is $k[a^3,a^2b, ab^2,b^3]$. Thus, we have established the isomorphism
$$
k[w,x,y,z]/I\cong k[a^3,a^2b, ab^2,b^3]\subset k[a,b]
$$
a subring of integral domain is always an integral domain.
\item Generalize this to $2\times n$ matrices
$$
\text{rank}
\begin{pmatrix}
x_0 & x_1 & \cdots &x_{n-1}\\
x_1 & x_2 & \cdots &x_{n}
\end{pmatrix}
\leq 1
$$
and the ideal $I$ is generated by the determinants of submatrices. 
$$
\begin{aligned}
k[x_0,...,x_{n}]/I&\cong k[a^n,a^{n-1}b,...,b^{n}]\\
x_i&\longmapsto a^{n-i}b^i
\end{aligned}
$$
\end{enumerate}
$A:=k[x_0,...,x_{n}]/I$ is again integral domain and hence $\spec A$ is irreducible

\end{proof}
\begin{exr}\label{chap3exr:affine_scheme_quasicompact}\ 
\begin{enumerate}[label=(\alph*)]
\item  Show that $\spec A$ is quasicompact.
\item Show that in general $\spec A$ can have nonquasicompact open sets.
\end{enumerate}
\end{exr}
\begin{proof}
\begin{enumerate}[label=(\alph*)]
\item
Consider an open cover of $\spec A=\cup_{i\in J}U_i$. Recall the definition of base of topology, each open set $U_i$ is a union of $\cup_{j\in J^{(i)}}D(f^{(i)}_j)$. 
$$
\spec A=\cup_{i\in J}\cup_{j\in J^{(i)}}D(f^{(i)}_j)
$$
by Exercise~\ref{exr:distinguished_open_cover}, there is a finite subcover
$$
\spec A=\cup_{(i,j)\in J'}D(f^{(i)}_j),
$$
where $J'$ is a finite index subset in $\cup_{i\in J}J^{(i)}$. Then there are finitely many $i\in J$ s.t., $(i,j)\in J'$, denote it by $I'$. Then because $U_i\supset D(f^{(i)}_j)$, we have a finite subcover of the initial cover
$$
\spec A=\cup_{i\in I'} U_i.
$$
\item Consider the ring $A:=k[x_1,x_2,...]$ and the maximal ideal $\scm=(x_1,x_2,...)$. There is an open subset 
$$
U=\spec A-V(\scm)
$$
with open cover
$$
U=\cup_{i\in\intg} D(x_i)=\{[\scp]\in\spec A\text{ s.t. at least one of $x_i\notin \scp$}\}
$$
There is no finite subcover because $[\scp_i]:=[(x_1,...,x_{i-1},x_{i+1},...)]\in U$ but $[\scp_i]\notin U-D(x_i)$.
\end{enumerate}
\end{proof}

\begin{exr}\ 
\begin{enumerate}[label=(\alph*)]
\item 
 If $X$ is a topological space that is a finite union of quasicompact spaces, show that $X$ is quasicompact.
\item  Show that every closed subset of a quasicompact topological space is quasicompact.
\end{enumerate}
\end{exr}
\begin{proof}
\begin{enumerate}[label=(\alph*)]
\item Assume $X=\cup_{0\leq i\leq n} X_i$, every open cover of $X$ is in particular an open cover of $X_i$. Then we can select finite subcovers such that they cover each $X_i$. It is also a finite sub cover that covers $X$.
\item $Y\subset X$ is closed subspace of X. Given any cover of $Y=\cup_i U_i$. For each $U_i$ we can find open subset $V_i\subset X$ such that $V_i\cap Y=U_i$, which is a consequence of subset topology. Then we adjoint the complement of $Y$. $\cup_{i} V_i\cup (X-Y)$ is an open cover of $X$, it has a finite subcover because $X$ is quasicompact. This finite subcover of $\cup_{i} V_i\cup (X-Y)$ induces a finite subcover of $\cup_{i} U_i$.
\end{enumerate}
\end{proof}

\begin{exr}
Show that the closed points of $\spec A$ correspond to the maximal ideals.
\end{exr}
\begin{proof}
$\{[\scp]\}$ is a closed subset of $\spec A$, which means $[\scp]=V(S)$, where $S$ is a subset of $A$.

$V(S)$ consists of only one point. By definition it is equivalent to ``$\scp$ is the only prime such that $\scp\supset S$''. $(S)$ the the ideal generated by $S$, we have $V(S)=V((S))$. Each ideal is contained in some maximal ideals. The only possibility is that $(S)$ is maximal and $\scp=(S)$.
\end{proof}

\begin{exr}\label{chap3exr:closed_point_dense}\ 
\begin{enumerate}[label=(\alph*)]
\item 
Suppose that $k$ is a field, and $A$ is a finitely generated $k$-algebra. Show that closed points of $\spec A $ are dense, by showing that if $f \in A$, and $D(f)$ is a nonempty (distinguished) open subset of $\spec A$, then $D(f)$ contains a closed point of $\spec A$. 
\item Show that if $A$ is a $k$-algebra that is not finitely generated the closed points need not be dense. 
\end{enumerate}
\end{exr}
\begin{proof}
\begin{enumerate}[label=(\alph*)]
\item 
We already showed that $\spec A_f\cong D(f)\subset \spec A$. A every nonzero ring has a maximal ideal, There is a closed point in $\spec A_f$, we have to prove that the pull back of this closed point is also a closed point in $\spec A$.

Assume $\scp\subset A$ such that $\scp A_f$ is  maximal ideal in $A_f$. $\scp A_f$ contains no units $f/1\notin\scp A_f\Lrta f\notin\scp$. We have the isomorphism 
$$
\frac{A_f}{\scp A_f}\cong (A/\scp)_{f}.
$$
Notice that $A_f$ is also finitely generated $k$-algebra, by Nullstellensatz, $(A_f/\scp_f)$ is the finite field extension of $k$, hence so is $(A/\scp)_f$.

The $k$-integral domain $A/\scp$ is a subring of $(A/\scp)_f$. (It is a $k$-vector space of a finite dimensional $k$-vector space). Then we know $k$-integral domain $A/\scp$ must be a finite dimensional $k$-vector space. Then it must be a field by exercise~\ref{exr:chap3_fin-dim_k_domain_is_field}.
$\scp$ is also a maximal ideal and $[\scp]\in D(f)$. Hence, the closed points is dense in $\spec A$.
\item Consider the ring $A:=k[x]_{(x)}$, it has only one maximal ideal $\scm=(x)k[x]_{(x)}$ and one generic point $[(0)]$. The closed point $[\scm]$ is not dense because ${[0]}$ is open. This ring is not finitely generated. Consider the set $1/(1+x), 1/(1+x+x^2),1/(1+x+x^2+x^3),...$, each can not be expressed as a polynomial. (In fact $k[x]_{(x)}$ is not even a finitely generated $k[x]$-algebra.)
\end{enumerate}
\end{proof}

\begin{exr}
Suppose $k$ is an algebraically closed field, and $A=k[x_1,...,x_n]/I$ is a finitely generated $k$-algebra with $\mathfrak{N}(A) = \{0\}$. Consider the set $X = \spec A$ as a subset of $\affn^n_k$ . The space $\affn^n_k$ contains the “classical” points $k^n$. Show that functions on $X$ are determined by their values on the closed points.
\end{exr}
\begin{proof}
Suppose $f,g$ are two distinct functions on $X$.

$f-g\neq 0$ , i.e., $f-g\notin \mathfrak{N}(A)=\{0\}$. functions are determined by their value at points. \underline{Want}: Functions are determined by their value at closed points.

\underline{Claim}: If $f$ and $g$ are distinct functions, then $f-g$ is nowhere zero on an open set $D(f-g)$.

By definition, $D(f-g)$ is an the set of points on which $f-g$ does not vanish. It is not empty because $f-g\neq 0$ implies $f-g$ is not nilpotents in our setting. 

In particular, $A$ is finitely generated $k$-algebra, the closed point is dense. $D(f-g)$ contains at least one closed point because closed points is dense in $\spec A$. Hence we know distinct function $f$, $g$ must have different values on at least one closed point.

Without the hypothesis $\mathfrak{N}(A)=\{0\}$, we can not argue like in topology `` functions agreeing on a dense subset must agree on the whole set''.
\end{proof}

\begin{exr}
If $ X = \spec A$, show that $[\scq]$ is a specialization of $[\scp]$ if and only if $\scp \subset \scq$. Hence show that $V(\scp) = \overline{\{[\scp]\}}$.
\end{exr}
\begin{proof}
By definition, ``$[\scq]$ is a specialization of $[\scp]$'' means

 $[\scq]\in\overline{\{[\scp]\}}$.

$\Llrta$ every open neighborhood of $[\scq]$ contains $[\scp]$.

$\Llrta$ every distinguished open $D(f)$ containing $[\scq]$ contains $[\scp]$.

$\Llrta$ every $f$ not vanishing on $[\scq]$ neither vanishes on $[\scp]$.

$\Llrta$ $A-\scq\subset A-\scp$.

$\Llrta$ $\scq\supset \scp$.

$\Llrta$ $[\scq]\in V(\scp)$.
\end{proof}

\begin{exr}
Verify that $[(y - x^2)]\in\affn^2$  is a generic point for $V(y - x^2)$
\end{exr}
\begin{proof}
$\overline{\{[(y-x^2)]\}}=V(y-x^2)$ as $(y-x^2)$ is prime ideal.
\end{proof}

\begin{exr}
Suppose $[\scp]$ is a generic point for the closed subset $K$. Show that it is ``near every point $[\scq]$ of $K$ ''(every neighborhood of $[\scq]$ contains $[\scp]$), and ``not near any point $[\scn]$ not in $K$'' (there is a neighborhood of $[\scn]$ not containing $[\scp]$).
\end{exr}
\begin{proof}
$\overline{\{[\scp]\}}=K$ topologically means every open neighborhood of $[\scq]\in K$ contains $[\scp]$. On the other hand assume, $[\scn]\notin K$ hence $[\scn]\notin \overline{\{[\scp]\}}$ topologically means there is an open neighborhood of $[\scn]$ that does not contain $[\scp]$.
\end{proof}

\begin{exr}\label{chap3exr:UNION_OF IRREDUCIBLE_COMPONENTS}
(EVERY TOPOLOGICAL SPACE IS THE UNION OF IRREDUCIBLE COMPONENTS). Show that every point $x$ of a topological space $X$ is contained in an irreducible component of $X$.
\end{exr}
\begin{proof}
The point $\{x\}$ as a subset is irreducible. 

\underline{Claim}: Every irreducible subset $Z\subset X$ is contained in an irreducible component.

Let $Z\subset X$ be irreducible. Consider the set $\Sigma$ of irreducible subsets $Z\subset Z_\alpha\subset X$. Note that $\Sigma$ is nonempty since $Z\in \Sigma$. There is a partial ordering on $\Sigma$ coming from inclusion: $\alpha\leq \alpha'\Llrta Z_\alpha\subset Z_{\alpha'}$. Choose a maximal totally ordered subset $\Sigma' \subset \Sigma$, and let $Z'=\cup_{\alpha\in \Sigma'}Z_\alpha$. We claim that $Z'$ is irreducible. Namely, suppose that $Z'=T_1\cup T_2$ is a union of two closed subsets of $Z'$. For each $\alpha\in \Sigma'$ we have either $Z_\alpha\subset T_1$ or $Z_\alpha\subset T_2$, by irreducibility of $Z_\alpha$. Suppose that for some $\alpha_0\in \Sigma'$ we have $Z_{\alpha_0}\subset T_2$ (say, if not we're done anyway). Then, since $\Sigma'$ is totally ordered we see immediately that $Z_\alpha\subset T_2$ for all $\alpha\in \Sigma'$. Hence $Z=T_2$.

Unlike connected component, a point $x$ does not determine a unique connected component.
\end{proof}

\begin{exr}
Show that $\affn^2_\cplx$ is a Noetherian topological space: any decreasing sequence of closed subsets of $\affn^2_\cplx = \spec \cplx[x, y]$ must eventually stabilize. Note that it can take arbitrarily long to stabilize.  Show that $\cplx^2$ with the classical topology is not a Noetherian topological space.
\end{exr}
\begin{proof}
Consider a chain of closed subsets $V(S_0)\supset V(S_1)\supset V(S_2)\supset \cdots $. If it does not stabilize, w.l.o.g, we can assume each inclusion is strict inclusion. Assume there is a strictly descending chain $V(I_0)\supsetneq V(I_1)\supsetneq V(I_2)\supsetneq \cdots$. W.l.o.g. we assume each $I_i$ is radical ideal.

$\forall \scq\supset I_1, \scq\supset I_0$, and $\exists \scp\supset I_0, \scp\notsupset I_1$

$
I_1=\sqrt{I_1}\supset I_0
$
and $I_1\notsubset \sqrt{I_0}=I_0$. $I_1\supsetneq I_0$. We have a strictly ascending chain of radical ideals.

But this is impossible for $\cplx[x,y]$. We intensionally prove this without using Noetherian property of $\cplx[x,y]$. Our proof is basically the special case of Hilbert Basis theorem.

\underline{Claim}: each ideal $\sca$ in $A:=\cplx[x,y]$ is finitely generated.

Pick an element $\sca\ni f=\sum a_{n,m} x^ny^m$. Denote the maximal power of $x$ by $N$ and maximal power of $y$ by $M$.

Construct an $A$-module 
$$
M:=\oplus_{i\leq N,j\leq M}Ax^i y^j
$$
it is a finite dimensional $\cplx$-vector space.
$$
M\cap \sca+(f)=\sca
$$
$\supset$ is trivial.
Consider an element $\sca\ni g=bx^{k_1}y^{l_1}+cx^{k_2}y^{l_2}...,b,c\neq 0$ with $k_i>N$ or $l_i>M$ and $k_1$ is maximal among $k_i$, $l_2$ is maximal among $l_i$
$$
g-\left(\frac{b}{a}x^{k_1-n}y^{l_1-m}+\frac{c}{a}x^{k_2-n}y^{l_2-m}\right)f=0+\text{terms with lower powers in $x$, $y$}.
$$
We can induct on $k_1$, $l_2$ until they finally goes down to $N,M$. $g\in (f)+M\cap\sca$.

Then we claim: if the each ideal is finitely generated, then every ascending chain of ideal stabilize

Let $I_0\subseteq I_1\subseteq I_2\subseteq ...$ \underline{Want:}show that $\exists n_0,I_n=I_{n_0}\forall n\geq n_0$.
Define $I':=\cup_n I_n$. We know that every ideal of ${A}$ is finitely generated. Then Assume $I'$ to be finitely generated by $r$ elements $\{x_1,...,x_r\}$, with $x_j\in I_{n_j}$. Choose $n_0=\text{max}\{n_1,...,n_r\}$, then we have $x_1,...,x_r\in I_{n_0}\Lrta I=I_{n_0}$. $\Lrta I_n=I_{n_0},\forall n\geq n_0$.

$\cplx^2$ under traditional topology is not Noetherian, we can consider the a chain of balls
$$
B_n:=\left\{|x^2|+|y^2|\leq \frac{1}{n^2}\right\},
$$
which gives a strictly descending chain of closed sets.
\end{proof}

\begin{exr}\label{chap3exr:topology_irreducible_components_connected_components}
Show that every connected component of a topological space $X$ is the union of irreducible components of $X$. Show that any subset of $X$ that is simultaneously open and closed must be the union of some of the connected components of $X$. If $X$ is a Noetherian topological space, show that the union of any subset of the connected components of $X$ is always open and closed in $X$.
\end{exr}
\begin{proof}
``Every connected component of a topological space $X$ is the union of irreducible components of $X$:'' By~\ref{chap3exr:irreducible_topological_space_connected}, a irreducible component is connected, hence it is contained in some connected component. But every point in $X$ is contained in some irreducible component. A point $x$ in a connected component $C\subset X$. Every $x$ is contained in some irreducible component $Z_x$~\ref{chap3exr:UNION_OF IRREDUCIBLE_COMPONENTS}. Every connected set containing $x$ is a subset of the connected component $C$. Hence, $Z_x\subset C$ and $C=\cup_{x\in C} Z_x$.

``A clopen subset in $X$ must be some union of connected component'': If $C$ is a clopen subset in $X$, then $C\subset \cup_{x\in C}C_x$, where $C_x$ is the connected component in $X$ that contains $x$. (This inclusion is true for arbitrary subset $C$). For the reverse inclusion, consider an element $y\in C_x$. If $y\notin C$, because $C$ is closed, there is an open neighborhood $U_y$ of $y$ such that $U_y\cap C=\emptyset$. $U_y\cap C_x$ is open, and $A:=\cup_y U_y\cap C_x$ is open, $B:=C_x\cap C$ is also open, and $A\cap B=\emptyset, A\cup B=C_x$, which contradicts the fact that $C_x$ is connected. Then we know every point $y$ of $C_x$ must be contained in $C$. A clopen set is the union of the connected components that intersect $C$.

The reverse statement is not true for general topological space, but when the topological space is in addition Noetherian, we have ``union of arbitrary connected components in $X$ is clopen'': Any union of open set is always open, we only need to prove that $\cup_i C_i$ is closed. Consider $X-C_0\supset X-C_0\cup C_1\supset...$, is a descending chain of closed subsets, it will stabilize after $n_0$, In other words, $\cup_i C_i=\cup_{0\leq i\leq n_0} C_i$. A finite union of closed sets in closed. Done.

\end{proof}

\begin{exr}
Show that a ring $A$ is Noetherian if and only if every ideal of $A$ is finitely generated.
\end{exr}
\begin{proof}
In fact we have,

\textbf{(Lemma)} The following characterizations are equivalent:
\begin{enumerate}[label=(\alph*)]
\item ${A}$ satisfies the \textbf{ascending chain condition on ideals (ACC)} (All the sequence $\sca_1\subseteq \sca_2\subseteq ...$ stabilizes, i.e. $\exists n_0$ s.t. $\sca_n=\sca_{n_0} \forall n\geq 0$)
\item Every ideal of ${A}$ is finitely generated.
\item $\{\text{ideals in ${A}$}\}$ satisfies the \textbf{maximal property}: i.e. Every subset contains a maximal element. That is : For any nonempty collection $S$ of ideals in ${A}$, $\exists \sca\in S$ s.t. $\forall \scb\in S\Lrta \scb\not \supset \sca$
\end{enumerate}

(a)$\Lrta$(b). Let $\sca$ be an ideal. we may assume that $\sca$ is \textbf{NOT} finitely generated. Inductively construct $x_1, x_2,x_3...\in \sca$ such that $(x_1)\neq 0$ and $\sca\supsetneqq(x_1,x_2)\supsetneqq(x_1)$ an also $\sca\supsetneqq(x_1,x_2,x_3)\supsetneqq(x_1,x_2)$, but then this sequence contradict the \textbf{ACC}.

(a)$\Lrta$(c). Let $\emptyset \neq S\subseteq \{\text{ideals in }{A}\}$. If $S$ violates the maximal property, then start from arbitrary ideal $\sca_1$, we can find $\sca_1\subsetneq \sca_2\in S$. Similarly, we can find $\sca_{j+1}\supsetneq \sca_j,\forall j\in\intg_{\geq 0}$ by the countable choice axiom. Then the ACC fails.

(c)$\Lrta $(a). If ACC fails, $\exists \sca_1\subsetneq \sca_2\subsetneq ...$ Take $S:=\{\sca_1,\sca_2,\sca_3...\}$. Then $S$ violates maximal property.

(b)$\Lrta$(a). Let $\sca_1\subseteq\sca_2\subseteq ...$ \underline{Want:}show that $\exists n_0,\sca_n=\sca_{n_0}\forall n\geq n_0$.
Define $\sca:=\cup_n \sca_n$. We know that every ideal of ${A}$ is finitely generated. Then $\sca$ is also finitely generated by assumption (b). Then Assume it to be finitely generated by $r$ elements $\{x_1,...,x_r\}$, with $x_j\in\sca_{n_j}$. Choose $n_0=\text{max}\{n_1,...,n_r\}$, then we have $x_1,...,x_r\in \sca_{n_0}\Lrta \sca=\sca_{n_0}$. $\Lrta \sca_n=\sca_{n_0},\forall n\geq n_0$.
\end{proof}

\begin{exr}\label{chap3exr:Noetherian_rings_Noetherian_topological_spaces}
If $A$ is Noetherian, show that $\spec A$ is a Noetherian topological space. Describe a ring $A$ such that $\spec A$ is not a Noetherian topological space.
\end{exr}
\begin{proof}
Let $(V_i)_{i\in \mathbb{N}}$ be a descending chain of closed subsets of $\spec A$. For every $i\in\mathbb{N}$, let $I_i\in A$ be an ideal such that $V(I_i)=V_i$. We conclude that $V(I_i)\supset V(I_{i+1})$ for all $i\in \mathbb{N}$. W.l.o.g. we assume each $I_i$ is radical ideal.

$\forall \scq\supset I_{i+1}, \scq\supset I_i$

$
I_{i+1}=\sqrt{I_{i+1}}\supset I_i.
$
We have an ascending chain of radical ideals.

By the Noetherian property of $A$, this chain of ideals must stabilize. Hence the descending chain of closed subsets also stabilize, therefore $\spec A$ is a Noetherian topological space.

The converse is not true. Consider the ring $A:=k[x_1,x_2,...]/(x^2_1,x^2_2,...)$ for a field $k$. Let $\scp\subset A$ be prime idal. Then the ideal $(x_1,x_2,...)$ is contained in $\scp$. But the ideal $(x_1,x_2,...)$ is already maximal. Hence $\spec A$ is only one point and trivially Noetherian topological space. But $A$ is not Noetherian because $(x_1,x_2,..)$ is not finitely generated.

As for an example of non-Noetherian topological space, we can consider the ring $A:=k[x_1,x_2,...]$. $A$ is not a Noetherian ring and there is a strictly ascending chain of ideals $(x_1)\subsetneq (x_1,x_2)\subsetneq ...$. It give rise to a strictly descending chain of closed subsets
$$
V(x_1)\supsetneq V(x_1,x_2)\supsetneq\cdots
$$
\end{proof}
\begin{exr}\label{chap3exr:open_subset_Noetherian_space_quasicompact}
Show that every open subset of a Noetherian topological space is quasicompact. Hence if $A$ is Noetherian, every open subset of $\spec A$ is quasicompact.
\end{exr}
\begin{proof}
Let $Y\subset X$ be an open subset of $X$ with the induced topology on $Y$, Then $Y$ is also Noetherian. Consider an open covering $\cup_{i\in I}V_i=Y$ of $Y$. {\color{red}is it true that every cover has countable subcover?}. Without loss of generality, we assume $I=\mathbb{N}$ and define
$$
U_n=\bigcup_{0\leq i\leq n} V_i.
$$
$U_n$ form an ascending chain of open subsets. Since Y is Noetherian, we conclude the chain stabilizes
and we find the finite subcover.
\end{proof}

\begin{remark}
A closed subset of Noetherian space is always quasicompact, because a Noetherian space is itself quasicompact.
\end{remark}

\begin{exr}\label{chap3exr:submodule_Notherian}
Show that if $ M$ is a Noetherian $A$-module, then any submodule of $M$ is a finitely generated $A$-module.
\end{exr}
\begin{proof}
The proof is just identical to the different characterization of Noetherian rings.

If the submodule $M'$ of $M$ is not finitely generated, then we can inductively construct $M'\supsetneq A x_1$, $M'\supsetneq A x_1\oplus Ax_2...$. But this violate the ascending chain condition on $M$.
\end{proof}

\begin{exr}\label{chap3exr:Noetherian_module_SES}
If $0\lrta M' \lrta M\lrta M'' \lrta 0$ is exact, show that $M'$ and $M''$ are Noetherian if and only if $M$ is Noetherian.
\end{exr}
\begin{proof}
``$\Lrta$'' Use ACC. Let $N_1\subseteq N_2\subseteq ...$ be submodules of $M$. \underline{Want:}  show that $\exists n_0 :(n\geq n_0)\Lrta N_n=N_{n_0}$.
Consider $N''_j:=$ Image of $N_j$ in $M''$. $ N_1''\subseteq N_2''\subseteq ...$ By ACC of $M''$, $N''_{n_0}=N''_{n}\forall n\geq n_0$. Do the same for $N'_j:=M'\cap N_j$ ($M'\inj M$)\\
Need: if $N_i\subseteq N_j\subseteq M$ and $N_i''=N''_j, N'_i=N'_j$, then $N_i=N_j$. (Five Lemma)
\begin{center}
\begin{tikzcd}
0\ar[r] & N_i'\ar[r]\ar[d,equal] & N_i\ar[r]\ar[d] & N_i''\ar[r]\ar[d,equal] & 0\\
0\ar[r] & N_j'\ar[r] & N_j\ar[r] & N_j''\ar[r] & 0
\end{tikzcd}
\end{center}

For the $\Llta$ direction, $M$ is Noetherian. Because every ascending chain in $M'$ can be identified as ascending chain in $M$, $M'$ is Noetherian. Also for every ascending chain of submodules in $M''$, it's preimage is an ascending chain of submodule in $M$, hence the preimage chain will stabilize, hence the chain in $M''$ would also stabilize.
\end{proof}

\begin{exr}
Show that if $A$ is a Noetherian ring, then $A^{\oplus n}$ is a Noetherian
$A$-module.
\end{exr}
\begin{proof}
When we regard $A$ as an $A$-module, the submodule of $A$ is just the ideals in $A$. Hence $A$ is a Noetherian ring iff $A$ is a Noetherian $A$-module. Then $A\oplus A$ is Noetherian $A$-module by~\ref{chap3exr:Noetherian_module_SES} because
$$
0\lrta A\lrta A\oplus A\lrta A\lrta 0
$$
and then we can induct on $n$, by 
$$
0\lrta A^{\oplus{n-1}}\lrta A^{\oplus n}\lrta A\lrta 0
$$
\end{proof}

\begin{exr}
Show that if $A$ is a Noetherian ring and M is a finitely generated $A$-module, then $M$ is a Noetherian module. Hence by Exercise~\ref{chap3exr:submodule_Notherian}, any submodule of a finitely generated module over a Noetherian ring is finitely generated.
\end{exr}
\begin{proof}
Suppose $M$ is generated by $\{x_1,...,x_n\}$
We always have a SES
$$
0\lrta Ker(\varphi)\lrta {A}^n\overset{\varphi}{\lrta} M\lrta 0,
$$
then apply~\ref{chap3exr:Noetherian_module_SES}.
\end{proof}
\section{The function $I(\cdot),$ taking subsets of $\spec A$ to ideals of $A$}

\begin{exr}
Let $A = k[x, y]$. If $S = \{[(y)], [(x, y - 1)]\} $ , then $I(S)$ consists of those polynomials vanishing on the $y$-axis, and at the point $(0, 1)$. Give generators for this ideal.
\end{exr}
\begin{proof}
By definition, $I(S)=\cap_{[\scp]\in S}\scp= (y)\cap (x,y-1)$. Hence $I(S)$ can be identified as polynomials in $k[x,y]$ that vanishes on the both $y$-axis and $(0,1)$. The generators of this ideal are $\{xy,y(y-1)\}$
\end{proof}

\begin{exr}
Suppose $S \subset \affn^3_\cplx$ is the union of the three axes. Give generators for the ideal $I(S)$. We will see in Chapter 12 that this ideal is not generated by less than three elements.
\end{exr}
\begin{proof}
$S$ geometrically is the union of three axes. $I(S)$ is the polynomial that vanishes on all the three axis. The points in $S$ are of the form $[(0)],[(x,y)],[(y,z)], [(x,z)],[(x-a,y,z)],[(x,y-b,z)],[(x,y,z-c)]$. The union of corresponding ideals would be $(x,y)\cap (y,z)\cap (x,z)$. The ideal is finitely generated by $(xy,yz,zx)$.
The inclusion $(x,y)\cap (y,z)\cap (x,z)\supset (xy,yz,zx)$ is trivial. For the reverse inclusion. Consider an element $g\in (x,y)\cap (y,z)\cap (x,z)$, $g$ can't contain terms like $x,y,z,x^n,y^n,z^n$ because they are not element in the intersection. Then the  lowest degree terms of $g$ should be $xy,yz,xz$, and all other terms can be generated by $xy,yz,zx$.
\end{proof}

\begin{exr}
Show that $V(I(S)) = \overline{S}$. Hence $V(I(S)) = S$ for a closed set $S$.
\end{exr}
\begin{proof}
``$V(I(S))\supset \overline{S}$'':
By definition, $V(I(S))$ is closed set. And $S\subset V(I(S))$, because $[\scp]\in S, I(S)=\cap_{[\scq]\in S}\scq$ $I(S)\mod \scp=0,\Lrta I(S)\subset \scp\Lrta [\scp]\in V(I(S))$. Closure $\overline{S}$ is the smallest closed set that contains $S$, therefore $\overline{S}\subset V(I(S))$.

 ``$V(I(S))\subset \overline{S}$'': We need to verify that open neighborhood of each point in $V(I(S))$ intersects with $S$. Consider a point $[\scq]\in V(I(S))$, we have $\scq\supset I(S)$. Assume  a non-empty distinguished open $D(f)$ contains $[\scq]\Llrta f\notin\scq \Lrta f\notin I(S)\Lrta$ $f $ does not vanish on every point in $S$, $\exists[\scn]\in S$, s.t. $f\notin [\scn]\Llrta$ $D(f)\ni [\scn]$, therefore $D(f)\cap V(I(S))\neq \emptyset$. Hence $V(I(S))\subset \overline{S}$.
\end{proof}

\begin{exr}
Prove that if $J \subset A$ is an ideal, then $I(V(J)) = \sqrt J$.
\end{exr}
\begin{proof}
Suppose $J\subset A$ is an ideal. By definition, $I(V(J))$ is the set of functions that vanish on $V(J)$, 
$
I(V(J))=\cap_{[\scp]\in V(J)}\scp.
$
Also recall that $[\scp]\in V(J)\Llrta\scp\supset J$, we have
$$
I(V(J))=\cap_{[\scp]\in V(J)}\scp=\cap_{\scp\supset J}\scp=\sqrt J.
$$
\end{proof}

\begin{exr}\label{chap3exr:bijection_closed_subsets_points_affine}
Show that $V(\cdot)$ and $I(\cdot)$ give a bijection between irreducible closed subsets of $\spec A$ and prime ideals of $A$. From this conclude that in $\spec A$ there is a bijection between points of $\spec A$ and irreducible closed subsets of $\spec A$ (where a point determines an irreducible closed subset by taking the closure). Hence each irreducible closed subset of $\spec A$ has precisely one generic point  any irreducible closed subset $Z$ can be written uniquely as $\overline{\{z\}}$.
\end{exr}
\begin{proof}
``$S$ irreducible closed $\Lrta$ $I(S)$ is prime'':

Assume $I(S)$ is not prime. $\Llrta \exists x,y\notin I(S)$ but $xy\in I(S)$, which means $x,y$ each does not vanish on every point in $S$ but $xy$ vanishes on each point in $S$.
$$
V((x))\notsupset S
$$
$$
V((y))\notsupset S
$$
but 
$$
V((xy))=V((x))\cup V((y))\supset S
$$
Then 
$$
S=(V((x))\cap S)\cup (V((y))\cap S),
$$
where $(V((x))\cap S)$ and $ (V((y))\cap S)$ are non-empty and closed. Contradiction.


``$\scp$ is prime $\Lrta V(\scp)$ is irreducible closed.''

Assume $V(\scp)$ is reducible, it can be written as union of two non-empty closed subsets, $V(\scp)=W\cup Z$. Assume $\sca,\scb$ radical ideals such that $W=V(\sca)$, $Z=V(\scb)$, we have $\sca\supsetneq \scp$ and $\scb\supsetneq \scp$
$$
V(\sca)\cup V(\scb)=V(\sca\scb)=V(\scp)
$$
By the theorem about bijection between closed sets and radical ideals in $A$.
$$
\sca\scb\subset\sca\cap\scb=\sqrt{\sca\scb}=\scp.
$$
But $\sca\scb\subset \scp$ implies $\sca\subset \scp$ and $\sca\subset\scp$. Then we find elements $x\in \sca\backslash\scp$, $y\in \scb\backslash\scp$, but $xy\in\scp$. Contradiction.


This bijection is also inclusion reversing (just consider them as ordinary closed subsets and radical ideals).


\end{proof}


\begin{exr}\label{chap3exr:irreducible_components_minimal_pirmes_bij}
A prime ideal of a ring $A$ is a \textbf{minimal prime ideal} (or more simply, \textbf{minimal prime}) if it is minimal with respect to inclusion. (For example, the only minimal prime of $k[x, y] $ is $(0).$) If $A$ is any ring, show that the irreducible components of $\spec A$ are in bijection with the minimal prime ideals of $A$. In particular, $\spec A$ is irreducible if and only if $A$ has only one minimal prime ideal; this generalizes Exercise~\ref{chap3exr:Integral_domain_irreducible}.
\end{exr}
\begin{proof}
$\scn$ is minimal prime\\
$\Llrta$ there is no strictly smaller prime $\scp\subset\scn$.\\
$\Llrta$ There is no strictly bigger irreducible closed subset that contains $V(\scn)$. By the bijection in~\ref{chap3exr:bijection_closed_subsets_points_affine}. \\
$\Llrta$ $V(\scn)$ is maximal irreducible closed subset (irreducible component)

In particular $\spec A$ is irreducible iff $\spec A$ is itself the unique irreducible component. Then $I(\spec A)$ is prime and contained in any other irreducible sets.
\end{proof}

\begin{exr}
What are the minimal prime ideals of $A:=k[x, y]/(xy)$ (where $k$ is a field)?
\end{exr}
\begin{proof}
Geometrically, three are two irreducible components of $\spec A$, the $x$-axis and $y$-axis, corresponding to $(y)$ and $(x)$ respectively.
\end{proof}

\chapterimage{structure_sheaf.jpg}
\chapter{The structure sheaf, definition of schemes}
\section{The structure sheaf of an affine scheme
}
\begin{definition}
Define $\calo_{\spec A}(D(f))$ to be localization of $A$ at the multiplicative set $S$, where
$$
S:=\{\text{All functions that do not vanish outside $V(f)$ (Do not vanish on $D(f))$}\}.
$$
(i.e., those $g\in A$ such that $V(g)\subset V(f)$ or equivalently $D(f)\subset D(g)$)
\end{definition}
In particular, $\calo_{\spec A}(\emptyset)=\{0\}$, where localize at the multiplicative set of functions $g$ such that $V(g)\subset \spec A$. This multiplicative set includes $0$, hence the localization is $\{0\}$ ring.

\begin{exr}\label{chap4exr:4.1.A}
Show that the natural map $A_f \lrta \calo_{\spec A}(D(f))$ is an isomorphism.
\end{exr}
\begin{proof}
In particular, $S_f:=\{1,f,f^2,...\}$ is a multiplicative subset of the multiplicative set $T$ in the definition of $\calo_{\spec A}(D(f))$, where 
$$
T:=\{\text{All functions that do not vanish outside $V(f)$ (Do not vanish on $D(f))$}\}.
$$
$S_f\subset T$.
There is a natural homomorphism 
$$
A\overset{S_f^{-1}}{\lrta}A_f\overset{\tilde{T}^{-1}}{\lrta}\calo_{\spec A}(D(f)), 
$$
where we have denoted the image of $T$ in $A_f$ by $\tilde{T}$. $g\in S\Llrta D(f)\subset D(g)$

$\Llrta T^{-1}g$ is invertible in $A_f$ by Exercise~\ref{chap3exr:inclusion_distinguished_open}.

$\Llrta$ $\tilde{T}\subseteq A_f^\times$

$\Llrta$ $\tilde{T}^{-1}$ is an isomorphism. $A_f\cong \calo_{\spec A}(D(f))$.
\end{proof}

\begin{exr}
Prove the base identity axiom for any distinguished open $D(f)$. 
\end{exr}
\begin{proof}
Consider the $D(f)=\cup_{i\in I}D(f_i)$. We already showed that $\spec A_f\cong D(f)$ as topological spaces~\ref{chap3exr:specA_f_and_specA/I}. If $D(f)=\cup_{i\in I}D(f_i)=\cup_{i\in I}D(f_i)\cap D(f)=\cup_{i\in I}D(f_if)$.

$D(f_i f)\cong\spec A_{ff_i}$ $A_{ff_i}$ is the localization of $A_f$ at the image of $f_i$. $D(f_i f)$ corresponds to the point $[\scq]\in \spec A_f$ such that $\scq\not\in \frac{f_i}{1}$.

Then $D(f)=\cup_{i\in I}D(f_i)\subset \spec A$ $\Llrta$ $\spec A_f=\cup_{i\in I} D(f_i/1)$

$\calo_{\spec A}(D(f))\cong A_f=\calo_{\spec A_f}(\spec A_f)$. The function restricts to $0$ on each $D(f_i)$ iff its restriction to $D(f_i/1)$ vanishes.

Then the problem reduces to the proved case $D(f)=\spec A$.
\end{proof}

\begin{exr}
Alter this argument appropriately to show base gluability for any distinguished open $D(f)$.
\end{exr}
\begin{proof}
Again, we regard $D(f)\cong \spec A_f$. 

Then $D(f)=\cup_{i\in I}D(f_i)\subset \spec A$ $\Llrta$ $\spec A_f=\cup_{i\in I} D(f_i/1)$.

$\calo_{\spec A}(D(f))\cong A_f=\calo_{\spec A_f}(\spec A_f)$.

The base gluability follows from the special case we have proved for $\spec A=D(f)$.
\end{proof}
\begin{exr}
Suppose $M$ is an $A$-module. Show that the following construction describes a sheaf $M$  on the distinguished base. Define $ \tilde{M}(D(f))$ to be the localization of $M$ at the multiplicative set of all functions that do not vanish outside of $ V(f)$. Define restriction maps $\res_{D(f),D(g)}$ in the analogous way to $\calo_{\spec A}$. Show that this defines a sheaf on the distinguished base, and hence a sheaf on $\spec A$. Then show that this is an $\calo_{\spec A}$-module.
\end{exr}
\begin{proof}
Define $\tilde{M}_{\spec A}(D(f))$ to be localization of $M$ at the multiplicative set $S$, where
$$
S:=\{\text{All functions that do not vanish outside $V(f)$ (Do not vanish on $D(f)$}\}.
$$

\underline{Claim}: $\tilde{M}(D(f))\cong M_f$.

 In particular, $S_f:=\{1,f,f^2,...\}$ is a multiplicative subset of the multiplicative set $S$.

 There is a natural homomorphism 
$$
M\overset{S_f^{-1}}{\lrta}M_f\overset{\tilde{S}^{-1}}{\lrta}\tilde{M}(D(f)), 
$$
where we have denoted the image of $S$ in $A_f$ by $\tilde{S}$. $g\in S\Llrta D(f)\subset D(g)$

$\Llrta T^{-1}g$ is invertible in $A_f$ by Exercise~\ref{chap3exr:inclusion_distinguished_open}.

$\Llrta$ $\tilde{S}\subseteq A_f^\times$

$\Llrta$ $\tilde{S}^{-1}$ is an isomorphism. $M_f\cong \tilde{M}(D(f))$.

Define the restriction map:
 $$
 \res_{D(f),D(g)}: \tilde{M}(D(f))\lrta \tilde{M}(D(g))
 $$
 it is the further localization of modules. The restriction does generate a presheaf on distinguished base because of the functorial property of localization.

 Base identity axiom: Consider the special case where $D(f)=\spec A$. Then $\spec A=\cup_{i\in I} D(f_i)$, there is a finite subcover, say $\cup_{i=1,...,n}D(f_i)=\spec A$, i.e., $(f_1,...,f_n)=A$.

 Suppose $m\in M=\tilde{M}(\spec A)$ such that $\res_{\spec A,D(f_i)} m=0$. Then $f_i^{l_i}m=0,\forall i$. We can choose $N\geq \max \{l_i\}$ such that $f_i^N m=0$. Also $(f_1^N,...,f_n^N)=A$ ($\cup_{i}D(f_i^N)=A$), then we know exist $r_i\in A$ s.t. $\sum_i r_i f_i^N=1$
 $$
 m=\sum_i(r_i f_i^N)m=0.
 $$ 
 The restriction map is injective, hence the base identity axiom holds for $D(f)=\spec A$.

 As for general $D(f)$, we can replace $A$ by $A_f$ so that $D(f)=\spec A_f$. The problem reduces to what we have proved.

 Base gluability axiom: Suppose again $\cup_{i\in I}D(f_i)=\spec A$. We assume in addition $I$ is a finite index set. We have elements
$$
\frac{m_i}{f_i^{l_i}}\in M_{f_i}=\tilde{M}(D(f_i))
$$
Letting $g_i=f_i^{l_i}$ to simplify the notation.
$m_i/g_i$ and $m_j/g_j$ agree  when restricted to intersection $D(f_i)\cap D(f_j)=D(g_i)\cap D(g_j)=D(g_ig_j)$, which means
$$
(g_ig_j)^{n_{ij}}(g_jm_i-g_im_j)=0\in A.
$$
By setting $N=\max\{n_{ij}\}$, we have
$$
(g_ig_j)^{N}(g_jm_i-g_im_j)=0\in A.
$$
Set $b_i=m_ig_i^N$ and $h_i=g_j^{N+1}$, the overlap condition simplifies to 
$$
b_ih_j-b_jh_i=0.
$$
Note that $\cup_i D(f_i)=\cup_i D(h_i)=\spec A$ implies that $\exists r_i\in A$ s.t., $\sum_i r_i h_i=1$. Define 
$$
r=\sum_i r_i b_i.
$$
Then 
$$
rh_j=\sum_i r_i b_i h_j=\sum_i r_i  h_i b_j=b_j
$$
$$
 g_j^{N}(g_jr-m_j)=g_j^{N}(f^{l_j}_jr-m_j)=0
$$
$$
\res_{\spec A,D(f_j)}r=\frac{r}{1}=\frac{m_j}{f_j^{l_j}}\in \tilde{M}(D(f_j)),\forall j
$$

For infinite index set $I$, choose a finite subset $\{1,..,n\}\subset I$ with $\spec A=\cup_i D(f_i)$, we can construct $r$ as above. We will show that for $z\in I-\{1,...,n\}$, $r$ restricts to desired element $m_z\in M_{f_z}$. Repeat the entire process for $\{1,...,n,z\}$ in place of $\{1,..,n\}$ and we get a $r'\in M$ which restricts to $m_i\in M_{f_i}$ for $i=1,..,n$. Then by identity axiom $r=r'$. Hence $r$ restricts to $m_z$ as desired.

We have checked the base gluability in the special case $D(f)=\spec A$. For general $D(f)$, we can replace $A$ with $A_f$ so that $D(f)=\spec A_f$ and the problem reduces to the special case we have solved above.

We have up to now proved that $\tilde{M}$ is a sheaf on the distinguished base. We can recover form it a sheaf on $\spec A$.

Observe that $\tilde{M}(D(f))=M_f$ which is automatically an $A_f$-module ($\calo_{\spec A}$-module). $\tilde{M}$ is $\calo_{\spec A}(D(f))$-module on the distinguished base. And this would also extend to a $\calo_{\spec A}$-module.
\end{proof}





\section{Visualizing schemes II: nilpotents}
\section{Definition of schemes
}

\begin{exr}
Describe a bijection between the isomorphisms $\spec A \lrta \spec A'$ and the ring isomorphisms $A' \lrta A$. Hint: the hardest part is to show that if an isomorphism $\pi: \spec A\lrta \spec A'$ induces an isomorphism $\pi^\#:A' \lrta A$, which in turn induces an isomorphism $\rho: \spec A \lrta  \spec A'$, then $\pi = \rho$. First show this on the level of points; this is (surprisingly) the trickiest part. Then show $\pi = \rho$ as maps of topological spaces. Finally, to show $\pi = \rho$ on the level of structure sheaves, use the distinguished base. Feel free to use insights from later in this section, but be careful to avoid circular arguments. Even struggling with this exercise and failing (until reading later sections) will be helpful.
\end{exr}

\begin{proof}
For simplicity, we replace the notion of topological space $\spec A,\spec A'$ by $X,X'$ respectively. 
Recall what it means by an isomorphism of schemes. An isomorphism of ringed space
$$
\pi:X\cong X' (\text{homeomorphism})
$$
and 
$$
\pi_*\calo_{X}\cong \calo_{X'} (\text{sheaf isomorphism}).
$$

$A=\calo_X(\pi^{-1}(X'))=\pi_*\calo_{X}(X')=\calo_{X'}(X')=A'$. It induces an morphism
$$
\pi^\#:A'\lrta A
$$
$$
\calo_{X'}(X')\ni a'\mapsto a'\in \pi_*\calo_X(X'). 
$$
This $\pi^\#$ in turn would induces a isomorphism $\rho:\spec A\lrta \spec A'$. 
$$
\rho:[\scp]\mapsto [(\pi^\#)^{-1}\scp]
$$

We need to show that $\pi=\rho$. 
$$
\pi:[\scp]\mapsto ?
$$
TO BE ADDED
\end{proof}

\begin{exr}\label{chap4exr:iso_scheme_sidginguished_open}
Suppose $f \in A$. Show that under the identification of $D(f)$ in $\spec A$ with $\spec A_f$, there is a natural isomorphism of ringed spaces $$(D(f), \calo_{\spec A}|_{D(f)}) \cong (\spec A_f, \calo_{\spec A_f} ).$$ Hint: notice that distinguished open sets of $\spec A_f$ are already distinguished open sets in $\spec A$.
\end{exr}
\begin{proof}
By exercise~\ref{chap3exr:specA_f_and_specA/I}, we know $\spec A_f\subset \spec A$ and the Zariski topology on $\spec A_f$ is the subspace topology induced by Zariski topology of $\spec A$. Hence $D(f)\cong \spec A_f$ as topological spaces. In particular, the distinguished base of $\spec A_f$ is induced by  distinguished base of $\spec A$ intersects with $D(f)$. Specifically,
$$
D(f)\cap D(f_i)\cong D(f_i/1)\subset \spec A_f
$$
 On the other hand, by exercise~\ref{chap4exr:4.1.A}, We know $A_f\cong \calo_{\spec A}(D(f)).$ We need to verify that $\calo_{\spec A_f}\cong \calo_{\spec A}|_{D(f)}$
 $$
 \begin{aligned}
 \calo_{\spec A}|_{D(f)}(D(f)\cap D(g))&=\calo_{\spec A}(D(f)\cap D(g))=\calo_{\spec A}(D(fg))\cong A_{fg}\\
 &=(A_f)_{g/1}=\calo_{\spec A_f}(D(g/1)).
 \end{aligned}
 $$
 Because $f$ is invertible in $A_f$. $D(g/f^n)=D(g/1)\cap D(1/f^n)=D(g/1)\cap D(1/1)=D(g/1)$. ( All distinguished open in $\spec A_f$ are of the form.) We know the two sheaves coincide on the distinguished base, which means they are isomorphic.
\end{proof}

\begin{exr}\label{chap4exr:open_subscheme}
If $X$ is a scheme, and $U$ is any open subset, prove that $(U, \calo_X|_U)$ is also a scheme.
\end{exr}
\begin{proof}
Recall the definition of scheme, at any point $x\in X$, there is an open neighborhood $V$ such that $(V,\calo_X|_V)$ is an affine scheme. Assume $(V,\calo_X|_V)\cong (\spec A,\calo_{\spec A})$. This isomorphism says the subset topology on $V$ coincides with the Zariski topology. $U\cap V$ is also an open set in $\spec A$. 

\underline{Claim}: There exists $x\in D(f)\subset U\cap V$ such that $(D(f),\calo_X|_{D(f)})$ is an affine scheme.

Such $D(f)$ always exists because distinguished open form a base of Zariski topology and $D(f)$ can be regarded as open set in $X$. On the other hand, $\calo_X|_{D(f)}=(\calo_X|_U)|_{D(f)}=(\calo_X|_V)|_{D(f)}\cong (\calo_{\spec A})|_{D(f)}\cong \calo_{\spec A_f}$. The last isomorphism by~\ref{chap4exr:iso_scheme_sidginguished_open} 

In $(U,\calo_X|_U)$, there is an open $U\supset D(f)\ni x$ such that
$$
(D(f),(\calo_{X}|_U)|_{D(f)})\cong (\spec A_f,\calo_{\spec A_f}).
$$
which means $(U,\calo_X|_U)$ is also a scheme.
\end{proof}

\begin{exr}
Show that if $X$ is a scheme, then the affine open sets form a base for the Zariski topology.
\end{exr}
\begin{proof}
$x\in X$, and $U$ is an open neighborhood of $x$. We need to verify that there is an affine open set $x\in W\subset U$. 

In Exercise~\ref{chap4exr:open_subscheme}, we proved that $(U,\calo_X(U))$ is open subscheme of $(X,\calo_X)$. Then for each $x\in U$, there is an affine open set $W$ \textbf{in $U$}. But the Zariski topology of $U$ is just the subset topology induced by the Zariski topology of $X$. We know $W$ is also a open set in $X$. Finally, $\calo_X|_W=(\calo_X|_U)|_W$. $W$ is an affine open in $X$ iff $W$ is an affine open in X.  
\end{proof}

\begin{exr}
The disjoint union of schemes is defined as you would expect: it is the disjoint union of sets, with the expected topology, with the expected sheaf.
\begin{enumerate}[label=(\alph*)]
\item Show that the disjoint union of a finite number of affine schemes is also an affine scheme.
\item (a first example of a non-affine scheme) Show that an infinite disjoint union of (nonempty) affine schemes is not an affine scheme. 
\end{enumerate}
\end{exr}
\begin{proof}
\begin{enumerate}[label=(\alph*)]
\item In Exercise~\ref{exr:nonconnected_scheme}, we see that for finite index set $I$:
$$
\coprod_{i\in I}\spec A_i\cong\spec \prod_i A_i
$$
and we need to describe the structure sheaf and verify that
$$
\calo_{\coprod_i \spec A_i}\cong \calo_{\spec \prod_i A_i}.
$$

 Consider the inclusion map $\iota_i: \spec A_i\inj \coprod_i\spec  A_i$
$$
\calo_{\coprod_i\spec  A_i}:= \prod_i (\iota_{i})_*\calo_{\spec A_i}
$$
For $U=\coprod_i U_i\subset \coprod \spec A_i$,
$$
\left(\prod_i (\iota_{i})_* \calo_{\spec A_i}\right)(U)=\prod_i \calo_{\spec A_i}(\iota_i^{-1}U)=\prod_{i}\calo_{\spec A_i}(U_i).
$$ 
On the other hand, the structure sheaf $\calo_{\spec \prod_i A_i}$ has the down to earth definition. We need to verify that they agree on $D(f)$ for $f=(f_1,..f_i,..)\in \prod_i A_i$. It is not hard to check 
$$
D(f)=\coprod_i D(f_i).
$$
$$
\calo_{\spec \prod_i A_i}(D(f))=(\prod_i A_i)_f\overset{?}{=}\prod_i (A_i)_{f_i}=\prod_i\calo_{\spec A_i}(D(f_i))=\calo_{\coprod_i \spec A_i}(D(f))
$$
The equality under question mark need to be checked. We just verify the universal property directly.
\begin{center}
\begin{tikzcd}
R &  &  \\
 & (A\times B)_{(f,g)} \arrow[lu, "\exists u", dashed] & B_g \arrow[l, "\beta"] \arrow[llu, "\phi"'] \\
 & A_f \arrow[u, "\alpha"'] \arrow[luu, "\psi"] & 
\end{tikzcd}
\end{center}
where 
$$
\alpha:\frac{a}{f^n}\mapsto \frac{(a,g^n)}{(f,g)^n}
$$
$$
\beta:\frac{b}{g^m}\mapsto \frac{(f^m,b)}{(f,g)^m}
$$
Given $\phi$ and $\psi$, there is a unique morphism $u$ to make diagram commute.
$$
u:\frac{(a,b)}{(f,g)^k}\mapsto \psi\left(\frac{a}{f^k}\right)\phi\left(\frac{b}{g^k}\right)
$$
\item Exercise~\ref{chap3exr:affine_scheme_quasicompact} shows that the topological space of affine scheme is quasicompact but the infinite disjoint union of affine scheme has infinite many connected components hence can't be quasicompact. For example the cover $\cup_i \spec A_i$ does not have finite subcover.
\end{enumerate}
\end{proof}

\begin{exr}
Show that the stalk of $\calo_{\spec A}$ at the point $[\scp]$ is the local ring $A_\scp$.
\end{exr}
\begin{proof}
$$
\calo_{\spec A,[\scp]}=\dirlim_{[\scp]\in U}\calo_{\spec A}(U)
$$
without loss of generality, we can assume each $U$ is a distinguished open $D(f)$ for some $f$. $D(f)\ni [\scp]$ iff $f\notin \scp$ or equivalently $f\in S_\scp=A-\scp$. 

\underline{Claim}: 
$$
\calo_{\spec A,[\scp]}=\dirlim_{[\scp]\in D(f)}\calo_{\spec A}(D(f))=A_\scp,
$$
The index category is now a subset of distinguished base which contain $[\scp]$ and the morphism are inclusions.
\begin{center}
\begin{tikzcd}
A_\scp \arrow[rrr, dashed] &  &  & R \\
 &  &  &  \\
 & \calo_{\spec A}(D(f))=A_f \arrow[r] \arrow[rruu, "\phi_f"] \arrow[luu, "l_f"] & \calo_{\spec A}(D(g))=A_g \arrow[ruu, "\phi_g"] \arrow[lluu, "l_g"'] & 
\end{tikzcd}
\end{center}
If there is a ring $R$ with all morphisms from $A_f$ to it. The image of all $f\in S_\scp$ is invertible in $R$, equivalently, $S_{\scp}\subset R^\times$. Then there is a morphism from $A$ to $R$ by universal property of localization.

Equivalently we have general formula
$$
S^{-1}M=\dirlim_{f\in S}M_f.
$$
The index category is now the partially ordered set $S$, with $f\geq g$ iff $f\in \sqrt{(g)}$, see~\ref{chap3exr:inclusion_distinguished_open} 
\end{proof}
\begin{exr}\label{chap3exr:locally_ringed_space_function_value}\ 
\begin{enumerate}[label=(\alph*)]
\item If $f$ is a function on a locally ringed space $X$, show that the subset of $X$ where $f$ vanishes is closed. (Hint: show that if $f$ is a function on a ringed space $X$, show that the subset of $X$ where the germ of $f$ is invertible is open.)
\item Show that if $f$ is a function on a locally ringed space that vanishes nowhere, then $f$ is invertible.
\end{enumerate}
\end{exr}
\begin{proof}
\begin{enumerate}[label=(\alph*)]
\item $f$ is a function on a ringed space $X$ ($f\in\Gamma(X,\calo_X)$). A germ of $f$ at $p$ is denoted $f_p$, assume $f_p$ is invertible. $\exists g_p\in\calo_{X,p}$ such that $f_pg_p=1\in\calo_{X,p}$, which means $\exists U\ni p$ open such that $(f|_U)(g|_U)=1\in \calo_X(U)$. Then all $f_q$ is invertible for $q\in U$. This means the set in $X$ where germ of $f$ is invertible is open.  In the case of locally ringed space, the complement of subset where $f$ vanishes is just the set where germ of $f$ is invertible. (We don't have good notion of function vanishing on general ringed space.)
\item $f$ vanishes nowhere, then the germ of $f$ is invertible everywhere. Notice that $\{f_p\}_{p\in X}$ consists of compatible germs, then their inverse $\{g_p\}_{p\in X}$ also consists of compatible germs. Check it by hand: for each $g_p,p\in X$, there is an open $V_p\subset X$, and  $\tilde{g}_p\in \calo_X(V_p)$ such that $g_q=\tilde{g}_p|q$. We can just choose $\tilde{g}_p$ to be the representative in the equality $f_pg_p=1$. And by Exercise~\ref{chap2:exr2_compatible germs_glue_to_section}. There is a unique global section to glue to and we name it $g\in\calo_X(X)$. Finally, $fg=1$, because $f_pg_p=1_p=1,\forall p\in X$ by~\ref{chap2exr2.4.A}.
\end{enumerate}
\end{proof}
\section{Three examples}
\begin{exr}
Show that you can glue an arbitrary collection of schemes together. Suppose we are given:
\begin{itemize}
\item  schemes $X_i$ (as $i$ runs over some index set $I$, not necessarily finite), 
\item open subschemes $X_{ij} \subset  X_i$ with $X_{ii} = X_i$,
\item isomorphisms $f_{ij} : X_{ij} \lrta X_{ji}$ with $f_{ii}$ the identity
such that
\item  (the cocycle condition) the isomorphisms ``agree on triple intersections'', i.e., $f_{ik}|_{X_{ij}\cap X_{ik}} = f_{jk}|_{X_{ji}\cap X_{jk}} \circ  f_{ij}|_{X_{ij}\cap X_{ik}}$ (so implicitly, to make sense of the right side, $f_{ij}(X_{ik} \cap X_{ij}) \subset  X_{jk}$).
\end{itemize}
(The cocycle condition ensures that $f_{ij}$ and $f_{ji}$ are inverses. In fact, the hypothesis that $f_{ii}$ is the identity also follows from the cocycle condition.) Show that there is a unique scheme $X$ (up to unique isomorphism) along with open subsets isomorphic to the $X_i$ respecting this gluing data in the obvious sense. (Hint: what is $X$ as a set? What is the topology on this set? In terms of your description of the open sets of $X$, what are the sections of this sheaf over each open set?)
\end{exr}
\begin{proof}
Confer Exercise~\ref{chap3exr:glue_sheaf} so see that we can glue together sheaves on an open cover. The part of structure sheaf on $X$ is constructed similarly. 

The isomorphism of schemes $f_{ij}:X_{ij}\lrta X_{ji}$ can be considered as a homeomorphism  of  topological spaces $f_{ij}:X_{ij}\lrta X_{ji}$ (abuse notion here) together with an isomorphism of sheaves: $h_{ij}:\calo_{X_j}|_{X_{ji}}\lrta (f_{ij})_*\calo_{X_i}|_{X_{ij}}$.

The cocycle condition on topological space still means
$$f_{ik}|_{X_{ij}\cap X_{ik}} = f_{jk}|_{X_{ji}\cap X_{jk}} \circ  f_{ij}|_{X_{ij}\cap X_{ik}}$$
and for structure sheaves it means
$$(f_{ik}^{-1})_*h_{ik}|_{X_{ij}\cap X_{ik}} = (f_{jk}^{-1})_*h_{jk}|_{X_{ji}\cap X_{jk}} \circ (f_{ij}^{-1})_* h_{ij}|_{X_{ij}\cap X_{ik}}$$

For the topological space
$$
X=\coprod_i X_i/\sim,
$$
where $x\sim y$ iff $x\in X_{ij}, y\in X_{ji}$ and $f_{ij}x=y$ for some $i,j$. $X$ is endowed with quotient topology. Denote the quotient map $\pi:\coprod_i X_i\lrta X$
$$
\tau_X=\{U\subset X:\pi^{-1}(U)\in\tau_{\coprod_i X_i}\}.
$$
Then $X_i\overset{\iota_i}{\inj} \coprod_j X_j\overset{\pi}{\lrta}X$ and we denote the composition $g_i:=\pi\circ \iota_i$. $U_i:=g_i(X_i)$ form an open cover of the topological space $X$. $g_i$ is a homeomorphism. The structure sheaf $\calo_{X_i}$ push forward to $\calo_{U_i}:=g_{i*}\calo_{X_i}$ on $U_i$.

Given $x\in X_{ij}$, $g_i(x)=g_j(f_{ij}x)$ because $x\sim f_{ij}x$. Then we can identify $g_i(X_{ij})$ and $g_j(X_{ji})$ as $U_{ij}=U_i\cap U_j$. The cocycle condition of topological homeomorphisms means simply that we can identify $g_{i}(X_{ij}\cap X_{ik})=g_{j}(X_{ji}\cap X_{jk})=g_{k}(X_{kj}\cap X_{ki})$ as $U_{ijk}=U_i\cap U_j\cap U_k$. The sheave isomorphism $h_{ij}$ induces isomorphism of sheaves on $U_{ij}$:
\begin{center}
\begin{tikzcd}
\calo_{X_j}|_{X_{ji}} \arrow[rr, "h_{ij}"] \arrow[dd, "(g_{j})_*"'] &  & (f_{ij})_*\calo_{X_i}|_{X_{ij}} \arrow[dd, "(g_j)_*"] \\
 &  &  \\
(g_j)_*\calo_{X_j}|_{X_{ji}}=\calo_{U_j}|_{U_{ij}} \arrow[rr, "\phi_{ji}"] &  &\calo_{U_i}|_{U_{ij}}=(g_{j})_*(f_{ij})_*\calo_{X_i}|_{X_{ij}},
\end{tikzcd}
\end{center}
where we have used $g_j\circ f_{ij}=g_i$. Also notice that $g_i$ are isomorphism on $X_i$. 
$$
\begin{aligned}
\phi_{ji}&=(g_j)_*h_{ij} (g_j^{-1})_*\\
&= (g_j)_* (f_{ij})_* (f_{ij}^{-1})_*h_{ij} (g_j^{-1})_*\\
&=(g_i)_* (f_{ij}^{-1})_*h_{ij} (g_j^{-1})_*\\
\end{aligned}
$$
Then the cocycle relation translate to the standard cocycle relation in~\ref{chap3exr:glue_sheaf}, $ \phi_{jk} \circ \phi_{ij} = \phi_{ik}$ on $U_i\cap U_j \cap U_k$. Then we can glue up and get a sheaf $\calo_X$ such that $\calo_X|_{U_i}=\calo_{U_i}=g_{i*}\calo_{U_i}$
\end{proof}

\begin{exr}
Show that affine line with doubled origin is not affine scheme.
\end{exr}
\begin{proof}
Denote the affine line with doubled origin by $X$. $X$ is constructed by gluing $X_1:=\affn^1_k=\spec k[u]$ and $X_2:=\affn^1_k=\spec k[t]$ along $X_{12}=\spec [u,1/u]$ and $X_{21}=\spec [t,1/t]$. They glue up by isomorphism $X_{12}\cong X_{21}$, $u\llrta t$. Denote the quotient map $\pi:X_1\coprod X_2\lrta X$.

For an open set $W\subset X$
$$
\calo_X(W)=\calo_{X_1}(\pi^{-1}(W)\cap X_1)\times_{\calo_{X_1}(\pi^{-1}(W)\cap X_1\cap X_2)=\calo_{X_2}(\pi^{-1}(W)\cap X_1\cap X_2)}\calo_{X_2}(\pi^{-1}(W)\cap X_2),
$$
where the notion means fibered product.

We have 
$$
\begin{aligned}
\Gamma(X,\calo_X)&=\calo_{X_1}(\pi^{-1}(X)\cap X_1)\times_{\calo_{X_1}(\pi^{-1}(X)\cap X_{12})=\calo_{X_2}(\pi^{-1}(X)\cap X_{21})}\calo_{X_2}(X\cap X_2)\\
&=\calo_{X_1}(X_1)\times_{\calo_{X_1}(X_{12})=\calo_{X_2}(X_{21})}\calo_{X_2}(X_2)\\
&=k[u]\times_{k[u,1/u]=k[t,1/t]}k[t]\\
&=k[t]
\end{aligned}
$$
The structure sheaf of $X$ is identical as a single affine line. Then $X$ can not be affine because, $X\not\cong \spec \calo_X(X)$.
\end{proof}
\begin{exr}\label{chap4exr:affine_doubled_origin}
Do the same construction with $\affn^1$ replaced by $\affn^2$. You will have defined the affine plane with doubled origin. Describe two affine open subsets of this scheme whose intersection is not an affine open subset.
\end{exr}
\begin{proof}
$X$ is constructed by gluing $X_1=\affn^2=\spec k[s,t]$ and $X_2=\affn^2=\spec k[u,v]$ along
$$
X_{12}=D(s)\cup D(t)=\spec k[s,t,1/s]\cup \spec k[s,t,1/t]
$$ and  
$$
X_{21}=D(u)\cup D(v)=\spec k[u,v,1/u]\cup \spec k[u,v,1/v]
$$ via $s\llrta u, t\llrta v$. (or we can exchange the role of $s,t$, which will not cause any change)
$$
\calo_X(W)=\calo_{X_1}(\pi^{-1}(W)\cap X_1)\times_{\calo_{X_1}(\pi^{-1}(W)\cap X_1\cap X_2)=\calo_{X_2}(\pi^{-1}(W)\cap X_1\cap X_2)}\calo_{X_2}(\pi^{-1}(W)\cap X_2),
$$
In particular,
$$
\begin{aligned}
\Gamma(X,\calo_X)&=\calo_{X_1}(\pi^{-1}(X)\cap X_1)\times_{\calo_{X_1}(\pi^{-1}(X)\cap X_{12})=\calo_{X_2}(\pi^{-1}(X)\cap X_{21})}\calo_{X_2}(X\cap X_2)\\
&=\calo_{X_1}(X_1)\times_{\calo_{X_1}(X_{12})=\calo_{X_2}(X_{21})}\calo_{X_2}(X_2)\\
&=k[s,t]\times_{k[s,t]=k[u,v]}k[u,v]\\
&=k[s,t],
\end{aligned}
$$
where we know $\calo_{X_1)}=k[s,t]$ from example 4.4.1

$U_1=\pi(X_1)$ and $U_2=\pi(X_2)$ are affine open but their intersection $U_{12}=\pi(X_{12})=\pi(X_{21})$ is the affine plain without origin, which is not affine.
\end{proof}

\begin{exr}
Check that the gluing of projective space $\proj^n$ satisfies the cocycle relation, as painlessly as possible. 
\end{exr}
\begin{proof}
The ``Charts'' of projective space
$$
X_i:=\spec k[x_{0/i},x_{1/i},...,x_{n/i}]/(x_{i/i}-1),
$$
$$
X_{ij}:=D(x_{j/i})\subset X_i
$$
$$
X_{ij}\cap X_{ik}=D(x_{j/i})\cap D(X_{k/i})=D(x_{j/i}x_{k/i})\subset X_i
$$
and we have the isomorphism
$$
\begin{aligned}
f_{ij}:X_{ij}&\lrta X_{ji}\\
x_{k/i}&\longmapsto x_{k/j}/x_{k/i}.
\end{aligned}
$$
The cocycle relation on triple intersection should be
$$
f_{ik}|_{X_{ij}\cap X_{ik}} = f_{jk}|_{X_{ji}\cap X_{jk}} \circ  f_{ij}|_{X_{ij}\cap X_{ik}}.
$$
Notice that 
$$
X_{ij}=\spec k[x_{0/i},x_{1/i},...,x_{n/i}, x_{j/i}^{-1}]/(x_{i/i}-1),
$$
and
$$
X_{ij}\cap X_{ik}=\spec k[x_{0/i},x_{1/i},...,x_{n/i}, x_{j/i}^{-1},x_{k/i}^{-1}]/(x_{i/i}-1)
$$
$$
\begin{aligned}
&f_{jk}\circ f_{ij}: \\
&\spec k[x_{0/i},x_{1/i},...,x_{n/i}, x_{j/i}^{-1},x_{k/i}^{-1}]/(x_{i/i}-1)\lrta \spec k[x_{0/k},x_{1/k},...,x_{n/k}, x_{j/k}^{-1},x_{i/k}^{-1}]/(x_{k/k}-1)\\
& f_{jk}\circ f_{ij}:x_{l/i}\mapsto x_{l/j}/_{i/j}\mapsto (x_{l/k}/x_{j/k})/(x_{i/k}/x_{j/k})=x_{l/k}/x_{i/k}=f_{ik}(x_{l/i})
\end{aligned}
$$
\end{proof}

\begin{exr}
Show that the only functions on $\proj^n_k$ are constants ($\Gamma(\proj^n_k,\calo)\cong k$), and hence that $\proj^n_k$ is not affine if $n > 0$.
\end{exr}
\begin{proof}
Consider two open sets $X_i$ and $X_j$ defined in last exercise. Also, we denote $\proj^n_k$ by $X$.
$$
\Gamma(X_i\cup X_j,\calo_X)=\calo_{X_i}(X_{i})\times_{\calo_X(X_{ij})=\calo_X(X_{ji})}\calo_{X_{j}}(X_{j})
$$
polynomials in $\calo_{X_i}(X_i)$ and $\calo_{X_j}(X_j)$ that agree on the intersection $X_{ij}$. 
$$
f\in k[x_{0/i},x_{1/i},...,x_{n/i}]/(x_{i/i}-1)
$$
and
$$
g\in k[x_{0/j},x_{1/j},...,x_{n/j}]/(x_{j/j}-1)
$$
agree when restricted to $X_{ij}$.
$$
f(x_{0/i},x_{1/i},...,x_{n/i})=g\left(\frac{x_{0/i}}{x_{j/i}},\frac{x_{1/j}}{x_{j/i}},...,\frac{x_{n/j}}{x_{j/i}}\right)\in \spec k[x_{0/i},x_{1/i},...,x_{n/i}, x_{j/i}^{-1}]/(x_{i/i}-1)
$$
Then $x_{j/i}$ can not appear in $f$.
$h\in\calo_X(X)$ can restrict to any pair of open sets and agree on the intersection, then must be constant. This solution is intuitive, but  a more rigorous statement is: Given $\omega$ be the base of topology by choosing the open sets contained in a single $X_i$.
$$
\calo_X(U)=\invlim_{B\in\omega, B\subset U} \calo_X(B)
$$
and then check the universal property of $k$.

\end{proof}

\begin{exr}
Show that if $k$ is algebraically closed, the closed points of $\proj^n_k$ may be interpreted in the traditional way: the points are of the form $[a_0,...,a_n]$, where the $a_i$ are not all zero, and $[a_0,...,a_n]$ is identified with $[\lambda a_0,...,\lambda a_n]$ where $\lambda \in k^\times$.
\end{exr}
\begin{proof}
The closed point in $X:=\proj^n_k$ is a point $p\in X$ which is also closed subset of $X$. $X$ is a scheme, for each $x$, we can find an open set $X_i\ni x$ and $X_i,\calo_X|_{X_i}$ is an affine scheme. We can choose $X_i$ as in the former exercise. The closed point in $X$ is also a closed point in $X_i$. By Nullstellensatz, the closed point (maximal ideal) in $k[t_1,...,t_n]$ is exactly of the form $(t_1-a_1,...,t_n-a_n)$, where $k$ is algebraically closed. In this case
$$
\scm=(x_{0/i}-a_0,..,x_{i-1/i}-a_{i-1},x_{i/i}-a_i,x_{i+1/i}-a_{i+1},...,x_{n/i}-a_{n}). \ \ x_{i/i}=a_i=1
$$ 
The same closed point can be identified in another open set $X_j$ as
$$
(x_{0/j}-b_0,..,x_{j-1/j}-b_{j-1},x_{j/j}-b_j,x_{j+1/j}-a_{j+1},...,x_{n/j}-a_{n}),  \ \ x_{j/j}=b_i=1,
$$
$(b_0,..,b_j=1,..,b_n)=\frac{1}{a_j}(a_0,...,a_i=1,...,a_n)$. This justifies our identification of closed point as $[a_0,...,a_n]$
\end{proof}


\section{Projective schemes, and the \textit{Proj} construction
}
\begin{exr}
Consider $\proj^2_k$, with projective coordinates $x_0, x_1$, and $x_2$. (The terminology ``projective coordinate'' will not be formally defined until \S 4.5.8, but you should be able to solve this problem anyway.) Think through how to define a scheme that should be interpreted as $x^2_0 + x^2_1 - x^2_2 = 0$ ``in $\proj^2_k$ ''.
\end{exr}
\begin{proof}
In the open set $X_2$, there is an affine scheme cut out by $x_{0/2}^2+x_{{}1/2}^2-1=0$. It is a closed subset in $\spec k[x_{0/2},x_{1/2}]$ and can be identified as $\spec k[x_{0/2},x_{1/2}]/(x_{0/2}^2+x_{1/2}^2-1)$ as shown in~\ref{chap3exr:specA_f_and_specA/I}.

In all three open charts, we can find
$$
V_0:=\spec k[x_{1/0},x_{2/0}]/(1+x_{1/0}^2-x_{2/0}^2),
$$
$$
V_1:=\spec k[x_{0/1},x_{2/1}]/(x_{0/1}^2+1-x_{2/1}^2),
$$
$$
V_2:=\spec k[x_{0/2},x_{1/2}]/(x_{0/2}^2+x_{1/2}^2-1)
$$
$$
V_{01}=\spec k[x_{1/0},x_{2/0},x_{1/0}^{-1}]/(1+x_{1/0}^2-x_{2/0}^2)
$$
$$
V_{10}=\spec k[x_{0/1},x_{2/1},x_{0/1}^{-1}]/(x_{0/1}^2+1-x_{2/1}^2)
$$
and so on.
We find that the gluing isomorphisms of $X_i$ induces well behaved isomorphisms on $V_i$
$$
g_{01}:V_{01}\lrta V_{10}
$$
induced by 
$$
f_{10}:x_{1/0}\mapsto x_{1/1}/x_{0/1}, x_{2/0}\mapsto x_{2/1}/x_{0/1},
$$
where $x_{i/i}$ are dummy variables that equals $1$. $f_{01}$ induces well-defined morphism of the quotient rings because 
$$
\begin{aligned}
f_{01}:(1+x_{1/0}^2-x_{2/0}^2)\lrta &(1+(x_{1/1}/x_{0/1})^2+(x_{2/1}/x_{0/1})^2)\\
&=(1/x_{0/1})^2(x_{0/1}^2+1-x_{2/1}^2)\\
&=(x_{0/1}^2+1-x_{2/1}^2)
\end{aligned}
$$
$g_{ij}$ also satisfies the cocycle relation because $f_{ij}$ does.

Hence $V_i$ glue up to a closed subscheme in $\proj^2_k$

\end{proof}

\begin{exr}
More generally, consider $\proj^n_A$, with projective coordinates $x_0,..., x_n$. Given a collection of homogeneous polynomials $f_i \in A[x_0,...,x_n]$, make sense of the scheme ``cut out in $\proj^n_A$ by the $f_i$.''
\end{exr}
\begin{proof}
Given a homogeneous polynomial $f_i\in A[x_0,...,x_n]$, we can consider map $x_j\rta x_{j/j}=1$ and $x_{k}\rta x_{k/j}$ so that we get a polynomial in $F_{i,(j)}\in A[x_{i/j},...]$ which no longer homogeneous. Then we can construct the closed subset in affine chart $X_i$ 
$$
U_{i,(j)}:=\spec A[x_{1/j},...,x_{n/j}]/(x_{j/j}-1, F_{i,(j)})
$$
and
$$
U_{i,(j);k,(l)}=\spec A[x_{1/j},...,x_{n/j},x_{l/j}^{-1}]/(x_{j/j}-1, F_{i,(j)}, F_{k,(j)})
$$
$$
U_{k,(l);i,(j)}=\spec A[x_{1/l},...,x_{n/l},x_{j/l}^{-1}]/(x_{l/l}-1, F_{i,(l)}, F_{k,(l)}).
$$
Again the gluing map of $X_i$ 
$$
h_{jl}:x_{m/j}\mapsto x_{m/l}/x_{j/l}
$$
would induce well-defined morphism on the quotient ring. Because 
$$
h_{jl}:F_{k,(j)}\mapsto \frac{1}{x_{j/l}^{n_{jl}}}F_{k,(l)},
$$
where $n_{jl}$ is the degree of $x_{l;j}$ in $F_{k,(j)}$. Denote the induced morphism $g_{i,(j);k,(l)}$. We can similarly verify that the induced morphism satisfies the cocycle relation because $h_{jl}$ does.
$$
g_{i,(j);k,(l)}=g_{p,(j);k,(l)}\circ g_{i,(j);p,(l)}
$$
 Hence we can glue the affine schemes to get a closed subscheme in $\proj^n_A$.
\end{proof}

\begin{exr}\label{chap4exr:homogeneous_ideal_characterization}\ 
\begin{enumerate}[label=(\alph*)]
\item Show that an ideal $I$ is homogeneous if and only if it contains the degree $n$ piece of each of its elements for each $n$.
\item Show that the set of homogeneous ideals of a given $\intg$-graded ring $S_\bullet$ is closed under sum, product, intersection, and radical.
\item 
Show that a homogeneous ideal $I \subset S_\bullet$ is prime if $I\neq S_\bullet$, and if for any homogeneous $a, b \in  S_\bullet$, if $ab \in I$, then $a \in I$ or $b \in I$.
\end{enumerate}
\end{exr}
\begin{proof}
\begin{enumerate}[label=(\alph*)]
\item Want:``$I$ is homogeneous ideal'' $\Llrta$ ``$I$ contains each degree $n$ piece of each element for each degree''

 $\Lrta$:  A homogeneous ideal $I$ is generated by homogeneous elements $\{a_i\in S_{n(i)}\}_{i\in I}$, where $n(i)$ is the degree of $a_i$. A general element looks like $b=\sum_i s_i a_i$. $s_i$ are not necessarily homogeneous, but we can decompose $s$ into homogeneous pieces and verify that homogeneous pieces of $b$ are still generated by $a_i$ which are contained in $I$.

 $\Llta$: $b=\sum_n b_n\in S$, where $b_n$ is homogeneous of degree $n$. We can simply choose the generating set to be the homogeneous pieces of each of the element in $I$. Then $I$ is generated by homogeneous elements. Then we can decompose $I$ as $\oplus_{n\in\intg}I_n$. And $S/I$ has a natural $\intg$-grading.
 \item $I,J$ are homogeneous ideals in $S_\bullet$. It is easy to verify that the $I+J$, $I\cdot J$ and $I\cap J$ are generated by homogeneous elements. As for radicals, 
 $$
 L=\sqrt{I}
 $$
 $x\in L\Llrta x^n\in L$. $x=\sum_{i\in\intg} x_i$ only finitely many of $x_i$ are not zero. Assume d the highest degree of $x$, $deg\  x_d$ is larger than any other homogeneous piece.  Specifically, $x_d^n\in I$ which means $x_i\in L$. Then $x-x_d\in L$, we can induct on this procedure until we find each homogeneous piece of $x$ is contained in $L$. Which means $I$ is homogeneous by $(a)$.
 \item The only if direction is trivial. We only need to check the if direction.

Consider general elements $x=\sum_i^d x_i$ and $y=\sum_i^{k} y_j$, they are finite sums by definition. We can check
$$
\sum_i^d x_i\cdot \sum_j^k y_j=\sum_n\sum_{i+j=n}x_i y_j.
$$
Assume neither $x$ nor $y$ is contained in $I$. Then at least one $x_d$ and $y_k$ are not in $I$. Assume $d,k$ to be maximal with this property.

We can subtract from $x$ the part $\sum_{n>d}x_n$ such that $\tilde{x}=x-\sum_{n>d}x_n\notin I$. Similarly, set $\tilde{y}=y-\sum_{m>k}y_m\notin I$. Then we have $\tilde{x}\tilde{y}\in I$ and the leading homogeneous piece being $x_dy_k$. Then $x_d y_k\in I$, $\Lrta x_d$ or $y_k\notin I$. Contradiction.
\end{enumerate}
\end{proof}

\begin{exr}\ 
\begin{enumerate}[label=(\alph*)]
\item
Show that a graded ring $S_\bullet$ over $A$ is a finitely generated graded ring (over $A$) if and only if $S_\bullet$ is a finitely generated graded $A$-algebra, i.e., generated over $A = S_0$ by a finite number of homogeneous elements of positive degree. (Hint for the forward implication: show that the generators of $S_+$ as an ideal are also generators of $S_\bullet$ as an algebra.)
\item Show that a graded ring $S_\bullet$ over $A$ is Noetherian if and only if $A = S_0$ is Noetherian and $S_\bullet$ is a finitely generated graded ring.
\end{enumerate}
\end{exr}
\begin{proof}
\begin{enumerate}[label=(\alph*)]
\item $\Lrta$: $S_\bullet$ is  a finitely generated graded ring over $A$, which means $S_0=A$ and the irrelevant ideal $S_+$ is finitely generated ideal. Such a graded ring is naturally an $A$-algebra. Assume $S_+$ is finitely generated by as $(s_{1,1},...,s_{n,m_{j}})$ each of degree higher than $s_{i,k}, k\leq m_j$ is of degree $i$. Then a general element in $S_1$ can be expressed as
$$
\sum_j^{m_1} a_j s_{1,j}.
$$
Each element in $S_2$ can be expressed as
$$
\sum_{j}^{m_2}a_{j}s_{2,j}+\sum_{j,k}^{m_1}b_{jk}s_{1,j}s_{1,k}.
$$
And it will also work for any degree, which means $S$ is generated as polynomial $A[s_{i,j}]$ (quotient some relations).

$\Llta$ If a graded ring $S_\bullet$ is finitely generated as an $A$-algebra ($S_0=A$), $S_+=S_\bullet -A$ is finitely generated.
\item $\Llta$: By (a), $S_\bullet$ being a finitely generated graded ring with $S_0=A$ means $S_\bullet$ is finitely generated $A$-algebra. Grading is only an extra structure on $S$, we can proceed by Hilbert basis theorem. If $A$ is Noetherian, the finitely generated $A$-algebra is Noetherian.

$\Lrta$: Graded ring $S_\bullet$ is Noetherian, then every ideal of $S_\bullet$ is finitely generated, specifically, $S_+$ is finitely generated. On the other hand $S_\bullet/S_+\cong A$, is a homomorphic image of a Noetherian ring. Hence $A$ must also be Noetherian.
\end{enumerate}
\end{proof}

\begin{exr}\label{chap4exr:interpreation_projS}
Suppose $f \in S_+$ is homogeneous.
\begin{enumerate}[label=(\alph*)]
\item Give a bijection between the prime ideals of $((S_\bullet)_f)_0$ and the homogeneous prime ideals of $(S_\bullet)_f$. Hint: Avoid notational confusion by proving instead that if $A$ is a $\intg$-graded ring with a homogeneous invertible element $f$ in positive degree, then there is a bijection between prime ideals of $A_0$ and homogeneous prime ideals of $A$. Using the ring map $A_0 \lrta A$, from each homogeneous prime ideal of $A$ we find a prime ideal of $A_0$. The reverse direction is the harder one. Given a prime ideal $P_0 \subset A_0$, define $P \subset A$ (a priori only a subset) as $\oplus Q_i$, where $Q_i \subset A_i$, and $a\in Q_i$ if and only if $a^{\deg f}/f^i\in P_0$ . Note that $Q_0 =P_0$. Show  that $a\in Q_i$ if and only if $a^2\in Q_{2i}$; show that if $a_1,a_2 \in Q_i$ then $a^2_1 +2a_1a_2 +a^2_2 \in Q_{2i}$ and hence $a_1 + a_2 \in Q_i$; then show that $P$ is a homogeneous ideal of $A$; then show that $P$ is prime.
\item Interpret the set of prime ideals of $((S_\bullet)_f)_0$ as a subset of $\Proj S_\bullet$.
\end{enumerate}
\end{exr}
\begin{proof}
\begin{enumerate}[label=(\alph*)]
\item Follow the hint, we avoid the notation trouble by working in  a $\intg$-graded ring $A$ with invertible  homogeneous element $f$ with $\deg f>0$.

\underline{Claim}: There is a bijection between prime ideals of $A_0$ and homogeneous prime ideal of $A$.

There is a natural inclusion $\iota:A_0\inj A$. For a given homogeneous prime ideal $P$ in $A$ we can find a prime ideal $\iota^* P\in A_0$ which is just $P\cap A_0$ which is a prime ideal in $A_0$.

The reverse direction is harder. Consider a prime ideal $P_0\in A_0$, define $P\subset A$ as $P=\oplus Q_i$, where $Q_i\subset A_i$ and a homogeneous element $a\in Q_i$ iff $a^{\deg f}/f^i\in P_0$. Note that $P_0=Q_0$.

\underline{Want}: $P$ is an ideal.
Indeed, for $a\in P$, $b\in A$. $a=\sum_i a_i$, $b=\sum_j b_j$, where each $a_i\in Q_i$ and $b_j\in A_j$. $a_i b_i\in A_{2i}$.
$$
\frac{(a_i b_j)^{\deg f}}{f^{i+j}}=\frac{a_i^{\deg f}}{f^i}\frac{b_i^{\deg f}}{f^j}\in P_0.
$$
Then product is contained in $P_0$ because $a_i^{\deg f}{f^i}\in P_0$ and $b_j^{\deg f}{f^j}\in A_0$ and $P_0$ is a prime ideal in $A_0$. Then we have 
$$
\frac{(a_i b_j)^{\deg f}}{f^{i+j}}\in Q_{i+j}.
$$
Similarly, we can prove $a\cdot b=\sum_i a_i\cdot \sum_i b_i\in P$. 

\underline{subclaim}: $a\in Q_i$ iff $a^2\in Q_i$. 
 one direction is clear and we only need to prove the only if part.
 $$
 \frac{(a^2)^{\deg f}}{f^{2i}}=\frac{(a)^{\deg f}}{f^{i}}\frac{(a)^{\deg f}}{f^{i}}\in P_0\Lrta \frac{(a)^{\deg f}}{f^{i}}\in P_0
 $$
 because $P_0$ is prime ideal.

On the other hand, given $a_1,a_2\in Q_i$
 We have $P$ is an ideal in $A$. We know $a_1^2+2a_1a_2+a_2^2\in Q_{2i}$. Hence by the subclaim, $a_1+a_2\in Q_i$.

Then we know $P$ is closed under summation and multiplication by $A$. ($P$ is an ideal.)

\underline{Want}: $P$ is homogeneous ideal. Recall~\ref{chap4exr:homogeneous_ideal_characterization}, part $(a)$, we only need to prove that $P$ contains  the homogeneous piece of each of its element of each degree. Assume $a\in P$ then $a=\sum_i a_i$ and $a_i\in Q_i$ by definition of $P$. We know $P$ is homogeneous ideal.

\underline{Want}: $P$ is a prime ideal in $A$. See~\ref{chap4exr:homogeneous_ideal_characterization}, part $(c)$. A homogeneous ideal is prime iff it is prime for homogeneous elements. Suppose $a\in A_i, b\in A_j$ and we have $ab\in Q_{ij}$, then by definition
$$
\frac{a^{\deg f}}{f^i}\frac{b^{\deg f}}{f^j}\in P_0
$$
then either $a\in Q_i$ or $b\in Q_j$ because $P_0$ is prime ideal.

We still need to prove that this indeed gives a bijection. The map $\alpha:I\lrta I\cap A_0$ as described above and the map $\beta: P_0\mapsto P$. $\alpha\circ \beta:P_0\mapsto P_0$. $\beta\circ \alpha:I\mapsto I\cap A_0\mapsto? I$.

$$
\beta(I\cap A_0)\supset I
$$
because $x\in I_j$ then $x^{\deg f}/f^j\in I\cap A_0$.

For the reverse inclusion, consider an element $y\in\beta(I\cap A_0)_j$ then by definition $y^{\deg f}/f^j\in I\cap A_0$, then $y^{\deg f}=y^{\deg f}/f^j\cdot f^j\in I_{j\deg f}$ $\Lrta y\in I_j$ because $I$ is prime ideal.

\item We have proved there is bijection between ``prime ideals in $((S_\bullet)_f)_0$'' and homogeneous prime ideals in $(S_\bullet)_f$. Recall the property of localization, there is a one to one homogeneous prime ideals in $(S_\bullet)_f$ and homogeneous prime ideals in $S_\bullet$ that does not intersect $\{f,f^2,...\}\subset S_+$. (Localization preserves the homogeneity of prime ideals.) In particular, these homogeneous prime ideals do not contain $S_+$, therefore they can be interpreted as subset in $\Proj S_\bullet$.
\end{enumerate}
\end{proof}
\begin{exr}\label{chap4exr:projective_affine_open_chart_localization}
Show that $D(f)$ ``is'' (or more precisely, ``corresponds to'') the subset $\spec((S_\bullet)_f)_0$ you described in Exercise~\ref{chap4exr:interpreation_projS}. For example, the $D(x_i)$ are the standard open sets covering projective space.
\end{exr}
\begin{proof}
$D(f)$ is the projective distinguished open set and
$$
\begin{aligned}
D(f)&=\Proj S_\bullet -V(f)\\
&=\{\text{Homogeneous prime ideals that do not contain $f$ or $S_+$}\}\\
&=\{\text{Homogeneous prime ideal that do not contain $f$}\} (f\in S_+)
\end{aligned}
$$
And we have the correspondences:
$$
\begin{aligned}
&\{\text{Homogeneous prime ideals in $S_\bullet$ that do not contain $f$}\}\\
&\llrta (\text{by localization  property})\\
&\{\text{Homogeneous prime ideals in $(S_\bullet)_f$}\}\\
&\llrta (\text{by Exercise~\ref{chap4exr:interpreation_projS}})\\
&\{\text{Prime ideals in $((S_\bullet)_f)_0$}\}\\
&=\spec ((S_\bullet)_f)_0
\end{aligned}
$$
\end{proof}

\begin{exr}
Verify that the projective distinguished open sets $D(f)$ (as $f$ runs through the homogeneous elements of $S_+$) form a base of the Zariski topology.
\end{exr}
\begin{proof}
As in the affine case, the closed set of Zariski's topology is of the form
$$
V(T)=V(\overline{T}).
$$
where $I:=\overline{T}$ is a homogeneous ideal generated by $T$. $\overline{T}\subset S_+$. 

It forms a topology
$$
\cap_i V(J_i)=V(\sum_i J_i)
$$
and 
$$
V(I)\cup V(J)=V(I\cdot J)
$$
and
$$
IV(J)=\sqrt{J}
$$
both are well-defined because \{homogeneous ideals\} is closed under intersection, product, addition and taking radicals.

\underline{Claim}: $\{D(f)\}$ where $f$ is homogeneous element in $S_+$ form a base of the topology.

Given a point $[P]$ in an open set $\Proj S_\bullet-V(I)$ , $P$ is a homogeneous prime ideal in $S_\bullet$ that does not contains $I$. We can find $f\in I\backslash P$. Then by definition 
$$
V(f)\supset V(I)
$$
and
$$
D(f)\subset \Proj S_\bullet-V(I).
$$
Also $P$ does not contain $f$, hence $D(f)\ni [P]$.

Which means $D(f)$ as $f$ ranges over $S_+$ form a base of the Zariski's topology.
\end{proof}

\begin{exr}\label{chap4exr:4.5.H}
Fix a graded ring $S_\bullet$.
\begin{enumerate}[label=(\alph*)]
\item Suppose $I$ is any homogeneous ideal of $S_\bullet$ contained in $S_+$, and $f$ is a homogeneous element of positive degree. Show that $f$ vanishes on $V(I)$ (i.e.,$V(I) \subset V(f)$) if and only if $f^n ? I$ for some $n$. (Hint: Mimic the affine case~\ref{exr:vanishing_set_radical},) In particular, as in the affine case (Exercise~\ref{chap3exr:inclusion_distinguished_open}), if $D(f) \subset D(g)$, then $f^n \in (g)$ for some $n$, and vice versa. (Here $g$ is also homogeneous of positive degree.)
\item If $Z \subset \Proj S_\bullet$, define $I(Z) \subset S_+$. Show that it is a homogeneous ideal of $S_\bullet$. For any two subsets, show that $I(Z_1 \cup Z_2) = I(Z_1) \cap I(Z_2)$.
\item For any subset $Z \subset \Proj S_\bullet$, show that $V(I(Z)) = \overline{Z}$.
\end{enumerate}
\end{exr}
\begin{proof}
\begin{enumerate}[label=(\alph*)]
\item 
If $f^n\in I $ for some $n$, then $f\in P$ for all $P$ homogeneous prime containing $I$ because $P$ is prime. Then $f$ vanishes on $V(I)$.

For the reverse direction, $f$ vanishes on $V(I)$, $V(I)\subset V(f)$. It means all homogeneous prime ideals that containing $I$ but not $S_+$ contains $f$:
$$
f\in \cap_{P\supset I, P\notsupset S_+} P.
$$
Notice when $f$ is homogeneous of positive degree, $f\in S_+$, $f$ is automatically contained in any homogeneous prime that contains $S_+$. the condition $f$ vanishes on $V(I)$ in fact is equivalent to 
$$
f\in\cap_{P\supset I}P.
$$

Note that $I$ itself is homogeneous.

\underline{Claim}: The radical of a homogeneous ideal $I$ is the intersection of all homogeneous prime ideal that contains $I$.

In general, we use $P$ to denote homogeneous prime
$$
\cap_{P\supset I}P\supset \cap_{\scp\supset I}\scp.
$$
For the reverse inclusion, consider $\scp^h$ to be the the homogeneous ideal generated by the homogeneous element of $\scp$. $\scp^h\subset \scp$ is prime on homogeneous elements then $\scp^h$ is homogeneous prime ideal by~\ref{chap4exr:homogeneous_ideal_characterization}. If $\scp$ is any prime containing $I$, then $\scp^h$ also contains $I$, because $\scp^h$ contains the homogeneous elements of $I$, and $I$ as a homogeneous ideal is generated by these homogeneous elements. 
We have
$$
\cap_{P\supset I} P\subset \cap_{\scp\supset I}\scp^h\subset \cap_{\scp\supset I}\scp .
$$
Then we know the intersection of homogeneous primes that containing a homogeneous ideal $I$ is just the radical of $I$.
$$
f\in\cap_{P\supset I}P=\cap_{\scp\supset I}\scp=\sqrt{I}.
$$

Hence exists $n\geq 0$ such that $f^n\in I$.
\item $I(Z)$ is defined to be the element of $S_\bullet$ that vanishes on all of $[P]\in Z$. $Z\subset \Proj S_\bullet$, means each $[P]\in Z$.
 $P$ is homogeneous prime that does not contain $S_+$.

  $f\in I(Z)$ $\Llrta$ $f\in \cap_{[P]\in Z}P$.

  Denote the degree $i$ piece of $f$ by $f_i$. $f\in\cap_{[P]\supset Z}P\Lrta f_i\in P$ $\forall [P]\in Z$, because each $P$ is homogeneous, this in turn means $f_i\in I(Z)$. $I(Z)$ contains each homogeneous piece of each element in it, which implies that $I(Z)$ is a homogeneous ideal. 

  For two subsets $Z_1$ and $Z_2$ in $\Proj S_\bullet$,
  $$
  f\in I(Z_1\cup Z_2)
  $$ 
  $$
  \Llrta
  f\in P,\forall [P]\in Z_1\cup Z_2
  $$
  $$
  \Llrta
  f\in P,\forall [P]\in Z_1\text{ AND }f\in Q,\forall [Q]\in Z_2 
  $$
  $$
  \Llrta f\in I(Z_1) \cap I(Z_2)
  $$
  \item ``$V(I(Z))\supset \overline{Z}$'':
By definition, $V(I(Z))$ is closed set. And $Z\subset V(I(Z))$, because $[P]\in Z, I(Z)=\cap_{[Q]\in Z}Q$ $I(Z)\mod P=0,\Lrta I(Z)\subset P\Lrta [P]\in V(I(Z))$. Closure $\overline{Z}$ is the smallest closed set that contains $Z$, therefore $\overline{Z}\subset V(I(Z))$.

 ``$V(I(Z))\subset \overline{Z}$'': We need to verify that open neighborhood of each point in $V(I(Z))$ intersects with $Z$. Consider a point $[Q]\in V(I(Z))$, we have $Q\supset I(Z)$. Assume  a non-empty distinguished open $D(f)$ contains $[Q]\Llrta f\notin Q \Lrta f\notin I(Z)\Lrta$ $f $ does not vanish on every point in $Z$, $\exists[P]\in Z$, s.t. $f\notin P\Llrta$ $D(f)\ni [P]$, therefore $D(f)\cap V(I(Z))\neq \emptyset$. Hence $V(I(Z))\subset \overline{Z}$.

\end{enumerate}
\end{proof}

\begin{exr}
Fix a graded ring $S_\bullet$, and a homogeneous ideal $I$. Show that the following are equivalent.
\begin{enumerate}[label=(\alph*)]
\item $V(I) = \emptyset$.
\item For any $f_i$ (as $i$ runs through some index set) generating $I$, $\cup_i D(f_i)=\Proj S_\bullet$. 
\item $\sqrt{I}\supset S_+$.
\end{enumerate}
\end{exr}
\begin{proof}
$(a)\Lrta (b)$: Assume $I$ is finitely generated by $\{f_i\}_{i\in J}$. $V(I)=\emptyset$, means $\forall [P]\in \Proj S_\bullet$, $P\notsupset I$. For each $[P]\in \Proj S_\bullet$, $P\not\ni f_i$ for some $f_i\in I$ $\Lrta [P]\in D(f_i)$ for some $f_i$, which means $\Proj S_\bullet \subset \cup_{i\in J} D(f_i)\subset \Proj S_\bullet$.

$(b)\Lrta (c)$: $\Proj S_\bullet = \cup_{i\in J} D(f_i)$ means each homogeneous prime ideal that does not contain $S_+$ in $S_\bullet$ does not contain $I$. The contrapositive says each homogeneous prime ideal contain $I$ would contain $S_+$:
$$
\cap_{P\supset I} P\supset S_+.
$$
By Exercise~\ref{chap4exr:4.5.H}, we know $\sqrt{I}=\cap_{P\supset I} P$ for homogeneous ideal $I$.

$(c)\Lrta (a)$, 
$$
\sqrt{I}=\cap_{P\supset I} P\supset S_+
$$
means for each point $[P]\in \Proj S_\bullet$, $P$ does not contain $I$. $\Lrta$ $[P]\notin V(I)$.

$\Lrta V(I)=\emptyset$.
\end{proof}

\begin{exr}
Suppose some homogeneous $f\in S_+$ is given. Via the inclusion
$$
D(f) = \spec((S_\bullet)_f)_0  \inj \Proj S_\bullet
$$
of Exercise~\ref{chap4exr:projective_affine_open_chart_localization}, show that the Zariski topology on $\Proj S_\bullet$ restricts to the Zariski topology on $\spec((S_\bullet)_f)_0$
\end{exr}
\begin{proof}
We can compare the closed subsets of these two topological space. Consider a subset $I\subset ((S_\bullet)_f)_0$
$$
\begin{aligned}
((S_\bullet)_f)_0\supset V'(I)&=\{ [\scp]\in\spec ((S_\bullet)_f)_0: I\subset \scp\}\\
&\llrta\{[P]\in \Proj (S_\bullet)_f: I^{h}\subset P\}\\
&(I^h\text{ is the homogeneous ideal generated by $I$.})\\
&\llrta\{[Q]\in \Proj S_\bullet: \iota^*(I^h)\subset Q, Q\not\ni f\}\\
&(\iota^*(I^h)\text{ denote the preimage of $I^h$ under localization})
\end{aligned}
$$
where we use $\scp$ to denote ordinary prime ideal and $P$ to mean homogeneous prime ideal. Also $\iota:S_\bullet\inj (S_\bullet)_f$ is the standard inclusion of localization.
$$
\begin{aligned}
((S_\bullet)_f)_0\supset V'(D)
&\llrta\{[Q]\in \Proj S_\bullet: \iota^*(I^h)\subset Q, Q\not\ni f\}\\
&\llrta \{[Q]\in \Proj S_\bullet: Q\supset \iota^*(I^h)\}\cap\{[Q]\in\Proj S_\bullet:Q\not\ni f\}\\
&=V(\iota^*(I^h))\cap D(f)
\end{aligned}
$$

For illustration, consider previous example
$S:=k[x_0,x_1,x_2]$, $S_{x_0}:=k[x_0,x_1,x_2,x_0^{-1}]$, $(S_{x_0})_0:=k[x_1/x_0,x_2/x_0]$. $I=(1+x_{1/0}^2-x_{2/0}^2)_0$ and $I^h=(1+x_{1/0}^2-x_{2/0}^2)$.
$$
V'(I)=V_0:=\spec k[x_{1/0},x_{2/0}]/(1+x_{1/0}^2-x_{2/0}^2),
$$
$\iota^*(I^h)=(x_0^2+x_1^2-x_2^2)$.
$$
V'(I)\llrta D(x_0)\cap V((x_0^2+x_1^2-x_2^2))
$$
\end{proof}


\begin{exr}
If $f, g \in S_+$ are homogeneous and nonzero, describe an isomorphism between $\spec((S_\bullet)_{fg})_0$ and the distinguished open subset $D(g^{\deg f}/f^{\deg g})$ of $\spec((S_\bullet)_f)_0$.
\end{exr}
\begin{proof}
Recall that 
$$
\spec((S_\bullet)_{fg})_0\cong D(fg)\subset \Proj S_\bullet
$$
Then problem reduce to describe the intersection of $D(fg)$ with $\spec ((S_\bullet)_f)_0$.
$$
\begin{aligned}
\spec ((S_\bullet)_{fg})_0&= \{[\scq]\in \spec ((S_\bullet)_{fg})_0:\}\\
&\llrta \{[Q]\in \Proj (S_\bullet)_{fg}\}\\
&\llrta \{[Q]\in\Proj S_\bullet:Q\not\ni fg\}\\
&\llrta \{[Q]\in\Proj S_\bullet:Q\not\ni f\} \cap \{[Q]\in\Proj S_\bullet: Q\not\ni g\}\\
&\llrta \{[P]\in \Proj (S_\bullet)_f: P\not\ni \frac{g}{1}\}\\
&\llrta\{ [\scp]\in\spec ((S_\bullet)_f)_0: \scp\not\ni\frac{g^{\deg f}}{f^{\deg g}}\}
\end{aligned}
$$
We only have to explain the last bijection. $\scp=P\cap ((S_\bullet)_f)_0$ and if $P\not\ni g/1$, we have $\scp\not\ni g^{\deg f}/f^{\deg g}$. For the reverse direction, we have to recall the construction in~\ref{chap4exr:interpreation_projS}, $g/1\notin P_{\deg g}$ iff $g^{\deg f}/f^{\deg g}\notin \scp=P_0$.
\end{proof}

Before finishing the definition, we have to collect some results about graded ring here in order to construct the isomorphism of sheaves on intersections.
\begin{theorem}
\label{lemma-standard-open}
Let $S_\bullet$ be a graded ring. Let $f \in S_\bullet$ homogeneous of positive degree.
\begin{enumerate}
\item If $g\in S$ homogeneous of positive degree
and $D(g) \subset D(f)$, then
\begin{enumerate}
\item $f$ is invertible in $(S_\bullet)_g$, and
$f^{\deg(g)}/g^{\deg(f)}$ is invertible in $S_{(g)}$,
\item $g^e = af$ for some $e \geq 1$ and $a \in S$ homogeneous,
\item there is a canonical $S_\bullet$-algebra map $(S_\bullet)_f \to (S_\bullet)_g$,
\item there is a canonical $(S_\bullet)_0$-algebra map $((S_\bullet)_{f})_0 \to ((S_\bullet)_{g})_0$
compatible with the map $(S_\bullet)_f \to (S_\bullet)_g$,
\item the map $((S_\bullet)_{f})_0 \to ((S_\bullet)_{g})_0$ induces an isomorphism
$$
(((S_\bullet)_{f})_0)_{g^{\deg(f)}/f^{\deg(g)}} \cong ((S_\bullet)_{g})_0,
$$
\item there are compatible canonical $S_\bullet)_{f}$ and $((S_\bullet)_{f})_0$-module
maps $(M_\bullet)_f \to (M_\bullet)_g$ and $((M_\bullet)_f)_0 \to ((M_\bullet)_g)_0$ for any graded $S_\bullet$-module $M_\bullet$,
and
\item the map $((M_\bullet)_{f})_0 \to((M_\bullet)_{g})_0$ induces an isomorphism
$$
(((M_\bullet)_{f})_0)_{g^{\deg(f)}/f^{\deg(g)}} \cong ((M_\bullet)_{g})_0.
$$
This is a horrible notation, where the first and third subscripts means grading while second and forth means localizations. 
\end{enumerate}
\item Any open covering of $D_{+}(f)$ can be refined to a finite
open covering of the form $D_{+}(f) = \bigcup_{i = 1}^n D_{+}(g_i)$.
\item Let $g_1, \ldots, g_n \in S$ be homogeneous of positive degree.
Then $D_{+}(f) \subset \bigcup D_{+}(g_i)$
if and only if
$g_1^{\deg(f)}/f^{\deg(g_1)}, \ldots, g_n^{\deg(f)}/f^{\deg(g_n)}$
generate the unit ideal in $S_{(f)}$.
\end{enumerate}
\end{theorem}\label{chap4thm:graded_ring_localization_iso}
\begin{proof}
We only prove $1.(a)-1.(e)$ and the remaining part can be found at \href{https://stacks.math.columbia.edu/tag/01M4}{Stack Project 26.8}.

\begin{enumerate}[label=(\alph*)]
\item
$f$ is is invertible in $(S_\bullet)_g$ by~\ref{chap3exr:inclusion_distinguished_open}. And $f^{\deg g}/g^{\deg f}$ is an invertible element an element in $((S_\bullet)_g)_0$ because it is not contained in any prime ideal in $((S_\bullet)_g)_0$. ($f/1$ is not contained in any prime ideal in $(S_\bullet)_g$ because $f/1$ is invertible.)
\item Then the inverse of $f$ in $(S_\bullet)_g$ can be written as $a'/g^d$ and $fa'/g^d=1/1$. We can replace $a'$ by its homogeneous part $a''$ and $(fa''-g^d)/1=0/1$ $g^k(fa-g^d)=0$. Then after  defining $a:=a''g^k$, we get $fa=g^e$, for $e=k+d$.
\item The morphism exists by universal property of localization. It maps $b/f^n$ to $a^nb/g^{ne}$. 

\item This clearly induce a map of the degree $0$ ring, $((S_\bullet)_f)_0\to ((S_\bullet)_g)_0$.
\item We need to look the ring morphism $((S_\bullet)_f)_0\to ((S_\bullet)_g)_0$ in detail. It is induced by 
$$
\phi:\left(\frac{b}{f^n}\right)_0\mapsto \left(\frac{a^n b}{g^{ne}}\right)_0
$$
In particular
$$
\frac{g^{\deg f}}{f^{\deg g}}\mapsto \frac{a^{\deg g}}{g^{e\deg g-\deg f}}
$$
Because $g^{\deg f}/f^{\deg g}$ is invertible in $(S_\bullet)_0 $, it induces a morphism
$$
\tilde{\phi}:(((S_\bullet)_f)_0)_{g^{\deg f}/f^{\deg g}}
\lrta ((S_\bullet)_g)_0
$$
$$
\begin{aligned}
\left(\frac{b}{f^n}\right)_0/\left(\frac{g^{\deg f}}{f^{\deg g}}\right)^m&\mapsto \phi\left(\frac{b}{f^n}\right)_0/\phi\left(\frac{g^{\deg f}}{f^{\deg g}}\right)^m\\
&=\left(\frac{\left(a^nb/g^{ne}\right)}{\left(a^{\deg g}/g^{e\deg g-\deg f}\right)^m}\right)_0\\
&=\left(\frac{a^{n-m\deg g}b}{g^{m\deg f-me\deg g+ne}}\right)_0\\
&=\left(\frac{b}{g^{m\deg f}f^{n-m\deg g}}\right)_0
\end{aligned}
$$
This morphism is surjective because $f$ is invertible in $(S_\bullet)_g$. Then we can look at the kernel of this ring morphism. If
$$
\left(\frac{b}{g^{m\deg f}f^{n-m\deg g}}\right)_0=0,
$$
It means exactly the degree $0$ part of $b/f^n$ is zero. Hence we have the isomorphism.
\end{enumerate}
\end{proof}

\begin{exr}
By checking that these gluings behave well on triple overlaps (see Exercise~\ref{chap3exr:glue_sheaf}), finish the definition of the scheme $\Proj S_\bullet$ .
\end{exr}
\begin{proof}
$D(f)$ give an open cover of the topological space of $X:=\Proj S_\bullet$. We identify the $\calo_{X,f}$ to be the structure sheaf of $\spec ((S_\bullet)_f)_0$: $\calo_{\spec ((S_\bullet)_f)_0}$.

$D(f)\cap D(g)=D(fg)$, we have the isomorphism of sheaves:
$$
\phi_{f,g}:\calo_{X,f}|_{D(fg)}\lrta \calo_{X,g}|_{D(fg)}.
$$
$$
\calo_{\spec ((S_\bullet)_f)_0}|_{D(g^{\deg f}/f^{\deg g})}\lrta \calo_{\spec ((S_\bullet)_g)_0}|_{D(f^{\deg g}/g^{\deg f})}
$$
We want to verify the cocycle relation on $D(fgh)$. We verify the cocycle relation on a base of topology. The distinguished base of $\spec (S_\bullet)_{fgh})_0$ is also induced by distinguished open in $\Proj S_\bullet$.

\underline{Want}:
$$
\phi_{f,h}(D(fghk))=\phi_{g,h}(D(fghk))\circ \phi_{f,g}(D(fghk)),\forall k\text{ homogeneous in $S_+$}.
$$
Notice that sheaf morphism is commutative with restriction maps. It suffices to check
$$
\phi_{f,h}(D(fgh))=\phi_{g,h}(D(fgh))\circ \phi_{f,g}(D(fgh)).
$$
$$
\begin{aligned}
\calo_{X,f}(D(fg))
&=\calo_{\spec ((S_\bullet)_f)_0}({D(g^{\deg f}/f^{\deg g})})\\
&=((S_\bullet)_f)_0)_{g^{\deg f}/f^{\deg g}},
\end{aligned}
$$
\underline{Claim}:
$$
((S_\bullet)_f)_0)_{g^{\deg f}/f^{\deg g}}\cong ((S_\bullet)_{fg})_0
$$
This isomorphism is already discussed in Theorem~\ref{chap4thm:graded_ring_localization_iso}.
We have the isomorphisms
$$
((S_\bullet)_f)_0)_{(gh)^{\deg f}/f^{\deg g+\deg h}}\cong ((S_\bullet)_{fgh})_0.
$$
Similarly, we have other two isomorphisms and the $\phi_{f,h}(D(fgh))$ is the composition of these ring isomorphism
\begin{center}
\tiny
\begin{tikzcd}
 &  &  &  & \calo_{X,h}(D(fgh))=((S_\bullet)_h)_0)_{(fg)^{\deg h}/h^{\deg f+\deg g}} \arrow[ldd] \\
 &  &  &  &  \\
\calo_{X,f}(D(fgh))=((S_\bullet)_f)_0)_{(gh)^{\deg f}/f^{\deg g+\deg h}} \arrow[rrr] \arrow[rrrrdd, "\phi_{fg}(D(fgh))"'] \arrow[rrrruu, "{\phi_{f,h}(D(fgh))}"] &  &  &  ((S_\bullet)_{fgh})_0 &  \\
 &  &  &  &  \\
 &  &  &  & \calo_{X,g}(D(fgh))=((S_\bullet)_g)_0)_{(fh)^{\deg g}/g^{\deg f+\deg h}} \arrow[luu] \arrow[uuuu, "{\phi_{g,h}(D(fgh))}"']
\end{tikzcd}
\end{center}

This means we can legally glue the structure sheaf of distinguished opens and there is a unique way to define the structure sheaf of $\Proj S_\bullet$.
\end{proof}

\begin{exr}
(Some will find this essential, others will prefer to ignore it.) (Re)interpret the structure sheaf of $\Proj S_\bullet$ in terms of compatible germs.
\end{exr}
\begin{proof}
We have to first discuss what the stalk of the structure sheaf $\calf:=\calo_{\Proj S_\bullet}$ is.

By definition
$$
\calf_{[P]}=\dirlim_{[P]\in U} \calf(U),
$$
where $[P]\in \Proj S_\bullet$ and $P$ is a homogeneous prime in $S_\bullet$.
Assume $[P]\in D(f)$ for a homogeneous element $f\in S_\bullet$, we define the structure sheaf on $D(f)$ to be the structure sheaf of $\spec ((S_\bullet)_f)_0$. $P$ corresponds to an ordinary prime $\scp\in ((S_\bullet)_f)_0$.

Recall the form of stalks of a affine scheme, we have the isomorphism
$$
\calf_{[P]}\cong (((S_\bullet)_f)_0)_\scp=:(S_{(f)})_\scp.
$$
This notation is terrifying, the second subscript means localize at an element and the forth subscript means localize at a prime $\scp$. And we will abbreviate it a little by writing $((S_\bullet)_f)_0$ as $S_{(f)}$, do not confuse it with localize at a prime ideal $(f)$. 

For an open set $U\subset\Proj S_\bullet$, $U$ is covered by $\{D(g)\}_{g\in G}$. Each stalk at $[P]$ should be interpreted as the equivalence class
$$
\calf_{[P]}:=\coprod_{(g,\scp)\in P}(S_{(g)})_{\scp_g}/\sim_{[P]}
$$
where $(g,\scp)\in P$ means $D(g)\ni [P]$ and $P\cap S_{(g)}=\scp$.
$$
(S_{(g)})_{\scp}\ni s_{g,\scp}\sim_{[P]} s_{f,\scq}(S_{(f)})_{\scq_f}\text{ iff } (g,\scp)\in P, (f,\scq)\in P\text{ and }
 s_{q,\scp}/1= s_{f,\scq}/1\in S_{(fg)}.
$$
Then we can interpret the sections in $\calf(U)$ via compatible germs
$$
\left\{\prod_{[P]}s_{[P]}\in\prod_{[P]\in U} \calf_{[P]}: \begin{matrix}
&\forall [Q]\in U, \exists D(f)\ni [Q] \text{ and }\tilde{s}\in S_{(f)}\text{ s.t. }\\
&\tilde{s}|_\scq/\sim_{[K]}=s_{[K]},\forall (f,\scq)\in K,\forall [K]\in D(g)\cap U
\end{matrix}
 \right\}
$$
where $\tilde{s}|_\scq$ means taking $\tilde{s}/1\in (S_{(f)})_\scq$.
\end{proof}

\begin{exr}
Check that defining $\proj^n_A$ as $\Proj A[x_1,...,x_n]$ agrees with our earlier construction by patching up specific affine schemes. (How do you know that the $D(x_i)$ cover $\Proj A[x_0,...,x_n]$?)

\end{exr}
\begin{proof}
We have already checked the new definition $\Proj A[x_1,..,x_n]$ is compatible with patching up all $D(f)$. $D(x_i)$ is only a special case. We only need to check $D(x_i)$ indeed covers $\Proj$
$V(x_i)=\Proj A[x_1,...,x_n] -D(x_i)$ and
$$
\cap_i^n V(x_i)=V\left(\sum_i^n (x_i)\right)=V((x_1,...,x_n))
$$
But $(x_1,...,x_n)$ is $S_+$ in the case $S_\bullet=A[x_1,..,x_n]$.
We have 
$$
\Proj A[x_1,..,x_n]-\cup_i D(x_i)=\cap_i V(x_i)=V(S_+)=\emptyset.
$$
\end{proof}

\begin{exr}
Suppose that $k$ is an algebraically closed field. We know from Exercise 4.4.F that the closed points of $\proj^n_k$ , are in bijection with the points of classical projective space. With our new definition of projective space, a point of classical projective space corresponds to a homogeneous prime ideal of $k[x_0,...,x_n]$. Which homogeneous prime ideal is it?
\end{exr}
\begin{proof} Assume $a_0\neq 0$,
$[a_0,...,a_n]$ corresponds to $(x_{1/0}-a_1/a_0,...,x_{n/0}-a_n/a_0)$ in the affine chart $D(x_0)$.

$(x_{1/0}-a_1/a_0,...,x_{n/0}-a_n/a_0)\subset (k[x_0,...,x_n]_{x_0})_0$ would corresponds to a homogeneous prime ideal in $k[x_0,...,x_n]$, as shown in~\ref{chap4exr:homogeneous_ideal_characterization}. We can construct the homogeneous prime by constructing it degree by degree and we omit the process and only list the final result (because we know it)
$$
\left(x_1-\frac{a_1}{a_0}x_0,...,x_n-\frac{a_n}{a_0}x_0\right)
$$
\end{proof}

\begin{exr}
If $S_\bullet$ is generated in degree $1$, and $f \in S_+$ is homogeneous, explain how to define $V(f)$ ``in'' $\Proj S_\bullet$ the vanishing scheme of $f$. 
(Warning: $f$ in general isn?t a function on $\Proj S_\bullet$ ). Hence define $V(I)$ for any homogeneous ideal $I$ of $S_+$.
\end{exr}
\begin{proof}
It is customary to define the topological space of $V(f)$ as the the homogeneous prime ideals containing $f$. The trouble is how define the structure sheaf on $V(f)$. 

In the affine case $V(I)=\spec A/I$. In the projective space assume $I$ or $f$ homogeneous, we can still define the vanishing scheme to be
$$
\Proj S_\bullet/I,\ \ \ \Proj S_\bullet/(f),
$$
which are topologically homeomorphic to $V(I)$ and $V(f)$ in $\proj S_\bullet$.


These two are well-defined schemes,
the subtlety lies in how to interpret them as subschemes of $\Proj S_\bullet$. These would be a long story and the key word is closed subscheme.
\end{proof}

\begin{exr}
Suppose $k$ is algebraically closed. Describe a natural bijection between one-dimensional subspaces of $V$ and the closed points of $\proj V$. Thus this construction canonically (in a basis-free manner) describes the one-dimensional subspaces of the vector space $V$.
\end{exr}
\begin{proof}
We define $\proj V$ to be $\Proj (\sym^\bullet V^\vee)$ and we already know $\sym^\bullet V^\vee\cong k[x_0,...,x_n]$, where $x_0,..,x_n$ form a basis  of $V^\vee$.

$\sym^\bullet V^\vee$ can be naturally regarded as functions on $V$ by sending $x_i$ to $x_i(v)$.

A closed point in $\Proj k[x_0,...,x_n]$ is of the form $[a_0,...,a_n]$ or $\left(x_1-\frac{a_1}{a_0}x_0,...,x_n-\frac{a_n}{a_0}x_0\right)$. We can define the corresponding one dimensional vector space by 
$$
\left\{v\in V: \left(x_1-\frac{a_1}{a_0}x_0,...,x_n-\frac{a_n}{a_0}x_0\right)v=0\right\}.
$$
In return, we can define the closed point to be the ideal of functions in $\sym^\bullet V^\vee$ that vanishes on the one dimensional linear space.
This gives a natural bijection.
\end{proof}

\chapter{Some properties of schemes}
\section{Topological properties}
\begin{exr}
Show that $\proj^n_k$ is irreducible.
\end{exr}
\begin{proof}
In this case $S_\bullet=k[x_1,..,x_n]$ and $S_+=(x_1,...,x_n)$.
Assume $\proj^n_k$ can be written as union of two proper closed subsets: $V(I_1)\cup V(I_2)$, where $I_1,I_2$ are two homogeneous ideal in $k[x_1,..,x_n]$, also we require $I_1,I_2$ not contain $S_+$.

$(0)$ is still a homogeneous prime ideal in $k[x_1,..,x_n]$ because it is integral domain and it does not contain the irrelevant ideal. W.l.o.g, assume $V(I_1)\ni [(0)]$, which means $I_1\subset (0)\Lrta I_1=(0)$. But $V(I_1)=V((0))=\Proj S_\bullet=\proj^n_k$, which contradicts the assumption that $V(I_1)$ is a proper closed subset.
\end{proof}

\begin{exr}\label{chap5exr:bijection_closed_generic_point_scheme}
Exercise~\ref{chap3exr:bijection_closed_subsets_points_affine} showed that there is a bijection between irreducible closed subsets and points for affine schemes (the map sending a point $p$ to the closed subset $\overline{\{p\}}$ is a bijection). Show that this is true of schemes in general.
\end{exr}
\begin{proof}
Given a scheme $X$, we want a bijection
$$
\begin{aligned}
X&\lrta \{Z\subset X:Z\text{ closed irreducible}\}\\
z&\longmapsto \overline{\{z\}}.
\end{aligned}
$$ 
We already showed this bijection exists for affine schemes. This map can be naturally extended to general scheme, with the Zariski topology.

We only need to show 
\begin{itemize}
\item $\overline{\{z\}}$ is irreducible in $X$, and
\item $\overline{\{z\}}\neq\overline{\{x\}}$ if $z\neq x$, and
\item Each irreducible closed subset in $X$ is of the form.
\end{itemize}

Each $\{z\}$ is irreducible in $X$, the closure of a irreducible subset is irreducible by~\ref{chap3exr:irreducible_subset_dense_open_irreducible_closure}.

If $\overline{\{z\}}=\overline{\{x\}}$, they are covered by some affine open subsets $U_i$. By definition of closure, each $U_i$ contain $z,x$. $\overline{\{z\}}\cap U_i=\overline{\{z\}}\cap U_i$ are the closure of $z,x$ in the subset topology of $U_i$, they are irreducible in $U_i$ because $\{z\},\{x\}$ are irreducible in $U_i$ with subset topology. Then because the bijection exists in affine schemes, there is a unique generic point corresponding to $\overline{\{z\}}\cap U_i$, $z=x$ in $U_i$, which means $z=x$ in $X$.

Given an irreducible subset $Z\subset X$. Let $U$ be an affine open subset such that $Z\cap U\neq \emptyset$. $Z\cap U$ is an open set in $Z$ with subspace topology on $Z$. By~\ref{chap3exr:irreducible_subset_dense_open_irreducible_closure}, $Z\cap U$ is dense and irreducible in $Z$, with subspace topology. Then it is also irreducible in $U$ with subspace topology on $U$. Hence it corresponds to a unique generic point in $U$. $\exists z\in U\cap Z$ s.t. $Z\cap U=\overline{\{z\}}'$ is the closure of $z$ in $Z\cap U$. However, we already know $Z\cap U$ is dense in $Z$, therefore $\overline{\{z\}}=Z$.

We have established to bijection between the irreducible closed subsets and generic points of the irreducible closed subsets.
\end{proof}


\begin{exr}
Prove that if $X$ is a scheme that has a finite cover $X = \cup_{i=1}^n \spec A_i$ where $A_i$ is Noetherian, then $X$is a Noetherian topological space (We will soon call a scheme with such a cover a \textbf{Noetherian scheme}) Hint: show that a topological space that is a finite union of Noetherian subspaces is itself Noetherian.
\end{exr}
\begin{proof}
We have shown in~\ref{chap3exr:Noetherian_rings_Noetherian_topological_spaces} that affine scheme corresponding to a Noetherian ring is Noetherian topological space.

It suffices to check finite union of Noetherian subspaces $X=\cup_i^nX_i$ is itself Noetherian.

A topological space is called Noetherian if it satisfies the descending chain condition for closed subsets. Given an descending chain of closed subsets in $X$.
$$
Z_1\supset Z_2\supset\cdots\supset Z_k\supset \cdots
$$
$Z_k=\cup_i^n(X_i\cap Z_k)$, $X_i\cap Z_k$ would give a descending chain of closed subsets in the subspace $X_i$ with subspace topology on $X_i$. Each would stabilize after $m_i$. Choose $m=\max\{m_i\}_{1\leq i\leq n}$, the descending chain $\{Z_k\}=\{\cup_i^n(X_i\cap Z_k)\}$ would also stabilize after $m$.
\end{proof}


\begin{exr}\label{chap5exr:quasicompact_scheme_union_affine_subschemes}
Show that a scheme $X$ is quasicompact if and only if it can be written as a finite union of affine open subschemes. (Hence $\proj^n_A$ is quasicompact for any ring $A$.)
\end{exr}
\begin{proof}
One direction is easy, if a scheme $X$ can be written as union of affine open subschemes $X_i$. Given a open cover of $X$, its intersection with affine open subschemes also given an open cover of $X_i$.
 Affine open subscheme is quasicompact, we can select a finite subcover of these intersection covers, and the union gives a finite subcover of the initial cover.

 For the reverse direction, consider a quasicompact scheme $X$. By definition, for each point $x$, there is an affine open subscheme $U_x$ contain $x$. These affine open subscheme give an open cover of $X$. $\cup_{x\in X}U_x$, we can select a finite subcover $\cup_{x\in I}U_x$ of $X$. Then the scheme $X$ is the union $\cup_{x\in I}U_x$. Then $X$ can be written as a finite union of affine open subschemes.

 In particular, $\proj^n_A$ is quasicompact.
\end{proof}

\begin{exr}\label{chap5exr:quasicompact_scheme_have_closed_points}
QUASICOMPACT SCHEMES HAVE CLOSED POINTS. Show that if $X$ is a quasicompact scheme, then every point has a closed point in its closure. Show that every nonempty closed subset of $X$ contains a closed point of $X$. In particular, every nonempty quasicompact scheme has a closed point.
\end{exr}
\begin{proof}
We already showed that $\overline{\{z\}}$ is a irreducible subset in $X$. Hence it suffices to show only every nonempty closed subset of a quasicompact scheme $X$ contains a closed point of $X$. 

Given a closed subset $Z\subset X$. Observe that a point $p\in Z$ is closed in $Z$ iff $p$ is closed in $X$. 

$X$ is a finite union of affine opens $\cup_i^nU_i$ by
~\ref{chap5exr:quasicompact_scheme_union_affine_subschemes}. Consider a finite sub-collection of $\{U_i\}_{1\leq i\leq n}$ that intersects with $Z$, w.l.o.g, we assume it is $\{U_i\}_{1\leq i\leq m}$. Then a point $p\in Z$ is closed iff $p$ is closed in each $\{U_i\}_{1\leq i\leq m}$ that contains $z$.

A closed point in $U_1$ corresponds to maximal ideals in $\spec A_1$. $Z\cap U_1$ is not empty, hence contains at least one $[\scp]$ and $\overline{\{[\scp]\}}\subset Z\cap U_1$. $\scp$ is contained in a maximal ideal $\scm$. $[\scm]\in \overline{\{[\scp]\}}$, which means there is at least one closed point in $Z\cap U_1$.

We can find a closed point (in $U_1$)  $z_1\in Z\cap U_1$. If $z_1$ is also closed in other $U_i$, we are done. 

If not, we can find in $\overline{\{z_1\}}\cap U_2$ a closed point $z_2$ such that $z_2$ is closed in $U_2$. $z_2\notin U_1$ because if $z_2\in U_1$, $\overline{\{z_1\}}\cap U_1\neq \emptyset$, which means $z_1$ is not closed point in $U_1$ in the first place. 

If $z_2$ is closed in all $U_i$ (besides $U_1$ ) that contains it, we are done. If not, assume $z_2$ is not closed point in $U_3$. $\overline{\{z_2\}}\subset \overline{\{z_1\}}$. We can find $z_3\in \overline{\{z_2\}}\cap U_3$, $z_3$ can't lie in $U_1$ or $U_2$, because it would imply $z_1$ not close point in $U_1$, $z_2$ not close point in $U_2$.

We can induct on these process until we find $z_m$ is closed in all $U_i$ that contains it. The process would terminate because we started with a finite cover with affine opens. 

Any closed subset in quasicompact scheme contains a closed point. In particular, a quasicompact scheme contains a closed point and the closure of each point in a quasicompact scheme contains a closed point.
\end{proof}

\begin{exr}\label{chap5exr:characterization_quasiseparatedness}
Show that a scheme is quasiseparated if and only if the intersection of any two affine open subsets is a finite union of affine open subsets.
\end{exr}
\begin{proof}
The $\Lrta$ direction is trivial, affine open subsets are quasicompact, if the scheme is quasiseparated, then the intersection of any two affine open subsets $A,B$ would be quasicompact and by~\ref{chap5exr:quasicompact_scheme_union_affine_subschemes} this intersection should be finite union of affine open subsets in $A\cap B$. Because $A\cap B$ is endowed with subspace topology of $X$, $A\cap B$ is also finite union of affine open subsets in $X$.

For the reverse direction. In a scheme any quasicompact subset is a finite union of affine open subsets, $Y=\cup_i^n U_i$ and $Z=\cup_{j}^m V_j$. Then the intersection of $Y$ and $Z$ should be 
$$
Y\cap Z=\cup_{1\leq i\leq n,1\leq j\leq m} U_i\cap V_j.
$$
If each $U_i\cap V_j$ can be written as finite union of affine open subsets, $Y\cap Z$ can be written as finite union of affine open subsets, which is equivalent to $Y\cap Z$ is quasicompact by~\ref{chap5exr:quasicompact_scheme_union_affine_subschemes}
\end{proof}

\begin{exr}
Show that affine schemes are quasiseparated.
\end{exr}
\begin{proof}
We have to firstly characterize a quasicompact open subset in $\spec A$.
(Do not try to characterize all affine opens in an affine scheme in general, it is difficult. The hint seems to be a trap)

Any open subset $U$, in $\spec A$ can written as a union of distinguished base. If in addition this set is quasicompact, we can write it as a finite union of $D(f_i)$.

$$
Y\cap Z=\cup_i^nD(f_i)\cap \cup_j^mD(g_j)=\cup_{1\leq i\leq n,1\leq j\leq m}D(f_i)\cap D(g_j)=\cup_{1\leq i\leq n,1\leq j\leq m}D(f_ig_j)
$$
which can be written as a finite union of affine opens. Then $Y\cap Z$ is quasicompact by~\ref{chap5exr:quasicompact_scheme_union_affine_subschemes}.

Then intersection of any two quasicompact opens is also quasicompact. Affine scheme is quasiseparated.
\end{proof}

\begin{exr}\label{chap5exr:characterization_quasicompact_And_quasiseparated}
Show that a scheme $X$ is quasicompact and quasiseparated if and only if $X$ can be covered by a finite number of affine open subsets, any two of which have intersection also covered by a finite number of affine open subsets.
\end{exr}
\begin{proof}
$\Lrta$ direction: $X$ can be written as a finite union of affine open subsets because $X$ is quasicompact by~\ref{chap5exr:quasicompact_scheme_union_affine_subschemes}. If in addition $X$ is quasiseparated, each intersection of these affine opens can be written as a finite union of affine open subsets by~\ref{chap5exr:characterization_quasiseparatedness}.

$\Llta$ direction: Assume $X$ can be written as a finite union of affine open subsets such that each pairwise intersection can be written as finite union of affine open subsets. Then $X$ is quasicompact by~\ref{chap5exr:quasicompact_scheme_union_affine_subschemes}. 

The quasiseparatedness part of the reverse direction is tricky. It can be cleanly proved using the diagonal morphism arguments, as shown in the first answer to this
\href{https://math.stackexchange.com/questions/9809/quasiseparated-if-finitely-covered-by-affines-in-appropriate-way}{StackExchange question}. We will give a more topological proof here.

Notice that we have proved affine schemes are quasicompact and quasiseparated. The specified affine opens are quasicompact and quasiseparated in the subspace topology. Given $Z$, $Y$ two quasicompact subsets in $X$. Each is covered by finitely many of specified affine opens.
$$
Z\cap Y=\cup_i(Z\cap Y\cap U_i)=\cup_i(Z\cap U_i)\cap (Y\cap U_i)
$$

\underline{Claim}: $Z\cap U_i$, $Y\cap U_i$ are quasicompact in $U_i$ for all $i$.

If we have the claim, we know $(Z\cap U_i)\cap (Y\cap U_i)$ is quasicompact because $U_i$ is quasiseparated. And a finite union of quasicompact space is always  quasicompact. Hence we are done.

Then the problem reduces to prove the claim:

\underline{proof of the claim}: $Z$ is quasicompact in $X$, therefore can be written as finite union of affine opens $Z=\cup_k Z_k$.
$$
Z=\cup_{k,i}(Z_k\cap U_i)
$$
$Z_k$ is affine open we can choose a distinguished base of $W_{k,i,\alpha}:=D(f_{k,i,\alpha})$, each $W_{k,i,\alpha}$ is quasicompact.

$Z_k\cap U_i$ is an open set in $Z_k$, hence can be written as union of $W_{k,\alpha}$ (not necessarily finite). We have
$$
Z=\cup_{k,i}(Z_k\cap U_i)=\cup_{k,i}\cup_{\alpha\in L_{ki}}W_{k,i,\alpha},
$$
where $L_{ki}$ is not finite in general. But because $Z$ is quasicompact, there is a finite subcover, we can choose a finite subset $F_{ki}$ for each $L_{ki}$ so that
$$
Z=\cup_{k,i}\cup_{\alpha\in F_{ki}}W_{k,i,\alpha}.
$$
(Now $\cup_{\alpha\in F_{ki}}W_{k,\alpha}$ no longer cover $Z_k\cap U_i$ in general)
Intersect it with $U_i$, we get
$$
Z\cap U_1=\cup_{k,i}\cup_{\alpha\in F_{ki}}W_{k,i,\alpha}\cap U_1.
$$
The condition that ``$U_i\cap U_j$ can be written as a finite union of affine opens $\cup_{l\in M_{ij}}V_{ij;l}$ '' would be used in the next step, where $M_{ij}$ is a finite index set depending on $i,j$.
$$
Z\cap U_1=\cup_{k,i}\cup_{\alpha\in F_{ki}}W_{k,i,\alpha}\cap U_1\overset{(*)}{=}\cup_{k,i}\cup_{\alpha\in F_{ki}}W_{k,i,\alpha}\cap U_1\cap U_i=\cup_{k,i}\cup_{\alpha\in F_{ki}}\cup_{l\in M_{1i}} W_{k,i,\alpha}\cap V_{1i;l}.
$$
where $(*)$ holds because each $W_{k.i.\alpha}\subset U_i$. $W_{k,i,\alpha}$ and $ V_{1i;l}$ are quasicompact subset in $U_i$ and their intersection is quasicompact because $U_i$ is quasiseparated. $Z\cap U_1$ is finite union of quasicompact subsets, therefore is compact. We proved the claim there fore the quasiseparatedness part of the reverse direction.
\end{proof}

\begin{exr}\label{chap5exr:Proj_A-scheme_quasicompact_quasiseparated}
Show that all projective $A$-schemes are quasicompact and quasiseparated. 
\end{exr}
\begin{proof}
By definition a projective $A$-scheme $X$ is defined to be $\Proj S_\bullet$, where $S_\bullet$ is a finitely generated graded ring over $A$. A graded ring is finitely generated iff $S_+$ is finitely generated. Assume $S_+=(x_1,...,x_n)$, where each $x_i$ is homogeneous, we can consider an open cover of $X$:
$$
D(x_i).
$$
We already showed that $D(x_i)$ is affine open subset that is isomorphic to $\spec ((S_\bullet)_{x_i})_0$.

$$
D(x_i)\cap D(x_j)=D(x_ix_j).
$$

In summary, $X=\Proj S_\bullet$ can be covered by a finite union of affine open subsets and each intersection of these affine open subsets can be written again as finite union of affine open subsets. Then we can conclude that $X=\Proj S_\bullet$ is quasicompact and quasiseparated by~\ref{chap5exr:characterization_quasicompact_And_quasiseparated}.
\end{proof}

\begin{exr}
Let $X = \spec k[x_1, x_2,...]$, and let $U$ be $X - [\scm]$ where $\scm$ is the maximal ideal $(x_1,x_2,...)$. Take two copies of $X$, glued along $U$ (``affine $\infty$-space with a doubled origin'', see Example 4.4.5 and Exercise~\ref{chap4exr:affine_doubled_origin}  for ``finite-dimensional'' versions). Show that the result is not quasiseparated.
\end{exr}
\begin{proof}
We glue two copies of $\spec k[x_1,x_2,...]$ $X,X'$along the open
$$
U:=\spec k[x_1,x_2,...]-[\scm],
$$
and denote the resulting scheme as $\affn^\infty_{(2)}$.
We showed in Exercise~\ref{chap3exr:affine_scheme_quasicompact}(b) that $U$ is not quasicompact in $\spec k[x_1,x_2,...]$.

Denote the the quotient map $\pi:X\coprod X'\lrta \affn^\infty_{(2)}$.
$\affn^\infty_{(2)}$ is equipped with quotient topology, which make the quotient map continuous. The images $\pi(X)$ and $\pi(X')$ in $\affn^\infty_{(2)}$ are quasicompact $\spec B$ is quasicompact for any ring $B$. And $\pi (U)$ is not quasicompact because $\cup_{i\in\intg}\pi(D(x_i))$ is an open cover of $\pi(U)$ with no finite subcover.

$\pi(X)\cap \pi(X')=\pi(U)$. Intersection of two quasicompact open subsets is not quasicompact, therefore $\affn^\infty_{(2)}$ is not quasiseparated.
\end{proof}

\section{Reducedness and integrality}

\begin{exr}\label{chap5exr:reducedness_stalk_local}
(REDUCEDNESS IS A stalk-local PROPERTY, I.E., CAN BE CHECKED AT STALKS). Show that a scheme is reduced if and only if none of the stalks have nonzero nilpotents. Hence show that if $f$ and $g$ are two functions (global sections of $\calo_X$) on a reduced scheme that agree at all points, then $f = g$.
\end{exr}
\begin{proof}
``$\Lrta$'': A scheme $X$ is reduced  then $\calo_X(U)$ is reduced for all open subsets $U\subset X$. Assume $[(f\in \calo_X(U);U)]\in\calo_{X,p}$ is nilpotent, which means $f|_V^n=0\in\calo_X(V)$ for some open neighborhood $V$ of $p$. But $\calo_X(V)$ is reduced ring, $\Lrta f|_V=0$ and $[(f\in \calo_X(U);U)]=0$, which implies there is no non-zero nilpotents in $\calo_{X,p}$.

``$\Llta$'': Given $f\in\calo_X(U)$, assume $f^n=0$. Then its image to stalks are also nilpotent. The condition that stalks are reduced implies the image of $f$ to each stalk is zero. By~\ref{chap2exr2.4.A}, we know the morphism 
$$
\calo_X(U)\lrta \prod_{p\in U}\calo_{X,p}
$$
is injective, therefore we know $f=0\in\calo_X(U)$.

We have to clarify what it means by ``agree on all points''. It means given any affine open subset $U$, $(U,\calo_X(U))\cong (\spec A,\calo_{\spec A})$, $f|_U\equiv g|_U\mod \scp,\forall [\scp]\in \spec A$. 
$$
f|_U\equiv g|_U\mod \scp
$$
means $f|_U\equiv g|_U$ divers only by a nilpotent. Because $\calo_X(U)$ is reduced, we have $f|_U=g|_U$. $f,g$ agree when restricted to a base of topology, then $f,g$ have same germs on each stalks which means $f=g$.
\end{proof}
\begin{remark}
We should have mentioned  this before. When $f=g+n$, where $n\in\sqrt{(0)}$, they are generally not equal when taking germs
$$
\frac{n}{1}\neq \frac{0}{1}\in A_\scp.
$$
For example, $A=k[x,\epsilon]/(\epsilon^2)$, $f=x^2,g=x^2+\epsilon$. $\epsilon/1\neq 0/1\in A_{(\epsilon)}$.
\end{remark}

\begin{exr}\label{chap5exr:reduced_ring_reduced_affine_scheme}
If $A$ is a reduced ring, show that $\spec A$ is reduced. Show that $\affn^n_k$ and $\proj^n_k$ are reduced.
\end{exr}
\begin{proof}
 
We state a stronger version: 

\textbf{$\spec A$ is reduced scheme iff $A$ is a reduced ring.}

One direction is already contained in the definition of reduced scheme.

For the other direction, we prove an algebraic fact. Any localization of a reduced ring is reduced.

 Given any multiplicative subset $S\subset A$. Assume $a/s\in S^{-1}A$ is nilpotent $(a/s)^n=0$, there exists $t\in S$ s.t. $ta^n=0\in A$, but then we have $(ta)^n=0\in A$ and $A$ being reduced, hence $ta=0$. $\Lrta a/s=0\in S^{-1}A$.

Exercise~\ref{chap5exr:reducedness_stalk_local} also translate to: $A$ is reduced iff $A_\scp$ is reduced for all $[\scp]\in\spec A$. We will give a algebraic proof here. Remember that taking radical is commutative with localization, $\mathfrak{N}(A)_\scp=\mathfrak{N}(A_\scp)$. Also notice that being zero is a local property for $A$-modules, each $A_\scp$ begin reduced implies $\mathfrak{N}(A)$  is zero.

We even have a stronger version: $A$ is reduced iff $A_\scm$ is reduced for all maximal ideals in $A$. (Because the ``being zero'' local property of $A$-module says an $A$ module $M$ is zero iff $M_\scm$ is zero module for all maximal ideals in $A$.)

But note that being reduced can not be checked only on closed points for general schemes, we need quasicompactness, see Exercise~\ref{chap5exr:quasicompact_scheme_reducedness_closed_point} below.


Then $\affn^n_k$ is reduced because $k[x_1,...,x_n]$ is reduced. On the other hand, $\proj^n_k$ is reduced because it can be covered by $\affn^n_k$ and hence each stalk is reduced.
\end{proof}



\begin{exr}
Show that  $\left(k[x, y]/(y^2, xy)\right)_x$  has no nonzero nilpotent elements.
(Possible hint: show that it is isomorphic to another ring, by considering the geometric picture. Exercise~\ref{chap3exr:3.2.L} may give another hint.) Show that the only point of $\spec k[x, y]/(y^2 , xy)$ with a nonreduced stalk is the origin.

\end{exr}
\begin{proof}
Geometrically $\spec \left(k[x, y]/(y^2, xy)\right)_{\bar{x}}=$ the distinguished open $D([x])$ in $\spec k[x, y]/(y^2, xy)$. It is equivalent to the non-vanishing loci of $x$ in $V(y^2,xy)$ in the affine plain, which is the affine line with origin removed. This geometric picture gives us some intuition

\underline{Claim}:
$$
\left(k[x, y]/(y^2, xy)\right)_{\bar{x}}=k[x]_{x},
$$
where $\bar{x}$ is the image of $x$ in the quotient ring. A general element in $k[x, y]/(y^2, xy)$ is of the form
$$
p(\bar{x})+a\bar{y},
$$
where $p$ is a polynomial in $k[x]$, $a$ is a constant in $k$ and $\bar{x}\bar{y}=0$ and $\bar{y}^2=0$.
A general element in $\spec \left(k[x, y]/(y^2, xy)\right)_{\bar{x}}$ is of the form
$$
\frac{p(\bar{x})}{\bar{x}^n}.
$$
$\bar{y}/1=0$ because $\bar{x}\bar{y}=0$.
 There is a well-define ring morphism
$$
\frac{p(\bar{x})}{\bar{x}^n}\mapsto \frac{p(x)}{x^n},
$$
Which is obviously an isomorphism.

$k[x]_x$ is reduced because $k[x]$ is reduced. We then know there is no nonzero nilpotents in $\left(k[x, y]/(y^2, xy)\right)_{\bar{x}}$.

Consider the ring of $\spec \left(k[x, y]/(y^2, xy)\right)$. The prime ideals of this ring corresponds to prime ideals in $k[x,y]$ that contains $(y^2,xy)$. Hence a general prime ideal in this ring is of the form $(\bar{x}-a, \bar{y})$ or $(\bar{y})$. 
When $a\neq 0$, $\bar{x}\notin (\bar{x}-a,\bar{y})$. $\bar{x}$ is invertible in $\left(k[x, y]/(y^2, xy)\right)_{(\bar{x}-a,\bar{y})}$, hence it is further localization of $\left(k[x, y]/(y^2, xy)\right)_{\bar{x}}$ and should be reduced.

When $a=0$, we focus on the ring
$$
\left(k[x, y]/(y^2, xy)\right)_{(\bar{x},\bar{y})}.
$$
The multiplicative set consists of elements like $b(\bar{x})+c\bar{y}$, where $b(\bar{x})$ is a polynomial of $\bar{x}$ and the constant term $b_0\neq 0$. $\bar{y}(b\bar{x}+c\bar{y})^n=b_0^n\bar{y}\neq 0$, therefore $\bar{y}/1\neq 0/1$. And $\bar{y}/1$ is nilpotent.

$\left(k[x, y]/(y^2, xy)\right)_{(\bar{x},\bar{y})}$ contains at least one nonzero nilpotent $\bar{y}/1$, hence is nonreduced.
\end{proof}


\begin{exr}\label{chap5exr:quasicompact_scheme_reducedness_closed_point}
If $X$ is a quasicompact scheme, show that it suffices to check reducedness at closed points. Hint: Do not try to show that reducedness is an open condition (see Remark 5.2.2). Instead show that any nonreduced point has a nonreduced closed point in its closure, using Exercise~\ref{chap5exr:quasicompact_global_section_vanishes_on_point_nilpotent}. (This result is interesting, but we won’t use it.)
\end{exr}
\begin{proof}
Follow the hint, we prove the contrapositive: If a quasicompact scheme is not reduced then at least one of its closed points has non-reduced stalk.

By definition, we know if $X$ is not reduced then at least one of its stalk is not reduced by~\ref{chap5exr:reducedness_stalk_local}. Denote this non-reduced point by $p$, then there is a closed point $z\in \overline{\{p\}}$ by~\ref{chap5exr:quasicompact_scheme_have_closed_points}. (I suppose the hint in the body of this problem is wrong). The stalk $\calo_{X,p}$ is a localization of the stalk $\calo_{X,z}$ at closed point. $\calo_{X,p}$ is not reduced, therefore $\calo_{X,z}$ is also non-reduced.
\end{proof}



\begin{exr}\label{chap5exr:quasicompact_global_section_vanishes_on_point_nilpotent}
Suppose $X$ is quasicompact, and $f$ is a function that vanishes at all points of $X$. Show that there is some $n$ such that $f^n = 0$. Show that this may fail if $X$ is not quasicompact. (This exercise is less important, but shows why we like quasicompactness, and gives a standard pathology when quasicompactness doesn’t hold.)
\end{exr}
\begin{proof}
$f$ is a global section on a quasicompact scheme $X$ that vanishes at each point. 

$X$ is quasicompact and can be written as a finite union of affine open subschemes $\cup_i U_i$, where each $U_i\cong \spec A_i$. We then have $f|_{U_i}\calo_X(U_i)=A_i$. $f|_{U_i}$ vanishes on each point $[\scp]\in \spec A_i$ and hence is an element in the nilpotent $\mathfrak{A_i}$, $\Lrta f_{U_i}^{n_i}=0$ for some $n_i\in\intg$. There are only finitely many such affine opens, we can choose $n:=\max_i\{n_i\}$ and then $f^n|_{U_i}=0\in \calo_X(U_i)$ and there is only one element $0$ to glue to in the global section.

For the counterexample when $X$ is not quasicompact, we follow the hint and consider the union of infinite disjoint union of $\spec A_n$, where $A_n=k[\epsilon]/(\epsilon^{n})$. Each $\spec A_n$ consists of only one point $[(\epsilon)]$. Consider the function $f=\coprod_n \epsilon$. We know $f$ vanishes at each each point but is not nilpotent in the global section.
\end{proof}

\begin{exr}\label{chap5exr:integral=irreducible+reduced}
Show that a scheme $X$ is integral if and only if it is irreducible and reduced. (Thus we picture integral schemes as: ``one piece, no fuzz'')
\end{exr}
\begin{proof}
In a scheme $X$ is integral, $\calo_X(U)$ is integral for all open subsets, hence $\calo_X(U)$ is also reduced because integral domain has no nonzero zero divisors.

An integral scheme should be irreducible. Assume contrarily $X$ is reducible, and can be written as union of two closed subsets $X=Y\cup Z$. Define the complements $U:=X-Y$ and $V=X-Z$, we know $U,V$ are nonempty opens and their have empty intersection. The structure sheaf $\calo_X(U\cup V)=\calo_X(U)\times \calo_X(V)$ which is not integral in general.


For the reverse direction, we would use the hint~\ref{chap3exr:locally_ringed_space_function_value}.

\underline{Claim}: Any open subset $U$ in a irreducible space $X$ is irreducible with the subset topology.

For any other nonempty open subset $U\subset X$. By~\ref{chap3exr:irreducible_subset_dense_open_irreducible_closure}, $U$ is dense in $X$. Assume $U$ is reducible and can be written as $(Y\cap U)\cup (Z\cap U)=(Y\cup Z)\cap U$, where $Y,Z$ are nonempty proper closed subsets in $X$. 
We also know $Y\cup Z\supset U$. Hence $\overline{Y\cup Z}=\overline{Y}\cup \overline{Z}=Y\cup Z\supset \overline{U}=X$, which means $X$ is reducible, contradiction. 

Given $f,g\in\calo_X(U)$ and $fg=0\in\calo_X(U)$. Recall~\ref{chap3exr:locally_ringed_space_function_value}, the set where $f,g$ vanishes is closed. (Notice here vanishes means $f_pg_p\in\scm_p$). $U$  could be covered by the closed set $V$ $f$ vanishes and the closed set $W$ where $g$ vanishes. Then because $X$ is irreducible, we know $W$ or $V$ must be the whole space $X$. Assume $f$ vanishes on every point in $X$, we know $f=0\in \calo_X(U)$ because $U$ is a reduced scheme~\ref{chap5exr:reducedness_stalk_local}.
\end{proof}

\begin{exr}\label{chap5exr:affine_integral_scheme}
Show that an affine scheme $\spec A$ is integral if and only if $A$ integral domain.
\end{exr}
\begin{proof}
The ``only if'' direction is obvious, because $\Gamma(\spec A,\calo_{\spec A})=A$.


Consider the ``if'' direction. We simply have $\spec A$ is irreducible if $A$ is an integral domain by~\ref{chap3exr:Integral_domain_irreducible}. Also an integral domain is always reduced ring, hence $\spec A$ is  reduced by~\ref{chap5exr:reduced_ring_reduced_affine_scheme}. We know $\spec A$ is irreducible and reduced if $A$ is integral domain, hence $\spec A$ is integral by~\ref{chap5exr:integral=irreducible+reduced}.
\end{proof}


\begin{exr}\label{chap5exr:generic_point_integral_scheme_function_field}
Suppose $X$ is an integral scheme. Then $X$ (being irreducible) has a generic point $\eta$. Suppose $\spec A$ is any nonempty affine open subset of $X$. Show that the stalk at $\eta$, $\calo_{X,\eta}$ is naturally identified with $K(A)$, the fraction field of $A$. This is called the \textbf{function field} $K(X)$ of $X$. It can be computed on any nonempty open set of $X$, as any such open set contains the generic point. The reason for the name: we will soon think of this as the field of \textit{rational functions} on $X$
\end{exr}
\begin{proof}
Suppose $X$ is an integral scheme. $X$ is itself irreducible closed and hence corresponds to a unique generic point by~\ref{chap5exr:bijection_closed_generic_point_scheme}. This generic point $\eta$ is contained in any open subset of $X$. Specifically, it is contained in affine open $U:=\spec A$ and corresponds to $[(0)]\in \spec A$, where $A$ is an integral domain. $\calo_{X,\eta}\cong A_{[(0)]}=K(A)$, where localization at the prime ideal $(0)$ is isomorphic to the fraction field $K(A)$.
\end{proof}

\begin{exr}
Suppose $X$ is an integral scheme. Show that the restriction maps $\res_{U,V}: \calo_X(U) \lrta \calo_X(V)$ are inclusions so long as $V\neq \emptyset$. Suppose $\spec A$ is any nonempty affine open subset of $X$ (so $A$ is an integral domain). Show that the natural map $\calo_X(U) \lrta  \calo_{X,\eta} = K(A)$ (where $U$ is any nonempty open subset) is an inclusion.
\end{exr}
\begin{proof}
In fact, the first question relies on the second. Because we have the composition of maps
$$
\calo_X(U)\overset{\res_{U,V}}{\lrta}\calo_X(V)\lrta\calo_{X,\eta},
$$
it suffices to prove $\calo_X(U)\lrta \calo_{X,\eta}$ is injective. ($\res_{U,V}$ has to be injective because its composition with an injection is again an injection.)

\underline{Claim}: $\calo_X(U)\lrta K(X)$ is injective.

Assume a section $f\in\calo_X(U)$ and $f_\eta=0$. \underline{Want}: $f=0\in \calo_X(U)$. It suffices to prove that $f_\eta\Lrta f|_W=0$ for all affine open subsets contained in $U$. 

For an affine open $W=\spec A$, the natural morphism
$$
\calo_X(W)\lrta \calo_{X,\eta}=K(A)
$$
as shown in~\ref{chap5exr:generic_point_integral_scheme_function_field}. This map must be the canonical inclusion because the stalk at $[(0)]$ is a further localization of any $A_f=\calo_{\spec  A}(D(f))$.

We have prove both questions.
\end{proof}

\section{Properties of schemes that can be checked ``affine-locally''}
\begin{exr}
Show that locally Noetherian schemes are quasiseparated.
\end{exr}
\begin{proof}
In a locally Noetherian scheme, each affine open is isomorphic to $\spec A$ for some Noetherian ring $A$. By~\ref{chap5exr:characterization_quasiseparatedness}, we can check whether any intersection of two affine open can be written as finite union of affine opens.

Given $U=\spec A$ and $V=\spec B$ for $A,B$ Noetherian rings, consider then intersection $U\cap V$.  By proposition 5.3.1, we know $U\cap V$ can be written as union of open sets such that are simultaneously distinguished open subschemes in $\spec A$ and $\spec B$. But notice by~\ref{chap3exr:open_subset_Noetherian_space_quasicompact} that open subset of Noetherian space is quasicompact, we can cover $U\cap V$ with finitely many simultaneous distinguished opens, which are affine opens.

We don't need Prop 5.3.1 to prove this statement but it is more clear to see which of the affine opens are selected.
\end{proof}


\begin{exr}\label{chap5exr:Noetherian_scheme_finite_irreducible_compoenents}
Show that a Noetherian scheme has a finite number of irreducible components. (Hint: Proposition 3.6.15.) Show that a Noetherian scheme has a finite number of connected components, each a finite union of irreducible components.
\end{exr}
\begin{proof}
\underline{Claim}: The underline topological space of a Noetherian scheme is Noetherian.

Note that a Noetherian scheme is locally Noetherian and quasicompact, we know a Noetherian scheme $X$ can be covered by finitely many Noetherian affine subschemes.

\underline{subclaim}: A finite union of Noetherian subspace is also Noetherian.

We can prove the subclaim directly by checking the descending chain condition on closed subsets. Given a chain $V_0\supset V_1\supset \cdots$, each $V_i$ can be covered by finitely many $U_j$, where $U_j$ is Noetherian subsapce. It renders a descending chain of closed subset in $U_j$ with subset topology, hence will stabilize after $n_j$. Choose $N=\max\{n_j\}$, we have $V_n=\cup_n (V_N\cap U_j)$ would stabilize when $n\geq N$.


Then in particular, we know a Noetherian scheme is Noetherian topological space and by Prop 3.6.15, we know it has only finitely many irreducible components.


Exercise~\ref{chap3exr:topology_irreducible_components_connected_components} already showed that connected components in topological space $X$ are unions of irreducible components of $X$. But in the case of Noetherian scheme, there are only finitely many irreducible components, hence $X$ is the union of finitely many connected components and each connected component is union of finitely many irreducible components.
\end{proof}

\begin{exr}\label{chap5exr:Noetherian_scheme_integral_stalk}
Show that a Noetherian scheme $X$ is integral if and only if $X$ is nonempty and connected and all stalks $\calo_{X,p}$ are integral domains. Thus in ``good situations'', integrality is the union of local (stalks are integral domains) and global (connected) conditions.
\end{exr}
\begin{proof}
We follow the hint and recall Exercise~\ref{chap5exr:integral=irreducible+reduced} says integral $=$ irreducible $+$ reduced. Note that ``being reduced'' is stalk-local~\ref{chap5exr:reducedness_stalk_local}. 

$\Lrta$ direction is easy:

``Irreducible'' alone guarantees that $X$ is connected. And also note that localization of integral domain is also integral domain because $\frac{a}{s}\frac{b}{t}=\frac{0}{1}\Lrta \exists r\in S, rab=0\Lrta ab=0$.

We want to prove the reverse direction under Noetherian hypothesis.

Each stalk being integral domain alone implies each stalk is reduced, hence the scheme $X$ is reduced.

\underline{Want}: under Noetherian hypothesis, connected $+$ stalks being integral domain $\Lrta$ $X$ is irreducible.

$X$ is Noetherian, we know each connected component is union of finitely many irreducible components. We know $X$ is  connected and hence $X$ is the union of finitely many irreducible components, and some of the pair-wise intersections of these irreducible components are nonempty to make sure their union is  connected.

Assume $Y,Z$ are two of these irreducible components and their intersection is non-empty. $p\in Y\cap Z$. \underline{Claim}: $\calo_{X,p}$ is not integral domain.

Choose an affine open $U\ni p$, $U=\spec A$ is the union of $U\cap Y$ and $U\cap Z$, $U$ is not irreducible and $A$ is not integral domain~\ref{chap3exr:Integral_domain_irreducible}. Assume $U\cap Y=V(\scq_1)$ and $U\cap Z=V(\scq_2)$, where $\scq_1,\scq_2$ are two prime ideals in $A$, they should be minimal prime because there is a bijection~\ref{chap3exr:irreducible_components_minimal_pirmes_bij} between the minimal primes and irreducible components in $\spec A$. $V(\scq_1)\cap V(\scq_2)=V(\scq_1+\scq_2)$. $p=[\scp]\in V(\scq_1+\scq_2)\Lrta \scq_1+\scq_2\subset \scp$. The stalk at $p$, is $\calo_{X,p}=A_\scp$. Then there are at least two minimal prime ideals in $A_\scp$, therefore $A_\scp$ can not be integral domain.

In short, we have
\textbf{the minimal primes in $\calo_{X,p}$ is in one to one correspondence with the irreducible components that pass through $p$}.

The only possibility is that $X$ is irreducible. Combined with our proof that $X$ is reduced, we know $X$ is integral scheme.
\end{proof}

\begin{exr}\label{chap5exr:quasiprojective=>finite_type;Open_subschem_of_Noetherian_projective_scheme_is_of_finite_type}\ 
\begin{enumerate}[label=(\alph*)]
\item (quasiprojective implies finite type)
If $X$ is a quasiprojective $A$-scheme (Definition 4.5.9), show that $X$ is of finite type over $A$. If $A$ is furthermore assumed to be Noetherian, show that $X$ is a Noetherian scheme, and hence has a finite number of irreducible components.
\item Suppose $U$ is an open subscheme of a projective $A$-scheme. Show that $U$ is locally of finite type over $A$. If $A$ is Noetherian, show that $U$ is quasicompact, and hence quasiprojective over $A$, and hence by (a) of finite type over $A$. Show this need not be true if $A$ is not Noetherian. Better: give an example of an open subscheme of a projective $A$-scheme that is not quasicompact, necessarily for some non-Noetherian $A$.
\end{enumerate}
\end{exr}
\begin{proof}\ 
\begin{enumerate}[label=(\alph*)]
\item a quasiprojective $A$-scheme $X$ is a quasicompact open subscheme of projective $A$-scheme $Y$.

A projective $A$-scheme is of the form $Y:=\Proj S_\bullet$, where $S_\bullet$ is finitely generated graded ring over $A$. It is locally of finite type over $A$. Because it is covered by affine open sets $\spec ((S_\bullet)_f)_0$, where $(S_\bullet)_f$ is finitely generated over $A$ because it is generated by $\{S_\bullet, 1/f\}$ hence its degree zero piece is finitely generated over $A$.

Recall the Affine Communication Lemma~5.3.2 and also check that ``being finite over $A$'' is indeed an affine-local property. Then $Y=\cup_{i\in I}\spec B_i$ and $\spec B_i$ is of finite generated algebra over $A$, $\Lrta$ any affine open $U\subset X$ is of finite type over $A$. In particular $X$ can be cover covered by affine open $\spec_{j\in J} B_j$, where $B_j$ is finitely generated $A$-algebra, therefore $X$ is of locally finite type over $A$. In addition $X$ is quasicompact, we know $X$ is of finite type over $A$.

If $A$ is furthermore assumed to be Noetherian, then each affine open of $X$ is isomorphic to $\spec B_i$, where $B_i$  is Noetherian because $B_i$ is finitely generated over a Noetherian ring $A$, hence is also Noetherian by Hilbert basis theorem. Then it has finite number of irreducible components by~\ref{chap5exr:Noetherian_scheme_finite_irreducible_compoenents}.

\item For same reasons, an open subscheme in a scheme of locally of finite type over $A$ is locally of finite type over $A$, therefore any open subscheme $U$ of a projective $A$-scheme $Y$ is locally of finite type over $A$.

If $A$ is Noetherian, we know $U$ is locally Noetherian. Also notice that a projective $A$-scheme is quasicompact and quasiseparated by~\ref{chap5exr:Proj_A-scheme_quasicompact_quasiseparated}. 

\textbf{A projective $A$-scheme over a Noetherian ring $A$ is a Noetherian scheme.} 


By~\ref{chap3exr:open_subset_Noetherian_space_quasicompact}, an open subset of a Noetherian topological space is quasicompact. We know $U$ is quasicompact and locally Noetherian, hence is Noetherian. $U$ is also quasiprojective and of finite type over $A$.

As for the counterexample. Consider the silly example 4.5.11, $\Proj A[T]\cong\spec A$, where $x_0=T$ is of degree $1$. Also recall~\ref{chap3exr:affine_scheme_quasicompact}, $\spec A$ can have non-quasicompact open subscheme.
\end{enumerate}
\end{proof}
\begin{remark}\label{chap5rmk:open_subscheme_Noetherian}
The part of (locally)Noetherian and quasicompactness does not depend on (Projective), hence we have in general:

\textbf{
Any open or closed subscheme of a (locally) Noetherian scheme is
(locally) Noetherian.}
\end{remark}

\begin{exr}\ 
\begin{enumerate}[label=(\alph*)]
\item
Showthat $\spec k[x_1,...,x_n]/I$ is an affine $k$-variety if and only if $I \subset k[x_1,...,x_n]$ is a radical ideal.
\item Suppose $ I \subset  k[x_0,...,x_n]$ is a radical graded ideal. Show that $\Proj k[x_0,...,x_n]/I$ is a projective $k$-variety. (Caution: The example of $ I = (x_0^2, x_0x_1,...,x_0x_n)$ shows that $\Proj k[x_0,...,x_n]/I$ can be a projective $k$-variety without $I$ being radical.)
\end{enumerate}
\end{exr}
\begin{proof}
\begin{enumerate}[label=(\alph*)]
\item $X:=\spec k[x_1,...,x_n]/I$ is of finite type over $k$.

$k$ is Noetherian, $\spec k[x_1,...,x_n]/I$ is a projective $k$-scheme. We know open subscheme of a projective $A$-scheme over a Noetherian ring is of finite type over $A$ by Exercise~\ref{chap5exr:quasiprojective=>finite_type;Open_subschem_of_Noetherian_projective_scheme_is_of_finite_type}.

\underline{Want}:
$\spec k[x_1,...,x_n]/I$ is reduced iff $I$ is radical ideal.

Recall~\ref{chap5exr:reduced_ring_reduced_affine_scheme}, an affine scheme $\spec A$ is reduced iff $A$ is reduced.
The nilradical of $k[x_1,...,x_n]/I$ is just $\sqrt{I}$. We know $k[x_1,...,x_n]/I$ is reduced iff $I$ is radical.
\item We define projective $k$-variety to be reduced $k$-scheme. $\Proj k[x_1,...,x_n]/I$ is automatically of finite type over $k$ by~\ref{chap5exr:quasiprojective=>finite_type;Open_subschem_of_Noetherian_projective_scheme_is_of_finite_type}. 

We know $I$ is radical homogeneous ideal, then $S_\bullet:=k[x_1,...,x_n]/I$ is a reduced graded ring, hence $((S_\bullet)_f)_0$ is reduced and so is any further localizations. We have checked any stalks of $\Proj k[x_1,...,x_n]/I$ is reduced ring, hence it is a reduced projective $k$-scheme, i.e., a projective $k$-variety.
\end{enumerate}
\end{proof}

\begin{exr}
Show that a point of a locally finite type $k$-scheme is a closed point if and only if the residue field of the stalk of the structure sheaf at that point is a finite extension of $k$. Show that the closed points are dense on such a scheme (even though it needn't be quasicompact, cf. Exercise~\ref{chap5exr:quasicompact_scheme_have_closed_points}).
\end{exr}
\begin{proof}
We have a stronger version, for a locally finite type $k$-scheme $X$, the followings are equivalent:
\begin{enumerate}[label=(\alph*)]
	\item the point $p\in X$ is closed point
	\item the field extension $k\inj K(\calo_{X,p})$  is finite
	\item the field extension $k\inj K(\calo_{X,p})$ is algebraic
\end{enumerate}
$p$ is contained in an affine open subset $U=\spec A$ of $X$, where $A$ is a $k$-algebra and it corresponds to a maximal ideal $\scm$ in $A$. By Hilbert Nullstellensatz, $K(\calo_{X,p})\cong A/\scm$ is a finite field extension of $k$.

A finite field extension is automatically an algebraic field extension. These proves $(a)\Lrta (b),(c)$.

For $(c)\Lrta (a)$, $p$ is a point in $\spec A$, and $\scp$ is the corresponding prime ideal. $K(\calo_{X,p})\cong Frac(A/\scp)$ is algebraic over $k$, (It is integral over $k$, hence the integral domain $A/\scp$ is integral over $k$.) Recall a lemma in commutative algebra, Lemma 1.9 in~\cite{gortz2010algebraic}. Or one can use directly \cite[\href{http://stacks.math.columbia.edu/tag/00GS}{Lemma 00GS}]{stacks-project}

$B$ is integral over $A$, and both of them are integral domains. Then
${A}$ is a field $\Llrta$ ${B}$ is a field.

We know then $A/\scp$ is a field, therefore $\scp$ is a maximal ideal ideal, which means $p$ is a closed point.


\underline{Claim}:Let $k$ be a field. Let $X$ be a $k$-scheme locally of finite type. Let $U$ be a non-empty affine open subset of $X$. Then any closed point of $U$ is a closed point of $X$.

In an affine open subset $U\subset X$, the residue field at a closed point $p$ is finite field extension of $k$, which means it is also a closed point in $X$.

Hence, there is a closed point in each  affine open, therefore the closed point is dense in $X$.
\end{proof}

\begin{exr}
Suppose $X$ is a reduced, finite type $\cplx$-scheme. Define the corresponding complex analytic prevariety $X_{an}$. (The definition of an analytic prevariety is the same as the definition of a variety without the Hausdorff condition.) 
\end{exr}
\begin{proof}
It would be a long story, \textbf{To Be Added}. The major reference should be Chapter$4-5$ of ~\cite{neeman2007algebraic}
\end{proof}


\begin{exr}
Finish the proof of Proposition 5.3.3(a).
\end{exr}
\begin{proof} $A=(f_1,...,f_n)$.
It remains to show that, given a strictly increasing chain of ideals $I_1\subsetneq I_2\subsetneq \cdots$ in $A$, we can construct a strictly increasing chain of ideals in $A_{f_i}$ by
$$
I_{i,1}\subset I_{i,2}\subset \cdots,
$$
where $I_{i,j}=I_j\otimes_A A_{f_i}$. It suffices to show that for each $j$,
$$
I_{i,j}\subsetneq I_{i,j+1},
$$
for some $i$. (Notice that this does mean there is a strictly increasing chain is some $A_{f_i}$)
But it means there is  at least a chain $I_{i,1}\subset I_{i,2}\subset \cdots,$ not stabilize in the Noetherian ring $A_{f_i}$, contradiction.

Each $I_{i,j}$ could be interpreted as localization of the $A$-module $I_j$ at $f_i$ by~\ref{chap1exr:localization_extension_of_scalar}

\underline{Claim}:
$A=(f_1,...,f_n)$. Consider an $A$-module $M$, it is zero iff $M_{f_i}=0$ for all $f_i$. 

We only need the part: $\forall f_i, M_{f_i}=0\Lrta M=0$. Assume $M_{f_i}=0$ for all $i$, choose $m\in M$, exists $n_i\geq 0$ s.t., $f_i^n m=0$. On the other hand $1=\sum_{i=1}^nr_i f_i$, choose $N\geq \sum_{i=1}^n n_i+1$. $1\cdot m=1^Nm=(\sum_i r_i f_i)^N=\sum_i a_i f_i^{n_i}m=0$. Hence $M=0$.

Or equivalently, we can prove it by observing that each $M_{\scp}$ is further localization of some $M_{f_i}$. All $M_{f_i}=0$ implies all $M_\scp=0$, which in return implies $M=0$~\cite[\href{http://stacks.math.columbia.edu/tag/00HN}{Lemma 00HN}]{stacks-project}.

We apply the above claim to $I_{j+1}/I_j\neq0$, therefore we have for each $j$,
$$
I_{i,j+1}/I_{i,j}\neq 0
$$
for some $i$.
\end{proof}

\begin{exr} In text, we want to prove the proposition 

``Given $A=(f_1,...,f_n)$, then $A$ is a finitely generated $B$-algebra iff each $A_{f_i}$ is a finitely generated $B$-algebra.''

with some ``partition of unity argument''.
Make this argument precise.
\end{exr}
\begin{proof}
One direction is clear, $\spec A_{f_i}$ is generated by $A,1/f_i$. If $A$ is a finitely generated $B$-algebra, so is $A_{f_i}$.

$A$ is generated by $f_i$. $1=\sum_i c_i f_i$ for $c_i\in A$. $A_{f_i}$ is finitely generated as $B$-algebra and we assume its generators are $\{r_{ij}/f_i^{k_j}\}$. 

\underline{Claim:}$\{f_i\}\cup \{c_i\}\cup \{r_{ij}\}$ generate $A$ as a $B$-algebra.

$r$ is an element in $A$ and restricts to $r/1\in A_{f_i}$, which could be written as
$$
\sum_j b_{ij} \frac{r_{ij}}{f_{i}^{k_{ij}}}=\frac{\sum_j b_{ij}r_{ij}f_i^{K_{ij}} }{f_i^{Q_{i}}}=:\frac{P_i(r_{ij},f_i)}{g_i},
$$
where  we define $P_{i}:=\sum_j b_{ij}r_{ij}f_i^{K_{ij}}\in B[r_{ij},f_i]$ and $g_i:=f_i^{Q_i}$ for simplicity. 

$P_i/g_i$ and $P_j/g_j$ would agree on $A_{g_ig_j}$, which means
$$
(g_i g_j)^{m_{ij}}(P_i g_j-P_j g_i)=0.
$$
By taking $m=\max\{m_{ij}\}$, (the indexing set is finite), we can simplify the notation
$$
(g_i g_j)^{m}(P_i g_j-P_j g_i)=0,
$$
for all $i,j$.

Let $T_i:=g_i^m P_i$ and $h_j:=g_j^{m+1}$, we have
$$
T_i h_j-T_j h_i=0.
$$

Now, because $D(f_i^n)=D(f_i)$, we have $\spec A=\cup_i D(h_i)$ and this implies $1=\sum_i s_i h_i$, where each $s_i$ can be written as a polynomial of $c_i,f_i$.
Define 
$$
r':=\sum_i s_i T_i.
$$
$r'$ restricts to each $T_i/h_i$, because 
$$
r'h_j=\sum_i {s_i T_i}h_j=\sum_i s_i h_i T_j=T_j.
$$
In other words, it restricts to each 
$$
\sum_j b_{ij} \frac{r_{ij}}{f_{i}^{k_{ij}}}
$$
and by the identity axiom $r=r'$, which can be expressed as a polynomial in $B[\{f_i\}\cup \{c_i\}\cup \{r_{ij}\}]$

\end{proof}


\section{Normality and factoriality}
\begin{exr}
Show that integrally closed domains behave well under localization: if $A$ is an integrally closed domain, and $S$ is a multiplicative subset not containing $0$, show that $S^{-1}A$ is an integrally closed domain.
\end{exr}
\begin{proof}
Assume $x\in K(A)$ is integral over $S^{-1}A$. It satisfies a monic polynomial
$$
x^n+\frac{a_{n-1}}{s_{n-1}}x^{n-1}+\cdots +\frac{a_{0}}{s_0}=0.
$$
Define $s:=\prod_i s_i$ and if multiply the above equation by $s^n$, we get a polynomial of $sx$ with coefficients in $A$. Because $A$ is integral closed, we know $sx\in A$, hence $x\in S^{-1}A$, which means $S^{-1}A$ is integral closed.
\end{proof}


\begin{exr}
Show that a Noetherian scheme is normal if and only if it is the finite disjoint union of integral Noetherian normal schemes.
\end{exr}
\begin{proof}
$\Lrta$ direction: A Noetherian scheme $X$ is normal at least implies each stalks of this scheme is integral. Recall each connected component is open subscheme and open subscheme of Noetherian scheme is Noetherian (Remark~\ref{chap5rmk:open_subscheme_Noetherian}).

 By~\ref{chap5exr:Noetherian_scheme_integral_stalk}, each connected component of $X_i$ is a integral scheme. Also by~\ref{chap5exr:Noetherian_scheme_finite_irreducible_compoenents}, a Noetherian scheme has only finitely many connected components.

Then we know $X$ is a disjoint finite union of integral Noetherian normal schemes.

$\Llta$ direction is clear. Because normality is stalk-local, we know a finite union of integral Noetherian normal schemes is normal and it is Noetherian because it can be covered by finite Noetherian affine opens. 
\end{proof}

\begin{exr}
If $A$ is an integral domain, show that $A = \cap A_\scm$, where the intersection runs over all maximal ideals of $A$.
\end{exr}
\begin{proof}
If we interpret them as subsets in $K(A)$, we have clearly
$$
A\subset \cap A_\scm.
$$
As for the reverse direction, we claim:
$$
K(A)\backslash A\subset K(A)\backslash \cap A_\scm.
$$
If nonempty, pick $s\in K(A)\backslash A$, we can construct the ideal of denominators of $s$
$$
I_s:=\{r\in A:rs\in A\}
$$
$I_s\neq A$ because $1\notin I_s$. Then $I_s\subset \scm$ for some maximal ideal. $s\notin A_\scm$.
\end{proof}

\begin{exr}\label{chap5exr:5.4.D}
One might naively hope from experience with unique factorization domains that the ideal of denominators is principal. This is not true. As a counterexample, consider our new friend $A = k[w, x, y, z]/(wz - xy)$ (which we first met in Example 4.4.12, and which we will later recognize as the cone over the quadric surface), and $w/y = x/z \in K(A)$. Show that the ideal of denominators of this element of $K(A)$ is $(y, z)$.
\end{exr}
\begin{proof}
Denote the element $s:=w/y=x/z$. we easily check that $(y,z)\subset I_s$.

Assume $t\in I_s$, we have 
$$
tw/y=tx/z\in A,
$$
hence $tw=ya$. Then $t$ can not contain monomials in $w,x$. Furthermore, $t$ can not contain terms with only $w,x$ because $w^nx^m$ can not be a multiple of $y$. As a result, $t\in (y,z)$.
\end{proof}

\begin{exr}\label{chap5exr:localization_of_UFD}
Show that any nonzero localization of a unique factorization domain is a unique factorization domain.
\end{exr}
\begin{proof}
Recall the definition of a UFD. A UFD is an integral domain in which every non-zero element $x$ of $R$ can be written as a product of irreducible elements and a unit and this product is unique up to order of irreducible elements and units. 

An equivalent characterization is ``an integral domain $R$ in which every non-zero element can be written as a product of unit and  prime elements of $R$''.

We prove the statement based on the second characterization. $S\subset R$ is a multiplicative set and let $T$ be the set of all prime elements that divides an element of $S$, and let $M$ be the set of all prime element not in $T$.

\underline{Claim}: $p\in T$ iff image of $p$ in $S^{-1}R$ is a unit.

Indeed, if $p\in T$, then there exists $s\in S$, such that $p|s$; Let $x\in R$ with $px=s$. Then
$$
\frac{p}{1}\cdot\frac{x}{s}=\frac{ps}{s}\cdot \frac{x}{s}=1_{S^{-1}R}.
$$

For the ``if'' part, consider $p/1$ a unit, there is an inverse $v/s$, $\exists t\in S$ s.t., $pvt=st$, $pv=s$ because $0\notin S$ and $R$ is integral domain.

\underline{Claim}: if $p\in M$, then $p/1\in S^{-1}R$ is prime.

Assume $p/1|(a/s)(b/t)$, $(p c)/(r)=(ab)/(st)$, because $R$ is iintegral domain, $pcst=ab r$. Then we know $p|ab$ or $p|r$ because $p$ is prime, but we assumed $p\nmid s\forall s\in S$. Then we have $p|a$ or $p|b$ because $p$ is prime, which means $p/1|a/s$ or $p/1|b/t$ in return.

Now given $a=up_1^{b_1}\cdots p_r^{b_r}q_1^{c_1}\cdots q_t^{c_t}$ be a prime factorization of $a$, where $p_i\in T$ and $q_j\in M$. We have
$$
\frac{a}{s}=\frac{u}{s}\left(\frac{p_1s}{s}\right)^{b_1}\cdots \left(\frac{p_rs}{s}\right)^{b_r}\left(\frac{q_1s}{s}\right)^{c_1}\cdots \left(\frac{q_ts}{s}\right)^{c_t}.
$$
Then we know each $a/s\in S^{-1}R$ can be factorized into product of a unit and prime elements.

This proof is a improved version of~\href{https://math.stackexchange.com/q/140596}{This Answer}, in that we don't need to verify the uniqueness of factorization.
\end{proof}

\begin{exr}
Show that unique factorization domains are integrally closed. Hence factorial schemes are normal, and if $A$ is a unique factorization domain, then $\spec A$ is normal. 
\end{exr}
\begin{proof} Assme $R$ is a UFD, denote the fraction field by $K(R)$.
Let $x\in K(R)^{\times}$, say $x=r/s$, $gcd(x,s)=1$, $r,s\in R, s\neq 0$. Suppose $\exists a_1,..,a_n\in R$ s.t.,
$$
x^n+a_1 x^{n-1}+...+a_n=0.
$$
Then after multiplying it by $s^n$, set
$$
r^n=-(a_1 r^{n-1}s+a_2 r^{n-2}s^2+...+a_n s^n)
$$
$\Lrta s|r^n$, $gcd(r^n,s)=1,\Lrta s\in R^\times\Lrta x\in R$.

Then we have ``all stalks of a scheme $X$ are UFD'' $\Lrta$ ``all stalks of a scheme $X$ are normal'', which means factorial schemes are normal. 

If $A$ is a UFD,  each stalk  is a localization of $A$ is  therefore UFD by~\ref{chap5exr:localization_of_UFD} . $\spec A$ is factorial $\Lrta$ $\spec A$ is normal.
\end{proof}

\begin{exr}
Show that the following schemes are normal: $\affn^n_k$ , $\proj_k^n$ , $\spec \intg$. (As usual, k is a field. Although it is true that if A is integrally closed then A[x] is as well, this is not an easy fact, so do not use it here.)
\end{exr}
\begin{proof}
$\spec \intg$ is normal because $\intg$ is UFD.

$\affn^n_k=\spec k[x_1,..,x_n]$, where the polynomial ring over a field $k$ is UFD.

$\proj^n_k=\Proj k[x_0,...,x_n]$ can covered by $k[x_{0/i},...,x_{n/i}](x_{i/i}-1)$, which are UFDs, there fore we know each stalk is UFD. Then we know $\proj^n_k$ is factorial and hence normal.
\end{proof}

\begin{exr}\label{chap5exr:handy_characterization_normal_scheme}
Suppose $A$ is a unique factorization domain with $2$ invertible, and $z^2 - f$ is irreducible in $A[z]$.
\begin{enumerate}[label=(\alph*)]
\item Show that if $f \in A$ has no repeated prime factors, then $\spec A[z]/(z^2 - f)$ is normal. Hint: $B := A[z]/(z^2 - f)$ is an integral domain, as $(z^2 - f)$ is prime in $A[z]$. Suppose we have monic $F(T)\in B[T]$ so that $F(T) = 0$ has a root $\alpha$ in $ K(B)\backslash \ K(A).$ Then by replacing $F(T)$ by $\bar{F}(T)F(T)$, we can assume $F(T) \in A[T]$. Also, $\alpha = g+hz$ where $g,h \in K(A).$ Now $\alpha$ is the root of $Q(T) = 0$ for monic $Q(T) = T^2 - 2gT + (g^2 - h^2f) \in K(A)[T]$, so we can factor $F(T) = P(T)Q(T)$ in $K(A)[T]$. By Gauss’s lemma, $2g, g^2 - h^2f \in A$. Say $g = r/2$, $h = s/t$ ($s$ and $t$ have no common factors, $r, s, t \in A$). Then $g^2 - h^2f = (r^2t^2 - 4s^2f)/4t^2$. Then $t$ is invertible.
\item Show that if $f \in A$ has repeated prime factors, then $\spec A[z]/(z^2 - f)$ is not normal.
\end{enumerate}
\end{exr}
\begin{proof}
\begin{enumerate}[label=(\alph*)]
\item In fact, the hint already gives the detailed solution. We only write here some further explanations.

$A[z]/(z^2-f)$ is integral domain because $z^2-f$ is irreducible and hence is prime element in UFD.

$B$ is equivalent to formally adjoining a pair of roots $\pm \sqrt{f}$. We can take the ``conjugate'' of a monic $F(T)$ be mapping the coefficients $a+b\sqrt{f}$ to $a-b\sqrt{f}$ and we denote the conjugate polynomial $\bar{F}(T)$. Then $\bar{F}(T)F(T)\in A[T]$.

Assume $\alpha\in K(B)\backslash K(A)$ is a root of $F(T)=0$ hence a root of $\bar{F}(T)F(T)=0$,  We replace $F(T)$ by $\bar{F}(T)F(T)$ from now on. \underline{Want}: $\alpha$ is contained in $B$.

The root $\alpha$ can be written as $g+hz$ or equivalently ($g+h\sqrt{f}$), where $g,h\in K(A)$.

Then $\alpha$ is the root of $(T-g-h\sqrt{f})(T-g+h\sqrt{f})=T^2-2gT+g^2-h^2f\in K(A)[T]$, so we can factor $F(T)$ as $P(T)Q(T)\in K(A)[T]$. Recall Gauss' lemma for UFD:

\textbf{Given $A$ is a UFD, and $K(A)$ its fraction field. $F(T)\in A[T]$, if it factors in $K(A)[T]$, then it factors in $A[T]$.} 

Then we know $2g, g^2-h^2f\in A$. Because $2$ is invertible in $A$, we know $g\in A$ hence $h^2f\in A$. Assume $h=r/s$ and $gcd(r,s)=1$, we have $r^2f/s^2\in A$, because $f$ is square free, the only possibility is $s\in A^\times$. Then we know $h\in A$, therefore $\alpha=g+hz=g+h\sqrt{f}\in B$.

We have proved that $B=A[z]/(z^2-f)$ is integral closed, therefore $\spec A[z]/(z^2-f)$ is normal.
\item Assume $f\in A$ has repeated prime factors. 

\underline{Want}: $B:=A[z]/(z^2-f)$ and $B_\scp$ is not integral closed for some $\scp\in \spec B$. By Proposition 5.4.2, it is equivalent to prove $B$ is not integral closed.

\underline{Want}:$F(T)\in B[T]$, find an root of $F(T)=0$, $\beta\in K(B)\backslash K(A)$ s.t., $\beta\notin B$.

Like in $(a)$, we can restrict to the case $F(T)\in A[T]$. Assume $p$ is the repeated prime element in $f$, i.e, $p^2|f$.

We can choose the polynomial $F(T):=T^2-1/p^2 f$, $\sqrt{f}/p$ is a root of $F(T)$ in $K(B)$ and is not an element in $B$.
\end{enumerate}
\end{proof}

\begin{exr}\label{chap5exr:normal_schemes_examples}
Show that the following schemes are normal:
\begin{enumerate}[label=(\alph*)]
\item  $\spec\intg[x]/(x^2-n)$ where $n$ is a square-free integer congruent to $3$ modulo $4$. Caution: the hypotheses of Exercise~\ref{chap5exr:handy_characterization_normal_scheme} do not apply, so you will have to do this directly. (Your argument may also show the result when $3$ is replaced by $2$. A similar argument shows that $\intg[(1+\sqrt{n})/2]$ is integrally closed if $n \equiv 1 (\mod 4)$ is square-free.)
\item
 $\spec k[x_1,...,x_n]/(x^2_1+x_2^2+\cdots+x^2_m)$ where $\char k\neq 2, n\geq m\geq 3$.
\item $\spec k[w, x, y, z]/(wz - xy)$ where $\char\ k\neq  2$. This is our cone over a quadric surface example from Example 4.4.12 and Exercise~\ref{chap5exr:5.4.D}. Hint: Exercise~\ref{chap5exr:quadratic_forms_diagonalization} may help. (The result also holds for $\char k = 2$, but don’t worry about this.)
\end{enumerate}
\end{exr}
\begin{proof}
\begin{enumerate}[label=(\alph*)]
\item $\intg[x]/(x^2-n)\cong B:=\intg[\sqrt{n}]$. Assume a monic polynomial $P(T)\in B[T]$, we can take the conjugate $\bar{P}(T)$ and $\bar{P}(T)P(T)\in \intg[T]$. W.l.o.g, we can set $P(T)\in \intg[T]$. Assume $\alpha=g+h\sqrt{n}\in K(B)\backslash \ratl$ is a root of $P(T)=0$, where $g,h\in \ratl$. $\alpha$ is also a root of $G(T)= T^2 - 2gT + (g^2 - h^2n) \in \ratl[T]$. $P(T)=R(T)G(T)$.

Again, by Gauss's lemma, we have $G(T)\in \intg[T]$.
$$
2g\in \intg
$$
and
$$
g^2-h^2n\in \intg
$$
Assume $g=k/2$ and $h=r/s$, $k,r,s\in\intg$, we have
$$
\frac{k^2}{4}-\frac{r^2n}{s^2}\in \intg
$$
$n=4e+3$ and $k=2q$ or $k=2q+1$. The later case is ruled out. If $k=2q+1$,  the equation reduces to
$$
\frac{1}{4}-\frac{r^2n}{s^2}=m\in \intg
$$
Then we know $r^2 n=s^2/4 +ms^2$ and
$gcd(r,s)=1$ and $n$ is square-free, we know $4|s^2$, assume $s^2=4t^2$, therefore
$$
3r^2\equiv t^2\ (\mod 4), 
$$
where $gcd(r,t)=1$. We know $r,t$ are coprime and have same parity. $r$ and $t$ can't be both even number and have to be both odd. Assume $r=2u+1$ and $t=2v+1$, then the equation reduces to
$$
3\equiv 1\ (\mod 4).
$$
We get the contradiction. As a result, $k=2q$ and $r^2n/s^2\in\intg$, which means $s=1$. $g,h\in \intg$ $\beta=g+h\sqrt{f}\in B$.

The identical argument works when $n$ congruent to $2$ modulo $4$.

\item $\char k\neq 2$, then $2$ is invertible in $k$. For example, in the case
$k[x_1,x_2,x_3]/(x_1^2+x_2^2+x_3^2)$, regard $k[x_1,x_2,x_3]$ as $A[x_1]$, where $A=k[x_2,x_3]$, $A$ is a UFD, $2$ is invertible in $A$, $x_2^2+x_3^2$ is square-free in $A$, and $x_1^2+(x_2^2+x_3^2)$ is irreducible in $A[x_1]$. We can now use~\ref{chap5exr:handy_characterization_normal_scheme} to prove that $k[x_1,x_2,x_3]/(x_1^2+x_2^2+x_3^2)$ is integrally closed. The identical argument can be used to prove for arbitrary $n\geq m\geq 3$.
\item We can diagonalize the quadratic form $wz-xy$.
$$
wz-xy=\left(\frac{w+z}{2}\right)^2-\left(\frac{w-z}{2}\right)^2-\left(\frac{x+y}{2}\right)^2+\left(\frac{x-y}{2}\right)^2.
$$
Then it reduces to a special case of $(b)$ above.

\end{enumerate}
\end{proof}
\begin{exr}\label{chap5exr:quadratic_forms_diagonalization}
Suppose $k$ is an algebraically closed field of characteristic not $2$. (The hypothesis that $k$ is algebraically closed is not necessary, so feel free to deal with this more general case.)
\begin{enumerate}[label=(\alph*)]
\item Show that any quadratic form in $n$ variables can be ``diagonalized'' by changing coordinates to be a sum of at most $n$ squares.
\item Show that the number of squares appearing depends only on the quadratic. For example, $x^2 + y^2 + z^2$ cannot be written as a sum of two squares.
\end{enumerate}
\end{exr}
\begin{proof}
In linear algebra courses, we have discussed the diagonalization of real symmetric matrices and Hermitian matrices. This exercise is a simple generalization to fields with non-2 characteristics.
\begin{enumerate}[label=(\alph*)]
\item Consider a quadratic form 
$$
Q_{n}(\mathbf{x}):=\sum_{1\leq i,j\leq n}a_{ij}x_ix_j,
$$
with coefficients in $k$.
Follow the hints, we can induct on $n$. We firstly prove for $n=2$,
$$
Q_2(\mathbf{x})=a_{11}x_1^2+2a_{12}x_1x_2+a_{22}x_2^2.
$$
Assume $a_{11},a_{22}=0$ and $2a_{12}\neq 0$ (which means $a_{12}\neq 0$ because $\char k\neq 2$), we have $Q_{2}(\mathbf{x})=a_{12}(x_1+x_2)^2-a_{12}x_1^2-a_{12}x_2^2$. Else, one of $a_{11},a_{22}\neq 0$, w.l.o.g, $a_{11}\neq 0$, we get $Q_2(\mathbf{x})=a_{11}(x_1-a_{12}a_{11}^{-1}x_2)^2+(a_{22}-a_{12}^2/a_{11})x_2^2$.

Assume now all $Q_n$ can be diagonalized, consider a quadratic form $Q_{n+1}$, it can be written as 
$$
Q_{n+1}(\mathbf{x},x_{n+1})=P_{n}(\mathbf{x})+2c_{1,n+1}x_1x_{n+1}+...+2c_{n,n+1}x_nx_{n+1}+c_{n+1,n+1}x_{n+1}^2.
$$
Up to reordering, we require $c_{n+1,n+1}\neq 0$.
By induction hypothesis, we can write $P_n(\mathbf{x})$ as $\sum_{i=1}^n b_{ii}y_i^2$, with with possibility all $b_{ii}=0$ and $\mathbf{y}=A^{-1}\mathbf{x}$, where $A$ is a linear transforamtion.
$$
\begin{aligned}
Q_{n+1}(\mathbf{x},x_{n+1})&=\mathbf{y}^{T}B\mathbf{y}+2\sum_{i=1}^n\sum_{j=1}^n A_{ij}c_{i,n+1}y_jx_{n+1}+c_{n+1,n+1}x_{n+1}^2\\
&=\mathbf{y}^{T}B\mathbf{y}+2\sum_{j=1}^n \tilde{c}_{j,n+1}y_jx_{n+1}+c_{n+1,n+1}x_{n+1}^2,
\end{aligned}
$$ 
where $\tilde{c}_{j,n+1}:=\sum_{i=1}^nA_{ij} c_{i,n+1}$.
$$
\begin{aligned}
Q_{n+1}(\mathbf{x},x_{n+1})
&=c_{n+1,n+1}\left(x_{n+1}^2+\sum_{j=1}^n\frac{\tilde{c}_{j,n+1}}{c_{n+1,n+1}}y_j\right)^2+\sum_{i=1}^n\left(b_{ii}-\frac{c_{j,n+1}^2}{c_{n+1,n+1}}\right) y_i^2.
\end{aligned}
$$ 
There are at most $n+1$ terms.
\item Given a quadratic form $Q(\mathbf{x})=\mathbf{x}^{T}M\mathbf{x}$, under a transformation of basis, $\tilde{\mathbf{x}}=A \mathbf{x}$, we can write $Q(\tilde{\mathbf{x}})=\tilde{\mathbf{x}}^T \tilde{M} \tilde{\mathbf{x}}$, where 
$$
\tilde{M}=A^{T}M A.
$$
Assume $\mathbf{v}$ is in the null space of $M$, then $A^{-1}v$ is  in the null space of $\tilde{M}$, because $A$ is a non-degenerate matrix, we know the dimension of  null space is invariant under the change of basis. Hence, the rank of matrix $M$ is preserved under the base change. For quadratic forms, the rank is just the number of independent square terms.
\end{enumerate}
\end{proof}

\begin{exr}
(RINGS CAN BE INTEGRALLY CLOSED BUT NOT UNIQUE FACTORIZATION DOMAINS, ARITHMETIC VERSION). Show that $\intg[\sqrt{-5}]$ is integrally closed but not a unique factorization domain.
\end{exr}
\begin{proof}
This is a classical counter-example.

 $A:=\intg[\sqrt{-5}]$ is integrally closed by~\ref{chap5exr:normal_schemes_examples}(a) because 
 $$-5\equiv 3\mod 4.
 $$

 It is not unique factorization domain because for example $6=2\times 3$ and $6=(1-\sqrt{-5})(1+\sqrt{-5})$.
\end{proof}

\begin{exr}
(RINGS CAN BE INTEGRALLY CLOSED BUT NOT UNIQUE FACTORIZATION DOMAINS, GEOMETRIC VERSION). Suppose $\char k \neq 2$. Let $A = k[w, x, y, z]/(wz - xy)$, so $\spec A$ is the cone over the smooth quadric surface (cf. Exercises 4.4.12 and~\ref{chap5exr:5.4.D}).
\begin{enumerate}[label=(\alph*)]
\item Show that $A$ is integrally closed. (Hint: Exercises~\ref{chap5exr:normal_schemes_examples}(c) and~\ref{chap5exr:quadratic_forms_diagonalization})
\item Show that $A$ is not a unique factorization domain. (Clearly $wz = xy$. But why are $w, x, y$, and $z$ irreducible? Hint: $A$ is a graded integral domain. Show that if a homogeneous element factors, the factors must be homogeneous.)
\end{enumerate}
\end{exr}
\begin{proof}
For (a), we can diagonalize the quadratic form $wz-xy$.
$$
wz-xy=\left(\frac{w+z}{2}\right)^2-\left(\frac{w-z}{2}\right)^2-\left(\frac{x+y}{2}\right)^2+\left(\frac{x-y}{2}\right)^2.
$$
We know $A\cong k[x_1,x_2,x_3,x_4]/(x_1^2-x_2^2-x_3^2+x_4^2)$, then we know it is integrally closed by~\ref{chap5exr:normal_schemes_examples}(b).

As for $(b)$, we clearly have $wz=xy$. Then it suffices to prove $w,x,y,z$ are irreducible.  

Indeed, we can regard $A$ as a graded integral domain because it is a quotient of the polynomial ring by a homogeneous prime ideal. 

\underline{Claim}: If a homogeneous element factors, the factors must be homogeneous. 

An element $f$ is of homogeneous degree $n$ iff each monomial of $f$ is of degree $n$.

Now assume $f=gh$, where $g,h$ are not homogeneous. Assume $g=g_{0}+g_1$, where $g_0$ is the highest degree term in $g$ and $g_1$ is the rest. The degree of monomials in $g_0$ is strictly higher than those in $g_1$. We have $f=g_0 h+g_1 h$, but it is impossible for the two summands to have same highest degrees.

Then we know $x$ is irreducible, because the only degree 
$1$ homogeneous terms are $x,y,z,w$ and we won't have the relation $w=\lambda x$. 
\end{proof}

\begin{exr}
Suppose $A$ is a $k$-algebra, and $l/k$ is a finite extension of fields. (Most likely your proof will not use finiteness; this hypothesis is included to avoid distraction by infinite-dimensional vector spaces.) Show that if $A\otimes_k l$ is a normal integral domain, then $A$ is a normal integral domain as well.
\end{exr}
\begin{proof}
Here the notion ``normal integral domain'' is the same as ``integrally closed''. $A$ can be identified as a subring of $A\otimes_k l$, therefore $A\otimes_k l$ is an integral domain implies that $A$ being an integral domain. 

The proof is sketched in the hint, we just follow it step by step.

 Fix a $k$-basis for $l$, $b_1=1,...,b_d$. There is a dual basis $v_1=1,...,v_d$ such that $trace_{l/k}(v_i b_j)=\delta_{ij}$. We can extend the trace to $Tr:A\otimes_k l\lrta A: a\otimes x\mapsto a\cdot trace_{l/k}(x)$. The map $Tr$ is nondegenerate $Tr(a\otimes x)=0\Lrta a=0$ or $x=0$.

 \underline{Claim}: $1\otimes b_1,...,1\otimes b_d$ forms a free $A$-basis for $A\otimes_k l$. 

 Assume an element $a\otimes x\in A\otimes_k l$, where $x=\sum_i x_i b_i$. We have 
 $a\otimes x=a\otimes \sum_i x_i b_i$ and it equals to $\sum_i (x_i a) (1\otimes b_i)$. It indeed forms an generating set of  $A$-module $A\otimes_k l$. And we can show that $A\otimes_k l$ is a free $A$-module with this basis. Assume $\sum_i r_i(1\otimes b_i)=0$, $\sum_i r_i\otimes b_i=0$. Consider $Tr((1\otimes v_j)\cdot (\sum_i r_i\otimes b_i))=r_j=0$. We have proved that $A\otimes_k l$ is a free $A$-module with the specified basis.

 \underline{Claim}: The diagram commutes and the composition map $A\lrta K(A)\otimes_k l$ is injective.
\[
\begin{tikzcd}
A \arrow[r] \arrow[d] & K(A) \arrow[d] \\
A\otimes_k l \arrow[r] & K(A)\otimes_k l
\end{tikzcd}
\]
$A\lrta K(A)$ is injective and $l$ is a flat $k$-module. Hence the composition map is injective.

\underline{Claim}: $K(A)\otimes_k l\cong K(A\otimes_k  l)$.

We know $A\otimes_k l\subset K(A)\otimes_k l\subset K(A\otimes_k l)$, then it suffices to prove only $K(A)\otimes_k l$ is a field.

A general term in $K(A)\otimes_k l$ is of the form $\frac{a}{c}\otimes x$, $1\otimes b_i$ gives a $K(A)$-basis of $K(A)\otimes_k l$. (it comes from our first claim). By~\ref{exr:chap3_fin-dim_k_domain_is_field}, a finite $k$-algebra which is an integral domain must be a field. We therefore have $K(A)\otimes_k l$ is a  field. 

\underline{Claim}: $K(A)\cap A\otimes_k l=A$, where we regard both $K(A)$ and $A\otimes_k l$ as subrings of $K(A\otimes_k l)$.

$A\subset K(A)\cap A\otimes_k l$ is automatically true. $K(A)\otimes_k l$ is a $K(A)$-vector space with basis $1\otimes b_i$. The subspace $K(A)\cong K(A)\otimes_k 1$ is generated by $1\otimes 1$. $K(A)\cap A\otimes_k l$ is an $A$-submodule of $A\otimes_k l$ generated by $1\otimes 1$, therefore is contained in $A$.

Now we have finished our preparation. Given $P(T)$ a monic polynomial in $A[T]$. We can regard $P(T)$ as an element in $K(A\otimes_k l)[T]$. By assumption, $A\otimes_k l$ is normal integral closed. The root of $P(T)=0$ in $K(A\otimes_k l)$, $\alpha\in A\otimes_k l$. Assume further more $\alpha$ is a root in $K(A)$, we then have $\alpha\in K(A)\cap A\otimes_k l=A$. Whence, we can conclude that $A$ is normal integral domain. 

\end{proof}
\begin{exr}
(UFD-NESS IS NOT AFFINE-LOCAL). Let $A = (\ratl[x, y]_{x^2+y^2})_0$ denote the homogeneous degree $0$ part of the ring $\ratl[x,y]_{x^2 +y^2}$ . In other words, it consists of quotients $f(x, y)/(x^2 + y^2)^n$ , where $f$ has pure degree $2n$. Show that the distinguished open sets $D\left( \frac{x^2}{x^2+y^2}\right)$ and $D\left( \frac{y^2}{x^2+y^2}\right)$ cover $\spec A$. Show that $A_{\frac{x^2}{x^2+y^2}}$ and $A_{\frac{y^2}{x^2+y^2}}$
are unique factorization domains. Finally, show that A is not a unique factorization domain.
\end{exr}
\begin{proof}
By~\ref{exr:open_cover_distinguished}, we know 
$ \frac{x^2}{x^2+y^2}+ \frac{y^2}{x^2+y^2}=1\Lrta D\left( \frac{x^2}{x^2+y^2}\right)\cup D\left( \frac{y^2}{x^2+y^2}\right)=\spec A$.

A general element of $A_{\frac{x^2}{x^2+y^2}}$ is of the form
$$
\frac{\frac{f(x,y)}{(x^2+y^2)^n}}{\frac{x^{2m}}{(x^2+y^2)^m}}=\frac{f(x,y)/x^{2n}}{(x^2+y^2)^{n-m} (x^{2})^{m-n}}=\frac{g(t)}{(1+t^2)^{n-m}}
$$
where $t=\frac{y}{x}$ and $g(t)=f(x,y)/x^{2n}$. Hence we know both $A_{x^2/(x^2+y^2)}$ and $A_{y^2/(x^2+y^2)}$ are isomorphic to $\ratl[t]_{t^2+1}$, which is a localization of UFD, by~\ref{chap5exr:localization_of_UFD}, they are also UFDs.

Finally, $A$ is not a UFD, for example
$$
\left(\frac{xy}{x^2+y^2}\right)^2=\frac{x^2}{x^2+y^2}\cdot \frac{y^2}{x^2+y^2},
$$
where both 
$$
\frac{xy}{x^2+y^2}
$$
and
$$
\frac{x^2}{x^2+y^2}
$$
are irreducible because $x/(x^2+y^2),y/(x^2+y^2)\notin A$ because they are not of degree $0$.
\end{proof}
\section{The crucial points of a scheme: Associated points and primes}
\begin{exr}
Suppose $f$ is a function on $\spec k[x,y]/(y^2,xy)$(i.e.,$f\in k[x,y]/(y^2,xy)$). Show that $\supp f$ is either the empty set, or the origin, or the entire space.
\end{exr}
\begin{proof}Denote the ring $k[x,y]/(xy, y^2)=:A$
$f$ is a global section on $\spec k[x,y]/(xy, y^2)$. By definition 
$$
\supp(f):=\{p\in X=\spec A:f_p\neq 0\text{ in } \calo_{X,p}\}
$$
$\supp f$ is the point $[\scp]$ where $f/1\neq 0$ in $A_\scp$. $f/1\neq 0$ means $fs\neq 0\forall s\in A-\scp$. 

A general element $f$ in $k[x, y]/(y^2, xy)$ is of the form
$$
q(\bar{x})+a\bar{y},
$$
where $q$ is a polynomial in $k[X]$, $a$ is a constant in $k$ and $\bar{x}\bar{y}=0$ and $\bar{y}^2=0$. 

1. Assume $q=0$ and $a=0$, then $f=0$. In this trivial case the support of $f$ is empty set because $fs=0$ for whatever $s\in A$.

2. In the case $f=q(\bar{x})+a \bar{y}$, $q\neq 0,a\neq 0$ 
$$
\bar{y}^2f=0
$$
which implies $\bar{y}^2\in\scp$ for $[\scp]\in \supp f$. But in fact, all the prime ideals in $A$  contains $\bar{y}$. 

\underline{Claim}: if $f=q(\bar{x})+a\bar{y}$ with $q\neq 0, a\neq 0$, then $\supp f=\spec A$.

Indeed the point of $\spec A$ is of the form $(\bar{x}-b,\bar{y}) $ or $(\bar{y})$, for any $\scp_b=(\bar{x}-b,\bar{y})$, the elements in multiplicative subset $S_b=A-\scp_b$ are of  the form $g(\bar{x})+c\bar{y}$ where $g(b)\neq 0$. $(g(\bar{x})+c\bar{y})(q(\bar{x})+a\bar{y})\neq 0$. For $\scp=(\bar{y})$, the multiplicative set are of the form $g(\bar{x})$ and $g(\bar{x})(q(\bar{x})+c\bar{y})\neq 0$.

3. Assume $f$ is $a\bar{y}$, $f$ is not supported on $(\bar{y})$ because $\bar{x}\in A-(\bar{y})$. $f$ is neither supported over $(\bar{x}-b,\bar{y})$ since $\bar{x}\in A-(\bar{x}-b,\bar{y})$.

$f$ is only supported on the origin $(\bar{x},\bar{y})$, where the multiplicative set are just the constants.

We have discussed all the possible cases thus finished the proof.
\end{proof}

\begin{exr}\label{chap5exr:integral_associated_point_generic_point}(ASSUMING \textbf{(A)})
Suppose $A$ is an integral domain. Show that the generic point $[0]$ is the only associated point of $\spec A$
\end{exr}
\begin{proof}
\textbf{Confusion: when we talk about the associated point of $\spec A$, does it mean we choose $M=A$?}
 
We use the definition \textbf{(A)} in page 167 in the case $M=A$ and always assume $A$ to be Noetherian when discussing associated points.

Given $A$  a Noetherian integral domain, $(\spec A,\calo_{\spec A})$ is a Noetherian integral scheme by~\ref{chap5exr:affine_integral_scheme}, which means it is irreducible and reduced by~\ref{chap5exr:integral=irreducible+reduced}. 

The global sections of $\spec A,\calo_{\spec A}$ are just elements in $A$. Now each section $a$ is supported on the whole scheme $\spec A$. The only irreducible component is $\spec A$ itself. Hence the generic point is the only associated point of $\spec A$.
\end{proof}

\begin{exr}(ASSUMING \textbf{(A)})
Show that if $A$ is reduced, $\spec A$ has no embedded points.
\end{exr}
\begin{proof}
What we want to show is 

``the associated points of a scheme $\spec A$ for a reduced ring $A$ are generic points of the irreducible components of $\spec A$''

Follow the hint.

(i): In the case where $A$ is integral domain, we know the only associated point is the generic point of $\spec A$ by Exercise~\ref{chap5exr:integral_associated_point_generic_point}.

(ii) For a general reduced ring $A$, we \underline{Claim}:

If $f$ is a nonzero function on a reduced affine scheme,  then $\supp f = \overline{D(f)}$

Indeed, $D(f)=\{[\scq]\in\spec A|f\notin \scq\}$, $[\scp]\in D(f)\Llrta$ $f\in A-\scp$, $f/1\neq 0\in A_\scp$ because $A-\scp$ is multiplicatively closed and $A-\scp$ does not contain $0$. Then we know $D(f)\subset \supp f$, therefore $\overline{D(f)}\subset\supp f$ because $\supp f$ is closed. 
For the reverse inclusion, consider $[\scp]\in\supp f$ $\Lrta fs\neq 0, \forall s\in A-\scp$. Suppose $[\scp]\in D(g)$ for some $g\in A$, $g\notin \scp$($g\in A-\scp$), $D(f)\cap D(g)=D(fg)$. Then we know $fg\neq 0$. $fq$ is not nilpotent because $A$ is a reduced ring, therefore $D(fg)\neq \emptyset$. Useful exercises are~\ref{chap3exr:3.5.D} and~\ref{chap3exr:3.5.F}

\underline{Claim}: $\overline{D(f)}$ is the union of irreducible component that intersects $D(f)$. 

By Proposition 3.6.15, we know in Noetherian topological space, a non-empty closed set can be expressed uniquely as a finite union of irreducible closed subsets.

 Apply it to $\overline{D(f)}$, it can be expressed as union of $\cup_i D_i$. What we need to prove is that the irreducible closed subset $D_i$ are in fact a irreducible components $X_i$ in $\spec A$. 
Assume by contradiction, $D_i\subsetneq Z_i\subset X_i$ is not a maximal irreducible subset, where $Z_i$ is a minimal irreducible closed subset that contains $D_i$. (i.e. there is no intermediate irreducible closed subset between $D_i$ and $Z_i$). By~\ref{chap3exr:bijection_closed_subsets_points_affine}, there is a reverse inclusion of generic point $\scp_d\supsetneq \scp_z$, where $\overline{\scp_d}=D_i$ and $\overline{\scp_z}=Z_i$. (This also means $\scp_d$ is one of  associated primes of $A$ corresponding to $f$)

 Consider the ring $A/\scp_z$, $\scp_d/\scp_z$ is a non-zero prime ideal in the quotient ring. 

The canonical projection $A\lrta A/\scp_z$ induces a continuous map $\spec A/\scp_z\lrta \spec A$, which can be identified as the inclusion of a closed subset.

We know $\overline{\scp_d/\scp_z}=\overline{\scp_d}\cap \spec A/\scp_z$. On the other hand
$\overline{\scp_d}\subset \overline{D(f)}$, $\overline{\scp_d/\scp_z}=\overline{\scp_d}\cap \spec A/\scp_z\cap \overline{D(f)}=\overline{\scp_d}\cap \overline{D([f])}$, where $[f]$ is the equivalence class of $f$ in $A/\scp_z$. Every subspace of Noetherian space is also Noetherian by~\ref{chap3exr:open_subset_Noetherian_space_quasicompact}. We know $\overline{D([f])}$ is a subspace of Noetherian space $\overline{D(f)}$. Still by Proposition 3.6.15, both $\overline{D([f])}$ and $\overline{D(f)}$
can expressed uniquely as a union of their irreducible components. Also because each irreducible closed subset is contained in an irreducible component. 

We then know
\underline{Lemma:}
\textbf{ the irreducible component $Y_i$ of a subspace $Y$ in Noetherian space $X$ is the intersection of $Y\cap X_i$ with the irreducible component of $X$}.

As a result, $\overline{\scp_d/\scp_z}$ is the irreducible component of $\overline{D([f])}=\supp [f]$.

$\scp_d/\scp_z$ is an associated point of $A/\scp_z$, where $A/\scp_z$ is an integral domain, we then apply $(i)$. $\scp_d/\scp_z$ can not be an embedded point ($\scp_d/\scp_z=0$) contradicting to our setting $\scp_d\subsetneq \scp_z$.

Now we are done.

 

\end{proof}

\begin{exr}
(ASSUMING \textbf{(A)}). Suppose $m \in M$. Show that $\supp m$ is the closure of those associated points of $M$ where $m$ has non-zero germ.(Hint: $\supp m$ is a closed set containing the points described, and thus their closure. Why does it contain no other points?)
\end{exr}
\begin{proof}
Recall what we mean by the support of a section $m$
$$
\supp(m)=\{\scp\in X:m_{\scp}\neq 0\}.
$$
By definition in $\mathbf{(A)}$, associated points are the generic points of irreducible components of $\supp(n)$ for some $n$. In particular the associated described in the problem is contained in $\supp(m)$ hence contains the closure of those associated point.

Denote the set of associated point where $m$ has nonzero germ by $S$. We want to prove $\supp(m)-\overline{S}=\emptyset$. Choose a point $\scq\in \supp(m)$, want $U\ni\scq\Lrta U\cap S\neq \emptyset$. Consider $D(f)\ni \scq\Lrta$ $D(f)\cap \supp(m)\neq \emptyset$. $D(f)$ intersects with as least one irreducible component $W$ of $\supp(m)$ by definition $(A)$, it corresponds to an associated point $\scp$ and $\overline{\scp}=W$. We know $D(f)\cap W\neq \emptyset,\Lrta D(f)\ni \scp$. $\Lrta D(f)\cap S\neq \emptyset$.
\end{proof}
\begin{exr}(Assuming $\mathbf{(A)}$ and $\mathbf{(B)}$)
Show that the locus on $\spec A$ of points $[\scp]$ where $\calo_{\spec A,[\scp]} = A_\scp$ is nonreduced is the closure of those associated points of $\spec A$ whose stalks are nonreduced. (Hint: why do points in the closure of these associated points all have nonreduced stalks? Why can’t any other point have a nonreduced stalk?)
\end{exr}
\begin{proof}
Follow the hint, we first show that the points in the closure of those specified associated point all have nonreduced stalk.

Let $\{\scn_i\}$ be the points with nonreduced stalks and let $\{\scp_j\}$ be the associated points with nonreduced stalks. What we want to show is basically the equality
$$
\overline{\cup_j [\scp_j]}=\{[\scn_i]\}
$$
By $\mathbf{(B)}$, the closure of finite union is the finite union of closures.
$$
\cup_j \overline{[\scp_j]}\overset{?}{=}\{[\scn_i]\}
$$

The containment $\supset$ follows from the following fact: if $f\in A_\scn$ is a nilpotent, consider $\supp(f)$.
$0\neq f\in Nil(A_\scn)$ means $f$ is contained in all prime ideals in $\scn$ and also $\scn\in \supp f$.

Choose $\scq$ to the associated point corresponding to the irreducible component containing $[\scn]$. Then $f/1$ is nilpotent in $A_\scq$. This is because $\scq\subset \scn$ and $$f\in Nil(A_\scn)=\cap_{prime\ \sca\in \scn}\sca\subset \cap_{prime\ \scb\in \scq}\scb=Nil(A_\scq).$$

The reverse inclusion is easier because assume
$\sca\in \cup\overline{[\scp_i]}$, this means some $A_\scp$ is a further localization of $A_\sca$, where $A_\scp$ is nonreduced. We then know $A_\sca$ is nonreduced because we have shown before that localization of a reduced ring is reduced in Exercise~\ref{chap5exr:reduced_ring_reduced_affine_scheme}.
\end{proof}
%---------------------------------------------------------------------------------------

\part{Morphisms}
%----------------------------------------------------------------------------------------
%	CHAPTER 6-10
%----------------------------------------------------------------------------------------
\chapter{Morphisms of schemes}
\section{Introductions}
\section{Morphism of ringed space}
\begin{definition}
A \textbf{morphism of ringed spaces $\pi:X\lrta Y$}  is a continuous map of topological spaces along with a map $\calo_Y\lrta \pi_*\calo_X$, which we think of as a ``pullback map''.
\end{definition}
\begin{exr}\label{chap6exr:glue_morphisms_of_ringed_spaces}
(Glue morphisms of ringed spaces). Suppose $(X,\calo_X)$ and $(Y,\calo_Y)$ are ringed spaces, $X=\cup_i U_i$ is an open cover of $X$, and we have morphisms of ringed spaces $\pi_i:U_i\lrta Y$ that ``agree on overlaps'', i.e. $\pi_i\mid_{U_i\cap U_j}=\pi_{j}|_{U_i\cap U_j}$. Show that there is a unique morphism of ringed space $\pi:X\lrta Y$ such that $\pi|_{U_i}=\pi_i$.
\end{exr}
\begin{proof}[Sol]
Denote the morphisms of ringed spaces as $(\pi_i,\tilde{\pi}_i)$
The topological part has been shown in Exercise~\ref{chap2exr:sheaf_of_continuous_maps}, which says continuous maps from a topological space $X$ to $Y$ forms sheaf.

Suppose the topological part glue to a unique continuous map $\pi$. We denote pushforward sheaves by $\pi_*\calo_X$. Given an open set $V\subset Y$, on the open set level
$$
\pi_{i,*}\calo_{U_i}(V)=\calo_{U_i}(\pi_i^{-1}(V))=\calo_X((\pi|_{U_i})^{-1}(V))=\calo_X(\pi^{-1}(V)\cap U_i)
$$
For each element $f\in \calo_Y(V)$
$\pi_i(V)(f)\in\calo_X(\pi^{-1}(V)\cap U_i)$ are compatible, thus glue to a unique element in $\calo_X(\pi^{-1}(V))=\pi_*\calo_X(V)$. As discussed before, we can value-wisely define the morphism $\tilde{\pi}(V)$ and it is easy to check thus defined $\tilde{\pi}(V)$ is well-defined ring morphism and that $\tilde{\pi}$ forms a sheaf morphism. 
$(\pi,\tilde{\pi})$ forms the desired morphism of ringed spaces.
\end{proof}
\begin{exr}
$\calo$-module push forward. Given a morphism of ringed spaces $\pi:X\lrta Y$, show that sheaf push forward induces a functor $\mathsf{Mod}_{\calo_X}\lrta \mathsf{Mod}_{\calo_Y}$.
\end{exr}
\begin{proof}
Given a $\calo_X$-module $\calf$, and  a morphism of ringed spaces $(\pi,\tilde{\pi}):X\lrta Y$. The $\calf$ as a sheaf of abelian groups automatically pushforward to $\pi_*\calf$. It suffices to check the $\calo_Y$-module structure on $\pi_*\calf$.
$$
\begin{tikzcd}
\calo_Y(U)\times\pi_*\calf(U) \arrow[d, "{\res_{U,V}\times\res_{U,V}}"] \arrow[r,"{\tilde{\pi}(U)\times id}"] & \pi_*\calo_{X}(U)\times\pi_*\calf(U) \arrow[d, "{\res_{U,V}\times\res_{U,V}}"]\arrow[r] & \pi_*\calf(U) \arrow[d,"{\res_{U,V}}"] \\
\calo_Y(V)\times\pi_*\calf(V) \arrow[r,"{\tilde{\pi}(V)\times id}"] & \pi_*\calo_{X}(U)\times\pi_*\calf(U) \arrow[r] & \pi_*\calf(U)
\end{tikzcd}
$$
The right square commutes because $\pi_*\calf$ is a $\pi_*\calo_X$-module. The right square commutes by the definition of morphism of ringed spaces, therefore the outer square commutes, which indicates the $\calo_Y$-module structure.

We still denote the functor $\pi_*$. It remains to check the commutative diagram
$$
\begin{tikzcd}
\calf \arrow[d, "\pi_*"] \arrow[r, "\phi"] & \calg \arrow[d, "\pi_*"] \arrow[r, "\psi"] & \calh \arrow[d, "\pi_*"] \\
\pi_*\calf \arrow[r, "\pi_*(\phi)"] & \pi_*\calg \arrow[r, "\pi_*(\psi)"] & \pi_*\calh,
\end{tikzcd}
$$
where $\pi_*(\phi)(V)$ is defined to be $\phi(\pi^{-1}(V))$ and $\pi_*(\psi\circ \phi)(V)=\psi\circ \phi(\pi^{-1}(V))=\psi(\pi^{-1}(V))\circ \phi(\pi^{-1}(V)) =\pi_*(\psi)\circ \pi_*(\phi) (V)$
This works for all opens,  we thus have $\pi_*(\psi\circ \phi)=\pi_*(\psi)\circ \pi_*(\phi)$.

It still remains to check $\pi_*$ is a well defined morphism of $\calo_Y$ modules, but the procedure is similar. We just omit it.
\end{proof}

\begin{exr}
Given a morphism of ringed spaces $\pi:X\lrta Y$ with  $\pi(p)=q$, show that there is a map of stalks $(\calo_Y)_q\lrta (\calo_X)_p$
\end{exr}
\begin{proof}[Sol]
A morphism of ringed spaces $\pi:X\lrta Y$ consists of the data $(\pi,\tilde{\pi})$ where
$$
\tilde{\pi}:\calo_Y\lrta \pi_*\calo_X
$$
Exercise~\ref{chap2exr:morphism_of_sheaves_induce_on_stalks} shows that it induces a morphism on the stalks $\calo_{Y,q}\lrta(\pi_*\calo_X)_{q}$.

On the other hand Exercise~\ref{chap2exr:Pushforward_induces_morphism_stalks} says the pushforward induces a morphism on the stalks $(\pi_{*}\calo_X)_{q}\lrta\calo_{X,p}$ if $\pi(p)=q$

The composition of the above two morphisms gives the desired map.
\end{proof}
\begin{exr}\label{chap6exr:ring_morphism_induces_morphisms_ringed_spaces}
Suppose $\pi^\#:B\lrta A$ is a morphism of rings. Define a morphism of ringed space $\pi:\spec A\lrta \spec B$ as follows. The map of topological spaces was given by Exercise~\ref{exr:spec_functor}. To describe a morphism of sheaves $\calo_{\spec B}\lrta \pi_*\calo_{\spec A}$ on $\spec B$, it suffices to describe a morphism of sheaves on the distinguished base of $\spec B$. On $D(f)\subset \spec B$, we define
$$
\calo_{\spec B}(D(g))\lrta \calo_{\spec A}(\pi^{-1}D(g))=\calo_{\spec A}(D(\pi^\#g))
$$
by $B_g\lrta A_{\pi^\# g}.$ verify that this makes sense (e.g., is independent of $g$), and that this describes a morphism of sheaves on the distinguished base. (This is the third in a series of exercises. We saw that a morphism of rings induces a map of sets in 3.2.9, a map of topological spaces in Exercise~\ref{exr:spec_functor}, and now a map of ringed spaces here. )
\end{exr}
\begin{proof}[Sol]
We have shown that $A_f\cong \calo_{\spec A}(D(f))$ in Exercise 4.1.A.
Choose another representative of $D(g)$, say $h$. Then we  know $\sqrt{(h)}=\sqrt{(g)}$, which means $h$ is invertible in $B_g$ and $g$ is invertible in $B_h$. i.e. 
$$
B_g\cong B_h
$$
It also means $\pi^\#(h)$ is invertible in $(\pi^\#(g))$ and $\pi^\#(g)$ is invertible in $(\pi^\#(h))$, i.e.
$$
A_{\pi^\#(g)}=A_{\pi^\#(h)}.
$$
Then this map is indeed independent on the choice of representative of $D(g)$.
$$
\calo_{\spec B}(D(g))=B_g\lrta \pi_*\calo_{\spec A}(D(g))=A_{\pi^\#(g)}
$$
gives a well-defined morphism of sheaves.
\end{proof}
\begin{exr} (Not all map of affine ringed spaces are induced by ring morphisms).
Recall (Exercise 3.4.K) that $\spec k[y]_{(y)}$ has two points, $[(0)]$ and $[(y)]$, where the second point is closed, and he first is not. Describe a map of ringed spaces $\spec k(x)\lrta \spec k[y]_{(y)}$ sending the unique point of $\spec k(x)$ to the closed point $[(y)]$, where pullback map on global sections sends $k$ to $k$ by the identity, and sends $y$ to $x$. Show that this map of ringed spaces is not of the form described in  Exercise~\ref{chap6exr:ring_morphism_induces_morphisms_ringed_spaces}.
\end{exr}
\begin{proof}[Sol]
Denote $B:=k[y]_{(y)}$ and $A:=k(x)$.
The exotic choice  topologically sends the unique point $\eta=[(0)]$ in $\spec k(x)$ to the closed point $[(y)]$. These topological space discrete, making it a continuous map.  The only open sets on $\spec A$ are $\{[(0)],[(y)]\},\{[(0)]\}$ and $\emptyset$.
$$
\phi:\spec A=\{\eta\}\lrta \spec B=\{[(0)],[(y)]\}
$$
$$
\phi_*\calo_{\spec A}(\{[(0)],[(y)]\})=\calo_{\spec A}(\eta)=A
$$
$$
\phi_*\calo_{\spec A}(\{[(0)]\})=\calo_{\spec A}(\emptyset)=\{0\}
$$
where $\{0\}$ is the final object in $\mathsf{Rings}$. The corresponding maps on
$$
\tilde{\phi}(\{[(0)],[(y)]\}):B\lrta A (k\mapsto k, y\mapsto x),
$$
$$
\tilde{\phi}(\{[(0)]\}):B\lrta 0 .
$$
This does not corresponds to a ring morphism as in Exercise~\ref{chap6exr:ring_morphism_induces_morphisms_ringed_spaces}. Assume by contradiction it is, then $\{[(0)]\}=D(y)$ and $\phi^\#$ has been given by the morphism on global sections. However $\phi^\#(y)=x$ and $A_{\pi^\#(y)}=A_x=A\neq 0$.
\end{proof}
\section{From locally ringed spaces to morphisms of schemes}
In this section we will use the notion $\pi^\#$ to denote the induced morphism on stalks.
\begin{exr}
Show that morphisms of locally ringed spaces glue (cf. Exercise~\ref{chap6exr:glue_morphisms_of_ringed_spaces}). (Hint: your solution to Exercise~\ref{chap6exr:glue_morphisms_of_ringed_spaces} may work without change.)
\end{exr}
\begin{proof}[Sol]
Without change of the proof, we only need to check the glued morphism of ringed spaces is in fact a morphism of locally ringed spaces.

This is easy because for each pair of points $p,q,\pi(p)=q$, we can check that
$$
\tilde{\pi}_{p}=\tilde{\pi_i}_{p}
$$
simply because $\tilde{\pi}$ restricts to $\tilde{\pi}_i$. Each $\pi_i$ induces local morphisms on stalks, therefore $\pi$ also induces local morphisms on stalks.
\end{proof}
\begin{exr}
\ \begin{enumerate}[label=(\alph*)]
\item Show that $\spec A$ is a locally ringed space.
\item Show that the morphism of ringed spaces $\pi:\spec A\lrta \spec B$ defined by a ring morphism $\pi^\#:B\lrta A$ is a morphism of locally ringed spaces.
\end{enumerate}
\end{exr}
\begin{proof}[Sol]
\begin{enumerate}[label=(\alph*)]
\item $\spec A$ is a locally ringed space because we have shown in Exercise~\ref{chap4exr:affine_scheme_stalk_localization} that
$$
\calo_{\spec A,\scp}=A_\scp
$$
and $A_\scp$ is a local ring with maximal ideal $\scp A_\scp$.
\item In Exercise~\ref{chap6exr:ring_morphism_induces_morphisms_ringed_spaces}, we have given explicitly on the distinguished base. We only need to show that the maps on stalks indeed send maximal ideals to maximal ideals.

Given two prime ideals $\scq\subset B$ and $\scp\subset A$, $\pi([\scp])=[\scq]$ means $(\pi^{\#})^{-1}(\scp)=\scq$, recall that pull back of ideals preserve prime ideals. This means the $\pi^\#(B-\scq)\subset A-\scp$ and hence composition map $B\lrta A_\scp$ sends $B-\scq$ into invertible element in $A_\scp$.  
$$
\begin{tikzcd}
B \arrow[d] \arrow[r] \arrow[rd] & A \arrow[d] \\
B_\scq \arrow[r, "\phi", dashed] & A_\scp
\end{tikzcd}
$$
There is a ring morphism from $B_\scq$ to $A_\scp$ by universal property of localization. 
$$
\phi:
\frac{b}{s}\mapsto \frac{\pi^\#(b)}{\pi^\#(s)}.
$$
We use $\phi$ to avoid abuse of notations.

We want to show that
$$
\phi(\scq B_\scq)\subset \scp A_\scp.
$$
This is true because 
$$
\phi\left(\sum_i q_i \frac{b_i}{s_i}\right)=\sum_i \frac{\pi^\#(q_i b_i)}{\pi^\#(s_i)}\in \scp A_\scp.
$$ 
\end{enumerate}
\end{proof}
\begin{proposition}
If $\pi:\spec A\lrta \spec B$ is a morphism of locally ringed spaces then it is the morphism induced by the map $\pi^\#:B=\Gamma(\spec B,\calo_{\spec B})\lrta \Gamma(\spec A,\calo_{\spec A})=A.$
\end{proposition}
\begin{exr}\label{chap6exr:equivalent_defs_morphisms_of_schemes}
Show that  a morphism of schemes $\pi:X\lrta Y$ is a morphism of ringed spaces that looks locally like morphisms of affine schemes. Precisely, if $\spec A$ is an affine open subset of $X$ and $\spec B$ is an affine open subset of $Y$, and $\pi(\spec A)\subset \spec B$, then the induced morphism of ringed spaces is a morphism of affine schemes. (In case it helps, note: if $W\subset X$ and $Y\subset Z$ are both open embeddings of ringed spaces, then any morphism of ringed spaces $X\lrta Y$ induces a morphism of ringed spaces $W\lrta Z$, by composition $W\lrta X\lrta Y\lrta Z$.) Show that it suffices to check on a set $(\spec A_i,\spec B_i)$ where the $\spec A_i$ form an open cover of $X$ and the $\spec B_i$ from an open cover of $Y$.
\end{exr}
\begin{proof}[Sol]
By affine open subset, we mean an open subscheme which is affine. There is an open embedding from $\spec A$ to $X$. In particular, the $\calo_X|_{\spec A}=\calo_{\spec A}$. By definition,  ``morphism of schemes'' is a morphism of locally ringed spaces. Restricting to the affine opens, the morphism of schemes induces a morphism of locally ringed spaces from $\spec A$ to $\spec B$. By proposition above, we know it is induced by a ring morphism $B\lrta A$.


On the other hand, given the set $(\spec A_i,\spec B_i)$, s.t. $\pi(\spec A_i)\subset\spec B_i$ and $A_i, B_i$ covers $X, Y$ respectively. 
The morphism of affine schemes $\spec A_i\lrta \spec B_i$ induces morphism of locally ringed spaces $\spec A_i\lrta Y$ by composition. 
We have checked that morphism of locally ringed spaces glue, hence there is a morphism of locally ringed spaces from $X\lrta Y$. This criterion is valid.
\end{proof}
\begin{exr}
Show that the category of rings and the opposite category of affine schemes are equivalent.
\end{exr}
\begin{proof}[Sol]
We give a pair of functors which are reciprocal quasi-inverse.
$$
\spec ,\Gamma,
$$
where $\Gamma$ is the global section functor, both are contravariant.
$$
\begin{tikzcd}
B \arrow[r, "\spec"] \arrow[d, "\pi'"] & \spec B \arrow[r, "\Gamma"] \arrow[d, "{(\pi,\tilde{\pi})}"] & B \arrow[d, "\pi^\#"] \\
A \arrow[r, "\spec"] & \spec A \arrow[r, "\Gamma"] & A
\end{tikzcd}
$$
The above diagram is easy to verify and wesee from it
$$
\Gamma\circ \spec =Id_{\mathsf{Rings}}
$$

$$
\begin{tikzcd}
(U,\calo_U)\cong \spec B \arrow[r, "\Gamma"] \arrow[d] & B \arrow[r, "\spec "] \arrow[d] & \spec B \arrow[d] \\
(V,\calo_V)\cong\spec A \arrow[r, "\Gamma"] & A \arrow[r, "\spec"] & \spec A
\end{tikzcd}
$$
This means
$$
\spec\circ \Gamma\cong Id_{\mathsf{Affn}}.
$$
Because both of them are contravariant functors, this indicates a equivalence of categories between $\mathsf{Rings}$ and $\mathsf{Affn}^{op}$.

\end{proof}
\begin{exr}
(This exercise can give you some practice with understanding morphisms of schemes by cutting up into affine open sets.) Make sense of the following sentence: ``$ \affn^{n+1}_k\backslash \{0\}\lrta \proj^n_k$ given by 
$$
(x_0,x_1,...,x_n)\mapsto[x_0,x_1,...,x_n]
$$
is a morphism of schemes.''
(Can you generalize to the case where $k$ is replaced by a general ring $B$?)
\end{exr}
\begin{proof}[Sol]
$\affn^{n+1}_k\backslash \{0\}$ can be covered by affine opens like
$$
\spec A_i\simeq D(x_i),\ \ \  A_i:= k[x_0,...,x_n]_{x_i}.
$$
On the other hand $\proj^{n}_{k}$ can be covered by affine opens
$$
\spec B_i\simeq \affn^{n}_{k}, B_i:=k[x_{0/i},...,x_{k/i},...,x_{n/i}]/(x_{i/i}-1)\cong k[y_0,...,y_{n-1}].
$$
We denote the morphism 
$$
f:(x_0,...,x_n)\mapsto [x_0,...,x_n]
$$
and observe that $f_i:=f|_{D(x_i)}$ maps $D(x_i)$ into $\spec B_i$. Locally, it is a morphism of affine scheme induced by 
$$
B_i\lrta A_i,
$$
$$
x_{k/i}\mapsto \frac{x_{k}}{x_i}, x_{i/i}\mapsto 1
$$
which is a well-defined ring morphism.

Then it suffice to check that behaves well when we glue up the affine charts. It is equivalent to check that
$$
f_i|_{D(x_i)\cap D(x_j)}=f_j|_{D(x_i)\cap D(x_j)}.
$$
Note that $D(x_i)\cap D(x_j)=D(x_i x_j)=\spec A_{ij}=\spec k[x..]_{x_i x_j}$ and both $f_i|_{D(x_i)\cap D(x_j)}$ and $f_j|_{D(x_i)\cap D(x_j)}$ map into $$\spec B_{ij}:=\spec k[x_{0/i},...,x_{n/i}, 1/x_{j/i}]/(x_{i/i}-1)$$ and $$\spec B_{ji}=\spec k[x_{0/j},...,x_{n/j},1/x_{i/j}]/(x_{j/j}-1)$$.
The above two are isomorphic by the map $x_{k/j}=x_{k/i}/x_{j/i}$.
$
f_i|_{D(x_i)\cap D(x_j)}
$
is induced by
$$
\spec B_{ij}\mapsto \spec A_{ij}
$$
$$
x_{k/i}\mapsto \frac{x_k x_j}{x_i x_j},1/x_{j/i}\mapsto \frac{x_i^2}{x_j x_i}
$$
Together with the isomorphism above, $f_i|_{D(x_i)\cap D(x_j)}$ can be identified with $f_j|_{D(x_i)\cap D(x_j)}$.
\end{proof}
\begin{exr}\label{chap6exr:morphism_of_schemes_bijection_morphisms_of_global_sections}
Show that morphism $X\lrta \spec A$ are in natural bijection with ring morphisms $A\lrta \Gamma(X,\calo_X)$.
\end{exr}
\begin{proof}[Sol.]
The bijection
$$
\hom_{\mathsf{Sch}}(X,\spec A)\simeq \hom_{\mathsf{Rings}}(A,\Gamma(X,\calo_X))
$$
is explicitly
$$
[f:X\lrta \spec A]\longmapsto [A=\calo_{\spec A}(\spec A)\lrta \calo_X(f^{-1}\spec A)=\calo_X(X)],
$$
where the later is induced by 
$$
f^*:\calo_X\lrta f_*\calo_X.
$$

Consider $X$ covered covered by a collection of affine opens $X=\cup_i U_i$ with each $U_i\simeq\spec A_i$
$$
\begin{aligned}
\hom_{\mathsf{Sch}}(X,\spec A)&=\{(f_i)\mid f_i:U_i\lrta \spec A\text{ s.t. }f_i|_{U_i\cap U_j}=f_j|_{U_i\cap U_j}\}\\
&=\{(g_i)\mid g_i: A\lrta \Gamma(U_i, \calo_X),\forall a\in A, g_i(a)|_{U_i\cap U_j}=g_j(a)|_{U_i\cap U_j}\}\\
&\simeq\{g|g:A\lrta \calo_X(X)\}\text{ by sheaf condition}\\
&=\hom_{\mathsf{Rings}}(A,\calo_X(X))
\end{aligned}
$$
\end{proof}
\begin{exr}
Show that this definition of $A$-scheme given in \S $6.3.7$ agrees with the earlier definition of \S $5.3.6$.
\end{exr}
\begin{proof}[Sol.]
Def 1: In \S 5.3.6, the $A$-scheme is defined to be scheme whose ring of sections over any  open sets is an $A$-algebra and the restriction maps are morphisms of $A$-algebras.

Def 2: In \S 6.3.7, the $A$-scheme is defined to be simply a scheme $X$ with  a morphism $X\lrta \spec A$. 

$Def 2\Lrta Def 1$: Consider an open subsets $U\subset V\subset X$, $(V,\calo_X|_V)\overset{\iota}{\inj}(X,\calo_X)$. Compositing with the morphism of schemes $X\lrta \spec A$, we know there is a morphism of schemes from $V$ to $\spec A$. By Exercise~\ref{chap6exr:morphism_of_schemes_bijection_morphisms_of_global_sections}, $\calo_X|_V(V)=\calo_X(V)$ is an $A$-algebra. 

Morphism $\calo_X(V)\lrta\calo_X(U)$ is a morphism of $A$-algebras because
$$
\begin{tikzcd}
U \arrow[r, hook] \arrow[rd] & V \arrow[d] & \calo_X(U) & \calo_X(V) \arrow[l, "\res"] \\
 & \spec A &  & A \arrow[u] \arrow[lu]
\end{tikzcd}
$$
the later diagram gives the correct meaning of morphism of $A$-algebras.

$Def1\lrta Def2$: BY definition 1, in particular $\Gamma(X,\calo_X)$ is an $A$-algebra, again by Exercise~\ref{chap6exr:morphism_of_schemes_bijection_morphisms_of_global_sections}, there is a morphism of schemes between $X$ and $\spec A$.
\end{proof}


\begin{exr}
If $S_\bullet$ is a finitely generated graded $A$-algebra, describe a natural ``structure morphism'' $\Proj S_\bullet\lrta \spec A$.
\end{exr}
\begin{proof}[Sol.]
By Exercise~\ref{chap6exr:morphism_of_schemes_bijection_morphisms_of_global_sections}, the structure morphism corresponds to the morphism of rings
$$
A\lrta S_\bullet : a\mapsto a\times 1.
$$

Recall that $\Proj S_\bullet$ is covered by  affine charts $D(f)=\spec ((S_\bullet)_f)_0$ for each homogeneous $f\in S_+$, where the $((S_\bullet)_f)_0$ is the zero piece of the localization $(S_\bullet)_f$. A morphism $A\lrta S_\bullet$ naturally induces a morphism $A\lrta ((S_\bullet)_f)_0$ and this gives a collection of compatible morphism from $\phi_f:A\lrta \Gamma(D_f, \calo_{\Proj S_\bullet})$. These $\phi_f$ in turn defines a compatible collection of morphism of schemes $\spec ((S_\bullet)_f)_0\lrta \spec A$. And they glue up to a morphism of schemes $\Proj S_\bullet\lrta \spec A$.

Sections on general open sets in $\Proj S_\bullet$ could be quite nasty, please refer to~\ref{chap4exr:structure_sheaf_of_projective_schemes} for a detailed description of structure sheaf on $\Proj S_\bullet$.
\end{proof}
\begin{exr}
Show that $\spec \intg$ is the final object in the category of schemes. In other words, if $X$ is any scheme, there exists a unique morphism to $\spec \intg$. (Hence the category of schemes is isomorphic to the category of $\intg$-schemes.) If $k$ is a field, show that $\spec k$ is the final object in the category of $k$-schemes.
\end{exr}
\begin{proof}[Sol.]
Simple words: $\intg$ is the initial object in the category $\mathsf{Rings}$ and $k$ it the initial object in the category in the category $k-\mathsf{Alg}$ and then argue by~\ref{chap6exr:morphism_of_schemes_bijection_morphisms_of_global_sections}.
\end{proof}

\begin{exr}\label{chap6exr:morphism_of_schemes_stalk_to_X}
\ \begin{enumerate}[label=(\alph*)]
\item Suppose $p$ is a point of a scheme $X$. Describe a canonical (choice-free) morphism $\spec \calo_{X,p}\lrta X$. (Hint: do this for affine $X$ first. But then for general $X$ be sure to show that your morphism is independent of choice.)
\item Define a canonical morphism $\spec \kappa(p)\lrta X$. This is often written $p\lrta X$; one gives $p$ the obvious interpretation as a scheme, $\spec \kappa(p)$
\end{enumerate}
\end{exr}
\begin{proof}[Sol.]
\begin{enumerate}[label=(\alph*)]
\item For affine $X=\spec A$, consider the point $\scp$. There is a morphism of rings 
$$
\calo_X(X)=A\lrta A_\scp=\calo_{X,\scp}
$$
which gives a morphism
$$
\spec \calo_{X,\scp}\lrta X.
$$

For general case, a point $p\in X$, we can find an affine open neighborhood $U_i$ of $p$ so that $U=\spec B$. There is a morphism
$$
\spec \calo_{X,p}\lrta \spec B\lrta X.
$$
We must show that the composition map $\spec \calo_{X,p}\lrta X$ is independent of the choice of affine open $U$. Consider the $V$ is another affine open that contains $p$, we can check that restriction of the injections $U\inj X$ and $V\inj X$ to $U\cap V\inj X$ would agree. For quasiseparated scheme, the intersection $U\cap V$ is finite union of affine opens again, and the proof is finished here by picking the affine open that contains $p$ and in $U\cap V$. 

For non-quasiseparated schemes, we can still check that $U\cap V$ is a union of affine open charts (not necessarily finite). The process is the same.
 
\item $\kappa(p)$ is a field and $\spec \kappa(p)$ is a one point set. Each residue field $\kappa(p):=\calo_{X,p}/\scm_p$, there is a canonical morphism of rings $\calo_{X,p}\lrta \kappa(p)$ and therefore induces a canonical morphism of affine schemes
$$
\spec \kappa(p)\lrta \spec \calo_{X,p}.
$$
Together with the morphism described in $(a)$, we have the canonical morphism
$$
\spec \kappa(p)\lrta X.
$$
\end{enumerate}
\end{proof}
\begin{exr}
Suppose $X$ is a scheme, and $(A,\scm)$ is a local ring. Suppose we have a scheme morphism $\pi:\spec A\lrta X$ sending $[\scm]$ to $p$. Show that any open set containing $p$ contains the image of $\pi$. Show that there is a bijection between $\mor(\spec A,X)$ and 
$$
\{p\in X,\text{ local homomorphisms }\calo_{X,p}\lrta A\}.
$$
\end{exr}
\begin{proof}[Sol.]
In order to show that the image lies in an affine open subscheme of $X$, we need to prove the following two topological facts:
\begin{enumerate}
\item Any continuous map preserves the relations of specialization/generalization
\item If $p\in X$ is a point. Then any open neighborhood $U\ni x$ contains all the generalization of $p$ in $X$.
\end{enumerate}
The first of which can be checked in \cite[Tag 0060]{stacks-project} (Also check the general definition of specialization and generalization in topological spaces). The first point implies that $p$ is a closed point in $X$.

The second point is easy to prove. Though $[\scm]$ is close point, it is not necessarily true that $\pi([\scm])$ is closed. In general we only know $\overline{\pi(\scm)}\supset \pi(\overline{[\scm]})$. Given point $p$ and an open neighborhood $U$ containing it. We recall the definition of $q\in X$ is a  generalization of $p$: $p\in \overline{\{q\}}$. $p$ is in the closure of $\{q\}$, thus any open neighborhood of $p$ would contain $q$. 

Because $[\scm]$ is the only closed point in $\spec A$, all point of $\spec A$ would be a generalization of $[\scm]$, which means $\pi(\spec A)$ is contained in any open neighborhood of $p=\pi([\scm])$.\footnote{We can not simply argue by ``$U$ would contain the image of generic point $\pi(\eta)$ of $\spec A$''. The subtle stuff here is that $U$ might do not contain $\overline{\{\pi(p)\}}$. In general for continuous map $\pi$, $\pi(\overline{S})\subset \overline{\pi(S)}$.}

The above proof finishes the first question.

This implies that the morphism $\pi:\spec A\lrta X$, with $(A,\scm)$ being local and $\pi([\scm])=p$, would factor through $\spec \calo_{X,p}\lrta X$.

Take an affine open neighborhood $X\supset \spec B$ of $p$. Thus $p$ can be identified as a prime ideal in $B$. We denote it by $p=[\scp]\in\spec B$. $\spec A\lrta X$ would factor as $\spec A\lrta\spec B$, which corresponds to a ring morphism $\phi:B\lrta A$. The preimage of $\scm$ is some the prime ideal $\scp\in B$. Recall that $A-\scm$ is the set of units in local ring.  $\phi$ sends the invertible elements into $A-\scp$, by the universal property of localization $B\lrta A$ factors through the localization $B_\scp$, where $B_\scp=\calo_{X,p}$ and the morphism $\calo_{X,p}\lrta A$ is local morphism. It does not depends on the choice of $\spec B$, because they are induced by the same morphism of ringed spaces cf~\ref{chap6exr:equivalent_defs_morphisms_of_schemes}. $\pi:\spec A\lrta X$, uniquely determines a data of $\{p\in X,\text{a local morphism } \calo_{X,p}\lrta A\}$. This given is injectivity in the second question. The surjectivity is derived from Exercise~\ref{chap6exr:morphism_of_schemes_stalk_to_X} which says the morphism $\spec \calo_{X,p}\lrta X$ is canonical.
\end{proof}

\begin{exr}
\ \begin{enumerate}[label=(\alph*)]
\item Show that a morphism of schemes $X\lrta Y$ induces a map of $Z$-valued points $X(Z)\lrta Y(Z)$.
\item Note that morphisms of schemes $X\lrta Y$ are not determined by their underlying map of points. (Give an example) Show that are determined by their induced maps of $Z$-valued points, as $Z$ varies over all schemes. 
\end{enumerate}
\end{exr}
\begin{proof}
\begin{enumerate}[label=(\alph*)]
\item 
A morphism of schemes $f:X\lrta Y$ would induce a map
$$
\begin{aligned}
X(Z):=\mor(Z,X)
\lrta \mor(Z,Y)=:Y(Z)\\
[Z\rta X]\longmapsto [Z\rta X\overset{f}{\rta} Y]
\end{aligned}
$$
\item For example choose $Z$ to the final object $\spec \intg$ in $\mathsf{Sch}$. Any morphism between $X,Y$ would induces a map between one point sets, which gives a non-example. 

Then we will reprove the Yoneda's Lemma in the context of schemes. Pick $Z=X$, $f:X\lrta Y$ would induce a map between $\mor(X,X)\lrta \mor(X,Y)$, where $id_X\in \mor(X,X)$ is mapped to $f\in \mor(X,Y)$. $f$ is determined by the induces map on Hom sets.
\end{enumerate}
\end{proof}

\begin{exr}
\ \begin{enumerate}[label=(\alph*)]
\item Suppose $B$ is a ring. If $X$ is a $B$-scheme, and $f_0,...,f_n$ are $n+1$ functions on $X$ with no common zeros, then show that $[f_0,..,f_n]$ gives a morphism of $B$-scheme $X\lrta \proj^n_B$.
\item Suppose $g$ is nowhere vanishing function on $X$, and $f_i$ are as in part $(a)$. Show that the morphisms $[f_0,...,f_n]$ and $[gf_0,...,g f_n]$ to $\proj^n_B$ are the same.
\end{enumerate}
\end{exr}
\begin{proof}[Sol.]
\ \begin{enumerate}[label=(\alph*)]
\item Recall that $\proj_B^n=\Proj B[x_0,...,x_n]$. We consider an open over of $X=\cup_{f_i} X_{f_i}$ where $X_{f_i}$ is open set of $X$ where $f_i\neq 0$ and $\proj^n_B$ is covered by the canonical affine charts $U_i:=spec B[x_{0/i},...,x_{n/i}]$. We hope to define a morphism $\pi:X\lrta\proj^n_B$ so that it induces a morphism of schemes
$$
X_{f_i}\lrta \spec B[x_{0/i},...,x_{n/i}]
$$
and then verify that they agree on the overlaps.
 We do not need $X_{f_i}$ to be affine. By Exercise~\ref{chap6exr:morphism_of_schemes_bijection_morphisms_of_global_sections}, it suffices to construct the morphism of rings
$$
B[x_{0/i},...,x_{n/i}]\lrta \Gamma(X_{f_i},\calo_{X}).
$$

One the obvious choice of $B$-algebra morphism is
$$
\pi_i:x_{k/i}\longmapsto \frac{f_k}{f_i}.
$$
{\color{red} A crucial problem here is how to define the the sheaf on $X_{f_i}$, which is not necessarily affine.}

They are compatible on intersections $X_{f_i}\cap X_{f_j}$ where both $f_i,f_j$ are non-zero. Recall on the intersection $U_i\cap U_j$, the gluing isomorphism is 
$$
\varphi_{ij}: U_i\cap U_j\lrta U_i\cap U_j
$$
$$
x_{k/i}\longmapsto x_{k/j}/x_{i/j}
$$
Compositing with the constructed map $x_{k/j}\lrta \frac{f_k}{f_j}$.
We
have
$$
x_{k/i}\longmapsto x_{k/j}/x_{i/j}\longmapsto \frac{\frac{f_k}{f_j}}{\frac{f_i}{f_j}}=\frac{f_k}{f_i}.
$$
This means they are compatible on the intersections, we argue by~\ref{chap6exr:glue_morphisms_of_ringed_spaces} that we can glue it up to a morphism of schemes $\pi:X\lrta\proj^n_B$.

\item Given the explicit form of $(a)$, $(b)$ is obvious.
\end{enumerate}
\end{proof}
\begin{exr}
Omitted.
\end{exr}

\section{Maps of graded rings and maps of projective schemes}
\begin{exr}
Suppose that $\phi:S_\bullet\lrta R_\bullet$ is a morhpism of $\intg^{\geq 0}$-graded rings. (By \textbf{ map of graded rings}, we mean a map of rings that perserves the grading as a map of ``graded semigroups''). In other words, there is a $d> 0$ such that $S_n$ maps to $R_{dn}$ for all $n$. Show that this induces  a morphism of schemes
$$
(\Proj R_\bullet)\backslash V(\phi(S_+))\lrta \Proj S_\bullet.
$$
Hint:
Define a map $D(\phi(f))\lrta D(f)$. Show that they glue together as $f$ runs over all homogeneous elements of $S_+$. Show that this deines a map from all of $\Proj R_\bullet\backslash V(\phi(S_+))$. In particular, if $V(\phi(S_+))=\emptyset$, then we have a morphism $\Proj R_\bullet\lrta \Proj S_\bullet$. 
\end{exr}
\begin{proof}[Sol.]
\end{proof}

\section{rational maps from reduced schemes}

\chapter{Useful classes of morphisms}
\section{An example of a reasonable class of morphisms: open embeddings}
\section{Algebraic interlude: Lying over and Nakayama}
\section{A gazillion finiteness conditions on morphisms}
\section{Images of morphisms: Chevalley's theorem and elimination theory}
\chapter{Closed embeddings and related notions}
\section{Closed embeddings and closed subschemes}
\section{More projective geometry}
\section{The (closed sub)scheme-theoretic image}
\section{Effective Cartier divisors, regular sequences and regular embeddings}
\chapter{Fibered product of schemes}
\section{They exist}
\section{Computing fibered product in practice}
\section{Interpretations: Pulling back families, and fibers of morphisms}
\section{Properties preserved by base change}
\section{Properties not preserved by base change and how to fix them}
\section{Products of projective schemes: The Segre embedding}
\section{Normalization}
\chapter{Separated and proper schemes, varieties}
\section{Separated morphisms}
\section{Rational maps to separated schemes}
\section{Proper morphism}
\part{``Geometric'' properties: Dimension and smoothness}
\chapter{Dimension}
\section{Dimension and codimension}
\section{Dimension, transcendence degree, and Noether normalization}
\section{Codimension one miracle, Krull's and Hartogs's Theorems}
\begin{exr}
Show that an irreducible homogeneous polynomial in $n+1$ variables over a field $k$ describes an integral scheme of dimension $n-1$ in $\proj^n_k$ 
\end{exr}
\begin{proof}[Sol.]
$S_\bullet=k[X_0,...,X_n]$ is an Noetherian integral domain. By Krull's principal ideal theorem. Suppose $f$ is an irreducible homogeneous polynomial in $S_\bullet$.
\end{proof}
\section{Dimensions of fibers of morphisms of Varieties}
\section{Proof of Krulls's Principal Ideal and Height Theorems}
\chapter{Regularity and smoothness}

\part{Quasicoherent sheaves}
\chapter{Quasicoherent and coherent sheaves}
\section{Vector bundles and locally free sheaves}
\begin{exr}
Show that over $U_i\cap U_j$, the vector-valued function $\vec{s}^i$ is related to ${s}^j$ by the (same) transition functions: $T_{ij} \vec{s}^{i}=\vec{s}^j$. (Make sure you get the correct superscripts)
\end{exr}
\begin{proof}[Sol.]
There is in fact some abuse of notions, where $\vec{s}^i$ should be interpreted as the restriction of $\vec{s}^i$ to $U_i\cap U_j$.
Here the superscripts are not for the indices of matrices but only denote the  charts.
$$
\tiny
\begin{tikzcd}
 & \calo(U_i)^d \arrow[rd] &  \\
\calf(U_i) \arrow[ru, "\phi_i(U_i)"] \arrow[rd, "\res"] &  & \calo(U_i\cap U_j)^d \arrow[dd, "T_{ij}"] \\
 & \calf(U_i\cap U_j) \arrow[rd, "\phi_j(U_i\cap U_j)"] \arrow[ru, "\phi_i(U_i\cap U_j)"] &  \\
\calf(U_j) \arrow[ru, "\res"] \arrow[rd, "\phi_j(U_j)"] &  & \calo(U_i\cap U_j)^d \\
 & \calo(U_j)^d \arrow[ru] & 
\end{tikzcd}
$$

Keep in mind that $\phi_i$ together with the open cover $U_i$ give the definition of a locally free sheaf. $T_{ij}$ is thus defined to be $\phi_j(U_i\cap U_j)\circ\phi_i(U_i\cap U_j)^{-1}$. 
\end{proof}
\begin{exr}
Suppose $\calf$ and $\calg$ are locally free sheaves on $X$ of rank $m$ and $n$ respectively. Show that $\phom_{\calo_X}(\calf , \calg)$ is a locally free sheaf of rank $mn$.
\end{exr}
\begin{proof}[Sol.]
Recall the definition of sheaf-hom
$$
\phom(\calf,\calg)(U):=\hom(\calf|_U,\calg|_U)
$$
(in fact, this notion has been implicitly extended to $\calo_X$-modules-hom)

On the other hand, 
$$
\calf|_U\simeq \calo_U^{\oplus m}
$$
and
$$
\calf|_U\simeq \calo_U^{\oplus n}
$$
We have
$$
\phom_{\calo_X}(\calf|_U,\calg|_U)\cong\phom_{\calo_X}(\calo_U^{\oplus m},\calo_U^{\oplus n}).
$$
The later is isomorphic to 
$$
M_{m\times n}(\calo_U)\cong\calo_U^{\oplus mn}
$$
(Note that the above isomorphism is in the perspective of $\calo_X$-module and we ignore the multiplication structure of matrix)
because for all open set $V\subset U$
$$
\phom_{\calo_X|_U}(\calo_U^{\oplus m},\calo_U^{\oplus n})(V)=\mor_{\calo_X|_V}(\calo_U^{\oplus m}|_V,\calo_U^{\oplus n}|_V)
$$
as $\calo_X$-modules, the RHS is isomorphic to
$$
M_{m\times n} (\calo_U|_V)\cong \calo_V^{\oplus mn}
$$
and this can be checked routinely via further checking the sections and compatibility with restrictions.

\end{proof}

\begin{exr}\label{chap13exr:dual_and_double_dual_of_locally_free sheaves}
If $\cale$ is a locally free sheaf on $X$ of (finite) rank $n$, the above exercise implies that $\cale^\vee:=\phom(\cale,\calo_X)$ is also a locally free sheaf of rank $n$. This is called the dual of $\cale$. Given transition functions for $\cale$, describe transition functions of $\cale^\vee$. (Note that if $\cale$ is of rank $1$, i.e., invertible, the transition functions of the dual are the inverse of the transition functions of the original.) Show that $\cale\cong\cale^{\vee\vee}$. (Caution: your argument showing that there is  a canonical isomorphism $(\calf^\vee)^\vee$ better not also show that there is an isomorphism $\calf^\vee\cong \calf$. We will see an example in \S 14.1 of a locally free $\calf$ that is not isomorphic to its dual: the invertible sheaf $\calo(1)$ on $\proj^n$.)
\end{exr}
\begin{proof}[Sol.]
The first question is to describe the transition function of the dual sheaf. Consider the following diagram
$$
\tiny
\begin{tikzcd}
 &  & \calo(U_i)^n \arrow[rrd, red] &  &  &  \\
\cale(U_i) \arrow[rru,red] \arrow[rrd,red] &  &  & (\calo^{n}|_{U_i})^{\vee} \arrow[rrd,blue] & \calo(U_i\cap U_j)^n \arrow[dd,red] &  \\
 & \cale^\vee(U_i) \arrow[rru,blue] \arrow[rrd,blue] & \cale(U_i\cap U_j) \arrow[rru,red] \arrow[rrd,red] &  &  & (\calo^n|_{U_i\cap U_j})^\vee \arrow[dd,blue] \\
\cale(U_j) \arrow[rru,red] \arrow[rrd,red] &  &  & \cale^\vee({U_j\cap U_j}) \arrow[rrd,blue] \arrow[rru,blue] & \calo(U_i\cap U_j)^n &  \\
 & \cale^\vee(U_j) \arrow[rru,blue] \arrow[rrd,blue] & \calo(U_j)^n \arrow[rru,red] &  &  & (\calo^n|_{U_i\cap U_j})^\vee \\
 &  &  & (\calo^n|_{U_j})^\vee \arrow[rru,blue] &  & 
\end{tikzcd}
$$
The locally free sheaf $\cale$ is defined by via an ``atlas'' $\{(U_i,\phi_i)_{i\in I}\}$, where $\phi_i:\cale|_{U_i}\cong \calo_X^{n}|_{U_i}$. 

If $U_i,\phi_i$ are already given, $\cale^\vee$ is built from the same open cover, but 
$$
\begin{aligned}
\phi_i^{\vee}:&\cale^\vee|_{U_i}\overset{\cong}{\lrta}\calo_{U_i}^n\\
\phi_i^{\vee}(U_i):&\hom(\cale|_{U_i},\calo_{U_i})\overset{\cong}{\lrta}\calo_{U_i}^n
\end{aligned}
$$
where
$$
\begin{aligned}
\hom(\cale|_{U_i},\calo_{U_i})&\cong \hom(\calo^n_{U_i},\calo_{U_i})\\
f\mapsto  f\circ \phi_i^{-1}
\end{aligned}
$$
therefore the canonical choice of $\phi_i^\vee(U_i)$ should be
$$
f\mapsto f\circ \phi_i^{-1}\mapsto (\iota_1\circ f\circ \phi_i^{-1}(e_1),...,\iota_n\circ f\circ \phi_i^{-1}(e_n))^{T}
$$
where we superessed all the possible choice of open subsets $V\subset U_i$ and notions like $f(V), \phi_i(V)$ and so on.
Similarly the inverse $\phi_i^\vee(U_i)^{-1}$ could be expressed as
$$
\calo_{U_i}^n\ni\alpha\mapsto \alpha\cdot\_\mapsto [\alpha\cdot\_]\circ \phi_i
$$


In this way
$$
\begin{aligned}
T^{\vee}_{ij}(\beta)&:= \phi_j^\vee(U_i\cap U_j)\circ \phi_i^\vee(U_i\cap U_j)^{-1}(\beta)\\
&=[g\_\mapsto (\iota_1\circ g\_\circ \phi_j^{-1}(U_i\cap U_j)(e_1),...,\iota_n\circ g\_\circ \phi_j^{-1}(U_i\cap U_j)(e_n))^{T}]\circ [\alpha\_\mapsto [\alpha\_\cdot\_]\circ \phi_i(U_i\cap U_j)](\beta)\\
&=[\alpha\_\mapsto (\alpha\_^{T}\circ \phi_i(U_i\cap U_j)\circ\phi_j^{-1}(U_i\cap U_j))^{T}](\beta) \\
&=(\phi_i(U_i\cap U_j)\circ\phi_j^{-1}(U_i\cap U_j))^T\beta\\
&=T_{ji}^{T}\beta
\end{aligned}
$$
Hence, $T_{ij}^\vee=T_{ji}^{T}=(T_{ij}^{-1})^{T}=T_{ij}^*$. 


The second question requires us to prove the canonical isomorphism $\cale\cong\cale^{\vee\vee}$.

This can be directly checked by the transition function argument $T_{ij}^{**}=T_{ij}$ and the fact that transition function determines the locally free sheaf.

Or we can directly construct the sheaf isomorphism $\varphi$, each $\varphi(U_i)$ is defined as follows, for all $v\in \cale(U_i)$
$$
v\longmapsto v^{\vee\vee}:=[\forall u^\vee\in \cale(U_i)^\vee, u^\vee\longmapsto u^\vee(v)].
$$
When they are free module of rank $n$, the map $v\mapsto v^{\vee\vee}$ is an isomorphism. This map is injective: Let $e_i^\vee$ be the dual basis of $M^\vee$. If $v\in \cale(U_i)$ and $v\neq 0$, the $v$ has a non-zero coordinate with respect to $e_i^{\vee}(v)\neq 0$ for some $i$, which means $v^\{\vee\vee\}(e_i^\vee)\neq 0$.

It is surjective. pick an $f\in \cale^{\vee\vee}(U_i)$ we want to find a $v\in \cale(U_i)$ s.t. $f(u^*)=u^*(v)$ for all $u^*\in \cale^\vee(U_i)$. We can set $u^*=e_i^\vee$ for all $i$ to get a collection of $\calo_X(U_i)$-linear equations. Let $a_i=f(e_i^\vee)\in \calo_X(U_i)$ and define $v=\sum^n_{i=1}a_i e_i\in M$. Then $f(e_i^\vee)=a_i=e_i^\vee(v)$ therefore $f=v^{\vee\vee}$.

And the compatibility with restriction and cocycle relation is just routine.
\end{proof}

\begin{exr}\label{chap13exr:tensor_product_of_invertible_sheaves}
If $\calf$ and $\calg$ are locally free sheaves, show that $\calf\otimes \calg$ is a locally free sheaf.  If $\calf$ is an invertible sheaf, show that $\calf\otimes \calf^\vee\cong \calo_X$.
\end{exr}
\begin{proof}[Sol.]
Note the the tensor of sheaves over abelian category requires sheafification. We can try proving this at the level of stalks. Recall Exercise 2.6.J that taking stalks commutes with tensoring, we have
$$
\calf_{x}\otimes\call_{x}\cong (\calf\otimes \calg)_{x},
$$
$$
\calf_{U,x}\cong\calo_{U,x}^n,
$$
and
$$
\calg_{U,x}\cong \calo_{U,x}^m,
$$
which altogether means 
$$
(\calf\otimes \calg)_{U,x}\cong \calo_{U,x}^n\otimes \calo_{U,x}^m\cong \calo_{U,x}^{mn}.
$$
This in turn means
$$
(\calf\otimes \calg)|_U\cong \calo_{U}^{mn} .
$$
If $\calf$ is a locally free sheaf of rank $1$, still we only need to check first that

\textbf{If $\call_1,\call_2$ are invertible so is their tensor product.} 

If $\call_1|_U\simeq \calo_X|_U$ and $\call_2|_U\simeq \calo_X|_U$,  
$$
(\call_1\otimes\call_2)|_U\simeq \calo_X|_U\otimes\calo_X|_U\simeq \calo_X|_U,
$$ where the LHS is sheaf associated to presheaf, (given $V\subset U$) $V\longmapsto\call_1(V)\otimes\call_2(V)$ and the $\calo_X(V)\otimes \calo_X(V)\simeq \calo_X(V)$. As a presheaf it is isomorphic to the sheaf $\calo_{X}(V)$. So the LHS is the sheaf $V\longmapsto \calo_X(V)$ and $\call_1|_U\otimes \call_2|_U\simeq \calo|_U$ and therefore $\call_1\otimes \call_2$ is invertible.

To prove $\call\otimes \call^\vee\cong \calo_X$, recall that in general the stalk of sheaf-hom is not the hom of stalks (Though in some coherent case this is true, but we don't want to use this result. See \href{https://mathoverflow.net/questions/642/stalks-of-sheaf-hom}{StackExchange} for a review).

Instead,
\underline{Claim}: For any $U\subset X$, we have an isomorphism of modules, $\call^\vee(U)\cong \hom_{\calo_X(U)}(\call(U),\calo_X(U))$ as $\calo_X(U)$-modules.
where 
$$
\begin{aligned}
\call^\vee(U):=\hom_{\calo_X|_U-mod}(\call|_U,\calo_X|_U)&\lrta \hom_{\calo_X(U)-mod}(\call(U),\calo_X(U))\\
\lambda&\longmapsto \lambda(U).
\end{aligned}
$$
We can construct the inverse by mapping $\phi\in \hom_{\calo_X(U)-mod}(\call(U),\calo_X(U))$ to $\res_{U,V}(\phi)$ and identify all the data $\res_{U,V}(\phi)$ as $\lambda(V)$ of a morphism of sheaves $\lambda$.

 If we further require $U$ is contained in some open set $V$, where $\call|_V\cong \calo_X|_V$
 $$
 \begin{aligned}
 (\call\odot \call^\vee)(U)&\cong \call(U)\otimes\call^{\vee}(U)\cong \call(U)\otimes\hom_{\calo_X(U)-mod}(\call(U),\calo_X(U))\overset{(*)}{\lrta}\calo_X(U)\\
  (s,\lambda)& \longmapsto (\lambda(U))(s),
 \end{aligned}
 $$
 The last isomorphism $(*)$ is due to fact $\call(U)\otimes_{\calo(U)}(\call(U))^\vee\cong End_{\calo_X(U)}(\call(U))=\calo_X(U)$ because $\call(U)\cong \calo_X(U).$  

 The above construction gives an isomorphism of presheaves on $U$. After sheafification, we get an morphism of sheaves on $U$, which induces isomorphism on stalks.

 In the end, we have a collection of isomorphism on a base $\{U_i\}$ of topology. If or choice of base refines the defining chart of locally free sheaves $\call$, we can check the isomorphisms on base are compatible on intersections and we can glue it to an isomorphism of sheaves on $X$.
\end{proof}

\begin{exr}
Recall that the tensor products tend to be only right-exact in general. Show that tensoring by a locally free sheaf is exact. More precisely, if $\calf$ is a locally free sheaf, and $\calg'\lrta \calg\lrta \calg''$ is an exact sequence of $\calo_X$-module, then so is $\calg'\otimes \calf\lrta \calg\otimes \calf\lrta\calg''\otimes \calf$.
\end{exr}
\begin{proof}[Sol.]
This theorem should intrinsically depends on the fact the free modules are flat. (in general projective module is flat).

The exactness of sequence of sheaves should be checked at the level of stalks. Given $\calg'\lrta\calg\lrta \calg''$ to be an exact sequence of $\calo_X$-modules and let $\calf$ be a locally free sheaf. Then for any $x\in X$, we have $\calf|_U\cong \calo_U^{I}$, for any open neighborhood of $x$, which would imply $\calf_x\cong \calo_x^{I}$. Hence, we have
$$
(\calg\otimes\calf)_x\cong \calg_x\otimes_{\calo_x}\calf_x.
$$
Also recall that the sequence 
$$
\calg_x'\lrta\calg_x\lrta\calg''_x
$$
is an exact sequence of $\calo_{X,x}$-modules. As we have shown, $\calf_x$ is a free $\calo_{X,x}$-module, tensoring with it would still give an exact sequence of $\calo_{X,x}$-module due to the commutative algebra facts.

The propose the sequence is exact at the level of stalks therefore is an exact sequence of sheaves.
\end{proof}
\begin{exr}
If $\cale$ is a locally free sheaf of finite rank, and $\calf$ and $\calg$ are $\calo_X$-modules, show that $\phom(\calf,\calg\otimes\cale)\cong \phom(\calf\otimes \cale^\vee,\calg)$.
\end{exr}
\begin{proof}[Sol.]
First, we consider the case where $\cale$ if a free sheaf of finite rank i.e. $\cale\cong \calo_X^n$, and check that in this way $\cale^\vee\cong\phom(\calo_X^n,\calo_X)$. We want to prove
$$
\phom(\calf,\calg\otimes\calo_X^n)\cong \phom(\calf\otimes \phom(\calo_X^n,\calo_X),\calg)
$$

Note that functors $\phom$ and $\otimes$ are additive, we can reduce the problem to the case $n=1$.
$$
\Lrta \phom(\calf,\oplus^n(\calg\otimes\calo_X))\overset{?}{\cong}\phom(\calf\otimes \oplus^n\phom(\calo_X,\calo_X),\calg)
$$
$$
\Lrta \bigoplus^n\phom(\calf,(\calg\otimes\calo_X))\overset{?}{\cong}\bigoplus^n\phom(\calf\otimes \phom(\calo_X,\calo_X),\calg)
$$

\underline{Claim}: 
$$
\phom(\calf,\calg\otimes\calo_X)\cong\phom(\calf\otimes \calo_X^\vee,\calg)
$$

This isomorphism is an obvious fact due to the adjointness of tensor and hom functors.

\underline{Subclaim}: 
$$
\calf\otimes\calo_X^\vee\cong \phom(\calo_X,\calf)
$$
 We first deal with the subclaim. We have shown that $\calo_X^\vee\cong\calo_X$ in Exercise~\ref{chap13exr:dual_and_double_dual_of_locally_free sheaves} and have verified $\phom(\calo_X,\calf)\cong\calf$ in $(c)$ of Exercise~\ref{chap2exr:OX-moduel-hom_from_special_Ox-modules}. Combining this two, we have the isomorphism $\calf\otimes\calo_X^\vee\cong\calf\cong\phom(\calo_X,\calf)$.

Then the claim can be re-expressed as
$$
\phom(\calf,\calg\otimes\calo_X)\cong\phom(\phom(\calo_X,\calf),\calg).
$$
This has noting to do with the proof of the claim, but we want give a relation similar but \textbf{not identical} to the adjointness of $\otimes \cale$ and $\phom(\cale,\cdot)$, see~\href{http://stacks.math.columbia.edu/tag/01CN}{Tag 01CN} of StackProject. 

The claim it self can be checked similarly to the subclaim:
$$
\phom(\calf,\calg\otimes\calo_X)\cong\phom(\calf,\calg)\cong\phom(\calf\otimes \calo_X^\vee,\calg).
$$
We have finished the proof in the case where $\cale$ is free. It is easy to generalize to the case of locally free sheaves.

For each point $x\in X$, there is an open neighborhood  $U\ni x$ such that 
$$
\cale|_U\cong \calo_U^n
$$
and we can check that
$$
\begin{aligned}
&\phom(\calf,\calg\otimes\cale)(U)\\
&=\mor_{\calo_U-Mod}(\calf|_U,(\calg\otimes_{\calo_X}\cale)|_U)\\
&=\mor_{\calo_U-Mod}(\calf|_U,\calg|_U\otimes_{\calo_U}\cale|_U)\\
&\cong \mor_{\calo_U-Mod}(\calf|_U,\calg|_U\otimes_{\calo_U}\calo_U^n)\\
&\cong \mor_{\calo_U-Mod}(\calf|_U\otimes\phom(\calo_U^n,\calo_U),\calg|_U)\\
\end{aligned}
$$
the last isomorphism is due to the free sheaf case we have proved above. 
Furthermore, we have
$$
\begin{aligned}
&\mor_{\calo_U-Mod}(\calf|_U\otimes\phom(\calo_U^n,\calo_U),\calg|_U)\\
&\cong \mor_{\calo_U-Mod}(\calf|_U\otimes\phom(\cale,\calo_X)|_U,\calg|_U)\\
&=\phom(\calf\otimes\phom(\cale,\calo_X),\calg)(U)
\end{aligned}
$$
and thus
$$
\phi(U):\phom(\calf,\calg\otimes\cale)(U)\cong\phom(\calf\otimes\cale^\vee,\calg)(U).
$$
The strategy of proving isomorphism of sheaves is always to firstly construct a morphism and check that it induces isomorphism on all stalks. What we have proved up to now means only the stalks are isomorphic. We have to prove that the isomorphisms of stalks are induced by a morphism of sheaves. (to check $\phi(U)$ is compatible with the restriction maps and hence renders a morphism of sheaves.) Explicitly, by~\href{http://stacks.math.columbia.edu/tag/01CN}{Tag 01CN} of StackProject, we only need to find an open cover of $X=\cup U_i$ and a collection of morphism of sheaves:
$$
\phi_i:\phom(\calf,\calg\otimes\cale)|_{U_i}\lrta \phom(\calf\otimes\cale^\vee,\calg)|_{U_i}.
$$
such that they are compatible on the intersections:
$$
\phi_i|_{U_i\cap U_j}=\phi_j|_{U_i\cap U_j}.
$$
If we choose the open cover refining the defining open cover of the locally free sheaf $\cale$, the compatibility condition is automatically satisfied.
\end{proof}

\begin{exr}
Show that the invertible sheaves on $X$ up to isomorphism form an abelian group under product. This is called the \textbf{Picard group} of $X$, and is denoted $\pic X$. 
\end{exr}
\begin{proof}[Sol.]
We only need to check 
\begin{enumerate}[label=(\alph*)]
\item $\call_1$ and $\call_2$ invertible sheaves $\Lrta \call_1\otimes\call_2$ is invertible.
\item $\call_1\otimes\calo_X\cong \call_1$
\item $\call_1\otimes\call_1^\vee\cong \calo_X$
\item $\call_1\otimes \call_2\cong\call_2\otimes\call_1$
\end{enumerate}
$(a)$: If $\call_1|_U\cong \calo_X|_U$ and $\call_2|_U\cong\calo_X|_U$, we want 
$$
(\call_1\otimes\call_2)|_U \cong \calo_X|_U.
$$
To avoid confusion, we use $\calf\odot\calg$ to denote the presheaf generated by $[V\longmapsto \calf(V)\otimes\calg(V)]$. In this way $\call_1\otimes\call_2=(\call_1\odot\call_2)^{sh}$.

Then we can check that as presheaves $(\call_1\odot\call_2)|_U\cong \calo_X|_U\odot \calo_X|_U\cong (\calo_X|_U)_{pre}$. Then the sheafification functor would send this isomorphism of presheaves to an isomorphism of sheaves
$$
(\call_1\otimes\call_2)|_U=((\call_1\odot\call_2)|_U)^{sh}\cong ((\calo_X|_U)_{pre})^{sh}\cong \calo_X|_U.
$$
$(b)$ For any $U\subset X$,
 $$
 \begin{aligned}
 \call(U)\otimes\calo_X(U)&\overset{\simeq}{\lrta}\call(U)\\
 (s,a)&\longmapsto as
 \end{aligned}
 $$
 this induces an isomorphism of presheaves
 $$
 \call\odot \calo_X\lrta \call_{pre}.
 $$
 This induces a morphism of sheaves after sheafification. At the level of stalks,  it induces an isomorphism
 $$
\call_x\otimes_{\calo_{X,x}}\calo_{X,x}\lrta\call_x.
 $$
 $(c)$ is already proved in~\ref{chap13exr:tensor_product_of_invertible_sheaves}.\\
$(d)$ We can check the isomorphism of the two presheaves, 
 $$
\call_1\odot \call_2 \cong \call_2\odot \call_1,
 $$
 where for any open  set $U\subset X$,
 $$
 (\call_1\odot \call_2) (U):=\call_1(U)\otimes\call_2(U) \cong \call_2(U)\otimes\call_1(U)=:  \call_2\odot \call_1(U).
 $$
We get the desired isomorphism of sheaves after sheafification.

\end{proof}
\begin{exr}
If $\pi:X\lrta Y$ is a morphism of ringed spaces, and $\calg$ is a locally free sheaf of rank $n$ on $Y$, figure out how to define a ``pull back'' locally free sheaf of rank $n$ on X, denoted $\pi^*\calg$. Then show that $\pic$ is a contravariant functor from the category of ringed spaces to the category $\mathsf{Ab}$ of abelian groups.
\end{exr}
\begin{proof}[Sol.]
In standard references, the pullback of quasi-coherent sheaf on ringed space is defined to be
$$
\pi^{-1}\calg\otimes_{\pi^{-1}\calo_Y}\calo_X.
$$
We first consider the simplest case where $\calg$ is chosen to be $\calo_Y$. Is $\pi^*\calo_Y$ a well defined locally free sheaf? In fact, we expect
$$
\pi^*\calo_Y\cong \calo_X.
$$

Recall the inverse image sheaf $\pi^{-1}\calo_Y$ on $X$
$$
\pi^{-1}_{pre}\calg (U):= \colim_{V\supset \pi(U)}\calg(V)
$$
and $\pi^{-1}\calg:=(\pi^{-1}_{pre}\calg)^{sh}$.

\underline{Warning}: We can not use directly the definition of inverse image sheaf, which is in general not an $\calo_X$-module.

\underline{Fact1}: $\calo_X$ is a $\pi^{-1}\calo_Y$-module.
Given the morphism of ringed spaces $\pi=(\pi,\pi^\#)$, where $\pi^\#:\calo_Y\lrta \pi_*\calo_X$. If we apply the inverse image functor on it, we get a morphism of sheaves
$$
\pi^{-1}\calo_Y\overset{\pi^{-1}(\pi^\#)}{\lrta}\pi^{-1}\pi_{*}\calo_X
$$
By the adjoint relation between $\pi^{-1}$ and $\pi_*$ (c.f.~\ref{chap2exr:adjoint_inverse_image_direct_image}, we have a morphism
$$
\pi^{-1}\pi_{*}\calo_X\lrta\calo_X.
$$
Composite the above two morphisms of sheaves together, we get a morphism
$$
\pi^{-1}\calo_Y\lrta\calo_X.
$$
It is  a morphism of sheaves of rings, we know $\calo_X$ can be regarded as a $\pi^{-1}\calo_Y$-module. 

On the other hand, for each point $y\in Y$, there is an open set $V\subset Y$ s.t. $\calg|_V\cong(\calo_Y|_V)^n$. Functors map isomorphisms to isomorphisms. We have
$$
(\pi^{-1}\calg)|_{\pi^{-1}(V)}\cong \left((\pi^{-1}\calo_Y)|_{\pi^{-1}(V)}\right)^n.
$$
Then, because of the algebraic fact
$$
B^n\otimes_B A\cong A^n
$$
where $A$ is a $B$-algebra.
We can establish an isomorphism of presheaves, for $U\subset \pi^{-1}V$,
$$
(\pi^{-1}\calg\odot_{\pi^{-1}\calo_Y}\calo_X)|_U\lrta (\calo_X|_U)^n.
$$

We need to check $\pi^*$ is compatible with isomorphism class of invertible sheaves.
If $\calg'\cong\calg$ are isomorphic sheaves on, we can check $\pi^*\calg'\cong\pi^*\calg$, which is guaranteed by the functorial property of $\pi^*$.
$$
\begin{tikzcd}
X \ar[r, "\pi"] \ar[d, "\pic"] & Y \arrow[d, "\pic"] \\
\text{Pic}\ X & \text{Pic}\ Y \arrow[l, "\pi^*"]
\end{tikzcd}
$$
\end{proof}
\begin{exr}
Suppose $s$ is a section of locally free sheaf $\calf$ on a scheme $X$. Define the notion of the subscheme cut out by $s=0$, denoted $V(s)$. Be sure to check that your definition is independent of choices.
\end{exr}
\begin{proof}[Sol.]
We have defined in Exercise~8.1.I what it means by a vanishing subscheme cut out be a global section in $\Gamma(X,\calo_X)$. In order to generalize this to the locally free sheaves, we don't have to use the affine chart, but only use the ``defining free chart'' of the locally free sheaf $\calf$. Take an open set  $U\subset X$ s.t. $\calf|_U\cong \calo_X^n|_U$. On $U$, $s$ can be expressed as a tuple of functions $s^i\in \Gamma(U,\calo_X|_U)$, we can then locally define $V_U(s):=\cap_i V_U(s^i) $ or the fibered product $V_U(s^1)\times_{U} V_U(s^2)\times_{U}\cdots\times_U V_U(s^n)$, where each $V_U(s^i)$ is already defined in Exercise~8.1.I. $V(s)$ is to be defined by gluing up $V_U(s)$. It is dependent of the choice of $U$, because for two intersecting free charts $x\in U_i\cap U_j$, $T_{ij}(x)s_i(x)=s_j(x)$, $s_i(x)=0\Llrta s_j(0)$, because $T_{ij}(x)\in GL(\calo_{X}(U_i\cap U_j)^n)$.

This proof works for locally free sheaves of finite rank. For infinite rank case, we can not argue as $GL(\calo_{X}(U_i\cap U_j)^\infty)$, but just remember $T_{ij}$ is still an isomorphism.

The coordinate-free proof is \textbf{TO BE ADDED}.
\end{proof}
\begin{exr}
Show that locally free sheaves on locally Noetherian normal schemes satisfy ``Hartog's Lemma''. sections defined away from a set of codimension at least 2 extends over that set. 
\end{exr}

\begin{exr}
Suppose $s$ is a nonzero rational section of an invertible sheaf on a locally Noetherian normal scheme. Show that if $s$ has no poles, then $s$ is regular.
\end{exr}
\begin{exr}
Omitted.
\end{exr}
\begin{exr}
Omitted.
\end{exr}

\section{Quasicoherent sheaves}
Given an $A$-module $M$, we define an $\calo_{\spec A}$-module $\tilde{M}$ on $\spec A$, where the sections over $D(f)$ were $M_f$.
\begin{definition}
If $X$ is a scheme, then an $\calo_X$-module $\calf$ is \textbf{quasicoherent} if for every affine open subset $\spec A\subset X$, $\calf|_{\spec A}\cong\tilde{M}$ for some $A$-module $M$. By theorem 13.2.1 It suffices to check on an affine open cover of $X$ due to Affine communication lemma.
\end{definition}
\begin{exr}
(NOT every $\calo_X$-module is quasicoherent)
\begin{enumerate}[label=(\alph*)]
\item Suppose $X=\spec k[t]$. Let $\calf$ be the skyscrper sheaf supported at the origin $[(t)]$, with group $k(t)$ and the usual $k[t]$-module structure. Show that this is an $\calo_X$-module  but not quasicoherent sheaf.
\item Suppose $X=\spec k[t]$. Let $\calf$ be the skyscraper sheaf supported at the generic point $[(0)]$, with group $k(t)$. Given this the structure of an $\calo_X$-module. Show that this is a quasicoherent sheaf. describe the restriction maps in the distinguished topology of $X$
\end{enumerate}
\end{exr}
\begin{proof}[Sol.]
\begin{enumerate}[label=(\alph*)]
\item 
\end{enumerate}
\end{proof}
\chapter{Line bundles: Invertible sheaves and divisors}
\section{Some line bundles on projective space}
\begin{exr}
Show that $\dim_K \Gamma(\proj^1,\calo(n))=n+1$ if $n\geq 0$ and $0$ otherwise.
\end{exr}
\begin{proof}[Sol.]
A global section would be a polynomial  $f(x_{1/0})\in K(x_{1/0})$ and a polynomial $g(x_{0/1})\in K[x_{0/1}]$ such that $f(1/x_{0/1})x_{0/1}^n=g(x_{0/1})$. $f$ must be a polynomial of degree less than $n$
$$
f(x_{1/0})=a_{0}x_{1/0}^n+a_{1}x_{1/0}^{n-1}+\cdots+a_n
$$
where $a_i\in K$. Therefore, we know 
$$
\dim_K\Gamma(\proj^1_K,\calo(n))=n+1
$$
for $n\geq 0$.

In the case where $n<0$, only such polynomial would be $f=g=0$. 
\end{proof}
\begin{exr}
Show that if $m\neq n$, then $\calo(m)\neq \calo(n)$. Hence conclude that we have an injection of groups $\intg\inj \pic\proj^1_K$ given by $n\mapsto \calo(n)$.
\end{exr}
\begin{proof}[Sol.]
For $m\neq n$,  we know $\calo(m)\not\cong \calo(n)$ because they have distinct transition functions.
\end{proof}
\begin{exr}
Show that $\dim_K\Gamma(\proj^m_K,\calo_{\proj^m_K}(n))={m+n\choose m}$.
\end{exr}
\begin{proof}[Sol.]
To find a global section of $\calo_{\proj^m(n)}$ means to find an $m+1$ polynomials
$f_i\in K[x_{0/i},...,x_{m/i}]/(x_{i/i}-1)$ such that
$$
f_i(x_{0/j}\times x_{i/j}^{-1} ,...,x_{m/j}\times x_{i/j}^{-1} )\times x_{i/j}^{n}=f_j(x_{0/j},...,x_{m/j}).
$$
This means each $f_i$ has degree less than or equal to $n$. $f_0$ must be of the form
$$
\sum_{|J|\leq n} a_J x_J
$$
where $J$ is a $\intg^{m}$-valued vector such that $|J|=\sum J_k\leq n$. The degree of freedom sums as
$$
{n+m-1\choose m-1}+{n+m-2\choose m-1}+\cdots+{m-1\choose m-1}
$$
which equals
$$
{n+m\choose m}.
$$
This combinatorial equality can be proved by induction.
\end{proof}
\begin{exr}
Show that every invertible sheaf on $\proj^1_K$ is of the form $\calo(n)$ for some $n$.
\end{exr}
\begin{proof}[Sol.]
Exercise~13.2.C said that all invertible sheaves on $\affn^n_K$ is trivial. We know any invertible sheaf on $\proj^1_K$ when restricted to $U_0\cong \affn^1_K$ is trivial. It suffices to show that the possible transition function is of the form in the definition of $\calo(n)$.
$U_0\cong \spec K[x_{1/0}]$ and $U_1\cong\spec K[x_{0/1}]$. There are only two open sets in the open cover, we hence know that the cocycle condition reduces to
$T_{ij}=T_{ji}^{-1}$. Here, we don't require $T_{ij}$ to be ring homomorphism, but only $K$-linear  transformation. It is determined by $f(x_{1/0})\mapsto f(x_{0/1})\times \lambda x_{0/1}^n$ for $\lambda \in K^\times$, with an obvious inverse. For the isomorphism class of invertible sheaves, $\lambda$ is irrelevant.
\end{proof}

\chapter{Quasicoherent sheaves and projective $A$-schemes}
\chapter{Pushforwards and pullbacks of quasicoherent sheaves}
\section{Introduction}
\section{Pushforward of quasicoherent sheaves}
\section{Pullbacks of quasicoherent sheaves}
\section{Line bundles and maps to projective spaces}
\section{The Curve-to-Projective Extension Theorem
}
\section{Ample and very ample line bundles}
\section{The Grassmannian as a moduli space}
\part{More}

%----------------------------------------------------------------------------------------
%	BIBLIOGRAPHY
%----------------------------------------------------------------------------------------

\chapter*{Bibliography}
\addcontentsline{toc}{chapter}{\textcolor{ocre}{Bibliography}}

%------------------------------------------------

\section*{Articles}
\addcontentsline{toc}{section}{Articles}
\printbibliography[heading=bibempty,type=article]

%------------------------------------------------

\section*{Books}
\addcontentsline{toc}{section}{Books}
\printbibliography[heading=bibempty,type=book]

%----------------------------------------------------------------------------------------
%	INDEX
%----------------------------------------------------------------------------------------

\cleardoublepage
\phantomsection
\setlength{\columnsep}{0.75cm}
\addcontentsline{toc}{chapter}{\textcolor{ocre}{Index}}
\printindex

%----------------------------------------------------------------------------------------

\end{document}
